% !TeX program = xelatex
% Standalone English volume. Use XeLaTeX until the references stabilize.
\documentclass[a4paper,11pt]{article}
\usepackage[T1]{fontenc}
\usepackage{lmodern}
\usepackage[english]{babel}
\usepackage[left=26mm,right=26mm,top=25mm,bottom=25mm,headheight=15pt,headsep=7mm,footskip=11mm]{geometry}
\usepackage{amsmath,amssymb,mathtools,amsthm}
\newtheorem{proposition}{Proposition}[section]
\newtheorem{lemma}[proposition]{Lemma}
\newtheorem{theorem}[proposition]{Theorem}
\newtheorem{corollary}[proposition]{Corollary}
\newtheorem{maintheorem}[proposition]{Main Theorem}
\theoremstyle{definition}
\newtheorem{definition}[proposition]{Definition}
\theoremstyle{remark}
\newtheorem{remark}[proposition]{Remark}
\makeatletter
\renewcommand*\l@section{\@dottedtocline{1}{0em}{2.4em}}
\renewcommand*\l@subsection{\@dottedtocline{2}{1.45em}{3.45em}}
\renewcommand*\l@subsubsection{\@dottedtocline{3}{3.8em}{4.55em}}
\makeatother
\usepackage{xcolor,fancyhdr,titlesec,needspace,xurl}
\usepackage{longtable,multicol}
\usepackage{hyperref}
\definecolor{ink}{HTML}{1A3542}
\definecolor{muted}{HTML}{4F606A}
\hypersetup{unicode=true,colorlinks=true,linkcolor=ink,filecolor=ink,urlcolor=ink,citecolor=ink,bookmarksnumbered=true,
 pdfdisplaydoctitle=true,pdflang={en},
 pdftitle={The generalized Riemann hypothesis for Dirichlet L-functions. Edition 930, Volume IX},
 pdfsubject={A natural-language reconstruction of the proposed proof},
 pdfauthor={A mathematical reconstruction}}
\titleformat{\section}{\Large\bfseries\color{ink}}{\thesection}{.65em}{}
\titlespacing*{\section}{0pt}{.8em}{.65em}
\setlength{\parindent}{1.2em}
\linespread{1.045}
\setlength{\parskip}{.22em plus .05em minus .03em}
\setlength{\emergencystretch}{2em}
\setcounter{topnumber}{2}
\setcounter{tocdepth}{3}\setcounter{secnumdepth}{3}
\widowpenalty=4500
\clubpenalty=4500
\pagestyle{fancy}\fancyhf{}
\fancyhead[L]{\small\color{muted}Dirichlet GRH \textperiodcentered\ Volume I}
\fancyhead[R]{\small\color{muted}Dirichlet zeros and power intervals}
\fancyfoot[C]{\small\thepage}
\renewcommand{\headrulewidth}{.3pt}
\newcommand{\lead}[1]{\par\smallskip\noindent{\bfseries\color{ink}#1.}\enspace}
\newcommand{\sour}[2]{\unskip\penalty150\hspace{.35em}{\small\hyperlink{ref#1}{[#2]}}}
\titleformat{\subsection}{\large\bfseries}{\thesubsection}{.6em}{}
\titlespacing*{\subsection}{0pt}{1.15em plus .15em}{.4em}
\setlength{\abovedisplayskip}{9pt plus 2pt minus 3pt}
\setlength{\belowdisplayskip}{9pt plus 2pt minus 3pt}
\newcommand{\RR}{\mathbb R}\newcommand{\CC}{\mathbb C}\newcommand{\QQ}{\mathbb Q}
\newcommand{\ZZ}{\mathbb Z}\newcommand{\ord}{\operatorname{ord}}
\newcommand{\Frob}{\operatorname{Frob}}\newcommand{\Gal}{\operatorname{Gal}}
\newcommand{\len}{\operatorname{length}}\newcommand{\Res}{\operatorname*{Res}}
\newcommand{\wmean}{\operatorname{av}}\newcommand{\Ker}{\operatorname{ker}}
\newcommand{\Lh}{\mathcal L}\newcommand{\Gg}{\mathcal G}
\newcommand{\real}{\operatorname{Re}}\newcommand{\imag}{\operatorname{Im}}
\newcommand{\abs}[1]{\left|#1\right|}\newcommand{\eps}{\varepsilon}
\newcommand{\anchoritem}[2]{\par\Needspace{2.8\baselineskip}\hypertarget{ref#1}{}\noindent{\bfseries\color{ink}#2.}\enspace}
\newcommand{\srcloc}[2]{{\small\nolinkurl{#1.lean}}, line~#2}
\numberwithin{equation}{section}
\raggedbottom

\newcommand{\Sel}{\mathcal S}
\newcommand{\Cl}{\operatorname{Cl}}
\newcommand{\Aut}{\operatorname{Aut}}
\newcommand{\GL}{\operatorname{GL}}
\newcommand{\OO}{\mathcal O}
\hypersetup{bookmarksdepth=3}


\addto\captionsenglish{\renewcommand{\partname}{Volume}}


\makeatletter
\expandafter\def\csname r@cor:actualmarks714\endcsname{{12.10}{38}{}{proposition.12.10}{Dirichlet-GRH-v930-Volume-II.pdf}}
\expandafter\def\csname r@cor:classfiber372\endcsname{{37.37}{70}{}{proposition.37.37}{Dirichlet-GRH-v930-Volume-VII.pdf}}
\expandafter\def\csname r@cor:coc145\endcsname{{25.15}{36}{}{proposition.25.15}{Dirichlet-GRH-v930-Volume-V.pdf}}
\expandafter\def\csname r@cor:dyadicorbit459\endcsname{{56.6}{53}{}{proposition.56.6}{Dirichlet-GRH-v930-Volume-XI.pdf}}
\expandafter\def\csname r@cor:eliminateupper472\endcsname{{57.3}{68}{}{proposition.57.3}{Dirichlet-GRH-v930-Volume-XI.pdf}}
\expandafter\def\csname r@cor:fieldrelations712\endcsname{{12.4}{32}{}{proposition.12.4}{Dirichlet-GRH-v930-Volume-II.pdf}}
\expandafter\def\csname r@cor:irreduciblelabel374\endcsname{{38.8}{75}{}{proposition.38.8}{Dirichlet-GRH-v930-Volume-VII.pdf}}
\expandafter\def\csname r@cor:leveltwentyfive467\endcsname{{56.15}{60}{}{proposition.56.15}{Dirichlet-GRH-v930-Volume-XI.pdf}}
\expandafter\def\csname r@cor:oddwindow677\endcsname{{11.30}{27}{}{proposition.11.30}{Dirichlet-GRH-v930-Volume-II.pdf}}
\expandafter\def\csname r@cor:projectsymbol483\endcsname{{68.17}{97}{}{proposition.68.17}{Dirichlet-GRH-v930-Volume-XIII.pdf}}
\expandafter\def\csname r@cor:scalarfrob713\endcsname{{12.7}{35}{}{proposition.12.7}{Dirichlet-GRH-v930-Volume-II.pdf}}
\expandafter\def\csname r@cor:selmergenerators410\endcsname{{42.11}{72}{}{proposition.42.11}{Dirichlet-GRH-v930-Volume-VIII.pdf}}
\expandafter\def\csname r@cor:unboundedprimes428\endcsname{{54.19}{18}{}{proposition.54.19}{Dirichlet-GRH-v930-Volume-XI.pdf}}
\expandafter\def\csname r@lem:collisionfamily464\endcsname{{56.7}{54}{}{proposition.56.7}{Dirichlet-GRH-v930-Volume-XI.pdf}}
\expandafter\def\csname r@lem:cyclointegral377\endcsname{{40.1}{2}{}{proposition.40.1}{Dirichlet-GRH-v930-Volume-VIII.pdf}}
\expandafter\def\csname r@lem:dual241\endcsname{{36.2}{2}{}{proposition.36.2}{Dirichlet-GRH-v930-Volume-VII.pdf}}
\expandafter\def\csname r@lem:finitech410\endcsname{{42.8}{70}{}{proposition.42.8}{Dirichlet-GRH-v930-Volume-VIII.pdf}}
\expandafter\def\csname r@lem:frameeuler365\endcsname{{37.7}{40}{}{proposition.37.7}{Dirichlet-GRH-v930-Volume-VII.pdf}}
\expandafter\def\csname r@lem:frobtest366\endcsname{{37.12}{45}{}{proposition.37.12}{Dirichlet-GRH-v930-Volume-VII.pdf}}
\expandafter\def\csname r@lem:imagespan374\endcsname{{38.5}{73}{}{proposition.38.5}{Dirichlet-GRH-v930-Volume-VII.pdf}}
\expandafter\def\csname r@lem:inertaux364\endcsname{{37.3}{35}{}{proposition.37.3}{Dirichlet-GRH-v930-Volume-VII.pdf}}
\expandafter\def\csname r@lem:localranks368\endcsname{{37.18}{53}{}{proposition.37.18}{Dirichlet-GRH-v930-Volume-VII.pdf}}
\expandafter\def\csname r@lem:parameterdepth409\endcsname{{42.5}{66}{}{proposition.42.5}{Dirichlet-GRH-v930-Volume-VIII.pdf}}
\expandafter\def\csname r@lem:rel70\endcsname{{32.12}{36}{}{proposition.32.12}{Dirichlet-GRH-v930-Volume-VI.pdf}}
\expandafter\def\csname r@lem:relativepro407\endcsname{{41.29}{57}{}{proposition.41.29}{Dirichlet-GRH-v930-Volume-VIII.pdf}}
\expandafter\def\csname r@lem:spandual367\endcsname{{37.14}{49}{}{proposition.37.14}{Dirichlet-GRH-v930-Volume-VII.pdf}}
\expandafter\def\csname r@lem:support35\endcsname{{31.4}{16}{}{proposition.31.4}{Dirichlet-GRH-v930-Volume-VI.pdf}}
\expandafter\def\csname r@lem:testpoints714\endcsname{{12.8}{36}{}{proposition.12.8}{Dirichlet-GRH-v930-Volume-II.pdf}}
\expandafter\def\csname r@prop:auxfield361\endcsname{{36.15}{15}{}{proposition.36.15}{Dirichlet-GRH-v930-Volume-VII.pdf}}
\expandafter\def\csname r@prop:classrank365\endcsname{{37.6}{39}{}{proposition.37.6}{Dirichlet-GRH-v930-Volume-VII.pdf}}
\expandafter\def\csname r@prop:curvelattice366\endcsname{{37.11}{44}{}{proposition.37.11}{Dirichlet-GRH-v930-Volume-VII.pdf}}
\expandafter\def\csname r@prop:finite15\endcsname{{11.4}{6}{}{proposition.11.4}{Dirichlet-GRH-v930-Volume-II.pdf}}
\expandafter\def\csname r@prop:localcomponent412\endcsname{{42.17}{77}{}{proposition.42.17}{Dirichlet-GRH-v930-Volume-VIII.pdf}}
\expandafter\def\csname r@prop:loss367\endcsname{{37.16}{50}{}{proposition.37.16}{Dirichlet-GRH-v930-Volume-VII.pdf}}
\expandafter\def\csname r@prop:mainselmer724\endcsname{{12.17}{43}{}{proposition.12.17}{Dirichlet-GRH-v930-Volume-II.pdf}}
\expandafter\def\csname r@prop:ordinarycorner414\endcsname{{42.23}{82}{}{proposition.42.23}{Dirichlet-GRH-v930-Volume-VIII.pdf}}
\expandafter\def\csname r@prop:orighecke374\endcsname{{38.7}{74}{}{proposition.38.7}{Dirichlet-GRH-v930-Volume-VII.pdf}}
\expandafter\def\csname r@prop:origideal373\endcsname{{38.2}{71}{}{proposition.38.2}{Dirichlet-GRH-v930-Volume-VII.pdf}}
\expandafter\def\csname r@prop:orthogonalblocks410\endcsname{{42.9}{70}{}{proposition.42.9}{Dirichlet-GRH-v930-Volume-VIII.pdf}}
\expandafter\def\csname r@prop:planelabels379\endcsname{{40.11}{8}{}{proposition.40.11}{Dirichlet-GRH-v930-Volume-VIII.pdf}}
\expandafter\def\csname r@prop:quadratic345\endcsname{{36.3}{3}{}{proposition.36.3}{Dirichlet-GRH-v930-Volume-VII.pdf}}
\expandafter\def\csname r@prop:strictclasses365\endcsname{{37.5}{38}{}{proposition.37.5}{Dirichlet-GRH-v930-Volume-VII.pdf}}
\expandafter\def\csname r@prop:tamecurve365\endcsname{{37.9}{42}{}{proposition.37.9}{Dirichlet-GRH-v930-Volume-VII.pdf}}
\expandafter\def\csname r@prop:tangentcurve409\endcsname{{42.6}{67}{}{proposition.42.6}{Dirichlet-GRH-v930-Volume-VIII.pdf}}
\expandafter\def\csname r@prop:testbounds368\endcsname{{37.19}{54}{}{proposition.37.19}{Dirichlet-GRH-v930-Volume-VII.pdf}}
\expandafter\def\csname r@prop:tracecrystal375\endcsname{{38.10}{76}{}{proposition.38.10}{Dirichlet-GRH-v930-Volume-VII.pdf}}
\expandafter\def\csname r@prop:upfamily372\endcsname{{37.34}{67}{}{proposition.37.34}{Dirichlet-GRH-v930-Volume-VII.pdf}}
\expandafter\def\csname r@sec:commonframe365\endcsname{{37.2}{38}{}{subsection.37.2}{Dirichlet-GRH-v930-Volume-VII.pdf}}
\expandafter\def\csname r@sec:evenroot358\endcsname{{36.3}{5}{}{subsection.36.3}{Dirichlet-GRH-v930-Volume-VII.pdf}}
\expandafter\def\csname r@sec:finiteexception467\endcsname{{56.5}{59}{}{subsection.56.5}{Dirichlet-GRH-v930-Volume-XI.pdf}}
\expandafter\def\csname r@sec:irrcurvesupport409\endcsname{{42.2}{66}{}{subsection.42.2}{Dirichlet-GRH-v930-Volume-VIII.pdf}}
\expandafter\def\csname r@sec:levelmass678\endcsname{{130.1}{3}{}{subsection.130.1}{Dirichlet-GRH-v930-Volume-XXV.pdf}}
\expandafter\def\csname r@sec:lower36\endcsname{{31.5}{20}{}{subsection.31.5}{Dirichlet-GRH-v930-Volume-VI.pdf}}
\expandafter\def\csname r@sec:markedcurve366\endcsname{{37.3}{43}{}{subsection.37.3}{Dirichlet-GRH-v930-Volume-VII.pdf}}
\expandafter\def\csname r@sec:missing\endcsname{{6.6}{32}{}{subsection.6.6}{Dirichlet-GRH-v930-Volume-I.pdf}}
\expandafter\def\csname r@sec:oddblock410\endcsname{{42.3}{69}{}{subsection.42.3}{Dirichlet-GRH-v930-Volume-VIII.pdf}}
\expandafter\def\csname r@sec:oddrelations411\endcsname{{42.4}{73}{}{subsection.42.4}{Dirichlet-GRH-v930-Volume-VIII.pdf}}
\expandafter\def\csname r@sec:reduciblepropagation412\endcsname{{42.5}{76}{}{subsection.42.5}{Dirichlet-GRH-v930-Volume-VIII.pdf}}
\expandafter\def\csname r@sec:seed361\endcsname{{36.7}{15}{}{subsection.36.7}{Dirichlet-GRH-v930-Volume-VII.pdf}}
\expandafter\def\csname r@sec:spacingdyadic464\endcsname{{56.3}{54}{}{subsection.56.3}{Dirichlet-GRH-v930-Volume-XI.pdf}}
\expandafter\def\csname r@sec:upper35\endcsname{{31.3}{15}{}{subsection.31.3}{Dirichlet-GRH-v930-Volume-VI.pdf}}
\expandafter\def\csname r@seriesguide\endcsname{{1.1}{4}{}{subsection.1.1}{Dirichlet-GRH-v930-Volume-I.pdf}}
\expandafter\def\csname r@thm:alllabels375\endcsname{{38.11}{77}{}{proposition.38.11}{Dirichlet-GRH-v930-Volume-VII.pdf}}
\expandafter\def\csname r@thm:auxdescent361\endcsname{{36.18}{18}{}{proposition.36.18}{Dirichlet-GRH-v930-Volume-VII.pdf}}
\expandafter\def\csname r@thm:auxtriangle713\endcsname{{12.6}{33}{}{proposition.12.6}{Dirichlet-GRH-v930-Volume-II.pdf}}
\expandafter\def\csname r@thm:commonframe365\endcsname{{37.8}{41}{}{proposition.37.8}{Dirichlet-GRH-v930-Volume-VII.pdf}}
\expandafter\def\csname r@thm:commonorbits459\endcsname{{56.5}{53}{}{proposition.56.5}{Dirichlet-GRH-v930-Volume-XI.pdf}}
\expandafter\def\csname r@thm:commonsplit483\endcsname{{68.15}{96}{}{proposition.68.15}{Dirichlet-GRH-v930-Volume-XIII.pdf}}
\expandafter\def\csname r@thm:conductor232\endcsname{{19.15}{12}{}{proposition.19.15}{Dirichlet-GRH-v930-Volume-IV.pdf}}
\expandafter\def\csname r@thm:derivation367\endcsname{{37.15}{50}{}{proposition.37.15}{Dirichlet-GRH-v930-Volume-VII.pdf}}
\expandafter\def\csname r@thm:end14\endcsname{{7.3}{90}{}{proposition.7.3}{Dirichlet-GRH-v930-Volume-I.pdf}}
\expandafter\def\csname r@thm:fourier145\endcsname{{25.14}{35}{}{proposition.25.14}{Dirichlet-GRH-v930-Volume-V.pdf}}
\expandafter\def\csname r@thm:full36\endcsname{{31.8}{20}{}{proposition.31.8}{Dirichlet-GRH-v930-Volume-VI.pdf}}
\expandafter\def\csname r@thm:heckedivisor345\endcsname{{36.4}{4}{}{proposition.36.4}{Dirichlet-GRH-v930-Volume-VII.pdf}}
\expandafter\def\csname r@thm:integer\endcsname{{30.11}{10}{}{proposition.30.11}{Dirichlet-GRH-v930-Volume-VI.pdf}}
\expandafter\def\csname r@thm:irrcurvesupport409\endcsname{{42.7}{68}{}{proposition.42.7}{Dirichlet-GRH-v930-Volume-VIII.pdf}}
\expandafter\def\csname r@thm:levelstrict345\endcsname{{22.17}{75}{}{proposition.22.17}{Dirichlet-GRH-v930-Volume-IV.pdf}}
\expandafter\def\csname r@thm:limitselmer407\endcsname{{41.32}{60}{}{proposition.41.32}{Dirichlet-GRH-v930-Volume-VIII.pdf}}
\expandafter\def\csname r@thm:markedcurve366\endcsname{{37.13}{46}{}{proposition.37.13}{Dirichlet-GRH-v930-Volume-VII.pdf}}
\expandafter\def\csname r@thm:modulepatch408\endcsname{{42.3}{64}{}{proposition.42.3}{Dirichlet-GRH-v930-Volume-VIII.pdf}}
\expandafter\def\csname r@thm:oddinteger677\endcsname{{11.29}{26}{}{proposition.11.29}{Dirichlet-GRH-v930-Volume-II.pdf}}
\expandafter\def\csname r@thm:oddrelations411\endcsname{{42.15}{76}{}{proposition.42.15}{Dirichlet-GRH-v930-Volume-VIII.pdf}}
\expandafter\def\csname r@thm:oddsupport410\endcsname{{42.10}{71}{}{proposition.42.10}{Dirichlet-GRH-v930-Volume-VIII.pdf}}
\expandafter\def\csname r@thm:pres70\endcsname{{32.13}{37}{}{proposition.32.13}{Dirichlet-GRH-v930-Volume-VI.pdf}}
\expandafter\def\csname r@thm:primes367\endcsname{{37.17}{51}{}{proposition.37.17}{Dirichlet-GRH-v930-Volume-VII.pdf}}
\expandafter\def\csname r@thm:reducibleminors378\endcsname{{40.7}{6}{}{proposition.40.7}{Dirichlet-GRH-v930-Volume-VIII.pdf}}
\expandafter\def\csname r@thm:reducibleorder379\endcsname{{40.12}{9}{}{proposition.40.12}{Dirichlet-GRH-v930-Volume-VIII.pdf}}
\expandafter\def\csname r@thm:reduciblepropagation412\endcsname{{42.18}{78}{}{proposition.42.18}{Dirichlet-GRH-v930-Volume-VIII.pdf}}
\expandafter\def\csname r@thm:reduciblerigid379\endcsname{{40.10}{8}{}{proposition.40.10}{Dirichlet-GRH-v930-Volume-VIII.pdf}}
\expandafter\def\csname r@thm:rigidtrace377\endcsname{{40.2}{2}{}{proposition.40.2}{Dirichlet-GRH-v930-Volume-VIII.pdf}}
\expandafter\def\csname r@thm:ring39\endcsname{{32.1}{23}{}{proposition.32.1}{Dirichlet-GRH-v930-Volume-VI.pdf}}
\expandafter\def\csname r@thm:scalarcost717\endcsname{{12.12}{39}{}{proposition.12.12}{Dirichlet-GRH-v930-Volume-II.pdf}}
\expandafter\def\csname r@thm:selmerlength368\endcsname{{37.20}{54}{}{proposition.37.20}{Dirichlet-GRH-v930-Volume-VII.pdf}}
\expandafter\def\csname r@thm:selmerlift369\endcsname{{37.25}{58}{}{proposition.37.25}{Dirichlet-GRH-v930-Volume-VII.pdf}}
\expandafter\def\csname r@thm:stableprime371\endcsname{{37.32}{65}{}{proposition.37.32}{Dirichlet-GRH-v930-Volume-VII.pdf}}
\expandafter\def\csname r@thm:trianglebudget714\endcsname{{12.9}{37}{}{proposition.12.9}{Dirichlet-GRH-v930-Volume-II.pdf}}
\expandafter\def\csname r@thm:unitray472\endcsname{{57.2}{67}{}{proposition.57.2}{Dirichlet-GRH-v930-Volume-XI.pdf}}
\expandafter\def\csname r@thm:upper35\endcsname{{31.5}{17}{}{proposition.31.5}{Dirichlet-GRH-v930-Volume-VI.pdf}}
\makeatother
\begin{document}
\setcounter{part}{8}
\fancyhead[L]{\small\color{muted}Dirichlet GRH \textperiodcentered\ Volume IX}
\fancyhead[R]{\small\color{muted}Hecke support and deformations}
\thispagestyle{plain}\begin{center}
{\LARGE\bfseries The generalized Riemann hypothesis\\[2pt]for Dirichlet \(L\)-functions}\\[6pt]
{\small 27 November 2026 \quad\textperiodcentered\quad Edition 930}\end{center}
\part{Fixed-weight Hecke support and symplectic deformations}\label{part:IX}
\begin{abstract}\noindent We study ordinary Hecke support at the prescribed arithmetic weight. Auxiliary polynomials and trace determinants relate strict cohomology to symplectic deformations and their ordinary equations.\end{abstract}
\setcounter{section}{43}
\noindent The \hyperref[seriesguide]{guide in Volume~I} lists all twenty-seven volumes.
\begingroup\small\setlength{\parskip}{0pt}\tableofcontents\endgroup
\clearpage

\section{Hecke support at the prescribed weight}\label{nav:chapter-fixed-weight-hecke}\label{nav:volume-IX}
\subsection{An integral module at the prescribed arithmetic weight}
\label{sec:fixedweightcorner415}

Fix a sufficiently distant interval \([V,V^{3/2}]\) of
Section~\ref{sec:missing} and a prime
\(\ell>\exp(V^\varepsilon/12)\) as in
Proposition~\ref{prop:auxfield361}. Use its common
abelian CM field \(L=L_a\) at the chosen depth. Let
\(E_x\) be an actual period field on the interval, let
\(1\le r<\lfloor x^{401}\rfloor\), and let \(\alpha\) be a
character of \(\operatorname{Gal}(E_x/\QQ)\) of sign
\((-1)^r\). Put \(\Psi=\alpha\kappa_\ell^r\) and
\(\vartheta=\Psi^{-1}\). Let \(O\) be the coefficient DVR
of the ordinary problem containing the values of \(\alpha\).
We describe a finite \(O\)-module for this very character
and weight, using the unlocalized finite algebra \(C_0\)
of Section~\ref{sec:reduciblepropagation412}.

\begin{lemma}[Separation from the fixed weight]\leavevmode
\label{lem:weightseparation415}
Let \(I_{\rm wt}\subset R_0\) impose the original ordered
inertial weights. At each selected curve prime \(\mathfrak p\)
one has
\[
         I_{\rm wt}(R_0)_{\mathfrak p}=(R_0)_{\mathfrak p},
               \qquad H_0\otimes_{R_0}(R_0/I_{\rm wt})=0 .
\]
\end{lemma}
\begin{proof}
At a chosen place above \(\ell\) the three curve characters
remain pairwise distinct on the kernel of the full cyclotomic
character in inertia by Proposition~\ref{prop:tamecurve365}.
Choose an element \(h\) of this kernel on which the first two
differ. In the original weights the two values coincide:
the algebraic weights are powers of \(\kappa_\ell\), the
conductor character is unramified there, and any common
normalizing twist cancels in their ratio. Thus
\(\psi_1(h)\psi_2(h)^{-1}-1\) belongs to \(I_{\rm wt}\)
and lies outside \(\mathfrak p\). It becomes a unit there.
The map to \(H_0\) factors through this localization,
proving both assertions.
\end{proof}

On \(G_L\) write the split diagonal characters at the original
weight in their given order as
\[
        (\chi_1,\chi_2,\chi_3)=\xi(\Psi,1,\Psi^{-1})|_{G_L},
\]
where \(\xi\) is the common cyclotomic normalization.
This diagonal determinant, with its ordered local characters,
defines a continuous map \(d:R_0\to O\). Indeed it has the
required residual characters, old inertia quotient, diagonal
ordinary flag and polarization (the outer lines pair and the
middle line pairs with itself). Write \(C_d=C_0\otimes_{R_0,d}O\).
Fix an embedding at \(\ell\) and choose
\(\gamma\in G_L\cap I_{\QQ_\ell}\) on which \(\bar\Psi\)
has order greater than two. Such an element exists because
\(\bar\kappa_\ell\) is onto on inertia of \(L\) and
\(0<2r<\ell-1\). Let \(e_i\) be the three ordered polynomial
projectors for \(\gamma\) in \(C_d\). Their denominators are
units in \(O\). We use \([g]\) for the group element in this
algebra and \(C_{ij}=e_i C_d e_j\).

\begin{proposition}[The ordinary corner and the split diagonal]\leavevmode
\label{prop:splitcorner415}
The projectors lift to \(C_0\); on the completion at a
selected two-constituent curve the corner from the first
projector to the second is isomorphic to \(N_0\) of
Proposition~\ref{prop:ordinarycorner414}. At the split
weight \(C_{ii}=Oe_i\). Set
\[
 B=C_{21}/(C_{23}C_{31}),\qquad
 u(g)=\chi_1(g)^{-1} e_2[g]e_1\pmod{C_{23}C_{31}}
                     \quad(g\in G_L),
\]
where products of submodules denote the spans of products.
Then \(B\) is finite over \(O\), the values of \(u\) generate
\(B\), \(u(\gamma)=0\), and
\[
                 u(gh)=u(g)+\vartheta(g)u(h).
\]
\end{proposition}
\begin{proof}
Over the complete local ring \(R_0\), the characteristic
polynomial of \(\gamma\) has three simple residual roots.
Hensel lifting and interpolation give orthogonal projectors
in the Cayley--Hamilton algebra whose sum is one. At the
curve the second lies on the character line: its reduction
in the stable lattice is the middle character, and the
other two residual roots are distinct from it. The first
idempotent projects to a primitive line in the plane.
In the semisimple residue algebra the two systems of matrix
projectors here and in Proposition~\ref{prop:orthogonalblocks410}
are conjugate after matching the labels. A unit realizing
this conjugacy lifts to a unit of the finite complete
algebra, since its kernel to the semisimple residue is the
Jacobson radical. Once two systems \((f_i),(f'_i)\) agree
modulo that radical, \(\sum_i f'_i f_i\) is a unit
intertwining them. This identifies their corresponding
corners over \(H_0\).

At \(d\) the split representation gives a surjection
\(C_d\to O^3\), as it contains all three coordinate
idempotents. Each primitive corner has determinant degree
one, by the same corner identity used in
Proposition~\ref{prop:orthogonalblocks410}. It is exactly
\(Oe_i\), and its scalar coordinate on \([g]\) is
\(\chi_i(g)\). For any \(a,a'\in C_d\), insert
\(1=e_1+e_2+e_3\) into \(e_2 a a'e_1\). The third product
vanishes in \(B\). The other two are the \(21\)-entry of
\(a\) times the first scalar of \(a'\), and the second
scalar of \(a\) times the \(21\)-entry of \(a'\).
Dividing this formula for group elements by \(\chi_1(gh)\)
proves the identity. Since \(\gamma\) commutes with its
projectors, \(u(\gamma)=0\). The original group elements
span \(C_d\) over \(O\) by the finiteness and closure
argument of Lemma~\ref{lem:finitech410}. Taking the corner
and the quotient proves the spanning and finiteness assertions.
The cocycle is continuous in the finite module topology.
\end{proof}

For \(g\in G_{\QQ}\) define an element of \(B\) by
\[
        v(g)=\frac{u(g\gamma g^{-1})}{1-\vartheta(\gamma)} .
\]
Normality of \(G_L\) makes the argument of \(u\) valid and
the denominator is a unit. Define \(B_{\alpha,r}^{\rm str}\)
to be the quotient of \(B\) by the submodule generated by
the following elements, for the indicated group elements:
\[
\begin{array}{ll}
 v(gh)-v(g)-\vartheta(g)v(h)             & g,h\in G_{\QQ},\\
 v(g)                                    & g\in G_{\QQ_\ell},\\
 v(g)                                    & g\in I_p\cap\ker\alpha,\quad p\ne\ell .
\end{array}
\]
For the last row fix one inertia group for each rational
prime. The first row ensures that the local conditions on
classes are independent of the chosen embeddings.
All the submodules in this quotient are finitely generated
since \(B\) is finite over the DVR. Its quotient topology
is Hausdorff. This construction keeps the rational
descent equations at primes where \(L/\QQ\) is ramified.

\begin{theorem}[The strict group at the original weight]\leavevmode
\label{thm:strictcorner415}
For every finite-length \(O\)-module \(W\), with the indicated
Galois action carried by its twist \(W(\vartheta)\), there
is a natural isomorphism
\[
  \operatorname{Hom}_O(B_{\alpha,r}^{\rm str},W)
   \ \simeq\
   \{z\in H^1(\QQ,W(\vartheta)):
          z_\ell=0,\ z|_{I_p}=0\ \text{for every }p\ne\ell\}.
\]
The equalities on the right are in local cohomology.
In particular \(B_{\alpha,r}^{\rm str}\) is the finite
\(O\)-linear dual, with values in \(\operatorname{Frac}(O)/O\),
of the strict group of Section~\ref{sec:evenroot358} after
extension of coefficients to \(O\).
\end{theorem}
\begin{proof}
A class on the right has a unique representative \(z\)
with \(z(\gamma)=0\), since \(\vartheta(\gamma)-1\) is a
unit. This representative is zero on the whole chosen
\(G_{\QQ_\ell}\): strictness makes it a coboundary there,
and its value at \(\gamma\) fixes the witness to zero.
At each \(p\ne\ell\) vanishing of the inertial class implies
\(z=0\) on \(I_p\cap\ker\alpha\), where the action is
trivial. Conversely this last condition makes the inertial
cocycle factor through \(\alpha(I_p)\), of order prime to
\(\ell\), so its cohomology class there is zero by averaging.

On \(G_L\) write \(d(g)=\operatorname{diag}(\chi_i(g))\).
Over the square-zero algebra \(O\oplus W\), form the matrices
\[
           \rho_z(g)=(1+z(g)E_{21})d(g).
\]
The ratio of the second diagonal character to the first
is \(\vartheta\), so the cocycle identity makes this a
representation. The determinant of every sum of these
group matrices, also after extending scalars, is the
product of its diagonal entries by triangularity. Thus
its determinant is exactly the split one as a polynomial
law. Away from \(\ell\) it kills the original specified
inertia kernels, since the finite part of the character
has prime-to-\(\ell\) order and reduces injectively on its
values. At all places of \(L\) above \(\ell\) the restricted
class is zero, because it comes from a rational strict
class. On each local group a change of basis by a matrix
\(1+t E_{21}\) therefore gives the split matrices. The ordinary matrix
words with the designated diagonal vanish in these
coordinates. The adjoints of those words vanish as well:
they are evaluated on the conjugate local groups, where
the class also splits, and the split diagonal itself
satisfies the adjoint words. The same reasoning applies
to adjoints of old inertia differences. These are all the
matrix conditions on \(C_0\). Continuity and the finite
coefficient reductions give a map from its continuous
group algebra quotient, using Lemma~\ref{lem:relativepro407}
for the relative pro-\(\ell\) quotient. With the determinant
fixed at \(d\) this is a map from \(C_d\).

The element \(\gamma\) maps to its diagonal matrix.
Consequently each \(e_i\) maps to the corresponding diagonal
matrix unit. Taking entry \(21\) defines an \(O\)-linear
functional on \(C_{21}\) with values in \(W\). It kills
\(C_{23}C_{31}\) because both off-diagonal entries in that
product vanish in these matrices. The induced functional
\(\varphi\) on \(B\) satisfies \(\varphi(u(g))=z(g)\)
for \(g\in G_L\). On the full rational group the cocycle
identity and \(z(\gamma)=0\) give
\[
        z(g\gamma g^{-1})=(1-\vartheta(\gamma))z(g).
\]
Thus \(\varphi(v(g))=z(g)\) for every \(g\in G_{\QQ}\).
The cocycle law and the two local conditions make
\(\varphi\) descend to \(B_{\alpha,r}^{\rm str}\).

Conversely any functional on this quotient gives the
continuous function \(g\mapsto\varphi(v(g))\). The defining
relations give a cocycle zero at the chosen coefficient
place and unramified elsewhere by the same averaging.
For \(g\in G_L\) the identity of the proposition applied
to \(g\gamma g^{-1}\) gives \(v(g)=u(g)\), so this cocycle
recovers the original functional on a spanning set of
\(B\). It is zero at \(\gamma\) and hence recovers the
normalized representative uniquely. The constructions
commute with coefficient maps and are \(O\)-linear by
their formulas on these generators.

Finally put \(D_O=\operatorname{Frac}(O)/O\). Every
continuous normalized cocycle with these discrete
coefficients has finite image by compactness. Its
strictness at \(\ell\) and the vanishing on the old inertia
kernels hold for its values inside a sufficiently large
finite submodule of \(D_O\), so the same construction
applies. Equivalently, one takes the direct limit of the
displayed identifications along \(D_O[\varpi^n]\). The
result is \(\operatorname{Hom}_O(B_{\alpha,r}^{\rm str},D_O)\).
The strict group over the character ring is finite by
Theorem~\ref{thm:levelstrict345}; extension of coefficients
commutes with these groups by its finite free basis argument.
A finite \(O\)-module with positive free rank would have
infinitely many homomorphisms into \(D_O\), so the
representing module is torsion. Duality on its finite
cyclic summands identifies it with the dual of the
stated strict group.
\end{proof}

\begin{corollary}[Integral minors on the actual period field]\leavevmode
\label{cor:strictminors415}
For each orbit of the indicated characters choose a finite
\(O\)-presentation \(O^m\xrightarrow{T_\alpha}O^s\to
B_{\alpha,r}^{\rm str}\to0\). Let \(c_\alpha\) be the
least normalized valuation of a nonzero \(s\)-rowed minor;
use zero when \(s=0\). Let \(O_\alpha\) be the character
DVR, \(f_\alpha\) its residue degree over \(\mathbf F_\ell\)
and \(e_\alpha'=e(O/O_\alpha)\). Then for the complementary
sign on \(E_x\) one has the exact identity
\[
 \operatorname{length}_{\ZZ_\ell}
     \mathcal S_{\ell,r}(E_x)^{c=(-1)^r}
       =\sum_{[\alpha]:\,\alpha(c)=(-1)^r}
                        \frac{f_\alpha}{e_\alpha'}\,c_\alpha .
\]
\end{corollary}
\begin{proof}
Since \(B_{\alpha,r}^{\rm str}\) is torsion, its presentation
has full row rank. Smith normal form gives its length
as \(c_\alpha\), also when the module is zero. By the theorem
this is the length of the corresponding strict group over
\(O\). Scalar extension of a cyclic module from
\(O_\alpha\) multiplies length by \(e_\alpha'\). On this
large-prime range \(\ell\nmid [E_x:\QQ]N_*\), so the
character projectors are integral and the unramified
conditions agree with the conductor conditions by
Theorem~\ref{thm:levelstrict345}. Its orbit formula now
gives the displayed equality.
\end{proof}

The matrices in this corollary present modules at the split
arithmetic weight, before curve localization. Section~\ref{sec:missing}
requires a uniform upper bound on the weighted sum of their
valuations as \(x,r,\ell\) vary and estimates for the signed
integral theta periods on the same interval. Both estimates
remain open in this manuscript.


\subsection{Exact descent of strict cohomology on a common auxiliary field}
\label{sec:strictcommonbase416}

Fix an interval \(V\le x\le V^{3/2}\) of Section~\ref{sec:seed361},
with largest period field \(E_*\), modulus \(N_*\) and
\(K=\QQ(\sqrt{-s})\) as there. Let \(\mathcal P\) be any finite set
of primes \(\ell>\exp(V^\varepsilon/12)\). Take \(V\) large enough
that \(\ell>sN_*\) and \(\ell>4t_*+5\) for all of them.
For each \(\ell\in\mathcal P\) take all the pairs \(1\le r<t_*\),
\(\alpha\in\widehat{\operatorname{Gal}(E_*/\QQ)}\) with
\(\alpha(c)=(-1)^r\). When a character is evaluated on a Galois
group it is understood to be inflated. Put
\[
 O_\alpha=\ZZ_\ell[\alpha],\quad
 k_\alpha=O_\alpha/\ell O_\alpha,\quad
 \vartheta=\alpha^{-1}\kappa_\ell^{-r},\qquad
 \delta=\vartheta^{-1}\kappa_\ell .
\]
Here \(O_\alpha\) is unramified since \(\ell\nmid[E_*:\QQ]\).
For a module \(U\) with Galois action write \(S_D(U)\) for the
classes in \(H^1(D,U)\) zero on decomposition groups above \(\ell\)
and zero on inertia at every other finite place. We use continuous
cohomology and impose these conditions on classes.

\begin{lemma}[One Frobenius for the dual classes and the auxiliary fields]\leavevmode
\label{lem:jointfrobenius416}
For every triple above one can choose a class
\(d_{\ell,r,\alpha}\in H^1(\QQ,O_\alpha(\delta))\) unramified
away from \(\ell\), and one can choose a single prime \(w\) with
all the following properties. It is one modulo \(sN_*\) and every
\(\ell\in\mathcal P\). For each such \(\ell\), the integer \(\ell\)
is a non-\(\ell\)-th power modulo \(w\). For every triple the
reduction of \(d_{\ell,r,\alpha}\) is nonzero on Frobenius at \(w\).
The primes satisfying these conditions have positive density.
\end{lemma}
\begin{proof}
Use the integral class of Lemma~\ref{lem:dual241}. Its proof
uses the oddness of \(\delta=\alpha\kappa_\ell^{r+1}\), positive
rank from the global Euler characteristic, and prime-to-\(\ell\)
inertia away from \(\ell\). These hypotheses hold for every pair
in the stated range, including pairs whose strict group vanishes.
Choose one class per orbit under residue-field Frobenius and
conjugate it for the other members. Its reduction is nonzero
globally. Let \(C=\QQ(\mu_{sN_*\prod_{\ell\in\mathcal P}\ell})\).
The reduction remains nonzero on \(G_C\). Indeed it is trivial
as a coefficient character on \(G_C\), and the kernel of
restriction in degree one is inflated from
\(\operatorname{Gal}(C/\QQ)\). In this abelian group complex
conjugation acts on the coefficients by \(-1\). Evaluating a
cocycle at \(cg=gc\) makes it a coboundary since \(2\) is a
unit.

For each \(\ell\) take also the homomorphism on \(G_C\) given
by its action on \(\ell^{1/\ell}\), with values in \(\mu_\ell\).
It is nonzero. Otherwise \(C\) would contain the real field
\(\QQ(\ell^{1/\ell})\). The polynomial \(X^\ell-\ell\) is
Eisenstein at \(\ell\); its real root generates a real field
missing its nonreal conjugates. Every subfield of the abelian
extension \(C/\QQ\) is Galois. Conjugation on this Kummer
homomorphism has character \(\bar\kappa_\ell\).

On \(G_C\) the reduced cohomology classes are additive
homomorphisms, with conjugation on their values by
\(\bar\delta\). For a fixed \(\ell\) each target \(k_\alpha\)
is simple as an \(\mathbf F_\ell[\operatorname{Gal}(C/\QQ)]\)-module.
In fact the values of \(\bar\delta\) generate \(k_\alpha\) over
\(\mathbf F_\ell\), since the cyclotomic factor is in the prime
field. The Kummer target and these targets for distinct orbits
are pairwise nonisomorphic. On inertia at \(\ell\) their
exponents are respectively \(1\) and \(r+1\) modulo \(\ell-1\).
These distinguish the Kummer factor and the different \(r\)'s,
even after residue-field conjugation. At fixed \(r\), isomorphism
of the simple modules would identify the reductions of the two
finite characters up to that conjugation, hence also their
prime-to-\(\ell\) lifts.

Consider the image of all the homomorphisms together. For fixed
\(\ell\) it surjects onto each simple target and is stable under
conjugation. The joint action on the targets factors through an
abelian group of order prime to \(\ell\), whose group algebra
over \(\mathbf F_\ell\) is a product of finite fields. Projecting
with its idempotents shows that the image is the whole direct
sum. The primary components for different \(\ell\)'s separate
in the finite additive image, so every set of coordinate values
occurs simultaneously. Choose an element with every coordinate
nonzero. The joint kernel is open in \(G_C\) and normal in
\(G_\QQ\). Chebotarev applied to the conjugacy class of the
chosen element in this finite Galois quotient gives a
positive-density set of primes splitting in \(C\). At each such
prime every reduced class evaluates nontrivially. The nonzero
Kummer coordinates mean that \(X^\ell-\ell\) has no root in
\(\mathbf F_w\), since \(w\) splits in \(\QQ(\mu_\ell)\).
Discard the finite set of ramified primes and of other forbidden
choices. For an empty list one needs only complete splitting in
\(\QQ(\mu_{sN_*})\).
\end{proof}

\begin{proposition}[Redundancy of the unramified condition at the chosen prime]\leavevmode
\label{prop:oneprimerelax416}
Let \(w\) be chosen by Lemma~\ref{lem:jointfrobenius416} and
let \(n\ge1\). Put \(M_n=(O_\alpha/\ell^n)(\vartheta)\).
Every class in \(H^1(\QQ,M_n)\) zero at \(\ell\) and unramified
at every finite prime outside \(\{\ell,w\}\) is unramified at
\(w\). The assertion holds also for
\(M_\infty=(\operatorname{Frac}(O_\alpha)/O_\alpha)(\vartheta)\).
\end{proposition}
\begin{proof}
Write \(T_n=(O_\alpha/\ell^n)(\delta)\). It is the Cartier dual
of \(M_n\) under the trace pairing, since \(O_\alpha\) is
unramified and \(\vartheta\delta=\kappa_\ell\).
The two unramified subgroups at \(w\) are exact annihilators
under local Tate duality. Their cup product factors through the
residue Galois group of cohomological dimension one. Each
unramified subgroup has order equal to the order of its
local invariants, by taking the cokernel of Frobenius minus one.
The local Euler characteristic formula at \(w\ne\ell\), together
with local duality in degree two, gives
\[
 |H^1(\QQ_w,M_n)|=
       |H^0(\QQ_w,M_n)|\,|H^0(\QQ_w,T_n)|.
\]
These equalities of orders give the asserted annihilators.

The unramified subgroup for \(T_n\) is the cyclic
\(O_\alpha\)-module \(T_n/(\delta(\operatorname{Frob}_w)-1)T_n\).
The reduction of the local value of \(d_{\ell,r,\alpha}\)
is nonzero by the lemma. Its value modulo \(\ell^n\) therefore
generates this whole cyclic module. Pair a global class as in
the proposition with each scalar multiple of this reduction
modulo \(\ell^n\). Global reciprocity makes the sum of the
local pairings zero. Every term away from \(w\) vanishes:
the class is zero at \(\ell\), at infinity the first cohomology
vanishes, and at the other finite places both classes are
unramified. Even at a prime ramifying in the coefficients,
their unramified cup product is inflated from the residue
Galois group using the inertia-invariant coefficient modules.
The pairing at \(w\) is consequently zero against the entire
unramified subgroup of \(T_n\). The annihilator calculation
gives the claim.~\cite[I, Sections~2 and~4]{IX-VII-V-milne98}

For divisible coefficients, pick a continuous cocycle representing
the class. Its image is finite by compactness. At \(\ell\) choose
a witness that it is a coboundary and include this witness in
a finite coefficient submodule containing the image. At primes
outside \(\{\ell,w\}\) where the coefficient is unramified the
inertial cocycle is identically zero. At the finitely many
conductor primes choose and include witnesses as well. The
cocycle and all its required conditions thus occur in one
\(M_n\), and the finite assertion applies.
\end{proof}

Choose \(v\equiv-1\bmod 8sN_*\prod_{\ell\in\mathcal P}\ell\)
prime, distinct from \(w\) and the primes already used.
For \(\ell\in\mathcal P\) let \(F_{\ell,w}\) be the degree
\(\ell\) subfield of \(\QQ(\mu_w)\) and \(F_{\ell,a}\) the
degree \(\ell^a\) subfield of \(\QQ(\mu_{\ell^{a+1}})\).
Given any depths \(a_\ell\ge0\), form the abelian CM field
\[
 A=K\,\QQ(\sqrt{sv})\prod_{\ell\in\mathcal P}
                       F_{\ell,w}F_{\ell,a_\ell}.
\]
For each \(\ell\in\mathcal P\) this contains the auxiliary field
\(L_{a_\ell}\) of Proposition~\ref{prop:auxfield361}. Frobenius
of \(\ell\) generates \(\operatorname{Gal}(F_{\ell,w}/\QQ)\)
by the Kummer condition. Thus the large local residue degree
required in Lemma~\ref{lem:inertaux364} is preserved. The prime
\(w\) was chosen before the depths.

\begin{theorem}[Exact strict descent on the auxiliary extension]\leavevmode
\label{thm:strictbasedescent416}
Let \(\ell\in\mathcal P\), take \(M=M_n\) or \(M_\infty\) as in
the proposition, and put \(\Delta=\operatorname{Gal}(A/\QQ)\).
Restriction gives an isomorphism
\[
                   S_\QQ(M)\ \simeq\ S_A(M)^\Delta .
\]
It also gives this isomorphism with \(A\) replaced by
\(L_{a_\ell}\). Both assertions hold simultaneously for every
indicated \(r,\alpha,n\), and for every choice of the depths.
\end{theorem}
\begin{proof}
At \(\ell\) the only ramified factor of \(A\) is
\(F_{\ell,a_\ell}\), which is totally ramified of degree
\(\ell^{a_\ell}\). A compositum with unramified local
extensions has the same ramification index. The index of
inertia for a completion of \(A\) inside rational inertia is
therefore a power of \(\ell\); its image under the residual
cyclotomic character is the whole \(\mathbf F_\ell^\times\).
Since \(1\le r<\ell-1\), an element acts on \(M\) by a scalar
whose difference from one is a unit. Thus \(M^{G_{A_u}}=0\)
at any \(u\mid\ell\), and also \(M^{G_A}=0\).
Inflation--restriction identifies the full global first
cohomology with the invariants upstairs. The local version
makes restriction at \(\ell\) injective.

Let \(z\) be the unique rational antecedent of an invariant
strict class. It is zero at \(\ell\). At any \(p\ne\ell,w\)
the index of inertia for \(A/\QQ\) is prime to \(\ell\).
Indeed only quadratic factors can ramify at \(s,v\), and
at any other prime in \(\mathcal P\) the ramification index
is a power of that prime. All remaining factors are unramified
there. Restriction followed by corestriction on inertia is
multiplication by its index, so \(z\) is unramified at \(p\).
Proposition~\ref{prop:oneprimerelax416} now gives the condition
at \(w\) too. Conversely, restriction preserves all the local
conditions. The same inertia and cohomology arguments apply
to the specified subfield \(L_{a_\ell}\), where the only
factor ramifying at \(\ell\) is again \(F_{\ell,a_\ell}\).
\end{proof}

\begin{corollary}[All odd coefficient primes on the period fields]\leavevmode
\label{cor:intervalbase416}
One can take a single \(A\) on the interval containing ordinary
auxiliary fields at prescribed depths for every contributing
large prime. Write \(H_x=E_xA\), \(\Delta_x=\operatorname{Gal}
(H_x/E_x)\) and \(e=(1+(-1)^rc)/2\), using complex conjugation
in the respective fields. For every \(x\) on the interval,
\(1\le r<\lfloor x^{401}\rfloor\), and every odd prime \(\ell\),
restriction gives
\[
       e\,\mathcal S_{\ell,r}(E_x)
          \ \simeq\ e\big(\mathcal S_{\ell,r}(H_x)^{\Delta_x}\big).
\]
Consequently the residual groups \(R_{\ell,r}\) in
Section~\ref{sec:missing} have exactly their original lengths
when their defining operators are applied to the invariant
groups on the right.
\end{corollary}
\begin{proof}
On the bounded interval only finitely many pairs \((E_x,r)\)
occur. For each pair only finitely many coefficient primes
have nonzero strict group by Proposition~\ref{prop:finite15}.
Include in \(\mathcal P\) all such primes above
\(\exp(V^\varepsilon/12)\). One may add any finite list of
other primes above that threshold. The field \(A\) ramifies only
at \(s,v,w\) and the primes in \(\mathcal P\), all prime to \(N_*\).
In particular \(A\cap E_*=\QQ\), since their intersection
would be unramified at every finite rational prime. This uses
Minkowski's discriminant bound. Therefore \(H_x/\QQ\) has the
product of the two Galois groups.

If \(\ell\notin\mathcal P\) is odd, the degree of \(H_x/E_x\)
is prime to \(\ell\): the degrees of the factors of \(A\) are
powers of two or of members of \(\mathcal P\). Global restriction
is an isomorphism onto invariants by inflation--restriction
and averaging. Local restriction is injective on decomposition
groups at \(\ell\) and on inertia away from \(\ell\) by
corestriction. This proves exact descent for the entire strict
group in this case.

If \(\ell\in\mathcal P\), then \([E_x:\QQ]\) is prime to
\(\ell\). Extend coefficients to an unramified ring splitting
its characters and take the \(\alpha\)-component for each
\(\alpha(c)=(-1)^r\). Over \(E_x\) this component is the
rational strict group with the additional twist \(\alpha^{-1}\).
Indeed after tensoring with that character it is the invariant
component, and restriction, both globally and locally, gives
prime-to-degree descent just as above. Applied to \(H_x/A\),
the same argument identifies its \(\alpha\)-component with
the strict group over \(A\) for this twisted module. These
identifications commute with \(\Delta_x\), since the two
Galois actions commute on cohomology. Tensoring with a finite
free coefficient ring commutes with cochains, kernels and
invariants by taking an integral basis. The theorem now gives
the isomorphism on every component; faithful scalar extension
gives it over \(\ZZ_\ell\). On the invariant group any two
lifts of \(c\) from \(E_x\) have the same action. Restriction
is equivariant for \(\operatorname{Gal}(E_x/\QQ)\), hence also
for the field averages in Section~\ref{sec:missing}. Applying
those averages and taking the prescribed torsion quotients
after taking invariants proves the assertion on \(R_{\ell,r}\).
\end{proof}

The corollary identifies the lengths on the original \(E_x\)
without a local correction at the auxiliary ramified prime.
It gives no bound for the size of the invariant groups. Such
a bound and the signed integral theta-period comparison on
these \(E_x\) are still required in Section~\ref{sec:missing}.


\subsection{Propagation along curves with two constituents}
\label{sec:twoblockchain417}

Fix one large-prime triple on the interval of Section~\ref{sec:seed361}
and its ordinary problem over \(L=L_a\), with real subfield \(F\).
Use the depth of Theorem~\ref{thm:commonframe365}, the original
ordinary determinant ring \(R_0\), its finite Galois algebra
\(C_0\), and the ordinary module \(M_K\) of
Theorem~\ref{thm:reduciblepropagation412}. The underlying
field may be taken inside the common field of
Corollary~\ref{cor:intervalbase416}. Write \(G=G_L\) and
\(G'=\ker(\kappa_\ell|_G)\). A prime of \(R_0\) is called
irreducible here if its determinant over a geometric residue
field is irreducible. The support of \(M_K\) consists exactly
of the primes containing the kernel of the good Hecke map.

Consider characteristic-\(\ell\) primes \(\mathfrak p\) of
\(R_0\) of quotient dimension one with the following concrete
matrix realizations. After a finite extension of the fraction
field of \(R_0/\mathfrak p\), take an ordinary representation
on a lattice over a complete DVR with finite residue field.
Its determinant and ordered local characters realize the
given point, its matrices satisfy the words of \(C_0\), and it
has the symmetric conjugate-dual polarization of
Proposition~\ref{prop:curvelattice366}. Over the fraction
field it is the sum of a line and an absolutely irreducible
plane, polarized separately. After removing the common
normalizing character the line reduces to \(1\) and the plane
has residual semisimplification \(\bar\Psi\oplus\bar\Psi^{-1}\).
At each chosen place above \(\ell\) the three characters on
its ordered flag are pairwise distinct on inertia in \(G'\).
These conditions are on the matrices of the curve. The selected
two-constituent curves of Theorem~\ref{thm:markedcurve366}
satisfy them.

\begin{lemma}[The projective image of the plane]\leavevmode
\label{lem:planecurve417}
At every such curve the plane remains absolutely irreducible
on \(G'\). In each nonempty relatively open subset of \(G'\)
there is an element having three distinct eigenvalues on the
full representation.
\end{lemma}
\begin{proof}
Work over an algebraic closure of the curve field. The two
local characters on the plane are among the characters of the
ordered flag. Indeed the characteristic polynomials on a
triangular local group have those three diagonal characters,
and semisimplification of its restriction to the two summands
respects this multiset. On inertia in \(G'\) the ratio of the
two plane characters takes a value different from one by
hypothesis. Its residue is one, since on that inertia all
three residual characters coincide. In characteristic \(\ell\)
a nonidentity element of \(1+z k[[z]]\) has infinite order:
raising it to an \(\ell\)-power raises the first nonzero term
of its difference from one to that power, and raising to an
integer prime to \(\ell\) leaves its order of vanishing unchanged.
Here \(z\) is a uniformizer, with the same argument after
finite extension. The projective image of the plane is
therefore infinite.

Its Zariski closure in \(\mathrm{PGL}_2\) is the whole group.
For the details, a proper positive-dimensional connected
subgroup is solvable. This follows also from its dimension
being at most two and the fact that a nontrivial connected
semisimple group has dimension at least three. Such a subgroup
preserves a line by triangularization. If it has nontrivial
unipotent radical, the common fixed line of that radical is
unique, hence is preserved by the full closure, contradicting
absolute irreducibility. Otherwise the identity component is a
torus. Its normalizer preserves its pair of eigenlines and
the full group must interchange them. Changing the sign on one
of those lines then intertwines the plane with its twist by a
nontrivial quadratic character \(\eta'\). Semisimplifying the
residual representations gives
\[
 \{\bar\Psi,\bar\Psi^{-1}\}
       =\{\bar\eta'\bar\Psi,\bar\eta'\bar\Psi^{-1}\}.
\]
The reduction of a nontrivial quadratic character is nontrivial
at odd \(\ell\). It must interchange the two distinct residual
characters, which would give \(\bar\Psi^4=1\). On inertia at
\(\ell\), \(\bar\kappa_\ell\) has full image and \(0<4r<\ell-1\).
This excludes the torus case.

The quotient \(G/G'\) is abelian. The projective image of
\(G'\) is therefore Zariski dense in \(\mathrm{PGL}_2\), since
it contains the image of the commutators and \(\mathrm{PGL}_2\)
is its own derived group. This gives the irreducibility
assertion. Every open subgroup \(U\subseteq G'\) has finite
index, so its projective image is also dense: finitely many
translates of a proper closed subset cannot cover the connected
irreducible group \(\mathrm{PGL}_2\). Divide the plane
representation by its line character. The closure of the
resulting image of \(U\) in \(\mathrm{GL}_2\) contains
\(\mathrm{SL}_2\). In fact the closure of its commutators lies
in \(\mathrm{SL}_2\), and its projective image is dense in
\(\mathrm{PGL}_2\), so its dimension is at least three.
Every translate of \(\mathrm{SL}_2\) contains a diagonal matrix
of its prescribed nonzero determinant having distinct entries
both different from one. Thus a coset of \(U\) cannot have
its image contained in the closed locus of matrices with a
repeated eigenvalue or eigenvalue one. The line in this
normalization has eigenvalue one. Every nonempty relatively
open subset of the profinite group \(G'\) contains a coset
of such a \(U\), as required.
\end{proof}

\begin{proposition}[One patched module for a finite list of curves]\leavevmode
\label{prop:jointpatch417}
For any finite list of the curves above, the auxiliary places,
the limiting group and the localized patched module can be
chosen simultaneously. Let \(q\) be the common number of
places and \(D=3([F:\QQ]+q)\). Adjoin a parameter \(s'\) as
in Section~\ref{sec:irrcurvesupport409}. Write \(T\) for the
patched ring with the trace relations, \(\mathcal P\) for the
finite patched module and \(P\) for the ordered splitting
factor over the augmented original ring. At an ordered curve
let \(H\) be the completed localization of \(T\) and
\(\mathcal P_H=H\otimes_T\mathcal P\).
If a prime in \(\operatorname{Supp}_H\mathcal P_H\) has
irreducible determinant, then there is a prime \(Z\) of \(H\)
in that support, with irreducible determinant, contained in
every prime where the determinant is irreducible or the
geometric fiber of the patched Galois algebra is nonsemisimple.
Here the curve parameter is sent to a uniformizer after a
finite extension and the ordering is compatible with the
residual roots.
\end{proposition}
\begin{proof}
We explain how the local calculations apply to this list.
At each curve the line and plane lattices can be taken as
the intersections with the given lattice. A coefficient-place
inertia element with three distinct residual eigenvalues gives
three integral eigenprojectors by Hensel lifting. Each
projector lies on one summand, so the lattice is their direct
sum, and the ordinary flag consists of sums of these
eigenlines at the chosen place. The same argument works
at the other coefficient places. The form on each summand is
symmetric and nondegenerate. All required changes of coordinates
over the fraction field and square roots for orthonormal
coordinates are defined over finite extensions.

On \(G'\) the two summands are absolutely irreducible and
inequivalent by the lemma and their dimensions. Hence the
block algebra with its square-zero copies of the off-diagonal
ordinary cohomology has full span on \(G'\) by the proof of
Lemma~\ref{lem:spandual367}. In that proof the only inputs for
the off-diagonal span are independence modulo coboundaries
and the vanishing of the intertwiners between the two summands
on \(G'\). Theorem~\ref{thm:derivation367} then detects its
twisted dual by the trace evaluations. Intersect the relatively
open dense regular loci from the lemma for the finitely many
curves, which preserves density. The finite matrix of cocycle
relations in Proposition~\ref{prop:loss367} and the simultaneous
Chebotarev argument of Theorem~\ref{thm:primes367} now give one
list of separating Frobenius elements and auxiliary levels for
all these curves. In this application the open groups are
pulled back to the common restricted-ramification group; the
maps on cyclotomic kernels are open surjections.

At the old places inertia has the same fixed prime-to-\(\ell\)
quotient. At the chosen place above \(p_0\) the local
intertwiner vanishes for every representation with the given
residual characters by Lemma~\ref{lem:frobtest366}. At
coefficient places the off-diagonal character ratios on
cyclotomic-kernel inertia are nontrivial by the assumed
separation. These are precisely the vanishing inputs of
Lemma~\ref{lem:localranks368}. The ranks of the upper
triangular spaces and the real skew-adjoint spaces depend only
on the two block sizes and the symmetric forms. Thus
Proposition~\ref{prop:testbounds368}, the local duality count
of Theorem~\ref{thm:selmerlength368}, and the lifting argument
of Theorem~\ref{thm:selmerlift369} apply with each curve's own
off-diagonal rank \(b\). In the lifting argument the projection
from the mixed upper triangular module onto the diagonal one
uses a local upper triangular representative for each chosen
ordinary class; these exist by the definition of its ordinary
subspace. The finite tests pass to the single limiting group
by the compactness proof of Theorem~\ref{thm:limitselmer407}.
This supplies the bound \(D\) on the paired ordinary rank,
the relations on diagonal cocycles and the identification of
off-diagonal cohomology used in Sections~\ref{sec:oddblock410}
and~\ref{sec:oddrelations411}.

Use these same auxiliary levels in
Theorem~\ref{thm:modulepatch408}. The comparison maps,
the monomial description of the graded module for the diamond
ideal and the trace relations are identities on the
ordinary modules at these levels. Their proofs use the residual
ordering and the discriminant element \(f\); they use a
classical family only in the separate assertion of nonzero
fiber at the curves selected there. We use the identities on
\(\mathcal P\) before deciding which fibers are nonzero.
In particular \(T\) maps onto \(P[1/f]\) with diamonds zero and
the diamond quotient of \(\mathcal P\) is
\((P\otimes_{R_0}M_K)[1/f]\). The discriminant \(f\) survives
at every curve in the list.

Now suppose the local module has an irreducible prime in its
support. It is nonzero and its residue fiber is nonzero by
Nakayama. Its depth is at least \(D+1\) by
Lemma~\ref{lem:parameterdepth409}: the diamond variables are
regular, and on their fiber the uniformizer and the translated
weight coordinates form a regular sequence. The coefficient
DVR \(\widehat{(O'[[s']])_{(\varpi)}}\), formed by completing
the localization, acts faithfully through \(H\). Indeed \(\varpi\)
is a nonzerodivisor on the nonzero diamond quotient, so no
power of \(\varpi\) vanishes in \(H\). Here \(O'\) is the
unramified extension of the original coefficient ring whose
residue field supplies the curve constants.

The block calculation of Proposition~\ref{prop:orthogonalblocks410}
and Theorem~\ref{thm:oddsupport410} now gives a finite off-diagonal
module \(N\) over \(H\). Its support is exactly the union of the
two loci in the proposition. In Corollary~\ref{cor:selmergenerators410}
the identification of its minimal generator space with the dual
of ordinary off-diagonal cohomology uses the matrix words and
the split local flags. The further conclusion \(b\ge1\) follows
here because the given irreducible prime lies in the support of
\(N\). With this observation and the coefficient DVR just
constructed, the relations of Theorem~\ref{thm:oddrelations411}
apply. Choose a minimal prime \(Z\) of the module support below
the given irreducible prime. Irreducibility persists under
generization: in the displayed block algebra it is the
nonvanishing of the ideal of off-diagonal products.
The prime \(Z\) is associated to \(\mathcal P_H\), hence
\(\dim H/Z\ge\operatorname{depth}_H\mathcal P_H\ge D+1\).
Proposition~\ref{prop:localcomponent412} puts \(Z\) below
every prime in \(\operatorname{Supp}_H N\), proving the claim.
\end{proof}

\begin{lemma}[Lifting a chain to the ordered curves]\leavevmode
\label{lem:chainorder417}
Let \(\mathfrak p_0,\ldots,\mathfrak p_h\) be a list as above.
Suppose \(\mathfrak e_i\subseteq
\mathfrak p_i\cap\mathfrak p_{i+1}\) are primes of \(R_0\)
for \(0\le i<h\). After a common extension of constants and
adjoining \(s'\), there are ordered curve primes
\(\mathfrak x_i\) of \(P[1/f]\) and primes
\(\mathfrak b_i\subseteq\mathfrak x_i\cap\mathfrak x_{i+1}\)
over \(\mathfrak e_i\). Every prime of \(R_0\) contained in an
endpoint \(\mathfrak p_i\) lifts to a prime below the designated
\(\mathfrak x_i\).
\end{lemma}
\begin{proof}
First record a lifting fact for a finite faithfully flat
integral algebra. Starting over \(\mathfrak p_0\), go down by
flatness to a prime over \(\mathfrak e_0\); go up by integrality
from that prime to one over \(\mathfrak p_1\). Alternating
these steps lifts the whole chain. Going down also lifts an
endpoint generization with the endpoint fixed.

Split the finitely many designated Frobenius polynomials and
realize the lattices and forms over finite extensions of the
curve fields. Choose a finite field \(k'\) containing all
their residue fields and let \(O'/O\) be the unramified
extension with residue field \(k'\). Apply the lifting fact
to \(R_0\to R'=O'\otimes_O R_0\). At a lifted curve the
residue field is a factor of
\(\kappa(\mathfrak p_i)\otimes_k k'\), where \(k\) is the
original residue field. This choice amounts to an embedding
of the constants. Extend the chosen curve field by \(k'\)
with that embedding. Its valuation ring is isomorphic to
\(k'[[z_i]]\) with a uniformizer \(z_i\). Over
\(B=R'[[s']]\) take the prime obtained by evaluating \(s'\)
at \(z_i\) in this ring. Evaluation is defined by convergence;
its image contains \(k'[[z_i]]\) from \(O'[[s']]\). The
residue field at the prime is therefore \(k'((z_i))\).
A lifted link \(\mathfrak e'_i\) extends to the prime
\(\mathfrak e'_i B\), because the quotient is the power
series ring over the domain \(R'/\mathfrak e'_i\); here
Noetherianity identifies the kernel of coefficient reduction
with the extended ideal. Its coefficients vanish on both
adjacent curves regardless of the two choices of uniformizer.
The same construction lifts a prime below either endpoint.

The splitting algebra for the ordered roots is finite free
over \(B\) by successive monic polynomial division. Its
nonzero residual factor \(P\) is a projective summand,
hence free of positive rank over the local ring \(B\).
Apply the lifting fact again. At every vertex the polynomials
split with simple roots in \(k'((z_i))\), so the residue
field of each lifted prime is exactly this field. Each
ordering has \(f\ne0\) there. The links and the endpoint
generizations survive in \(P[1/f]\), proving the lemma.
The corresponding lifting by flatness and integrality occurs
in \srcloc{Development/AutomorphicArithmetic/AutomorphicVarieties/PatchingModules/PatchingBaseChange/ChainLift}{64}.
\end{proof}

\begin{theorem}[Hecke support along a finite chain]\leavevmode
\label{thm:chaintransfer417}
Let \(\mathfrak p_0,\ldots,\mathfrak p_h\) be curves as above.
For each successive pair suppose there is a prime
\(\mathfrak e_i\) below both, at which either the determinant
is irreducible or the geometric fiber of \(C_0\) is
nonsemisimple. Assume one irreducible prime below
\(\mathfrak p_0\) belongs to \(\operatorname{Supp}_{R_0} M_K\).
Then every prime below \(\mathfrak p_h\) satisfying either of
these two conditions belongs to \(\operatorname{Supp}_{R_0} M_K\).
In particular the initial hypothesis holds when the first
curve is a selected two-constituent curve and the initial
prime is supplied by Theorem~\ref{thm:stableprime371}.
\end{theorem}
\begin{proof}
Use the single \(T,\mathcal P\) of
Proposition~\ref{prop:jointpatch417}. In \(T\) take the
preimages \(\mathfrak t_i\) of the ordered vertices under
\(T\to P[1/f]\); let \(H_i=\widehat{T_{\mathfrak t_i}}\).
Lift the initial supported prime below \(\mathfrak x_0\)
by the lemma. Its preimage in \(T\) contains the diamonds,
and its residue fiber of \(\mathcal P\) is the extension of
the nonzero original residue fiber of \(M_K\), by the
diamond quotient formula. Lift the prime to \(H_0\) using
faithful flatness of completion of the local ring. Its
residue fiber remains nonzero and its determinant remains
geometrically irreducible. The proposition therefore
supplies a supported irreducible prime \(Z_0\) of \(H_0\)
below all the indicated primes there.

Suppose \(Z_i\) has been obtained and \(i<h\). Lift the
preimage in \(T\) of \(\mathfrak b_i\) to a prime \(Q\) of
\(H_i\) by the same faithful flatness. If the original link
was irreducible, \(Q\) remains irreducible. In the other
case its geometric Galois algebra surjects onto the geometric
fiber of \(C_0\) after passage to a common algebraically closed
field. Indeed the diamonds vanish at this prime. The base
matrices satisfy the first auxiliary word by Cayley--Hamilton
and the second by triviality of the new inertia; lifting group
elements through the surjection of limiting groups generates
the target. This is the fiber map proved in
Theorem~\ref{thm:reduciblepropagation412}. A nonsemisimple
geometric fiber remains so under algebraically closed field
extension, and a quotient of a semisimple algebra is
semisimple. Thus \(Q\) meets one of the two conditions, and
\(Z_i\subseteq Q\).

Contract \(Z_i\) all the way to the common ring \(T\), obtaining
a prime \(z_i\). It lies below the preimage of \(\mathfrak b_i\),
hence below \(\mathfrak t_{i+1}\). Its determinant is
geometrically irreducible: the determinant at \(Z_i\) is its
scalar extension. It lies in the support of \(\mathcal P\).
In fact the fiber of a finite module after base change at
a prime is its original residue fiber tensored with the
new residue field. Nonvanishing is equivalent in the two
fields. Lift \(z_i\) to a prime of \(H_{i+1}\). The same
fiber calculation gives a supported irreducible prime
there. Applying the proposition at \(i+1\) gives \(Z_{i+1}\).
This step transports the irreducible generization through
the common patched ring; merely knowing the link itself to
be supported would not provide that generization.

Finally lift the proposed endpoint prime below \(\mathfrak x_h\)
and then to \(H_h\). Its lift satisfies the same alternative
as \(Q\), so contains \(Z_h\) and is supported. It lies over
the diamond quotient and its fiber is a scalar extension of
the original \(M_K\)-fiber. That original fiber is nonzero.
The argument includes \(h=0\). The contraction used at the
induction step is visible in
\srcloc{Development/AutomorphicArithmetic/AutomorphicVarieties/PatchingModules/AnnihilatorBaseChange/ChainAbs}{38}
and in
\srcloc{Development/AutomorphicArithmetic/AutomorphicVarieties/PatchingModules/AnnihilatorBaseChange/PmMove}{7}.
\end{proof}

Using this theorem to bound the minors of
Corollary~\ref{cor:strictminors415} still requires connecting
the components that occur at the original weights to the
supported ones and controlling their integral multiplicities.
The estimates in Section~\ref{sec:missing} also involve the
signed period on the original field.


\subsection{Complete intersections at irreducible ordinary curves}
\label{sec:irreduciblechain418}

Keep the field, prime and ordinary ring of
Section~\ref{sec:twoblockchain417}. Consider a characteristic-\(\ell\)
curve \(\mathfrak p\) of quotient dimension one, with an integral
matrix realization after finite extension of its curve field.
Require the same ordinary flags, matrix words and symmetric
conjugate-dual polarization as there. In this section suppose
the generic representation \(U\) on three variables is absolutely
irreducible, its residual semisimplification is
\(\xi(\bar\Psi\oplus1\oplus\bar\Psi^{-1})\) for the common
normalizing character \(\xi\), and its three ordinary characters
on inertia in \(G'\) are pairwise distinct at each chosen place.
We treat these curves even when they first occur in an
intersection with a component already carrying Hecke support.

\begin{lemma}[Irreducibility on the connected monodromy group]\leavevmode
\label{lem:connectedthree418}
Over an algebraic closure of the curve field let \(H\) be the
reduced Zariski closure of \(U(G_L)\). Its identity component
\(H^\circ\) and its connected derived group \(D\) both act
irreducibly. The group \(D\) is semisimple.
\end{lemma}
\begin{proof}
The restriction to the normal subgroup \(H^\circ\) is
semisimple. Indeed the sum of its simple submodules is nonzero
and is preserved by \(H\), so it is the entire representation.
The isotypic spaces are permuted transitively for the same
reason. Dimension three leaves just two possibilities for a
reducible restriction: a scalar action or three distinct
character lines permuted transitively.

In the scalar case the projective image is finite since an
algebraic group has finitely many components. At a chosen
coefficient place the local eigencharacters on inertia in
\(G'\) have ratios with residue one. Each ratio of two distinct
characters takes a nonidentity value there, of infinite order
by the calculation in Lemma~\ref{lem:planecurve417}. This
excludes a finite projective image.

In the three-line case the permutation action on the lines
is continuous. They are defined over a finite extension, and
a finite set of distinct points of projective space has
discrete induced valuation topology. The residual representation
can be put in upper triangular form over a finite residue-field
extension. Its diagonal reduction gives an abelian quotient
of the compact matrix image whose kernel is pro-\(\ell\):
reduction maps that kernel into a finite unipotent group,
and the principal congruence group is pro-\(\ell\).
Since \(\ell>3\), its image in \(S_3\) is trivial. The
permutation quotient is therefore transitive abelian, hence
cyclic of order three. In the basis of the permuted lines,
conjugation by the diagonal matrix with entries
\(1,\zeta_3,\zeta_3^2\) in cyclic order gives a self-twist by
a character of order three. Its reduction still has order
three. Removing \(\xi\), multiplication by this reduced
character must permute the multiset
\(\{1,\bar\Psi,\bar\Psi^{-1}\}\).
The orbit of \(1\) is the whole set and consists of characters
of order dividing three. But \(\bar\Psi^3\ne1\) on inertia
at \(\ell\), since \(0<3r<\ell-1\). This proves the first
irreducibility assertion.

The unipotent radical of \(H^\circ\) has a fixed vector by
triangularization; normality makes its fixed space
\(H^\circ\)-invariant. Hence it acts trivially and is the
identity subgroup of this matrix group. The connected
reductive group \(H^\circ\) is the product of its connected
central torus and its connected semisimple derived group.
Its central torus acts by scalars by Schur's lemma, so \(D\)
acts irreducibly.~\cite{IX-VII-VI-milne122}
\end{proof}

\begin{proposition}[Separable Frobenius tests on every open coset]\leavevmode
\label{prop:regularthree418}
The restriction of \(U\) to \(G'\) is absolutely irreducible.
Every nonempty relatively open subset of \(G'\) contains
an element with three distinct eigenvalues.
\end{proposition}
\begin{proof}
We first check that every element of \(H\) is a scalar times
an element of \(D\), over the algebraic closure. As a
nontrivial semisimple subgroup of \(\mathrm{SL}_3\), \(D\)
has rank one or two. If its rank is two and the root system
is irreducible, it has at least six roots. Its dimension is
at least \(2+6=8\), giving \(D=\mathrm{SL}_3\). In this case
a cube root of the determinant writes any invertible matrix
as a scalar times an element of \(D\). A reducible root
system of rank two would give two commuting connected
semisimple factors with finite central intersection. The
restriction of the irreducible representation to either
factor is isotypic: its simple socle and its isotypic
summands are preserved by the other factor, which commutes
with it. Since three is prime, a non-scalar first factor
then acts irreducibly, forcing the second to be scalar by
Schur's lemma. A positive-dimensional semisimple factor cannot
be a scalar subgroup. This excludes that case.

In rank one the based root datum of \(D\) has trivial
automorphism group. Indeed an automorphism must fix the sole
simple root, which spans the rational character space, hence
fixes the character lattice. By the isomorphism theorem for
pinned reductive groups every algebraic automorphism of \(D\)
over the algebraically closed field is inner. Equivalently
one can move the image of a chosen pinning back by inner
conjugation, after which the induced automorphism of the
based root datum determines the pinned automorphism uniquely.
Conjugation by any \(g\in H\) preserves \(D\). Remove an
element of \(D\) inducing that conjugation. The remaining
element commutes with \(D\), hence is scalar by
Lemma~\ref{lem:connectedthree418}. This proves the assertion
in rank one as well.~\cite{IX-VII-VI-milne122}

The closure of \(U(G')\) contains \(D\). It is normal in
\(H\), and the quotient is abelian because \(G_L/G'\) is
abelian and the image of \(G_L\) is dense. This gives
irreducibility. Also \(D\) contains a matrix with three
distinct eigenvalues. In fact on the specified local
inertia in \(G'\) each of the three character ratios has
infinite image. Its kernel is a closed subgroup of infinite
index and has empty interior. The complements of the three
kernels are open dense sets and have nonempty intersection.
An element there has three different local eigenvalues.
It is a scalar times a member of \(D\) by the preceding
paragraphs.

For an open subgroup \(U_0\subseteq G'\) the closure of
\(U(U_0)\) contains the identity component of the closure
of \(U(G')\): the latter group is a finite union of cosets of
the former closed subgroup. In particular it contains \(D\).
For each \(g\in G'\) the closure of \(U(gU_0)\) contains
\(U(g)D\), a scalar translate of \(D\). It therefore contains
a matrix with distinct eigenvalues. The discriminant is a
regular function, so some matrix in the coset image already
has nonzero discriminant. This proves the required open-set
assertion.
\end{proof}

\begin{theorem}[The local ordinary algebra at an irreducible curve]\leavevmode
\label{thm:completehecke418}
Let \(\mathfrak p\) be an irreducible curve as above in the
support of \(M_K\), and put \(\mathfrak l=\mathfrak p\cap\Lambda\)
in its principal weight algebra. Then
\[
       (R_0)_{\mathfrak p}\ \simeq\ (B_K)_{\bar{\mathfrak p}},
       \qquad \dim (R_0)_{\mathfrak p}=3[F:\QQ].
\]
This is a local complete intersection ring, faithfully flat
over \(\Lambda_{\mathfrak l}\); its closed fiber over this
weight ring is an Artinian complete intersection. Moreover
\((M_K)_{\mathfrak p}\) is free of positive rank over the
displayed ring. In particular every prime below
\(\mathfrak p\) belongs to the Hecke support.
\end{theorem}
\begin{proof}
Use the augmented and ordered construction of
Lemma~\ref{lem:chainorder417} for this curve. The auxiliary
level argument in Proposition~\ref{prop:jointpatch417} applies
with one constituent. Here the diagonal block is the full
matrix algebra, by irreducibility on \(G'\), and there are
no off-diagonal constituent blocks. The preceding proposition
gives the dense regular set on which to place all the dual
trace tests. The ratios of distinct characters give the local
invariant vanishing at the coefficient places, and
Lemma~\ref{lem:frobtest366} supplies it at \(p_0\).
Thus the first assertion of Theorem~\ref{thm:limitselmer407}
bounds the rank of the ordinary polarized endomorphism
cohomology by \(D_0=3([F:\QQ]+q)\), where \(q\) is the
number of auxiliary places. The patched module has the
same \(3q\) diamond coordinates as before.

At an ordered prime \(\mathfrak x\) above the augmented curve,
write \(U_*=T_{\mathfrak P}\) for the patched local ring,
where \(\mathfrak P\) is the preimage in \(T\), and put
\(Y=(Y_1,\ldots,Y_{3q})\) for the diamond ideal.
The base fiber over \(\kappa(\mathfrak x)\) is nonzero:
it is a scalar extension of the residue fiber of \(M_K\)
at \(\mathfrak p\). Hence the localized finite patched module
has depth at least \(D_0+1\) by
Lemma~\ref{lem:parameterdepth409}. On the other hand a
first-order lift fixing the augmented coefficient parameter
determines an ordinary endomorphism class injectively by
Proposition~\ref{prop:tangentcurve409}. To recall the
injectivity, the trace pairing on the full matrix algebra
recovers the lifted representation up to strict conjugacy.
A coboundary fixes every trace; the ordinary projectors
then fix the local characters, the distinct Frobenius roots
fix the ordered coordinates, and the inertia word fixes
the diamonds. These are all the generators. The rank bound
and the relative tangent calculation of
Lemma~\ref{lem:parameterdepth409} give embedding dimension
at most \(D_0+1\).

Depth of a nonzero finite module is at most the Krull
dimension, which in turn is at most embedding dimension.
Thus \(U_*\) is regular local of dimension \(D_0+1\), and
the Auslander--Buchsbaum formula gives
\(\mathcal P_{\mathfrak P}\simeq U_*^m\) for some \(m>0\),
as in Theorem~\ref{thm:irrcurvesupport409}.
In particular the diamonds, already regular on this module,
form a regular sequence in \(U_*\) itself. The natural
surjection \(U_*/Y\to P_{\mathfrak x}\) is injective.
Indeed the diamond quotient is \((U_*/Y)^m\), and its action
factors through \(P_{\mathfrak x}\) by the comparison
\[
       \mathcal P_{\mathfrak P}/Y\mathcal P_{\mathfrak P}
                     \simeq P_{\mathfrak x}\otimes_{R_0}M_K .
\]
A free module of positive rank is faithful. We conclude that
\(P_{\mathfrak x}\) is a quotient of a regular local ring
by the \(3q\) indicated regular elements, and is a complete
intersection of dimension \(3[F:\QQ]+1\).

The map \((R_0)_{\mathfrak p}\to P_{\mathfrak x}\) is flat
and local. Indeed the unramified extension of constants,
adjoining a formal parameter over a Noetherian ring, the
finite free ordered splitting factor and the localizations
are all flat. The module \((M_K)_{\mathfrak p}\) becomes
free of rank \(m\) there. A basis of its residue space lifts
to \(m\) generators by Nakayama. The kernel of this
presentation vanishes after faithfully flat extension, where
the surjection between free modules of equal rank is an
isomorphism. It vanishes already over \((R_0)_{\mathfrak p}\).
Its annihilator is consequently zero, proving the displayed
Hecke isomorphism. The complete intersection property
descends through this flat local map by the complete
intersection descent theorem for Noetherian local rings.
\href{https://stacks.math.columbia.edu/tag/09Q7}
{The Stacks Project, Lemma 23.8.9.}

Finally \(M_K\) is finite free over the principal weight
algebra \(\Lambda\), a power series ring in \(3[F:\QQ]\)
variables over the coefficient DVR on the fixed residual
factor. Localizing \(M_K\) at \(\mathfrak p\) preserves
flatness over \(\Lambda\): this localization is a filtered
direct limit of copies of the flat module with transition
maps given by multiplication by elements of \(R_0\).
The free-module formula just proved then gives flatness of
\((R_0)_{\mathfrak p}\) over \(\Lambda_{\mathfrak l}\);
a flat local map is faithfully flat. The Hecke image is
finite over \(\Lambda\), since it is a submodule of
\(\operatorname{End}_\Lambda M_K\). The integral extension
\(\Lambda/\mathfrak l\subseteq R_0/\mathfrak p
  =B_K/\bar{\mathfrak p}\) gives
\(\dim\Lambda/\mathfrak l=1\). In the regular complete
local ring \(\Lambda\) this gives
\(\operatorname{ht}\mathfrak l=3[F:\QQ]\).
Incomparability in the finite Hecke algebra gives this as
an upper bound for the dimension of its local ring, and
going down by flatness gives the matching lower bound.
The images of a regular system of parameters for
\(\Lambda_{\mathfrak l}\) form a regular sequence there
by flatness. The quotient is zero-dimensional by finiteness
of the Hecke fibers, hence is an Artinian complete
intersection. Further localization preserves the positive-rank
freeness of \(M_K\), proving the support assertion.
\end{proof}

\begin{corollary}[Chains containing both kinds of ordinary curves]\leavevmode
\label{cor:mixedchain418}
In Theorem~\ref{thm:chaintransfer417} the vertices may also
include the irreducible curves of this section. The same
initial hypothesis and the same conclusion hold at endpoints
of either type. If the first vertex is irreducible, it
suffices that this vertex itself be in the support of \(M_K\).
\end{corollary}
\begin{proof}
Treat each maximal run of two-constituent vertices by
Theorem~\ref{thm:chaintransfer417}. A link below an
irreducible vertex has irreducible determinant. To see this
directly, choose nine group matrices at the vertex whose
trace pairing is nonsingular. Its determinant remains
nonzero at every prime below it. Over an algebraically
closed field a reducible semisimple representation of
dimension three has trace pairing of rank at most five on
its group algebra, from the two possible block sizes \(1,2\),
or at most three for three characters. This excludes a
reducible determinant at the link. Thus upon leaving a
run of two-constituent vertices one obtains a supported
irreducible link, and the next irreducible vertex is
supported by specialization. At this vertex the theorem
makes every generization supported. Its next link gives
a supported irreducible starting prime for the following
run. Induction over the finite list proves the result,
including lists of just one type. The assertion for the
first irreducible vertex uses that vertex itself as its
irreducible starting prime.
\end{proof}

The statements apply to any exhibited chain with these
inertia and matrix conditions. A global argument that the
components relevant to the fixed weights are so reached,
and a quantitative comparison of the resulting integral
orders with the signed theta periods, are still needed in
Section~\ref{sec:missing}.


\subsection{Dimensions of intersections that transfer Hecke support}
\label{sec:weightintersections419}

Use the ordinary ring \(R_0\), the Galois algebra \(C_0\), the
module \(M_K\) and the fields \(L/F\) of
Section~\ref{sec:twoblockchain417}. Write \(I_K\) for the
kernel of \(R_0\to B_K\), and \(\varpi\) for a uniformizer of
the coefficient ring \(O\). We measure dimensions on the
full three-character weight space of the same arithmetic
problem. Set
\[
 d=[F:\QQ],\qquad d_u=[F_u:\QQ_\ell],\qquad
 b=\min\bigl(\{d-1\}\cup\{d_u-1:u\mid\ell\}\bigr),
 \qquad m=3d-b .
\]
Here \(b\ge5\) by Theorem~\ref{thm:commonframe365}. At each
\(u\mid\ell\) fix one of the two places of \(L\) over \(u\).
Let \(H_u\) be the principal units of \(F_u\), and let
\(H=\prod_{u\mid\ell}H_u\). On the chosen finite character
factor the weight algebra and its special fiber are
\(\Lambda=O[[H^3]]\) and \(\bar\Lambda=k[[H^3]]\),
where \(k=O/(\varpi)\). Write \(\lambda_{u,i}\) for the
\(i\)-th universal character on \(H_u\), \(1\le i\le3\).
A geometric determinant at a prime means its associated
semisimple representation over an algebraic closure of the
residue field. The fixed cyclotomic twist may be removed
in statements about conjugate duality. In characteristic
\(\ell\) it acts trivially on principal units.

\begin{proposition}[Two exceptional sets of weights]\leavevmode
\label{prop:badweights419}
There is a closed subset \(W\subset\operatorname{Spec}\bar\Lambda\)
of dimension at most \(m\) with the following property. At a
characteristic-\(\ell\) prime of \(R_0\) of quotient dimension
one whose weight is outside \(W\), the three ordinary characters
are pairwise distinct on inertia in the full cyclotomic kernel
at each chosen place. Its geometric determinant is irreducible
or is a sum of a character and an absolutely irreducible plane.
\end{proposition}
\begin{proof}
Put \(H_u^0=\ker(N_{F_u/\QQ_\ell}|_{H_u})\). This is a free
direct summand of rank \(d_u-1\), by the local unit calculation
in Theorem~\ref{thm:commonframe365}. For each \(u,i\ne j\),
impose \(\lambda_{u,i}(h)=\lambda_{u,j}(h)\) for all
\(h\in H_u^0\). The equations kill the copy of \(H_u^0\)
embedded in \(H^3\) with coordinates \(h,h^{-1}\) in
positions \(i,j\) at \(u\). Projection to position \(i\)
and a projection of \(H_u\) onto \(H_u^0\) split this
embedding. In power series coordinates these equations
therefore cut out codimension \(d_u-1\). Local reciprocity
identifies the cyclotomic character on inertia with the
local norm on units, up to inversion. In particular norm-one
principal units lift to inertia in its kernel. Equality of
the characters on that inertia forces these equations.

For the second exceptional set let
\(\mathcal E=\{e\in\mathcal O_F^\times:e\in H_u
\text{ for all }u\mid\ell\}\). It has finite index in the
unit group. For each choice of two distinct row indices
\(i_u,j_u\) at every \(u\), impose
\begin{equation}
       \prod_{u\mid\ell}
           \lambda_{u,i_u}(e)\lambda_{u,j_u}(e)^{-1}=1
                              \qquad(e\in\mathcal E).
       \label{eq:realunits419}
\end{equation}
The diagonal closure of \(\mathcal E\) in \(H\) has
\(\ZZ_\ell\)-rank \(d-1\). This is Leopoldt's theorem
for the totally real abelian field \(F\). The map
\(H\to H^3\) putting the \(u\)-coordinate with exponents
\(1,-1\) in these two rows admits a left inverse obtained
by reading row \(i_u\) for each \(u\).
In Smith normal form the embedded closed subgroup is
generated by \(\ell^{a_i}\) times distinct basis vectors
of \(H^3\), \(1\le i\le d-1\). In characteristic \(\ell\)
its equations have radical \((X_1,\ldots,X_{d-1})\):
indeed \((1+X_i)^{\ell^{a_i}}-1=X_i^{\ell^{a_i}}\).
Equations for a dense subgroup give the same closed ideal
as those for its closure by completeness and closedness of
ideals in the Noetherian complete local ring.
Thus \eqref{eq:realunits419} has codimension \(d-1\).
\cite{IX-VII-brumer330}

Suppose now that the curve determinant splits into three
characters. They can be taken over a finite extension of
the curve field. The integral closure of the curve ring in that
extension is a complete DVR. The finite \(R_0\)-algebra \(C_0\)
acts on the three simple summands through its specialization.
The values of every group element and its inverse on a
character are integral over the curve ring, so are units of
the valuation ring. The characters are continuous: the map
from the finite algebra to that valuation ring is a linear
map of finite modules over the curve ring, and the original
group map is continuous. After removing the common twist,
order the characters by the three distinct residual
characters \(\bar\Psi,1,\bar\Psi^{-1}\). The conjugate inverse
of the first must be the third, by the polarization and its
reduction. At the chosen local place these two characters
occupy some distinct rows \(i_u,j_u\). Indeed the ordinary
polynomial has the prescribed local characters as its three
roots on every local group element, hence these are the
three local constituents by uniqueness of semisimplification.
One can also recover their row indices by the separating
residual inertia element.

Divide the first global character by its residual character,
lifted into the finite constant field. The quotient takes
values in principal units of the curve DVR. This group is
torsion-free and pro-\(\ell\). At a finite place of residue
characteristic different from \(\ell\), local units are an
extension of a finite group by a pro-group of that residue
characteristic. Their image under the quotient character is
therefore trivial. At a chosen place above \(u\) its value
on \(e\in\mathcal E\) is \(\lambda_{u,i_u}(e)\); at the
conjugate place it is \(\lambda_{u,j_u}(e)^{-1}\). This uses
that \(e\) is real and that residual characters on principal
units are trivial. Complex places contribute trivially.
Global reciprocity on the principal idele of \(e\) gives
\eqref{eq:realunits419}. Taking the finite union of the two
sets of loci just constructed proves both assertions.
\end{proof}

\begin{lemma}[Specialization of the polarized matrices]\leavevmode
\label{lem:positivefibers419}
Let \(J\subseteq\mathfrak p\) be primes of \(R_0\) with
geometrically irreducible determinant at \(J\). Suppose
\(\mathfrak p\) is a characteristic-\(\ell\) curve of
quotient dimension one with weight outside \(W\). It has
an integral ordinary, symmetrically polarized realization of
one of the two kinds in Sections~\ref{sec:twoblockchain417}
and~\ref{sec:irreduciblechain418}.
\end{lemma}
\begin{proof}
Work in the normalization with the common twist removed,
so that \(g^\star=cg^{-1}c\) induces the adjoint and fixes
the trace. Over a geometric field the semisimple representation
realizing the determinant factors through \(C_0\): its
algebra is a product of full matrix algebras, the trace
pairing there is nondegenerate, and the trace of each defining
word against every group element is zero. Here the residual
characters are distinct, so there are no repeated constituents
at the curve. More generally even multiplicities at most three
are units in the present characteristics. The group matrices
span this semisimple algebra. Trace duality shows that
\(\star\) descends to it: if \(x\) is in the kernel, then
\(t(hx^\star)=t(xh^\star)=0\) for every group element \(h\).
At \(J\) the resulting anti-involution of \(M_3\) is adjunction
for a symmetric form. Indeed an anti-involution of a split
matrix algebra is transposition followed by inner conjugation.
The form matrix has transpose equal to itself up to sign;
a nondegenerate alternating form has even dimension.

We check the sign on a plane after specializing.
Take the separating element \(\gamma\) of
Theorem~\ref{thm:commonframe365}, set
\(Z=\gamma-\gamma^\star\), let \(a=t(Z^2)/2\), and write
\(E=Z^2/a\) in the group algebra over \(R_0/J\).
The residue of \(a\) is
\((\bar\Psi(\gamma)-\bar\Psi(\gamma)^{-1})^2\ne0\), so
\(a\) is a unit. At the generic point of \(R_0/J\), skew
adjunction gives the eigenvalues \(0,\sqrt a,-\sqrt a\)
of \(Z\). Thus \(E\) is the orthogonal projector onto a
nondegenerate symmetric plane. Skew-adjoint endomorphisms
on that plane form a one-dimensional space, identified
with alternating bilinear forms. For any three group
elements \(g,h,k_0\) the identity
\[
 t\!\left([\,E(g-g^\star)E,\,E(h-h^\star)E\,]k_0\right)=0
\]
holds at that generic point. Its left side is an element
of the domain \(R_0/J\), so it vanishes at \(\mathfrak p\).

If the curve determinant is irreducible the odd-dimensional
argument already gives its symmetric form. Otherwise the
preceding proposition gives the algebra \(M_2\times M_1\).
The involution fixes both factors since they have different
dimensions. It is the identity on the line. On the plane
\(Z\) is traceless with trace of its square \(2a\), so
\(Z^2=a\) there. Consequently \(E\) projects to the plane.
The displayed identity and the nondegenerate trace pairing
show that any two skew-adjoint endomorphisms of this plane
commute. Indeed every such endomorphism is half its difference
from its adjoint, and the group matrices span the algebra.
For an alternating form in dimension two, both off-diagonal
matrix units in a symplectic basis are skew adjoint and their
commutator is nonzero. This excludes alternating sign.
The line is self-dual for \(\star\), so its residual character
is \(1\): the two outer residual characters are interchanged
by conjugate inversion and are distinct.

A parameter \(T\) of the complete one-dimensional domain
\(R_0/\mathfrak p\) makes it finite over \(k[[T]]\). To see
finiteness, lift a basis of its quotient by \(T\) and expand
successive remainders using completeness. Its fraction field
is a finite extension of \(k((T))\). The simple matrices,
block decomposition and forms are defined over a further
finite extension since their finitely many structure
constants belong to the algebraic closure. The span of the
translates of any basis over the valuation ring is a stable
lattice: it is contained in the module obtained by applying
a finite list of generators of \(C_0\) to that basis and
taking their integral span. This also realizes the action
continuously on the lattice. The matrix words of \(C_0\)
hold there.

For the local flags, take \(\gamma_u\) on inertia with
three distinct residual eigenvalues. Take its eigenbasis
in the order of the three local characters. In the ordered
word \(\prod_{i=1}^3(\rho(h_i)-\psi_i(h_i))=0\), set
\(h_i=\gamma_u\) except for \(h_k=g\). At row \(i\ge k\)
and column \(j\le k\) the factors to the left and right of
the \(k\)-th factor have invertible entries on the indicated
row and column respectively. Thus the \(ij\)-entry of
\(\rho(g)-\psi_k(g)\mathrm{id}\) vanishes for these indices.
The resulting flag is upper triangular with the prescribed
diagonal. Finally apply the lattice construction in
Proposition~\ref{prop:curvelattice366}. It uses only this
flag, a stable lattice, the symmetric form and the distinct
residual eigenvalues of \(Z\): rescaling its eigenlines
first splits the reduction and then averaging the exponents
with those of the dual lattice makes the lattice self-dual.
Intersecting with each local flag gives saturated steps.
Restoring the common twist gives the required realization.
\end{proof}

\begin{theorem}[Hecke support across a large intersection]\leavevmode
\label{thm:largeintersection419}
Let \(J\) be a prime of \(R_0\) containing \(I_K\) with
geometrically irreducible determinant. Suppose \(P\) is
a prime containing \(J+(\varpi)\) and \(\dim R_0/P>m\).
Then every prime \(Q\subseteq P\) at which the determinant
is geometrically irreducible or the geometric algebra
\(C_0\otimes_{R_0}\overline{\kappa(Q)}\) is nonsemisimple
contains \(I_K\). If the determinant at \(P\) is geometrically
irreducible, the map from \((R_0)_P\) to the local Hecke
algebra is an isomorphism, \((R_0)_P\) is a complete
intersection, and \((M_K)_P\) is free of positive rank.
\end{theorem}
\begin{proof}
The domain \(D=R_0/P\) is a quotient of \(B_K\) and hence
finite over \(\bar\Lambda\). The inverse image of \(W\) in
\(\operatorname{Spec}D\) has dimension at most \(m\).
Indeed on each component of this inverse image the domain
is finite integral over the quotient by its contraction,
and their dimensions agree. Hence the complement is a
nonempty open set. If \(P\) is geometrically irreducible,
choose nine group elements for which the trace Gram
determinant is nonzero modulo \(P\), and also require
that this determinant remain nonzero. The two open
conditions still meet because \(D\) is a domain.

Every nonempty open set in this spectrum contains a prime
of quotient dimension one. Choose a nonzero \(f\in D\) such
that the principal open where \(f\ne0\) is contained in
the prescribed open. We may multiply by a nonzero nonunit
to make \(f\) a nonunit. Write \(n=\dim D\). If \(n=1\)
use the zero prime. Otherwise \(D\) is a complete local
domain, hence catenary and equidimensional, and each minimal
prime over \(f\) has height one. Thus \(f\) extends to a
system of parameters \(f,u_2,\ldots,u_n\). A minimal prime
\(q\) over \((u_2,\ldots,u_n)\) has height at most \(n-1\)
and cannot contain \(f\), since the full parameter ideal
is primary to the maximal ideal. In the local domain \(D/q\)
the nonzero image of \(f\) generates an ideal with maximal
radical. The principal ideal theorem gives \(\dim D/q=1\).

Lift this prime to \(\mathfrak p\subset R_0\). It contains
\(P\), has the required weight, and lies above the supported
irreducible prime \(J\). The preceding lemma supplies its
matrix realization. If the curve is irreducible,
Theorem~\ref{thm:completehecke418} puts every prime below
it in the support. If the curve has two constituents, apply
Theorem~\ref{thm:chaintransfer417} to this single curve
with starting prime \(J\). Its conclusion applies to each
of the stated \(Q\). If \(P\) is irreducible, the extra Gram
condition forces the curve to be irreducible as well: nine
independent matrices are impossible for a sum with a proper
constituent. Localize the complete intersection and the
positive-rank free module of Theorem~\ref{thm:completehecke418}
further at \(P\). Localization preserves the complete
intersection property and freeness; the localized Hecke
kernel remains zero. This proves the last assertions.
\end{proof}

\begin{corollary}[The component through the cuspidal adjoint]\leavevmode
\label{cor:componentgraph419}
Let \(J_0\) be the kernel of the full component map through
the cuspidal adjoint in Proposition~\ref{prop:upfamily372}.
Every minimal prime over \(J_0+(\varpi)\) satisfies the
dimension hypothesis of Theorem~\ref{thm:largeintersection419}.
More generally let \(J_0,J_1,\ldots,J_s\) be primes of
\(R_0\) with geometrically irreducible determinants and
\[
       \dim R_0/(J_i+J_{i+1}+(\varpi))>3d-b
                       \qquad (0\le i<s).
\]
Then all \(J_i\) contain \(I_K\). The displayed bound
holds in particular if \(\dim R_0/(J_i+J_{i+1})\ge 3d-b+2\).
\end{corollary}
\begin{proof}
The full adjoint component is a finite domain containing
\(\Lambda\) and is a quotient of \(B_K\). It has dimension
\(3d+1\). Each minimal prime over \(\varpi\) has height
one by the principal ideal theorem; catenarity of the
complete local domain gives quotient dimension \(3d\).
For a successive pair, choose a minimal prime \(P\) over
the displayed sum of ideals attaining its quotient dimension.
With \(J=J_i\) and \(Q=J_{i+1}\) the theorem transfers the
Hecke support, starting with the given quotient at \(J_0\).
The final sufficient bound follows since quotienting a
Noetherian local ring by one element of its maximal ideal
lowers its dimension by at most one.
\end{proof}

The corresponding choice of a curve above a large intersection
in marked representation coordinates occurs in
\srcloc{Development/AutomorphicArithmetic/AutomorphicVarieties/CurveDVRSpecialization/ExistsNicePrimeAboveOfLargeDim}{7}.
The corollary gives the transfer on the original determinant
ring whenever the indicated intersections exist. Establishing
coverage of the components at the arithmetic weights, and
bounding the integral multiplicities there, are separate
steps still needed for Section~\ref{sec:missing}.


\subsection{Intersections in the marked ordinary deformation ring}
\label{sec:markedintersections420}

Fix the same \(L/F,\ell,\bar\Psi\) and \(d=[F:\QQ]\) as in
Section~\ref{sec:weightintersections419}. We use
\(\Gamma=\Gamma_a\), \(H=\Gamma_a^0\), and a complex conjugation
\(c\) as in Section~\ref{sec:commonframe365}. Remove the common
cyclotomic twist and let \(J\) be the antidiagonal matrix with
entries one on its antidiagonal. Let \(k'\) be a finite extension
of the residual coefficient field and \(O'/O\) its unramified
lift. Fix a continuous matrix representation
\(\rho_*:H\to {\rm GL}_3(k')\) with semisimplification
\(\bar\Psi\oplus1\oplus\bar\Psi^{-1}\), satisfying the ordinary
conditions of Section~\ref{sec:markedcurve366}. Assume it has
the specified pairing \(J\) and marking
\[
 \rho_*(\gamma)-\rho_*(c\gamma^{-1}c)
       =\delta\,\operatorname{diag}(1,0,-1),\qquad
 \delta=\bar\Psi(\gamma)-\bar\Psi(\gamma)^{-1}\ne0 .
\]
This includes the split residual representation. Write \(A_*\)
for the complete local \(O'\)-algebra of its marked ordinary
matrix lifts on \(H\), with all three ordinary weights variable.
The marking at a lift has the same shape with a scalar lifting
\(\delta\). Matrices in the chosen basis are part of the data.
The presentation by generator matrices and local words in
Theorem~\ref{thm:markedcurve366} works for this residual point:
the kernel of its finite image has finitely generated maximal
pro-\(\ell\) quotient by Theorem~\ref{thm:commonframe365}.
On the finite old inertia quotients the residual type lifts
uniquely up to conjugacy. Indeed their order is prime to
\(\ell\), and averaging an intertwiner lifting the identity
makes it equivariant and still invertible at every small
extension. The trace and the ordered characters, with the
common twist restored, give a map \(R_0\to A_*\).

\begin{lemma}[Tangent space and obstruction for the marked matrices]\leavevmode
\label{lem:markedcomplex420}
Let \(W=M_3(k')\), with \(H\) acting by conjugation through
\(\rho_*\) and \(c\) acting by \(X\mapsto-JX^{\mathsf T}J\).
At each \(u\mid\ell\) let \(b_u\subset W\) preserve the residual
ordinary flag at the chosen place. Let \(H^i_C\) be the cohomology
of the mapping fiber of the continuous restriction of cochains
on \(\Gamma\) with values in \(W\) to the direct sum of the
local cochains with values in \(W/b_u\); put \(h_i=\dim_{k'}H^i_C\).
Then the relative tangent dimension of \(A_*/O'\) is
\(h_1+1-h_0\). For a small extension with kernel \(V\)
killed by the maximal ideal, there is a functorial obstruction
in \(H_C^2\otimes_{k'}V\); it vanishes exactly when the marked
ordinary lift exists.
\end{lemma}
\begin{proof}
First allow change of basis by matrices reducing to the identity.
Put \(\iota(g)=J(g^{\mathsf T})^{-1}J\). The normalized polarized
problem on \(H\) extends to \(\Gamma=H\rtimes\langle c\rangle\)
with target \({\rm GL}_3\rtimes\langle\iota\rangle\) and the
prescribed projection to the group of order two. In the marked
basis \(c\) has image \(\iota\). The derivative of this action
is the indicated action on \(W\). For each local group the flag
with the given ordering is unique over every coefficient ring.
Indeed the three distinct residual eigenvalues of a separating
inertia element lift to distinct eigenvalues with unit differences.
Its three projectors give eigenlines. The ordered words force the
upper triangular flag in their prescribed order by the row and
column calculation of Lemma~\ref{lem:positivefibers419}. Conversely
a stable ordered flag gives the words and its diagonal characters.

For a lift to the dual numbers write its variation as \(z(g)\)
in the Lie algebra, and write the variation of each flag as
\(-X_u\in W/b_u\). The representation equation is \(dz=0\);
flag stability is \(z|_{G_{F_u}}=dX_u\bmod b_u\). Changing
basis changes this pair by the coboundary of a single \(X\in W\).
These are exactly the degree-one cocycles and coboundaries of
the mapping fiber. Its degree zero is the space of infinitesimal
automorphisms commuting with the extended action and preserving
the flags.

We record the lifting calculation. Given a small extension
\(D'\to D\) as in the statement, lift the global matrices as
a continuous set map with the prescribed projection, and lift
each local flag by lifting an adapted basis. The multiplicative
defect of the lifted global matrices is a 2-cocycle in \(W\otimes V\):
write it as \(\widetilde\rho(g)\widetilde\rho(h)
\widetilde\rho(gh)^{-1}-1\) in the identity component and use
associativity and \(V^2=0\). In bases for the lifted flags the
local matrices reduce to upper triangular matrices over \(D\).
Lift those upper triangular matrices in the flag stabilizer
by a set map. The discrepancy from the global lift gives a
local 1-cochain modulo \(b_u\). Its differential is the
restriction of the global defect modulo \(b_u\), because the
multiplicative defect of the stabilizer lifts is in \(b_u\otimes V\).
Changing matrix lifts changes this pair by a fiber coboundary;
changing flag lifts changes the local cochains by local
coboundaries. The resulting class in degree two vanishes
precisely if one can subtract a global 1-cochain to kill the
multiplicative defect and move the local flags so that all local
discrepancies vanish modulo their stabilizers. This yields the
required representation with stable flags. All operations
commute with maps of the kernel \(V\). For Artinian coefficient
rings the set lifts can be chosen on finite images. The same
argument applies to complete coefficient rings with finite kernel:
a sufficiently deep open ideal of \(D'\) has zero intersection
with \(V\), and maps homeomorphically onto an open additive
subgroup of \(D\). Choosing representatives on its finitely many
cosets gives a continuous section, usable for each matrix entry.

It remains to enforce the marking, both in this lifting test
and in the tangent count. Suppose the lift over \(D\) is already
marked. If the image of \(c\) upstairs is \(a\iota\), its square
being one says that \(aJ\) is symmetric. It reduces to \(J\)
over \(D\). Conjugate by \(1-\tfrac12(aJ-J)J\) to put it back
to \(J\): under this change \(aJ\) transforms by congruence,
and products of two kernel entries vanish. Once \(c\) is fixed,
the operator \(Z=\rho(\gamma)-\rho(c\gamma^{-1}c)\) is skew
adjoint. Its three eigenvalues lift the unit-separated roots
on the base; they are \(z,0,-z\). Choose eigenvectors lifting
the original coordinate vectors. The two outer vectors pair to
a unit and the middle vector pairs with itself to a unit, both
equal to one over \(D\); every other pairing vanishes by skew
adjunction. Rescale one outer vector and the middle vector to
make these pairings one, using the square root congruent to one
for the middle vector. This is an orthogonal change of basis
that reduces to the identity over \(D\). It gives the marked
lift, preserving the flags under transport.

On dual numbers every change between two marked representatives
fixes \(c\) and the ordered eigenspaces of \(Z\). Its infinitesimal
generator is therefore diagonal and skew adjoint for \(J\), hence
a multiple of \(\operatorname{diag}(1,0,-1)\). Conversely each
such generator acts by change of marked basis. Its stabilizer is
exactly \(H_C^0\), by the description of infinitesimal automorphisms
above. Passing from unmarked tangent classes to marked matrices
thus adds \(1-h_0\) dimensions. This proves the lemma.
\end{proof}

\begin{theorem}[Number of relations in the marked ordinary ring]\leavevmode
\label{thm:markedpresentation420}
There is a presentation
\[
 A_*=O'[[X_1,\ldots,X_g]]/(f_1,\ldots,f_s),
                \qquad g-s\ge 3d+1.
\]
It holds at every finite residual matrix point specified above,
including those with nonsplit residual matrices.
\end{theorem}
\begin{proof}
The module \(W^\vee(1)\) has no global invariants. In fact
Theorem~\ref{thm:commonframe365} gives an element of \(H\) with
\(\bar\Psi=1,\ \bar\kappa_\ell=-1\). All eigenvalues of its
conjugation on \(W\) are one, even if its action is unipotent.
All eigenvalues on \(W^\vee(1)\) are therefore minus one.
In Lemma~\ref{lem:frameeuler365} the dimensions of the local
Borel spaces are six. At each real place the conjugation has
image \(a\iota\) with \(aJ\) symmetric and nonsingular. Its
invariants in \(W\) are the skew-adjoint matrices for that
form, of dimension three. The Euler formula gives
\[
                    h_1-h_0-h_2=6d-3d=3d .
\]

Use \(g=h_1+1-h_0\) relative tangent generators to obtain a
surjection \(P=O'[[X_1,\ldots,X_g]]\to A_*\) by completeness and
Nakayama's lemma. For its kernel \(D\), minimality of the relative
tangent generators gives
\(D\subseteq\varpi P+\mathfrak m_P^2\). We claim that the minimal
number of generators of \(D\) is at most \(h_2\).
Apply the obstruction calculation to
\(P/\mathfrak m_P D\to A_*\), whose kernel is
\(V=D/\mathfrak m_P D\), and to the universal marked matrices
over \(A_*\). Contracting its obstruction with
\(\lambda\in V^\vee\) gives a linear map \(V^\vee\to H_C^2\).
Suppose a nonzero \(\lambda\) maps to zero. Quotient the extension
by its kernel in \(V\). Functoriality and the lifting test give
a lift of the universal marked matrices to the resulting ring
\(D'\), hence a section \(A_*\to D'\) of \(D'\to A_*\).
The induced map \(P\to A_*\to D'\) and the canonical \(P\to D'\)
agree modulo the one-dimensional small kernel and on \(O'\).
Their difference is a derivation into that kernel for the
residual action. It kills \(\varpi P+\mathfrak m_P^2\), and so
kills \(D\). But the first map sends \(D\) to zero and the second
maps it onto the nonzero kernel, a contradiction. Hence the linear
map is injective. Taking \(s=\dim V\) and lifting generators, we
obtain \(g-s\ge h_1+1-h_0-h_2=3d+1\). The corresponding presentation
in the source is \srcloc{Development/AutomorphicArithmetic/AutomorphicVarieties/FrobeniusSmoothVectors/HasObs}{26}.
\end{proof}

\begin{proposition}[Dimensions of marked component intersections]\leavevmode
\label{prop:dimensionchain420}
Each irreducible component of \(\operatorname{Spec} A_*\) has
dimension at least \(3d+2\). Its components form a connected
graph when two different components are joined if their
intersection has dimension at least \(3d+1\). In the special
fiber \(A_*/\varpi\), the corresponding bounds are \(3d+1\)
for components and \(3d\) for their intersections.
\end{proposition}
\begin{proof}
Let \(P'\) be a complete regular local ring of dimension \(N\)
and \(I\) an ideal generated by \(s\) elements in its maximal
ideal. Minimal primes of \(I\) have height at most \(s\), so
their quotient dimensions are at least \(N-s\). For the
intersection bound, suppose the minimal primes admit a partition
into two nonempty lists with every cross intersection of
dimension less than \(N-s-1\). Let \(U,V\) be the intersections
of the primes on the respective lists. At a minimal prime
\(\mathfrak q\) of \(U+V\) one has
\(h=\operatorname{ht}\mathfrak q\ge s+2\), since \(\mathfrak q\)
contains a prime on each list and the regular local ring is
catenary. Complete \((P')_{\mathfrak q}\), obtaining a regular
complete local domain of dimension \(h\). Each of the quotients
by \(U\) and \(V\) still has positive dimension there: below
\(\mathfrak q\) it has a prime of height at most \(s\), and
completion preserves quotient dimension. Hartshorne--Lichtenbaum
vanishing makes both top local cohomology groups with these
supports zero. The group in degree \(h\) supported on \(U+V\)
is nonzero by nonvanishing for the maximal ideal. Mayer--Vietoris
now requires nonzero local cohomology of \(U\cap V\) in degree
\(h-1\). This is impossible: \(U\cap V=\sqrt I\) before
localization, flatness preserves intersections, and the extended
ideals still have the same radical. The \v Cech complex on the
\(s\) generators of \(I\) vanishes above degree \(s\).
This proves the graph statement for \(P'/I\).
\href{https://stacks.math.columbia.edu/tag/0EB0}{The Stacks Project,
Section 51.16.}

Apply this calculation to the power series presentation over
\(O'\), where \(N=g+1\), and to its reduction over \(k'\),
where \(N=g\). The number of generators in the reduction is
at most \(s\). The theorem gives the claimed bounds in both rings.
\end{proof}

\begin{theorem}[Loss of dimension on passage to traces]\leavevmode
\label{thm:tracefiber420}
Let \(P_*\) be a prime of \(A_*\) and let \(p=P_*\cap R_0\).
If the geometric determinant at \(p\) is irreducible or is
a sum of a character and an irreducible plane, then
\[
                         \dim A_*/P_*\le \dim R_0/p+1.
\]
\end{theorem}
\begin{proof}
We first realize the marked problem inside an affine scheme of
finite type over \(R_0\otimes_O O'\). Take finitely many module
generators of the finite algebra \(C_0\). Give each generator
nine matrix coordinates and impose its module relations,
multiplication table and identity. For the normalized group
matrices impose \(J\)-adjunction, the traces of the universal
determinant, and the prescribed diagonal form of \(Z\).
Its positive labeled eigenvalue is specified on the base:
it is the root of \(z^2=t(Z^2)/2\) with residue \(\delta\),
which exists uniquely by Hensel's lemma. Here the trace and \(Z\)
are computed after twisting by the fixed character. All these
are equations in the finite list of coordinates; arbitrary
families of equations define a quotient of the same polynomial
algebra. Denote the resulting algebra by \(S\). The point of
\(S\) for \(\rho_*\) is closed with residue field \(k'\).
The completion of its local ring is \(A_*\). Indeed any lift
of the matrices gives its determinant and ordered characters,
hence its base map and a representation of \(C_0\). Conversely
such a matrix realization satisfies the defining words of
\(C_0\), hence the ordered flag by the separating local
projectors. Over a local Artinian coefficient ring the group
action is continuous by finiteness and continuity of \(C_0\).
These two constructions are inverse on all Artinian test rings.

Consider the fiber over an algebraic closure of \(\kappa(p)\).
Each marked realization there is semisimple under the stated
hypothesis. This is immediate for an irreducible determinant.
In the second case both distinct constituents are preserved by
conjugate duality, since their dimensions differ. If a simple
submodule \(U\) had degenerate restriction of the pairing,
then \(U\cap U^\perp\), being invariant, would be all of \(U\).
The quotient by \(U^\perp\) would be its conjugate dual,
isomorphic to \(U\), giving a second copy among the composition
factors. This contradicts multiplicity one. Hence \(U\) has
an invariant orthogonal complement, proving semisimplicity.

Two matrices of representations in this geometric fiber are
intertwined because their semisimple determinants agree. Their
pairings are also intertwined after adjusting the isomorphism:
on each self-dual simple summand the two invariant forms differ
by a nonzero scalar by Schur's lemma, and one may rescale by a
square root. Distinct summands are orthogonal because they are
inequivalent and separately self-dual. The resulting intertwiner
is an isometry of \(J\) commuting with the fixed diagonal matrix
\(Z\), whose three eigenvalues are distinct. Such isometries
are \(\operatorname{diag}(s,e,s^{-1})\) with \(s\ne0\) and
\(e=\pm1\). Conjugation by this one-dimensional group therefore
accounts for all geometric matrix points for each choice of
the coefficient embedding. There are finitely many such
embeddings. The finite-type geometric fiber has dimension at
most one: its closed points are contained in the union of these
orbit images, whose closures have dimension at most one.

Contract \(P_*\) to \(S\), obtaining \(P_S\). Write
\(D=R_0/p\). The finite-type domain \(S/P_S\) over \(D\)
has generic transcendence degree at most one by the fiber
calculation. To bound its dimension, present it as a quotient
of \(D[X_1,\ldots,X_n]\) by a prime. After localization over
\(\operatorname{Frac}D\) this prime has height at least \(n-1\),
so it already has height at least \(n-1\) over \(D\). The polynomial
ring has dimension \(\dim D+n\), giving quotient dimension at
most \(\dim D+1\). Localization at the residual matrix point
and completion of that local quotient preserve the upper bound.
The ring \(A_*/P_*\) is a quotient of this completion. This
also treats primes in formal fibers, for which a direct count
of parameters over the residue field would be insufficient.
\end{proof}

\begin{corollary}[Transferring support along marked components]\leavevmode
\label{cor:markedtransfer420}
Let \(P_0,\ldots,P_t\) be components of \(\operatorname{Spec}A_*\),
written as minimal primes, whose contractions \(J_i\) to \(R_0\)
are geometrically irreducible. Suppose \(I_K\subseteq J_0\).
For each successive pair suppose that some prime \(Q\) above
\(P_i+P_{i+1}\) has \(\dim A_*/Q\ge3d+1\) and its geometric
determinant has at most two constituents. Then \(I_K\subseteq J_t\).
\end{corollary}
\begin{proof}
For \(q=Q\cap R_0\) the preceding theorem gives
\(\dim R_0/q\ge3d\). The special fiber of \(R_0/(J_i+J_{i+1})\)
therefore has dimension at least \(3d-1\), by the dimension
inequality for a single element. Since \(b\ge5\), this exceeds
\(3d-b\). Apply the induction step of
Corollary~\ref{cor:componentgraph419} with initial supported
prime \(J_0\). The connectedness assertion above supplies the
dimensions of intersections in the marked ring. To use an entire
connected component graph one must also treat intersections
whose generic determinants split into three characters and
components with reducible generic determinant.
\end{proof}


\subsection{The triangular locus and support at adjacent lower entries}
\label{sec:lowerchartcoverage421}

Use \(A_*,R_0,L/F\) and the normalized matrices of
Section~\ref{sec:markedintersections420}. In this section the
finite point \(\rho_*\) is lower triangular, with diagonal
\((\bar\Psi,1,\bar\Psi^{-1})\), and some strictly lower entry
of some group element is nonzero. Denote by \(A_{\rm low}\)
the quotient on which all the matrices are lower triangular
in the marked basis. Let \(h\) be the common upper bound for
the \(\ell\)-rank of the class groups of
\(T_a=L_aE_*(\mu_\ell)\) in
Proposition~\ref{prop:classrank365}. Here \(d=[F:\QQ]>3h+7\)
by the choice of depth in Theorem~\ref{thm:commonframe365}.

\begin{lemma}[The basis at a reducible determinant]\leavevmode
\label{lem:splitorder421}
At any prime of \(A_*\) whose geometric determinant is a sum
of three characters all matrices are lower triangular. At a
geometric determinant with one line and one irreducible plane
all matrices preserve the middle coordinate line and the plane
of the two outer coordinates.
\end{lemma}
\begin{proof}
In the first case there is an invariant full flag over the
algebraic closure. Every step is stable under
\(Z=\rho(\gamma)-\rho(c\gamma^{-1}c)\). Since \(Z\) is diagonal
with three distinct eigenvalues, each step is a coordinate
subspace. Thus there is an ordering of the coordinates in
which all group matrices are triangular. Draw an arrow
\(j\to i\) for \(i\ne j\) if some group element has nonzero
\(ij\)-entry. The triangular flag precludes directed cycles.
Adjunction by \(J\) pairs the entries \(21,32\) and the entries
\(12,23\). Whenever one of a pair occurs the other occurs,
possibly on another group element. A nonzero lower entry at
the residual point is a unit in \(A_*\) and persists. If a
lower adjacent entry occurs, the path \(1\to2\to3\) prohibits
every upper entry: an upper adjacent entry would give a
two-cycle, and an upper \(13\) would close the path. If the
lower \(31\) occurs, an upper \(13\) would reverse its arrow;
the pair of upper adjacent entries would give the return path
\(3\to2\to1\). Thus no upper entry occurs in either case.
In the second case Theorem~\ref{thm:tracefiber420} shows that
the representation is an orthogonal direct sum. On the self-dual
character line \(Z\) is zero by adjunction, so the line is the
zero eigenspace, namely the middle coordinate. Its orthogonal
complement is the outer plane.
\end{proof}

\begin{proposition}[Dimension of the lower triangular locus]\leavevmode
\label{prop:lowerdimension421}
On the given field,
\[
              \dim A_{\rm low}\le 2d+2+3h,\qquad
              \dim A_{\rm low}/\varpi\le 2d+1+3h .
\]
\end{proposition}
\begin{proof}
For lower triangular lifts the outer diagonal characters are
conjugate inverses and the middle diagonal character is its
own conjugate inverse. Remove the finite residual characters
by their Teichm\"uller lifts. The remaining variations take
values in principal units and are unramified away from
\(\ell\): at the old places the inertia quotient in \(H\) has
order prime to \(\ell\). Let \(X_L\) be the maximal abelian
pro-\(\ell\) Galois group of \(L\) unramified outside \(\ell\).
Use multiplicative notation for characters and additive notation
for \(X_L\). The first diagonal variation factors through \(X_L\),
and the middle one through its minus direct summand \(X_L^-\)
for complex conjugation, since \(2\) is invertible.

The ray class exact sequence, followed by the inverse limit
over powers of the primes above \(\ell\), identifies the
image of the product of local principal units in \(X_L\)
with its quotient by the closure of the projected global
units. Its cokernel is the \(\ell\)-part of the ideal class
group. The local unit product has ranks \(d,d\) in its plus
and minus summands: its factors occur in conjugate pairs of
equal local degree. The diagonal image of the real units has
rank \(d-1\) by Leopoldt for the abelian field \(F\). The
units of \(L\) have the same rank as those of \(F\), by the
unit theorem, so adjoining them changes this closure by at
most finite index. Consequently \(X_L\) has rank \(d+1\)
and \(X_L^-\) has rank \(d\). Both are finitely generated
\(\ZZ_\ell\)-modules. The complete diagonal character algebra
\[
                      D=O'[[X_L\times X_L^-]]
\]
is finite over a power series ring on \(2d+1\) variables
over \(O'\); its extra generators come from finite torsion
groups. It has dimension \(2d+2\), and its special fiber
has dimension \(2d+1\). The universal triangular diagonals
give a map \(D\to A_{\rm low}\).
\cite{IX-VII-brumer330}

We bound the relative tangent dimension at this map. Let
\(\mathfrak u\) be the three-dimensional strictly lower triangular
Lie algebra over \(k'\) with action of \(H\) by conjugation
through \(\rho_*\). A triangular infinitesimal lift with
fixed diagonal differs from \(\rho_*\) by a cocycle in
\(Z^1(H,\mathfrak u)\). This cocycle is a local coboundary at
every place over \(\ell\). To see this at a chosen place, use
the separating inertia matrix \(\rho(\gamma_u)\).
A lower triangular matrix with distinct diagonal entries
having unit differences can be diagonalized by a unique
lower unitriangular matrix of eigenvectors: solve for each
column starting with its first nonzero coordinate equal to
one and then solve successively below it. The ordinary flag
is the increasing flag of these columns, since its ordered
residual characters agree with the three labeled diagonal
entries. The lower triangular action also preserves the
opposite coordinate flag, which is unchanged by this lower
unitriangular change of basis. The full local action therefore
preserves both opposite flags and is diagonal in this basis.
For an infinitesimal variation with fixed diagonal the change
between the two diagonalizing bases is lower unitriangular and
reduces to the identity. This gives the local coboundary in
\(\mathfrak u\). Conjugate places satisfy the same assertion
by polarization, since inverse transposition followed by
reversal preserves the lower unitriangular group.

The map from this relative tangent space to
\(H^1_{\rm str}(H,\mathfrak u)\) is injective. If the cocycle
is the global coboundary of \(X\in\mathfrak u\), the two
representations are strictly conjugate by \(1+\epsilon X\)
up to changing the sign of \(X\). Their \(Z\)'s are diagonal.
Hence \([X,Z_*]\) is diagonal, forcing \(X=0\) by distinctness
of the three eigenvalues and strict lower triangularity.
The local character variations are then zero as well, since
the stable ordered flags are determined by their projectors.

The center \(k'E_{31}\) of \(\mathfrak u\) has character
\(\bar\Psi^{-2}\), and its quotient has two copies of
\(\bar\Psi^{-1}\). For either character the strict first
cohomology on \(H\) has dimension at most \(h\). Indeed
restriction to \(T_a\) is injective, since \([T_a:L]\) is
prime to \(\ell\). There the coefficient character is trivial
and the strict class is an everywhere unramified additive
homomorphism. It factors through the ideal class group,
whose \(\ell\)-rank is at most \(h\). This proves the assertion
for both powers. The quotient of \(\mathfrak u\) by its center
has zero local invariants above \(\ell\) by the residual inertia
characters. Its local exact sequence therefore injects the
first cohomology of the center into that of \(\mathfrak u\).
A strict class on \(\mathfrak u\) whose image in the quotient
vanishes globally lifts by exactness to a class on the center,
and local injectivity makes that lift strict. The total strict
dimension is thus at most \(3h\).

Finally lift generators of the relative cotangent space to
the maximal ideal of \(A_{\rm low}\). They give a surjection
from a power series ring on at most \(3h\) variables over \(D\):
together with the image of the maximal ideal of \(D\) they
generate the maximal ideal modulo its square, and successive
approximation gives surjectivity by completeness. Taking the
two quotient dimensions proves the proposition.
\end{proof}

\begin{lemma}[Choosing a curve at a marked intersection]\leavevmode
\label{lem:markedcurves421}
Let \(Q\subset A_*\) be prime of quotient dimension at least
\(3d\), and suppose \(I_K\subseteq Q\cap R_0\).
There is a characteristic-\(\ell\) curve \(\mathfrak p\) in
\(R_0\), containing \(Q\cap R_0\), with one of the two matrix
realizations in Sections~\ref{sec:twoblockchain417}
and~\ref{sec:irreduciblechain418}. It is the contraction of a
one-dimensional marked specialization above \(Q\).
\end{lemma}
\begin{proof}
Cut \(A_*/Q\) by \(\varpi\) if necessary and choose a prime
\(Q'\) of the resulting special fiber attaining its dimension.
Then \(\dim A_*/Q'\ge3d-1\). This is greater than
\(2d+1+3h\). The lemma on the marked basis and the proposition
show that the generic determinant at \(Q'\) is not a sum of
three characters. By Theorem~\ref{thm:tracefiber420}, for
\(q'=Q'\cap R_0\) we have \(\dim R_0/q'\ge3d-2\).
This exceeds \(3d-b\) of Section~\ref{sec:weightintersections419}.
Since \(I_K\subseteq q'\), \(R_0/q'\) is finite over the weight
algebra. Its generic point avoids the exceptional weights of
Proposition~\ref{prop:badweights419}, by the finite-map dimension
argument of Theorem~\ref{thm:largeintersection419}.

Pull this open condition back to \(\operatorname{Spec}(A_*/Q')\).
It is a nonempty open in a complete local domain. The system of
parameters argument of Theorem~\ref{thm:largeintersection419}
gives a one-dimensional quotient there still satisfying the
condition. Embed it in its normalization \(k''[[z]]\). The
contraction to \(R_0\) has quotient dimension exactly one.
In fact the weight is outside the exceptional set, so some
principal unit weight ratio differs from one. Its difference
from one gives an element \(0\ne f\in (z)\) of the image
of \(R_0\). This image contains \(k[[f]]\) by completeness
and continuity, with \(k\) the residue field of \(O\).
The ring \(k''[[z]]\) is finite over \(k[[f]]\) by expansion
from a basis modulo \(f\). The intermediate image is a finite
module over this Noetherian DVR and has dimension one.

The marked matrices over the normalization are already
polarized and ordinary. The determinant outside the exceptional
set has at most two constituents. In the two-constituent
case their representations split orthogonally as in
Theorem~\ref{thm:tracefiber420}, the line being the middle
one. These statements remain true over a finite extension
of the contracted curve field: the extension to the marked
normalization is finite by the same argument with \(k[[f]]\).
Use Proposition~\ref{prop:curvelattice366} to obtain a marked
self-dual lattice with split reduction after finite extension.
The norm-one local characters are separated by construction,
as required.
\end{proof}

\begin{theorem}[All components at a nonzero adjacent lower entry]\leavevmode
\label{thm:adjacentchart421}
Suppose some matrix of \(\rho_*\) has a nonzero \(21\)- or
\(32\)-entry. If a prime \(P_s\subset A_*\) with
\(\dim A_*/P_s\ge3d\) has Hecke-supported contraction, then
\(I_K A_*\) is nilpotent. In particular every marked
specialization at this residual point has Hecke-supported
determinant.
\end{theorem}
\begin{proof}
The nonzero entry is a unit throughout the chart. A geometric
determinant with one line and a plane is therefore impossible
by Lemma~\ref{lem:splitorder421}. At every prime of quotient
dimension at least \(3d\) the determinant must be irreducible:
the other possible case would lie in the lower triangular
quotient, of dimension at most \(2d+2+3h<3d\).
In particular this applies to \(P_s\), to every minimal prime
of \(A_*\), and to primes of maximal dimension on the successive
intersections of Proposition~\ref{prop:dimensionchain420}.

Take a minimal prime \(P_0\subseteq P_s\). Apply
Lemma~\ref{lem:markedcurves421} with \(Q=P_s\). Its curve is
irreducible since a line and plane cannot occur anywhere in
the marked chart. It is supported by specialization from
\(P_s\cap R_0\). Thus Theorem~\ref{thm:completehecke418}
makes every prime below the contracted curve supported,
including \(P_0\cap R_0\). For an adjacent pair of components
in Proposition~\ref{prop:dimensionchain420}, take a prime of
maximal quotient dimension in their intersection. Once the
first component is supported, this prime has supported
contraction and quotient dimension at least \(3d+1\). The
same choice of curve and the same local theorem give support
at the second component. Connectedness propagates this conclusion
to all minimal primes. Their intersection is the nilradical,
which is nilpotent in the Noetherian ring \(A_*\); this gives
the stated ideal containment and its consequence.
\end{proof}

\begin{corollary}[Adjacent lower entries on the arithmetic fields]\leavevmode
\label{cor:lowerpoints421}
For every large-prime pair on the interval, finite residual
matrix points with a nonzero adjacent lower entry as in the
theorem exist on the layer \(L/F\).
\end{corollary}
\begin{proof}
Use the nonzero strict class of
Proposition~\ref{prop:strictclasses365} with character
\(\bar\Psi^{-1}\eta\) on \(\Gamma\), where \(\eta\) cuts out
\(L/F\). Choose a representative \(z\) with \(z(c)=0\);
this is possible by a coboundary since the character takes
value \(-1\) at \(c\). On \(H\) it has coefficient \(\bar\Psi^{-1}\)
and satisfies \(z(cgc)=-z(g)\). Its class on \(H\) remains
nonzero by restriction and corestriction. Put \(N=E_{21}+E_{32}\)
and \(n(x)=1+xN+x^2N^2/2\). Then
\[
   g\longmapsto n(z(g))\operatorname{diag}
               (\bar\Psi(g),1,\bar\Psi(g)^{-1})
\]
is a representation of \(H\), since \(n(x)n(y)=n(x+y)\) and
the diagonal conjugates \(N\) by \(\bar\Psi(g)^{-1}\).
It has the required polarization, since \(JN^{\mathsf T}J=N\)
and inverse transposition followed by reversal sends \(n(x)\)
to \(n(-x)\). It is locally conjugate to the diagonal at the
coefficient places by strictness, giving the prescribed ordered
flags. The finite coefficient field gives a finite image.
To mark it, take the eigenvectors of its lower triangular skew
matrix \(Z\), labeled by the three distinct diagonal entries.
They can be chosen as a lower triangular basis, then rescaled
to have form matrix \(J\) over a finite field extension. This
only rescales the \(21\) extension class and changes its cocycle
by a coboundary, as one sees on the quotient by the last coordinate
line. The class is nonzero, so some adjacent entry persists.
\end{proof}

The last corollary constructs residual points. Hecke support for
a prime of the size required in Theorem~\ref{thm:adjacentchart421}
is an additional assertion. A residual point whose only nonzero
strictly lower entries are in position \(31\) can also admit
one-dimensional plus two-dimensional constituents; its component
intersections require the corresponding two-constituent argument.


\subsection{Support through the chart with one long lower entry}
\label{sec:longrootcoverage422}

Continue on the fields and at the depth of
Section~\ref{sec:lowerchartcoverage421}. Suppose all strictly
lower entries at the finite residual point \(\rho_*\) vanish
except for the \(31\)-entry, which is nonzero on some group
element. Then the residual matrices preserve the middle line
and the outer plane. Let
\(D_0=\operatorname{diag}(-1,1,-1)\). Conjugation by \(D_0\)
fixes the residual matrices, preserves the form \(J\) and
the marking \(Z\), and induces an involution \(\sigma\) of
\(A_*\) over \(R_0\). A plus or minus sign below denotes the
eigenvalue of this involution. Both eigenspace projections
are available since \(2\) is invertible. A prime of \(R_0\)
will meet the hypothesis on a link in
Theorem~\ref{thm:chaintransfer417} if its determinant is
geometrically irreducible or its geometric \(C_0\)-fiber
is nonsemisimple.

\begin{proposition}[Even and odd relations in the marked ring]\leavevmode
\label{prop:paritypresentation422}
There is a presentation of \(A_*\) as a quotient of
\[
 P=O'[[x_1,\ldots,x_a,y_1,\ldots,y_g]],\qquad
                   \sigma x_i=x_i,\quad \sigma y_j=-y_j ,
\]
by \(b\) even relations \(F_i\) and \(g\) odd relations
\((My)_j\), where \(y\) is the column of the \(y_j\) and
\(M\) is a square matrix of even elements of \(P\).
Here \(a-b\ge3d+1\). The closed set \(y=0\) is exactly the
locus on which the universal matrices preserve the middle
line and the outer plane.
\end{proposition}
\begin{proof}
The residual adjoint of Lemma~\ref{lem:markedcomplex420},
its local flag stabilizers and their mapping fiber decompose
under \(D_0\). In fact \(D_0\) commutes with the residual
local action. The three separating eigenlines at the local
places therefore lie individually in the line or the plane,
and every ordered flag is preserved. Write \(h_{i,\pm}\)
for the two cohomology dimensions. The odd adjoint consists
of the two off-diagonal blocks between a plane and a line.
The odd dimension of a local flag stabilizer is two: among
the three eigenlines there are two unordered pairs with
opposite signs, and each contributes one upper triangular
entry. At a real place \(c_\tau=h c\) the residual matrix
of \(h\) preserves the blocks. The symmetric form defining
adjunction there makes these blocks orthogonal; on the two
off-diagonal blocks a skew-adjoint endomorphism is determined
by one map from the line to the plane, of dimension two.
The odd contribution at infinity is therefore two as well.
The dual invariant vanishing used in
Theorem~\ref{thm:markedpresentation420} persists for this
summand. Lemma~\ref{lem:frameeuler365} gives
\[
             h_{1,-}-h_{0,-}-h_{2,-}=2d-2d=0 .
\]
In degree zero the marked automorphisms are contained in the
diagonal torus Lie algebra from Lemma~\ref{lem:markedcomplex420},
which is even. Thus \(h_{0,-}=0\). The extra change-of-marking
tangent directions there are also even. The relative odd tangent
dimension is \(g=h_{1,-}=h_{2,-}\), and the even tangent
dimension \(a\) satisfies \(a-h_{2,+}=3d+1\) by the full Euler
formula.

Choose the generators of a minimal power series presentation
with their parities, using the projections \((1\pm\sigma)/2\).
The odd generators can be chosen from off-diagonal block entries
of the universal group matrices. Indeed the complete local
algebra is generated by matrix entries and local characters
as in Theorem~\ref{thm:markedcurve366}. The characters and
diagonal block entries are even and the off-diagonal block
entries are odd. Therefore the indicated entries span the odd
relative cotangent space. Every odd series belongs to \((y)\),
since its monomials have odd total degree in these variables.
Consequently \(y=0\) is exactly the stated block locus.

Let \(J_P\) be the kernel of the presentation. The injection
\((J_P/\mathfrak m_P J_P)^\vee\to H_C^2\) in
Theorem~\ref{thm:markedpresentation420} is equivariant:
conjugating the universal matrices and lifted flags by
\(D_0\) conjugates both the multiplicative defect and its
local projections, and these compute the obstruction class.
Thus generators for \(J_P\) can be chosen with at most
\(h_{2,+}\) even and at most \(g\) odd ones. Lift eigenbases
of its quotient by \(\mathfrak m_P J_P\) with the prescribed
parities and use Nakayama's lemma. Append zero odd relations
to give exactly \(g\) of them. In each odd series factor the
first \(y_j\) occurring in each monomial, in a fixed order;
the remaining factor has even parity. Grouping these formal
series gives \(My\) as asserted.
\end{proof}

\begin{lemma}[A nonsemisimple fiber on the determinantal locus]\leavevmode
\label{lem:singularblock422}
Assume \(g>0\) and put \(A^\dagger=A_*/(\det M)\), with
\(\det M\) evaluated in \(A_*\). Every prime of \(A_*\) with
geometrically irreducible determinant contains \(\det M\).
At every prime \(P_*\) of \(A^\dagger\) of quotient dimension
at least \(3d\), the contraction to \(R_0\) meets the hypothesis
on a link.
\end{lemma}
\begin{proof}
Multiplication by the adjugate in \(My=0\) gives
\((\det M)y_j=0\) for all \(j\). At an irreducible determinant
some \(y_j\) is nonzero because a block-preserving representation
is reducible. This proves the first assertion. At a prime of
the indicated dimension the determinant cannot split into three
characters, by Lemma~\ref{lem:splitorder421} and
Proposition~\ref{prop:lowerdimension421}. If it has a line
and plane, the same lemma gives \(y=0\) at the prime.

Take an algebraic closure \(K\) of its residue field, and
evaluate the presentation at the resulting point. As
\(\det M=0\), there is a nonzero \(v\in K^g\) with \(Mv=0\)
there. There is a lift of the evaluation to \(K[\epsilon]\),
\(\epsilon^2=0\), given by keeping the \(x_i\) and sending
\(y_j\) to \(\epsilon v_j\). This substitution on formal series
is defined by first reducing modulo \((y)^2\), then applying
the given evaluation on \(O'[[x]]\) to each coefficient.
The even relations have no linear \(y\)-term, and the odd
ones have first variation \(Mv=0\). The resulting map is
defined on \(A_*\). Its restriction to \(R_0\) is the constant
evaluation, because \(\sigma\) fixes \(R_0\) and every invariant
series has zero first variation at \(y=0\). An invariant
element of the quotient has an invariant lift by averaging.

Let \(q=P_*\cap R_0\). The universal group matrices now give
a representation of \(C_0\otimes_{R_0} K\) over \(K[\epsilon]\)
with a nonzero variation in its off-diagonal blocks; one of
the chosen \(y_j\) is such an entry. Suppose the geometric
fiber of \(C_0\) at \(q\) were semisimple. Its extension to
\(K\) would be a product of split matrix algebras. Every
first-order variation of a representation of this product,
with the algebra held fixed, is inner. Explicitly, in each
matrix factor take \(\sum_i E_{i1}\otimes E_{1i}\), and sum
these tensors over the factors. Writing the result as
\(\sum a_i\otimes b_i\), we have \(\sum a_i b_i=1\) and
\(\sum xa_i\otimes b_i=\sum a_i\otimes b_i x\).
For a derivation \(\partial\) relative to the given representation
\(\rho\), applying \(\partial\) to the first component of this
identity and multiplying shows that \(U=\sum\partial(a_i)\rho(b_i)\)
satisfies \(\partial(x)=U\rho(x)-\rho(x)U\).
This is the asserted infinitesimal conjugation.

The marking \(Z\) has zero variation: its diagonal scalar is
determined by the base trace and has its fixed choice of square
root. Thus \(U\) commutes with \(Z\) and is diagonal, since
the three eigenvalues are distinct. Commuting a diagonal
matrix with a matrix that preserves the two blocks gives
zero in the off-diagonal blocks. This contradicts the
variation just constructed. The geometric fiber at \(q\)
is therefore nonsemisimple, as required.
\end{proof}

\begin{proposition}[Intersections after imposing the determinant]\leavevmode
\label{prop:incidencegraph422}
If an irreducible prime occurs in \(A_*\), in the sense of its
geometric determinant, then \(g>0\), \(A^\dagger\) is nonzero
and every component of \(\operatorname{Spec}A^\dagger\) has
dimension at least \(3d+1\). These components form a connected
graph under intersections of dimension at least \(3d\).
\end{proposition}
\begin{proof}
For \(g=0\) the block locus is the entire ring, so an
irreducible determinant forces \(g>0\). The lemma then
puts \(\det M\) in the corresponding prime. Use the
\(a+g\) variables in the presentation over \(O'\) and the
\(b+g+1\) relations \(F,My,\det M\). The proof of
Proposition~\ref{prop:dimensionchain420}, with
\(N=1+a+g\) and \(s=b+g+1\), gives component dimension at
least \(a-b\) and the connected graph with intersection
dimension at least \(a-b-1\). Now \(a-b\ge3d+1\).
\end{proof}

\begin{theorem}[Hecke support of irreducible determinants at the long entry]\leavevmode
\label{thm:longrootchart422}
Suppose \(P_s\subset A_*\) is prime of quotient dimension
at least \(3d\) and its contraction to \(R_0\) is both
Hecke-supported and geometrically irreducible. Then
\(I_K A^\dagger\) is nilpotent. Every prime of \(A_*\)
with geometrically irreducible determinant has Hecke-supported
contraction.
\end{theorem}
\begin{proof}
By Lemma~\ref{lem:singularblock422} the prime \(P_s\) lies
in the closed subscheme \(\operatorname{Spec} A^\dagger\).
Choose a minimal prime \(P_0\) of this subscheme below \(P_s\).
To reach any other minimal prime, take a chain
\(P_0,P_1,\ldots,P_t\) supplied by the proposition, viewed
as primes of \(A_*\). Write \(J_i=P_i\cap R_0\). Every \(J_i\)
meets the link hypothesis by Lemma~\ref{lem:singularblock422}.
Choose a marked curve above \(P_s\) using
Lemma~\ref{lem:markedcurves421}; write \(\mathfrak p_0\)
for its contraction to \(R_0\). It lies above the given
supported irreducible prime \(P_s\cap R_0\). By
Corollary~\ref{cor:mixedchain418} for this one curve,
its indicated generizations belong to Hecke support, in
particular \(J_0\).

Suppose curves \(\mathfrak p_0,\ldots,\mathfrak p_i\) have
been obtained with \(J_i\subseteq\mathfrak p_i\) supported,
and with the two successive curves at each earlier step
both above the corresponding \(J_j\). In the intersection
defined by \(P_i+P_{i+1}\), choose a prime \(Q\) of quotient
dimension at least \(3d\). Its contraction contains \(J_i\)
and is therefore supported. Lemma~\ref{lem:markedcurves421}
now provides the next curve \(\mathfrak p_{i+1}\) above
\(Q\cap R_0\). Both \(\mathfrak p_i\) and
\(\mathfrak p_{i+1}\) contain \(J_i\), which is an allowed
link even when its geometric determinant has two constituents.
The prime \(J_{i+1}\) is an allowed endpoint below the new
curve. Apply Corollary~\ref{cor:mixedchain418} to the whole
finite chain, using the original supported irreducible
starting prime below \(\mathfrak p_0\). Its conclusion
gives \(I_K\subseteq J_{i+1}\). Using the whole chain keeps
the irreducible starting hypothesis available at stages where
\(J_i\) has reducible determinant.

Induction proves the containment for every minimal prime of
\(A^\dagger\). Its nilradical is nilpotent by Noetherianity.
The first assertion of Lemma~\ref{lem:singularblock422} then
proves the conclusion for every irreducible determinant of
\(A_*\). The related determinantal construction in the source
appears in \srcloc{Development/AutomorphicArithmetic/AutomorphicVarieties/FrobeniusSmoothVectors/ChainedOfComap}{93}.
\end{proof}

\begin{corollary}[Local propagation at the nonzero lower points]\leavevmode
\label{cor:alllowercharts422}
At any marked lower triangular residual point with a nonzero
strictly lower entry, a supported geometrically irreducible
prime of marked quotient dimension at least \(3d\) gives
Hecke support at every geometrically irreducible prime in
that chart. Both adjacent-entry and pure \(31\)-entry residual
points exist on the fields of Section~\ref{sec:commonframe365}.
\end{corollary}
\begin{proof}
If an adjacent entry is nonzero, apply
Theorem~\ref{thm:adjacentchart421}; existence in this case
is Corollary~\ref{cor:lowerpoints421}. Otherwise apply the
theorem just proved. To construct a pure \(31\)-entry point,
use the strict class with character \(\bar\Psi^{-2}\eta\) in
Proposition~\ref{prop:strictclasses365}, choose \(z(c)=0\),
and restrict to \(H\). Its restriction remains nonzero.
Use \(N=E_{31}\) and \(n(x)=1+xN\) in the proof of
Corollary~\ref{cor:lowerpoints421}. Now \(N^2=0\), conjugation
by the residual diagonal has character \(\bar\Psi^{-2}\)
on \(N\), and \(JN^{\mathsf T}J=N\), giving the same local
splitting and polarization. The skew marking is diagonalized
within the outer plane and the middle vector is retained.
On the plane the \(31\) extension class is nonzero and is
only rescaled or changed by a coboundary in this basis. Thus
the long entry occurs and both adjacent entries are zero.
\end{proof}

This proves local propagation at all the nonzero lower residual
points, with the same group and local conditions at each step.
Passing from a supported arithmetic family to all of these
charts, and estimating the lengths at the prescribed integral
weights, remain necessary to obtain the bound of
Section~\ref{sec:missing}.


\subsection{Global Hecke support for irreducible ordinary determinants}
\label{sec:globalheckesupport423}

Fix one of the large-prime triples and its layer \(L/F\) from
Section~\ref{sec:commonframe365}. Write \(X=\operatorname{Spec} S\)
for the finite-type scheme over \(R_0\) constructed in
Theorem~\ref{thm:tracefiber420}, using the original coefficients
\(O\). Thus a point specifies matrices of \(C_0\) in dimension
three with the base determinant, pairing \(J\) and diagonal
marking \(Z\), after removing the common cyclotomic twist.
Let \(k\) be the finite residue field of \(R_0\) and \(X_k\)
the fiber at its closed point. The torus \(\mathbf G_m\)
acts on \(X\) by conjugation by
\(\operatorname{diag}(t,1,t^{-1})\). Let \(U\) be the open in
\(X_k\) where some strictly lower entry of a group matrix is
nonzero. Within it let \(U_1\) require some adjacent lower
entry to be nonzero.

\begin{lemma}[Marked points and the determinant map]\leavevmode
\label{lem:matrices423}
Every geometrically irreducible prime of \(R_0\) has a nonempty
geometric fiber in \(X\). Its geometric points form a single
one-dimensional torus orbit. Every irreducible component of
\(X\) is stable under the torus. At a closed point \(x\in U\)
with finite residue field \(k'\), the completed local ring is
the marked ring \(A_*\) for the matrices of that point over the
unramified coefficient extension with residue field \(k'\).
\end{lemma}
\begin{proof}
Over a geometric field the irreducible representation realizing
the determinant factors through \(C_0\), by the word identities
and the trace pairing as in Lemma~\ref{lem:positivefibers419}.
Adjunction for its normalized group matrices is given by a
nondegenerate symmetric form, by irreducibility and odd
dimension. The skew operator \(Z\) has eigenvalues \(z,0,-z\),
with \(z\) the specified square root of \(t(Z^2)/2\) in the
base. An eigenbasis has a nonzero pairing of the outer vectors
and a nonzero norm of its middle vector; rescale them to
obtain \(J\). This gives a point. As in
Theorem~\ref{thm:tracefiber420} any two such points are related
by \(\operatorname{diag}(s,e,s^{-1})\), \(e=\pm1\). If \(e=-1\)
multiply by the scalar \(-1\) without changing conjugation.
This is the given torus action. Its stabilizer on an irreducible
representation is trivial: a commuting matrix is scalar and
the middle entry of the torus element is one. Thus the orbit
has dimension one. Every irreducible component \(C\) of \(X\)
is stable because \(\mathbf G_m\times C\) is irreducible
and its image contains \(C\) at \(t=1\).

At a geometric point of \(X_k\) the determinant has three
character constituents. Its invariant flag is a coordinate flag
by the three eigenvalues of \(Z\). The directed-cycle argument
of Lemma~\ref{lem:splitorder421} then shows that every point of
\(U\) is lower triangular. The diagonal entries at any such point
are the prescribed \(\bar\Psi,1,\bar\Psi^{-1}\) on the group.
Indeed for the \(i\)-th polynomial projector \(e_i\) of \(Z\)
the entry is \(t(e_i[g])\), whose value is fixed by the split
trace on the closed fiber. The local conditions and marking
are part of \(X\). To pass to the completion at \(x\), the
unramified coefficient extension \(O'/O\) with residue field
\(k'\) maps uniquely to \(\widehat{\mathcal O}_{X,x}\) lifting
the specified residue-field map, by Hensel's lemma. Equivalently,
after this finite \'{e}tale coefficient extension the completed
factor at the chosen lift of \(x\) is the same ring. The
Artinian test-ring identification in
Theorem~\ref{thm:tracefiber420} gives exactly \(A_*\).
\end{proof}

\begin{proposition}[The lower triangular points in the closed fiber]\leavevmode
\label{prop:lowercore423}
The space \(U\) is nonempty and connected, and \(U_1\) is
dense in \(U\).
\end{proposition}
\begin{proof}
Work first over \(\bar k\). The constructions of
Corollary~\ref{cor:alllowercharts422} give a fixed pure
\(31\)-entry point and a fixed adjacent-entry point. Write
\(d(g)=\operatorname{diag}(\bar\Psi(g),1,\bar\Psi(g)^{-1})\).
At a pure point the matrices are \(d(g)+b(g)E_{31}\);
denote the function for the fixed pure point by \(b_0\).
At any lower point write \(a_{ij}(g)\) for its strictly lower
entries. For an indeterminate \(T\) define
\[
 \rho_T(g)=d(g)+T a_{21}(g)E_{21}+T a_{32}(g)E_{32}
               +\bigl(T^2a_{31}(g)+(1-T)b_0(g)\bigr)E_{31}.
\]
This is a representation: scaling a lower entry by
\(T^{i-j}\) respects matrix multiplication, and the additional
\(31\) term satisfies the homogeneous equation
\(b_0(gh)=b_0(g)d_{11}(h)+d_{33}(g)b_0(h)\). There are no
other products with that entry in the strictly lower part.
Adjunction pairs the adjacent entries and preserves the long
position, giving the original adjunction identity. The
off-diagonal entries of \(Z\) are zero for both input matrices
and remain zero under this formula.

The ordered local flags persist over \(\bar k[T]\). At each
chosen place the original local action is conjugate to \(d\)
by a lower unitriangular matrix \(V\), by the proof of
Proposition~\ref{prop:lowerdimension421}. Scaling the entries
below the diagonal of \(V\) by \(T^{i-j}\) gives a lower
unitriangular polynomial matrix effecting the same scaling
on the conjugated diagonal. The pure \(31\) function at this
place has the form \(b_0(g)=c_u(d_{11}(g)-d_{33}(g))\);
one may diagonalize its outer plane directly. Conjugating
the scaled matrices by \(1+(1-T)c_uE_{31}\) adds exactly the
prescribed final term. Thus the words hold with the original
local diagonal characters. The global determinant is constant
since the matrices are triangular. This gives a morphism
\(\mathbb A^1_{\bar k}\to X_k\). To check factorization through
the finite Galois algebra, the original group functions factor
through finite quotients: they are specializations of the
group map into the finite algebra \(C_0/\mathfrak m_{R_0}C_0\)
over a finite field. The new group action factors through their
common finite quotient. Every group algebra element has the
prescribed characteristic polynomial, by triangularity, so
Cayley--Hamilton, the old inertia relations and the verified
ordinary words give the required factorization.

At \(T=0\) this is the fixed pure point; at \(T=1\) it is
the chosen lower point. If the latter belongs to \(U\), the
inverse image of \(U\) is a nonempty open subset of an affine
line containing both endpoints. It is connected. Thus every
closed point of \(U_{\bar k}\) is connected in this space to
the fixed one. Any nonempty open and closed summand contains
a closed point, by finite type over \(\bar k\); this proves
connectedness. For density take any pure point and use its
\(b\) in place of \(b_0\), with the fixed adjacent point
as the other input. The value at zero is the chosen pure
point and all nonzero values have an adjacent entry. Thus
every closed point of \(U_{\bar k}\) is in the closure of
\((U_1)_{\bar k}\). Closed points are dense in every closed
subset, so this is the whole space. Surjectivity of the
base-change map gives connectedness and density over \(k\).
The polynomial deformations are also used in
\srcloc{Development/AutomorphicArithmetic/AutomorphicVarieties/FrobeniusSmoothVectors/ExistsNonzeroLowCocycle}{144}.
\end{proof}

\begin{lemma}[Specialization to a nonzero lower entry]\leavevmode
\label{lem:specializemarks423}
Every torus-stable closed irreducible subset \(C\subseteq X\)
with geometrically irreducible generic determinant contains a
closed point of \(U\).
\end{lemma}
\begin{proof}
Let \(p\) be its generic contraction to \(R_0\) and put
\(D=R_0/p\). It is a positive-dimensional complete local
domain: the closed determinant has three distinct character
constituents. The map \(C\to\operatorname{Spec}D\) is dominant
and of finite type, so its image contains a nonempty open by
Chevalley's theorem. Intersect this open with the locus where
a nine-element trace Gram determinant that is nonzero at \(p\)
remains nonzero. By the parameter argument of
Theorem~\ref{thm:largeintersection419}, which applies in either
characteristic, this open contains a prime giving a
one-dimensional quotient of \(D\). At its curve field choose
a point of \(C\) over a finite extension of the field, using
a closed point of the nonempty finite-type fiber. The
representation there is absolutely irreducible. The normalized
curve ring is a complete DVR with finite residue field: complete
Noetherian local domains here are excellent, so normalization
in dimension one is finite, and the Henselian property makes
the normalization in a finite field extension local.

Over this DVR the group matrices have a stable lattice since
the group algebra \(C_0\) is finite over the base. By the
construction in Proposition~\ref{prop:curvelattice366}, after
a finite extension we can take a marked self-dual basis with
split reduction. In detail the unit-separated eigenvalues of
\(Z\) split the lattice, a small spread of exponents splits
the residual flag, and averaging with the dual exponents makes
the lattice self-dual. The final unit rescalings restore \(J\).
That argument uses only the DVR, the pairing, and the distinct
residual eigenvalues, and hence works in mixed characteristic
as well. The normalized point in this basis stays in \(C\):
it differs from the initial point by the torus action by
Lemma~\ref{lem:matrices423}.

Write \(v\) for the discrete valuation. All nonzero strictly
lower entries now have positive valuation, and at least one
exists, since matrices all upper triangular would contradict
irreducibility. Set
\[
       e=\min_{\substack{i>j,\ g\in H\\ \rho(g)_{ij}\ne0}}
                         \frac{v(\rho(g)_{ij})}{i-j}>0 .
\]
The minimum is attained since the value group is discrete
and there are only three positions. After a finite ramified
extension take \(t\) of valuation \(e\), with the valuation
extended in the same normalization. Torus conjugation scales
the \(ij\)-entry by \(t^{j-i}\), so all entries stay integral,
every upper entry still has positive valuation, and some lower
entry is a unit. The resulting matrices give a point over the
valuation ring: the algebra is spanned by group elements and
the scalar relations and local characters keep their values.
Its generic point is in \(C\) by stability; its closed point
is in \(C\) by closedness and in \(U\) by the unit just found.
Its residue field is finite, so the image is a closed point
of \(X_k\), and of \(X\).
\end{proof}

\begin{theorem}[Hecke support at every irreducible ordinary determinant]\leavevmode
\label{thm:globalhecke423}
For the ordinary determinant ring \(R_0\) over the field \(L\)
and at the level of Section~\ref{sec:twoblockchain417}, every
prime with geometrically irreducible determinant contains the
kernel \(I_K\) of the good Hecke map.
\end{theorem}
\begin{proof}
Let a component of \(X\) be called irreducible at its generic
determinant if its generic point has that property. The locus
of such determinants in \(X\) is open: a nonzero determinant
of the trace pairing on nine group elements detects it, and
one takes the union of these opens. Among components with
this generic property, consider those on which the Hecke
kernel vanishes in the reduced coordinate ring. We first
exhibit one.

Take the full Hecke component through the cuspidal adjoint in
Proposition~\ref{prop:upfamily372}. Its domain \(B\) is a
quotient of \(R_0/I_K\), has geometrically irreducible generic
determinant, and is finite over and contains \(\Lambda\).
Thus \(\dim B=3d+1\). In \(X\) take the reduced closure \(Y\)
of the orbit over this generic point. The geometric fiber has
one orbit of dimension one by Lemma~\ref{lem:matrices423};
its reduced generic fiber is irreducible and its function
field has transcendence degree one over \(\operatorname{Frac}B\).
The closure is irreducible, torus-stable, dominates \(B\),
and is supported by \(I_K\). By the preceding lemma choose
a finite closed \(x\in Y\cap U\).

At this \(x\) the dimension of \(\mathcal O_{Y,x}\) is \(3d+2\).
Here is the dimension calculation. Present the coordinate ring
of \(Y\) as a domain quotient of \(B[T_1,\ldots,T_n]\) by a
prime contracting to zero. Its height is \(n-1\) by the
transcendence degree over the fraction field of \(B\): every
prime below it contracts to zero as well. The inverse image
of \(x\) has height \(\dim B+n\), by a chain to
\(\mathfrak m_B B[T]\) followed by a height-\(n\) closed point
in the residue polynomial ring, and the dimension upper bound
for a polynomial ring. The polynomial ring over the complete
local domain \(B\) is catenary, so the difference gives the
required local dimension. Completion preserves this dimension.
Choose a minimal prime in \(\widehat{\mathcal O}_{Y,x}\)
attaining it. This prime contracts to zero in
\(\mathcal O_{Y,x}\): completion is flat, and a nonzero element
of the domain remains a nonzero divisor and is in no minimal
prime. Its inverse image in
\(A_*=\widehat{\mathcal O}_{X,x}\) is therefore supported,
has geometrically irreducible determinant, and has quotient
dimension \(3d+2\).

Apply Corollary~\ref{cor:alllowercharts422}. It puts every
geometrically irreducible determinant of this chart in Hecke
support. Some global component containing \(Y\) has geometrically
irreducible generic determinant, by the Gram open. Every such
component through \(x\) has a branch in the completion with the
same generic contraction: by faithful flatness lift its local
minimal prime and then take a minimal prime below the lift.
It still contracts to that local minimal prime. Geometric
irreducibility persists under residue-field extension. Thus
the Hecke kernel vanishes at the generic point of this global
component. This supplies the promised component.

The same argument transfers support at every finite closed
point \(z\in U\) shared by two components having geometrically
irreducible generic determinants. If the first is supported,
a minimal prime of \(\widehat{\mathcal O}_{X,z}\) above its
generic prime is a supported irreducible prime of quotient
dimension at least \(3d+2\), by
Proposition~\ref{prop:dimensionchain420}. Applying
Corollary~\ref{cor:alllowercharts422} and contracting a branch
of the second component proves that it is supported too.
Moreover at each finite closed point of \(U_1\) every component
through the point has geometrically irreducible generic
determinant. For any minimal branch its dimension is at least
\(3d+2\); the arguments of Lemma~\ref{lem:splitorder421} and
Proposition~\ref{prop:lowerdimension421} exclude both a sum of
three characters and a line with a plane at its generic point.

There are only finitely many components. Their subunion with
geometrically irreducible generic determinant thus meets \(U\)
in a closed subset containing all closed points of \(U_1\);
it contains \(U\) by the proposition and density of closed
points. Intersect \(U\) with the union of the supported
components in that list, and with the union of the others.
These are two closed sets covering \(U\). Their intersection
is empty: if nonempty it has a finite closed point, since \(U\)
is of finite type over a finite field, contradicting the
transfer just proved. The first set is nonempty by the
construction from the adjoint. Connectedness of \(U\) makes
the second empty. Every component of the list meets \(U\)
by Lemma~\ref{lem:specializemarks423}, so every one is supported.

Finally take any geometrically irreducible prime of \(R_0\).
A geometric point above it exists by Lemma~\ref{lem:matrices423}.
Any component of \(X\) through this point has generically
irreducible determinant by openness, hence is supported. The
base prime contains \(I_K\). The corresponding component theorem
in the source is
\srcloc{Development/AutomorphicArithmetic/AutomorphicVarieties/CurveDVRSpecialization/PropagationTheoremGL3}{213}.
\end{proof}

\begin{corollary}[The irreducible closure at the prescribed weights]\leavevmode
\label{cor:irreducibleweight423}
Let \(R_{\rm irr}\) be the quotient of \(R_0\) by the
intersection of all its geometrically irreducible primes.
It is a finite \(\Lambda\)-algebra. At any original integral
arithmetic weight \(\Lambda\to O_w\), its base change is
finite over \(O_w\). Every geometric characteristic-zero
point there occurs as an eigencharacter on the ordinary module
at that weight, and their number is at most
\(\operatorname{rank}_{\Lambda}M_K\).
\end{corollary}
\begin{proof}
The theorem makes \(R_{\rm irr}\) a quotient of \(B_K\), which
is finite over \(\Lambda\). Finiteness survives base change.
A field point of this quotient has nonzero \(M_K\)-fiber, since
\(M_K\) is finite and faithful over \(B_K\). Base change to the
chosen geometric weight identifies this fiber with the residue
fiber of the specialized ordinary module. Over an algebraically
closed characteristic-zero field the specialized commutative
finite algebra is a product of Artinian local algebras. Each
such point gives a distinct nonzero summand of the module and
an eigenvector there, by the nilpotence of its maximal ideal.
The full dimension of the specialized module equals its stated
rank by weight freeness.
\end{proof}

This finiteness holds at the original arithmetic weight even
when the determinant there is reducible but lies in the indicated
closure. The integral lengths in Corollary~\ref{cor:strictminors415}
involve a matrix corner at the split determinant; finiteness of
the reduced trace closure alone gives no numerical bound for
those lengths or for the signed projected periods.


\subsection{The entire trace ring at the original arithmetic weight}
\label{sec:originalfixedhecke424}

Fix \(V\le x\le V^{3/2}\) on an interval of
Proposition~\ref{prop:quadratic345}. Put \(K=\QQ(\sqrt{-s})\).
Take \(\ell>P(x)\), \(1\le r<t(x)\), and a character \(\alpha\)
of \(E_x\) with \(\alpha(c)=(-1)^r\). Write \(O\) for the
unramified coefficient ring of this character, \(J=\operatorname{Frac}O\),
and \(T\) for the complete trace ring of the original marked
fixed-weight problem on \(G_K\). Its augmentation sends the
universal trace \(t\) to the trace of
\(\Psi\oplus1\oplus\Psi^{-1}\), where \(\Psi=\alpha\kappa_\ell^r\).
Here \(K\) is independent of the interval, \(x,r,\alpha\) and \(\ell\).

\begin{theorem}[Finiteness at the original weight, including torsion]\leavevmode
\label{thm:originalfinite424}
The ring \(T\) is finite over \(O\).
\end{theorem}
\begin{proof}
Pass first to the coefficient extension of Proposition~\ref{prop:orighecke374}
and denote the extended rings by the same letters. At the depth
of Section~\ref{sec:twoblockchain417} that proposition gives a
surjection \(R_0\to T\). It sends the weight algebra \(\Lambda\)
to the scalars by the original three ordered characters, with
the fixed common normalization. Thus \(T/I_K T\) is a quotient
of the specialization of \(B_K=R_0/I_K\) at this weight.
It is finite over the coefficient DVR, since \(B_K\) is finite
over \(\Lambda\).

Take a prime \(p\) of \(T\) over the uniformizer. Over an
algebraic closure of its residue field, the trace pairing on
the image algebra of the semisimple representation is
nondegenerate. On a simple matrix factor it is the matrix
trace multiplied by the multiplicity, which is at most three
and is invertible here. Consequently \(I_9(G_K)\not\subseteq p\)
precisely when the representation is absolutely irreducible.
Lemma~\ref{lem:imagespan374} gives
\(I_9(G_K)=I_9(G_L)T\). If \(I_9(G_K)\not\subseteq p\), the
preimage of \(p\) in \(R_0\) therefore has geometrically
irreducible determinant. Theorem~\ref{thm:globalhecke423} puts
\(I_K T\) inside \(p\). In the other case \(p\) contains
\(I_9(G_K)\), and \(T/I_9(G_K)\) is finite over the DVR by
Theorem~\ref{thm:reducibleminors378}. In either case \(p\)
is maximal since it lies in the special fiber of a finite
algebra. The Noetherian special fiber of \(T\) is therefore
Artinian.

Choose lifts of a residue-field basis of that special fiber.
Their \(O\)-span surjects modulo the uniformizer \(\varpi\).
Lift successive remainders after division by \(\varpi\), obtaining
coefficients in \(O\) converging in the \(\varpi\)-adic topology.
The errors lie in \(\varpi^n T\subseteq\mathfrak m_T^n\),
so separatedness gives equality in \(T\). Thus the lifts generate
the whole ring as a module. Finally the coefficient extension is
finite faithfully flat and formation of the trace ring commutes
with it by Proposition~\ref{prop:orighecke374}. Finiteness
descends. For instance the original \(T\) embeds in the extended
ring, which is a finite module over the original Noetherian DVR.
\end{proof}

\begin{proposition}[Classical labels for every nonsplit trace]\leavevmode
\label{prop:alloriginalpoints424}
Each geometrically irreducible characteristic-zero point of \(T\)
is represented, with the normalization of
Corollary~\ref{cor:irreduciblelabel374}, by a conjugate-self-dual
cuspidal representation of \(\mathrm{GL}_3/K\) at level \(N_*\).
Each geometrically reducible point other than the augmentation is
the sum of the trivial character and an irreducible plane. It has
the primitive modular form label of Proposition~\ref{prop:planelabels379}
at a level dividing \(sN_*\). For fixed \(\ell,r,\alpha\), a label
with its embedding and Frobenius convention determines at most
one trace point.
\end{proposition}
\begin{proof}
By finiteness, a geometric point is defined over a finite extension
of \(J\). It maps \(T\) into the integers there. This is a
continuous local map: the inverse image of the maximal ideal
is the unique maximal ideal over the uniformizer and its powers
map into the corresponding powers. At an irreducible point
the nine-by-nine minor argument in the preceding proof works
also in characteristic zero. Hence restriction to \(G_L\)
gives an irreducible eigencharacter of \(B_K\) at the original
weight. Corollary~\ref{cor:irreducibleweight423} supplies the
eigencharacter on the ordinary module. Ordinary control at
the dominant algebraic weight identifies its characteristic-zero
specialization with classical definite automorphic forms, exactly
as in the proof of Corollary~\ref{cor:classfiber372}. At \(r=1\)
the coefficient weight in this normalization is zero and the
same control theorem applies.~\cite{IX-VII-geraghty363}

In the transfer of this system to \(\mathrm{GL}_3/L\), the
attached semisimple Galois representation is the direct sum
from its isobaric parameters (with their prescribed twists).
Compatibility of good Frobenius traces and Chebotarev identify
this representation with the normalized restriction of the point.
Irreducibility forces a single rank-three cuspidal parameter.
Descend it along a tower of cyclic prime-degree steps from \(L\)
to \(K\), as in Theorem~\ref{thm:auxdescent361}. At each step
the original Galois representation makes the upper system
invariant, by strong multiplicity one. The possible descents
differ by characters of the cyclic group. The normalized central
character singles out one, since cubing is bijective on this
group of order prime to three. Uniqueness makes this descent
conjugate-self-dual. On the Galois side the two extensions of the
same absolutely irreducible restriction differ by a character by
Schur's lemma. Their determinants fix its cube, so this character
is trivial. Restriction is irreducible at each step because
it is irreducible already on \(G_L\). At the bottom
Proposition~\ref{prop:tracecrystal375} gives crystallinity and
the prescribed finite inertia. The level and weight computation
of Theorem~\ref{thm:auxdescent361} now places the descent in the
stated list.~\cite{IX-VII-I-ac123,IX-VII-lgp361,IX-VII-V-shin}

At a reducible point all nine-by-nine minors vanish. Three
character summands give the augmentation by
Proposition~\ref{prop:origideal373}. In the other case
Theorem~\ref{thm:alllabels375} identifies the line as the
trivial character. The point factors through \(T/I_9(G_K)\),
so Proposition~\ref{prop:planelabels379} applies. The label of
the plane determines its trace on \(G_K\). For a rank-three
label the good Frobenius traces determine the normalized trace
on \(G_K\) by continuity and Chebotarev. In both cases group
traces topologically generate \(T\), which proves the last
assertion.
\end{proof}

\begin{corollary}[The generic trace algebra]\leavevmode
\label{cor:originaletale424}
The algebra \(T[1/\ell]\) is finite \'{e}tale over \(J\).
\end{corollary}
\begin{proof}
At every reducible point Theorem~\ref{thm:reduciblerigid379}
already computes the local ring of the full trace algebra.
At an irreducible point the argument of Theorem~\ref{thm:rigidtrace377}
now applies to the classical system in the proposition.
More explicitly, a dual-number deformation of its determinant
gives a representation lift by absolute irreducibility. Its
continuous cocycle on \(G_K\) is invariant for the polarized
action of complex conjugation, hence descends to \(\QQ\)
by inflation--restriction. Its restriction to inertia away
from \(\ell\) is a coboundary by averaging over the fixed
finite inertia. Above \(\ell\) its self-extension is crystalline
by Proposition~\ref{prop:tracecrystal375}. These local conditions
descend by restriction and corestriction. Lemma~\ref{lem:cyclointegral377}
gives enormous image on the cyclotomic kernel. Theorem~A of
Newton and Thorne for the polarized adjoint of this cuspidal
parameter makes its finite Selmer group zero. The lift is
therefore conjugate to the constant lift and all its traces
are constant. After a finite splitting extension each local
factor of \(T[1/\ell]\) is Artinian and has zero cotangent
space, so its maximal ideal vanishes by Nakayama's lemma.
This is geometric reducedness of the finite algebra.~\cite{IX-VII-VI-nt119}
\end{proof}

\begin{theorem}[One integer for the torsion-free original trace orders]\leavevmode
\label{thm:originalorder424}
Fix the split prime \(u=u_1\) and a place of \(K\) above it
as in Proposition~\ref{prop:quadratic345}. For each triple put
\[
 T^{\mathrm{tf}}=T/T[\ell^\infty],\qquad
 I=\ker(T^{\mathrm{tf}}\longrightarrow O),\qquad
 h=t(\operatorname{Frob}_u)-
       (1+u^r\alpha(u)+u^{-r}\alpha(u)^{-1}).
\]
Write \(b_\alpha\) for the size of the residue orbit
\([\alpha]\). There is an integer \(D(V)\ge1\), independent of
\(x\) on the interval, such that
\[
 \sum_{\ell>P(x)}(\log\ell)
    \sum_{1\le r<t(x)}\sum_{[\alpha]} b_\alpha
            \operatorname{length}_O(I/hI)
       \ \le\ \log D(V)\ \ll_{\psi,\varepsilon} V^C.
\]
The orbits here have the specified complementary sign.
\end{theorem}
\begin{proof}
Take \(D(V)=D_2(V)\prod_{r,H}|I_{u,H,r}|\), where \(D_2(V)\)
is the integer of Theorem~\ref{thm:reducibleorder379}, the
other factors are those of Theorem~\ref{thm:integer}, and the
product is over \(1\le r<t_*\) and the character kernels on
\(E_*\) of the required sign as in Theorem~\ref{thm:heckedivisor345}.
These are the integers for the chosen prime \(u\) itself.
Theorem~\ref{thm:integer} bounds their logarithms by a
polynomial in \(N_*(t_*+2)\) times the corresponding orbit
sizes. The sum of those sizes for any \(r\) is at most
\([E_*:\QQ]\). Theorem~\ref{thm:reducibleorder379} bounds the
first factor and the interval bounds all these parameters by
powers of \(V\). All factors are nonzero.

The finite torsion-free \(O\)-module \(I\) is free. Its geometric
generic fiber has one line at each nonsplit point and no others
by Corollary~\ref{cor:originaletale424}. At all the corresponding
labels the Frobenius trace has complex absolute value at most
three. For planes this follows directly from the formula
\(1+u^{-r}\alpha(u)a_u(f)\) and Deligne's bound.
The split trace has absolute value at least \(u^r-1-u^{-r}>3\).
Thus multiplication by \(h\) is injective on \(I\). Its
determinant on the geometric generic fiber is the product
of the trace differences over these nonsplit points. The
Smith normal form over the unramified \(O\) gives
\[
 \operatorname{length}_O(I/hI)
       =\sum_{\sigma\ne{\rm split}} v_\ell(\sigma(h)),
                  \qquad v_\ell(\ell)=1 .
\]
For each \(\ell\) choose a common coefficient embedding.
The rank-three labels in the proposition occur with their signs
and embeddings in the product defining \(I_{u,H,r}\). The
plane labels occur separately in \(D_2(V)\), where each
difference is multiplied by \(u^r\). The denominator-clearing
powers of \(u\) in both products are \(\ell\)-adic units.
For fixed \(\alpha,r\), different points use different global
labels, even if their values at \(u\) coincide. The \(b_\alpha\)
local conjugates use different character labels with equal
valuations. Inflation \(E_x\subseteq E_*\) preserves these
labels. Hence every valuation counted at this \(\ell\) occurs
in \(v_\ell(D(V))\), and all unused factors are integral at \(\ell\).
Summing over primes proves the theorem.
\end{proof}

This estimate concerns the lattice \(I\). To use it for the
negative-weight cotangent module in Section~\ref{sec:missing},
one must also control the odd-weight directions of the marked
matrix problem at the same trace points.

\section{Auxiliary Hecke polynomials and trace determinants}\label{nav:hecke-polynomials}

\subsection{Spherical Hecke polynomials at the auxiliary primes}
\label{sec:heckechains716}

Let \(L\) be a nonarchimedean local field, with uniformizer
\(\varpi\), residue field of cardinality \(q\) and valuation
ring \(\mathcal O_L\). Put \(G=\mathrm{GL}_n(L)\) and
\(K=\mathrm{GL}_n(\mathcal O_L)\). Let \(\pi\) be an irreducible
unitarizable spherical representation of \(G\). We use the
arithmetic normalization of the Hecke polynomial \(P(X)\).
If \(t_i\) is the eigenvalue on the \(K\)-fixed line of the
double coset with \(i\) diagonal entries \(\varpi\) and the rest
\(1\), using counting measure on right cosets, then
\[
           P(X)=\sum_{i=0}^n(-1)^i q^{i(i-1)/2}t_iX^{n-i}.
\]
The Satake isomorphism says that its roots are the normalized
Satake parameters multiplied by \(q^{(n-1)/2}\). Indeed the
eigenvalue of that minuscule double coset is the \(i\)-th
elementary symmetric polynomial in the normalized parameters
times \(q^{i(n-i)/2}\). Adding the two exponents gives
\(i(n-1)/2\).

\begin{lemma}[A root of specified absolute value]\leavevmode
\label{lem:sphericalchain716}
Suppose \(0<d<n\), \(k=n-d\), and \(P(C)=0\) with
\(|C|^2=q^d\). The roots \(C,Cq,\ldots,Cq^{k-1}\) occur in
a single block of the spherical unitary classification.
In particular the monic polynomial
\[
             P^\natural(X)=C^{-n}P(CX)
                =\prod_{h=0}^{k-1}(X-q^h)\, B(X)
\]
has a monic quotient \(B\).
\end{lemma}
\begin{proof}
By Tadi\'c's theorem, \(\pi\) is a normalized parabolic
induction of characters \(\chi_j\circ\det\) on blocks of
dimensions \(d_j\) with \(\sum d_j=n\). Write
\(\chi_j(\varpi)=\omega_j q^{\epsilon_j}\), \(|\omega_j|=1\).
Here \(\epsilon_j=0\) for a unitary character block, and the
other blocks come in pairs with shifts \(\epsilon_j=\pm a\),
\(0<a<1/2\). The normalized Satake parameters of one determinant
character on a block are its value on \(\varpi\) times the
successive powers of \(q\) from \(-(d_j-1)/2\) to \((d_j-1)/2\).
This follows from the Borel modulus in normalized induction
of a character, or by applying the Satake isomorphism to the
one-dimensional representation. Normalized parabolic induction
takes the union of these multisets. Thus the roots of \(P\) on
this block are
\[
        \chi_j(\varpi)\,q^{(n-d_j)/2+h},\qquad 0\le h<d_j .
\]
These are the parameters of the spherical unitary classification
in the normalization fixed above.~\cite[Theorem 3.2]{IX-VII-VII-tadicsph716}

The base-\(q\) logarithm of the absolute value of
\(C/q^{(n-1)/2}\) is \(-(k-1)/2\), a half-integer. It cannot
occur in a block with \(0<|\epsilon_j|<1/2\), because every
unshifted exponent is a half-integer. In the block in which
it occurs, the exponents are consequently centered at zero
and all roots have the same complex phase. The block contains
both \(-(k-1)/2\) and \((k-1)/2\) and all unit steps between
them. This proves the assertion, including the phases.
\end{proof}

Let \(I\subset K\) be the inverse image of the upper triangular
Borel subgroup and \(I_1\subset I\) the inverse image of its
unipotent radical. Put \(a=\operatorname{diag}(\varpi,1,\ldots,1)\).
Write \(U\) for the unnormalized double coset operator of \(a\)
at either of these levels, on the corresponding fixed vectors.

\begin{proposition}[The roots retained by the Iwahori operator]
\leavevmode\label{prop:toproots716}
On \(\pi^I\) the operator \(U\) is killed by
\[
       \prod_j(X-\beta_j),\qquad
       \beta_j=\chi_j(\varpi)q^{(n+d_j-2)/2}.
\]
There is one factor for each inducing block, with repetitions.
Under the hypotheses of Lemma~\ref{lem:sphericalchain716}, if
\(D(X)\) is any multiple of \(B(X)\) over \(\CC\), then for any
spherical vector \(x\) and any \(y\in\pi^{I_1}\) one has
\[
       \big\langle x,(C^{-1}U-q^{k-1})D(C^{-1}U)y\big\rangle=0 .
\]
\end{proposition}
\begin{proof}
Use the standard parabolic of type \((1,n-1)\), with Levi
\(M\) and unipotent radical \(N\), the additive group of
upper first rows of length \(n-1\). The unnormalized Jacquet
module \(\pi_N\) is the quotient by the span of
\(\pi(v)z-z\) for \(v\in N\). The map \(\pi^I\to\pi_N\)
is injective. To verify this, a vector in its kernel is a
finite sum of such differences. Its Haar average over a
compact subgroup \(\varpi^{-h}N(\mathcal O_L)\) of rows is
zero for all sufficiently large \(h\), since this subgroup
contains each of the finitely many \(v\)'s in the sum.
On an \(I\)-fixed vector \(z\),
\[
    U^h z=\sum_{b\in(\mathcal O_L/\varpi^h)^{n-1}}
                   \pi(n(b)a^h)z .
\]
For \(h=1\), the representatives follow from reduction of the
upper first row in \(I/(I\cap aIa^{-1})\). Equivalently, use
the first-row unipotent factor and the lower-right invertible
block to clear that row modulo \(\varpi\). Iteration proves the
formula, since conjugation by \(a\) multiplies rows by
\(\varpi\). The sum equals \(q^{h(n-1)}\pi(a^h)\) times the
preceding average of \(z\). The Iwahori double coset of \(a\)
is invertible on \(\pi^I\). In the Iwahori--Matsumoto
presentation its characteristic function factors, along a
reduced expression, into simple generators and a length-zero
unit. Each simple generator satisfies
\(T_s^2=(q-1)T_s+q\) and is invertible. Thus \(U^h z=0\)
implies \(z=0\), proving injectivity.~\cite{IX-VII-VII-im716}

Under projection to the unnormalized Jacquet module, \(U\)
acts as \(q^{n-1}\pi_N(a)\) because the row factors act
trivially. On the normalized Jacquet module this operator
is \(q^{(n-1)/2}r_P(\pi)(a)\), since the parabolic modulus
at \(a\) is \(q^{-(n-1)}\). The geometric lemma gives a
filtration of \(r_P(\pi)\) with one subquotient for each
inducing block. Indeed, the relevant Weyl double cosets
record from which block the first coordinate comes.
For the block \(j\) that coordinate receives the normalized
Jacquet character of \(\chi_j\circ\det\) on the split-off
line, namely \(\chi_j(t)|t|^{-(d_j-1)/2}\). All other
coordinates lie in the \(\mathrm{GL}_{n-1}\) factor and
are evaluated at the identity. Hence the central Levi
element \(q^{(n-1)/2}r_P(\pi)(a)\) acts on that subquotient
by \(\beta_j\). Successively multiplying its differences
from these scalars kills the whole filtered module, including
when two scalars agree. Injectivity proves the first
assertion.~\cite[\S2.12]{IX-VII-VII-bzgeom716}

The chosen block of Lemma~\ref{lem:sphericalchain716} contains
every root \(Cq^h\), \(0\le h<k\). Delete from that block a
single occurrence of each root with \(0\le h<k-1\). Its top
root survives because \(Cq^{k-1}\) belongs to it. Every
other block retains its own top root, with multiplicities
counted by blocks. Dividing the remaining roots by \(C\)
shows that \(\prod_j(X-\beta_j/C)\) divides
\((X-q^{k-1})B(X)\). The asserted operator therefore kills
\(\pi^I\). The quotient \(I/I_1\) has diagonal unit
representatives, which normalize \(I_1\) and commute with
\(a\). Its averaging projection on \(\pi^{I_1}\) commutes
with \(U\). On its image the two versions of \(U\) agree,
using the same first-row representatives at both levels.
Since \(x\) is \(I\)-fixed, averaging \(y\) preserves its
pairing with \(x\) after applying this operator. The
orthogonality follows.
\end{proof}

\begin{theorem}[A common integral polynomial at each auxiliary prime]
\leavevmode\label{thm:commonhecke716}
Use the primes of Corollary~\ref{cor:markeddet715} for a main
summand \(T=\theta R\) of dimension \(d\) on the shorter
family, with \(k=n-d\ge2\) in the permitted weight range.
Use geometric Frobenius \(\phi_w\), a fixed coefficient prime
\(\ell\), and an identification of the complex and algebraically
closed \(\ell\)-adic coefficient fields. At each chosen place
let \(C_w\) be the prescribed normalization with
\(|C_w|^2=q^d\) at the complex place and \(C_w\) an
\(\ell\)-adic unit, \(q=Nw\). For any finite family of unitary spherical Hecke systems
there, retain those whose polynomials \(P_x\) satisfy
\[
 \begin{gathered}
 P_x(C_w)=0,\qquad P_x^\natural(X)=C_w^{-n}P_x(C_wX)
                       \text{ has integral coefficients},\\
 \overline{P_x^\natural}(X)=
       \det(XI-\overline{T}(\phi_w))
                     \prod_{h=0}^{k-1}(X-\overline q^{\,h}).
 \end{gathered}
\]
The bar means reduction at the maximal ideal, in a single
residue algebraic closure. There is one \(\ell\)-integral
polynomial \(D_w(X)\) for all these systems with \(D_w(1)=1\)
such that
\[
       (C_w^{-1}U_w-q^{k-1})D_w(C_w^{-1}U_w)
\]
has zero pairing with every one of their spherical vectors,
on every \(I_1\)-fixed input. This is an integral operator
on forms with the usual \(\ell\)-integral lattice away from
\(\ell\).
\end{theorem}
\begin{proof}
For each system let \(B_x\) be the monic quotient of
Lemma~\ref{lem:sphericalchain716}. Division by a monic
polynomial over the valuation ring preserves integrality
of coefficients. Cancel the common monic factor from the
residual polynomial identity in the residue polynomial ring.
This gives
\(\overline B_x(X)=\det(XI-\overline T(\phi_w))\), even when
\(q\equiv1\bmod\ell^\nu\) and all the cyclotomic roots coincide
in the residue field. The determinant condition on the selected
prime now makes \(B_x(1)\) a unit. In a common coefficient
field take
\[
                 D_w(X)=\prod_x\frac{B_x(X)}{B_x(1)}
\]
over precisely the indicated systems. An empty product is
\(1\). This polynomial is integral, takes value \(1\) at \(1\),
and is divisible by each quotient over the coefficient field.
The preceding proposition gives orthogonality in the complex
model. Finally \(U_w\) is a sum of right translations away
from \(\ell\). Each term preserves integrality of the values
of a form in its coefficient lattice, and \(C_w^{-1},q\)
and the coefficients of \(D_w\) are integral. This proves the
lattice assertion.
\end{proof}

For clarity, the local hypotheses can be checked using a
global representation attached to a Hecke system whenever
the following compatibility is available. Suppose the
normalized representation \(\rho_x\) has eigenvalue \(1\)
on every Galois element (the weak-lift condition). At a
good split place assume \(P_x^\natural\) annihilates
\(\rho_x(\phi_w)\) and agrees modulo the maximal ideal with
its characteristic polynomial. Then \(P_x^\natural(1)=0\),
by evaluation on a fixed eigenvector, and its coefficients
are integral if those of the characteristic polynomial are.
If the latter polynomials agree residually, on all Galois
elements, with those of
\(T\oplus\bigoplus_{0\le h<k}\kappa^{-h}\), evaluation at
geometric Frobenius gives the required residual identity
since \(\kappa^{-1}(\phi_w)=q\). An exact trace identity
at a characteristic-zero reference system gives this same
factorization there by applying Newton's identities to the
traces of the powers of each element. For finitely many
systems any finite exceptional sets for these compatibilities
can be omitted in the choice of the auxiliary primes.

This calculation supplies a polynomial at the very places
where the even cocycles are tested. It treats the weak
systems satisfying the indicated compatibility and residual
congruence. Passing from this local annihilation to the
denominator of the global orthogonal projection requires
the integral congruence relations among the different systems.
\clearpage


\subsection{A Fitting determinant for the scalar auxiliary relations}
\label{sec:scalarfitting718}

We pass from equalities of cocycle values to an integral
coefficient of a determinant. Write \(\mathcal O\) for a
discrete valuation ring with uniformizer \(\pi\) and valuation
\(v(\pi)=1\). Fix \(0\ne a\in\pi\mathcal O\), set
\(A=\mathcal O/a\mathcal O\), and use \(M^\vee=\operatorname{Hom}
_{\mathcal O}(M,A)\) for modules of finite length killed by \(a\).
This duality is exact and preserves lengths. Indeed \(A\) embeds
in \(\operatorname{Frac}(\mathcal O)/\mathcal O\) by division by
\(a\), with image its \(a\)-torsion. The fraction-field quotient
is divisible, hence injective over the principal ideal domain
\(\mathcal O\). Equivalently, a map from an ideal \((d)\) extends
to \(\mathcal O\) by dividing the image of \(d\) by \(d\).
The elementary divisor decomposition then identifies \(M\) with
its double dual via evaluation.

Let \(J\) be a finite ordered set, \(I\subseteq J\), and
\(B=\mathcal O[X_j:j\in J]\). For a module \(P=B^s/N\), write
\(\operatorname{Fitt}(P)\) for its zeroth Fitting ideal, generated
by the determinants of \(s\) rows in \(N\). Give \(P\) a family
of elements \(y(g)\), indexed by a set \(G\). Assume their
\(\mathcal O\)-span \(M\) in \(P\) has finite length and is killed
by \(a\). Assume each of the \(s\) standard generators of \(P\)
belongs to the family. Write \(Z\) for the functions on \(G\)
of the form \(g\mapsto\lambda(y(g))\) with \(\lambda\in M^\vee\).
Since the \(y(g)\) span \(M\), this identifies \(Z\) with \(M^\vee\).

\begin{theorem}[The determinant coefficient from paired relations]
\leavevmode\label{thm:fittingcost718}
Suppose that \(G\) contains marked elements \(\tau_j,\tau'_j,
\phi_j,\phi'_j\) and in \(P\) one has
\[
       y(\tau_j)=X_jy(\phi_j),\qquad
       y(\tau'_j)=X_jy(\phi'_j).
\]
Assume \(X_i P=0\) whenever \(i<j\in I\) and \(i\notin I\).
Choose \(L\subseteq G\), put \(Z_L=\{z\in Z:z|_L=0\}\), and
let \(C\in\mathcal O\) be nonzero. Suppose the following two
bounds hold. For \(j\in I\), a function \(z\in Z_L\) that agrees
on all earlier \(\tau_i,\tau'_i\) and on earlier
\(\phi_i,\phi'_i\) for \(i\in I\) satisfies
\(C(z(\tau_j)-z(\tau'_j))=0\). For any submodule \(S\) of
\(Z_L\) whose functions vanish at all \(\tau_j,\tau'_j\) and
agree at \(\phi_j,\phi'_j\) for \(j\in I\), suppose that
\[
                  \len(Z/Z_L)+\len(CS)\le b .
\]
Then for some \(F\in B\) one has
\[
       F\prod_{j\in I}X_j\in\operatorname{Fitt}(P),\qquad
       v(F(0))\le b+(|I|+s)v(C).
\]
\end{theorem}
\begin{proof}
We first record the cost of inserting a single new relation.
For \(w\in P\) killed by \(d\in B\), one has
\[
        d\,\operatorname{Fitt}(P/Bw)\subseteq\operatorname{Fitt}(P).
\]
Choose a lift of \(w\) to \(B^s\). Expand a determinant of rows
in \(N+B w\) by multilinearity, using this lift. Terms using it
twice vanish by alternation. In a term using it once, multiply
that row by \(d\), which puts it in \(N\). Terms using it zero
times already belong to the desired ideal. The same conclusion
holds for \(s=0\), by the empty determinant convention.

Let \(U\subset M\) be the \(\mathcal O\)-span of \(y(L)\), of length
\(h\). Exactness of duality gives \(h=\len(Z/Z_L)\). Take a
composition series of \(U\) and a lift of a generator of each
successive factor. When their \(B\)-spans are killed in this
order, each added vector is killed by \(\pi\) in the preceding
quotient. Thus passage from \(P\) to \(P/BU\) costs the factor
\(\pi^h\) in the displayed Fitting inclusion.

Now kill \(v_j=y(\phi_j)-y(\phi'_j)\) in increasing order of
\(j\in I\). In the quotient just before \(j\), write \(M'\) for
the image of \(M\). Each member of \((M')^\vee\) gives a
function in \(Z_L\) agreeing on the earlier tested Frobenius
pairs. On the earlier inertia pairs it also agrees. Indeed
\(X_i v_i=0\) for earlier \(i\in I\), since \(v_i\) was killed,
and \(X_iP=0\) at every other earlier index. The first hypothesis
therefore says that every such functional kills
\(C(y(\tau_j)-y(\tau'_j))\). This vector is in \(M'\).
Finite duality detects it, hence it is zero in the quotient.
By the two relations this says \(CX_jv_j=0\).
Inserting \(v_j\) consequently costs the single factor \(CX_j\).
Denote the resulting quotient by \(P'\).

Put \(P_0=P'/(X_j:j\in J)P'\). It is generated over
\(\mathcal O\) by the images of the original \(s\) generators,
so \(M\to P_0\) is surjective. In particular \(P_0\) is a
finite length module killed by \(a\). Its dual identifies with
a submodule \(S\) of \(Z_L\). Setting all variables to zero
kills both tame values for every index, and killing \(v_j\)
imposes the two Frobenius values' equality for \(j\in I\).
Hence the second hypothesis applies to this \(S\). The
elementary divisors of \(P_0\) and \(S\) agree. It follows that
\[
            \len P_0=\len(CS)+\len(P_0/CP_0)
                      \le b-h+s\,v(C).
\]
The last bound uses the surjection from
\((\mathcal O/C\mathcal O)^s\) onto \(P_0/CP_0\).
Over \(\mathcal O\), Smith normal form gives
\(\operatorname{Fitt}(P_0)=(\pi^{\len P_0})\).
Specialization at zero sends \(\operatorname{Fitt}(P')\)
onto this ideal, since it specializes the rows and their
determinants and every relation in the specialized presentation
comes from a relation before specialization. Lift the indicated
generator to \(D\in\operatorname{Fitt}(P')\).
Multiplying back the factors introduced above and taking
\(F=\pi^h C^{|I|}D\) proves both assertions.
\end{proof}

Here is the application to actual finite scalar cohomology.
Choose one node and the primes of
Theorem~\ref{thm:trianglebudget714}, once for all \(0\le e<t\).
Put \(a=\ell^\nu\), \(C=\ell^{b_\ell}\), and write \(d\) for
the size of the set consisting of \(S_\ell,\gamma\), the one
extra coset representative if needed at two, and all the
\(\tau_j,\tau'_j\). Repetitions are counted once. Choose
\(Z_e\) to be the module of all \(A(\kappa^{-e})\)-valued
continuous cocycles unramified outside \(\ell,w_j,\bar w_j\)
that vanish on \(L_e\) and at \(\gamma\). Write \(y_e(g)\)
for evaluation at \(g\) in \(M_e=Z_e^\vee\). Set
\[
 P_e=(B\otimes_{\mathcal O}M_e)\Big/
   \Big\langle y_e(\tau_j)-X_jy_e(\phi_j),\
               y_e(\tau'_j)-X_jy_e(\phi'_j):j\in J\Big\rangle_B .
\]
We identify \(y_e(g)\) here with \(1\otimes y_e(g)\).
For \(m=\prod_{j\in J}X_j\), write
\(\mathfrak n=(X_j^2:j\in J)\).

\begin{theorem}[One coefficient for all scalar weights]
\leavevmode\label{thm:allscalarfit718}
The ideal \(\operatorname{Fitt}(\bigoplus_{e=0}^{t-1}P_e)\)
contains a polynomial congruent to \(\alpha m\) modulo
\(\mathfrak n\), where \(0\ne\alpha\in\mathcal O\) satisfies
\[
 \frac{v(\alpha)}{v(\ell)}
       \le v_\ell(h_E\mathcal J_{N,t})+(2t+1)c_\ell
        \qquad\bigl(E\subseteq\QQ(\mu_N)\bigr).
\]
The bound holds on the period intervals of
Corollary~\ref{cor:oddwindow677} at every precision. On the
shorter family the auxiliary contribution is at most
\((2t+1)Y^6v_\ell(G_0)\) with \(Y,G_0\) as in
Corollary~\ref{cor:actualmarks714}.
\end{theorem}
\begin{proof}
The \(d\) evaluations just chosen generate \(M_e\) over
\(\mathcal O\). Indeed a cocycle vanishing on them is unramified
outside \(\ell\) by the tame-inertia argument in
Theorem~\ref{thm:scalarcost717}, and then is zero by
Lemma~\ref{lem:testpoints714}. The annihilator of their span
in \(M_e^\vee=Z_e\) is zero, so the span is all of \(M_e\).
Thus \(P_e\) has a presentation with \(d\) generators of the
kind used in the preceding theorem.

Set \(m_e=\prod_{j\in I_e}X_j\), \(m_{<e}=\prod_{i<e}m_i\)
and \(\mathfrak q_e=(\mathfrak n:m_{<e})\).
For a variable of earlier weight \(\mathfrak q_e\) contains
\(X_j\), since \(X_j m_{<e}\) has a square factor. Apply
Theorem~\ref{thm:fittingcost718} to \(P_e/\mathfrak q_eP_e\)
and \(I=I_e\). Its functions are a submodule of \(Z_e\) by
finite duality on the coefficient image, so they already
vanish on \(L_e\) and satisfy the normalization at \(\gamma\).
The proof of Theorem~\ref{thm:scalarcost717} gives both
hypotheses with \(b=\len O_e\). In particular its bound on
\(CS\) only uses \(S\subseteq Z_L\) and the normalization.
We obtain an \(F_e\) of constant-term valuation at most
\(\len O_e+(|I_e|+d)v(C)\) and
\[
           F_e m_e=d_e+\zeta_e,\qquad
           d_e\in\operatorname{Fitt}(P_e),\quad
           \zeta_e\in\mathfrak q_e .
\]
Indeed, after reduction modulo \(\mathfrak q_e\) both presentations
give the same module over \(B/\mathfrak q_e\), so their
maximal minors generate identical ideals there. This proves the
claimed lifting of \(d_e\).

Multiply these equations in increasing order of \(e\) modulo
\(\mathfrak n\). If the preceding product of the \(d_i\) equals
\(m_{<e}\prod_{i<e}F_i\) in this quotient, it kills
\(\zeta_e\) by the definition of \(\mathfrak q_e\).
Induction gives \(\prod d_e\equiv m\prod F_e\).
Since every variable kills \(m\) in the quotient, the right
side equals \(m\prod F_e(0)\). The product of the Fitting
ideals belongs to the Fitting ideal of the direct sum, by
taking block diagonal determinants. Set \(\alpha=\prod F_e(0)\).

Finally \(|J|=(t-1)|S_\ell|\),
\(d\le 2+|S_\ell|+2|J|\), and \(v(C)=b_\ell v(\ell)\).
In particular \((|J|+td)b_\ell\le(2t+1)c_\ell\) by the
definition of \(c_\ell\). The weight sum of
\(\len O_e/v(\ell)\) is at most \(v_\ell(h_E\mathcal J_{N,t})\)
by Theorem~\ref{thm:oddinteger677}. This proves the valuation
bound and its stated uniformity.
\end{proof}

These modules use finite cocycles on the fields on which the
periods are to be taken. To obtain a determinant in the Hecke
trace algebra one must also realize the scalar determinants
there together with the determinant on the higher-rank block
and control the cost of its local conditions.
\clearpage


\subsection{Integral relation matrices on the main and character summands}
\label{sec:relationmatrices719}

Use the auxiliary primes and both elements at each pair from
Theorem~\ref{thm:scalarcost717}. Fix a finite family of continuous
representations \(u_x:G_E\to\mathrm{GL}_n(\mathcal O)\).
In addition to \(\ell\) and the auxiliary pairs, allow ramification
at a finite list of old places away from \(\ell\). Include every
old place in the list, including the conjugates that occur.
For each, choose a lift \(\delta_w\) in inertia of a generator of
its pro-\(\ell\) tame quotient. Put \(a=\ell^\nu\), and keep
\(B,J,C=\ell^{b_\ell}\) from Section~\ref{sec:scalarfitting718}.
Write \(G=G_E\), \(\Lambda=B[G]\), and let \(I\) be the augmentation
ideal of \(\Lambda\). Fix an inertia element \(\eta\) away from
\(\ell\). Suppose a rank \(r\) representation \(T\) satisfies
\[
       T(\eta)=y\,1_r,\qquad y-1\in\mathcal O^\times,\qquad
       \operatorname{span}_{\mathcal O}T(G)=M_r(\mathcal O).
\]
These properties for a main summand of the marked type follow
from Theorem~\ref{thm:markedmatrix715}. Put
\[
       c_1=(y-1)^{-1},\qquad H=c_1(\eta-1),\qquad
       R_0=H+(1-H)(\gamma-1),
\]
where \(\gamma\) is the cyclotomic element of
Lemma~\ref{lem:testpoints714}. Fix \(\epsilon_j\in\mathcal O[Z]\)
with \(\epsilon_j(1)\equiv1\bmod a\). One may use the polynomials
of Theorem~\ref{thm:commonhecke716} for its compatible local data.

Denote by \(K\) the common kernel of the representations
\(u_x:\Lambda\to M_n(B)\), extended linearly. Define a left ideal
\(\mathcal D\) of \(\Lambda\) by adding \(K,aI,I R_0\) and the
left ideals generated by
\[
         b-1-X_j\,\epsilon_j(f)(f-1),
      \qquad (b,f)=(\tau_j,\phi_j),(\tau'_j,\phi'_j).
\]
The finite vector of matrices \(D(g)=(u_x(g)-1)_x\) spans a finite
free \(\mathcal O\)-module as \(g\) varies. Choose a list \(g_i\)
in three groups of indices \(W,H_0,D_0\). In \(W\) put all
\(\delta_w\). In \(D_0\) put the \(d\) tests of
Theorem~\ref{thm:allscalarfit718}. Complete these by a finite
list \(H_0\) so that the \(D(g_i)\) generate the entire span.
Such a finite sublist exists by the Noetherian property.
Index groups are kept disjoint even if two group elements agree.
Choose an \(\mathcal O\)-linear section of
\((s_i)\mapsto\sum_i s_iD(g_i)\) onto that span, and denote
the coordinates of \(D(g)\) by the row \(v(g)\).
The coordinates are continuous in \(g\). Indeed a finite free
submodule of a finite matrix lattice has its lattice topology
as subspace topology, by the elementary divisor theorem.
Define the module of rows
\[
       \mathcal R=\left\{p\in\Lambda^{k}:
                   \sum_i p_i(g_i-1)\in\mathcal D\right\},
       \qquad k=|W|+|H_0|+d .
\]
It is closed under addition and multiplication by group-algebra
elements on the left.

\begin{lemma}[Relations read by a representation]
\leavevmode\label{lem:readrows719}
The following rows belong to \(\mathcal R\), where \(e_i\) is
the \(i\)-th standard row and all scalar coordinates of \(v\)
are included in \(\Lambda\).
\[
\begin{gathered}
   v(g_i)-e_i,\qquad v(gh)-v(g)-g v(h),\qquad a e_i,\\
   v(b)-X_j\epsilon_j(f)v(f),\\
   (g-1)\bigl(c_1v(\eta)+(1-H)v(\gamma)\bigr).
\end{gathered}
\]
For the middle row use either pair \((b,f)\) above.
An integral relation among the \(D(g)\) is also a relation
among the coordinate rows \(v(g)\).
\end{lemma}
\begin{proof}
The value \(\sum_i v(g)_i(g_i-1)\) differs from \(g-1\) by an
element of \(K\), by the definition of the section. For the
first line use the identity \(gh-1=(g-1)+g(h-1)\) and the
inclusion \(a(g_i-1)\in aI\). The second line reads the
displayed generator of \(\mathcal D\) modulo \(K\).
The parenthesis in the last line reads \(R_0\) modulo \(K\),
and \((g-1)R_0\in I R_0\). The final assertion follows by
linearity of the section.
\end{proof}

For \(F=\operatorname{Frac}\mathcal O\), let \(\mathcal V_T\)
be the continuous cocycles \(G\to(F/\mathcal O)^r\) with action
\(T\) that vanish at \(\eta\) and all \(\tau_j,\tau'_j\), and
respect every integral linear relation among the \(D(g)\).
Write \(\lambda_T\) for its length, allowing infinity.

\begin{proposition}[The determinant on the main summand]
\leavevmode\label{prop:blockrows719}
There is a \(k\)-square matrix \(B_0\) of rows in \(\mathcal R\)
whose \(kr\)-square matrix of blocks, read under \(T\) at
\(X_j=0\), has nonzero determinant of valuation at most
\(\lambda_T\).
\end{proposition}
\begin{proof}
Write \(L\subset\mathcal O^{kr}\) for the span of all individual
component rows obtained by reading \(\mathcal R\) under \(T\)
at zero. Since \(a e_i\in\mathcal R\), the module \(L\) contains
\(a\mathcal O^{kr}\). Consider a column \(w=(w_i)_i\), each
\(w_i\in(F/\mathcal O)^r\), annihilated by \(L\).
The function \(z(g)=\sum_i v(g)_i w_i\) is continuous, because
the column is killed by \(a\) and the coordinates are continuous.
Lemma~\ref{lem:readrows719} gives \(z(g_i)=w_i\), the cocycle
identity for \(T\), all the linear relations, and \(z(b)=0\).
In the last relation take \(g=\eta\). Since \(T(H)=1_r\) and
\((y-1)c_1=1\), it gives \(z(\eta)=0\). Thus \(w\mapsto z\)
embeds the annihilator of \(L\) in \(\mathcal V_T\).
Finite duality identifies the length of this annihilator with
\(\len(\mathcal O^{kr}/L)\).

Choose a basis of \(L\) giving determinant valuation equal to
that length, by Smith normal form. Each set of \(r\) successive
basis rows can be realized as the component rows of a single
member of \(\mathcal R\). In fact multiplication of a row on
the left by a lift of a matrix unit \(E_{ij}\) copies its
\(j\)-th component into position \(i\) and kills the others.
The full-span assumption gives every such lift in
\(\mathcal O[G]\). Adding the resulting rows realizes arbitrary
\(\mathcal O\)-linear combinations in each position independently.
Taking the \(k\) sets of \(r\) rows gives \(B_0\).
\end{proof}

For \(0\le e<t\), put \(\chi_e=\kappa^{-e}\). Let
\(\mathcal Z_e\) be the module of all continuous cocycles
\(G\to(\mathcal O/a)(\chi_e)\) respecting every integral linear
relation among the \(D(g)\). Write \(l_e\) for the colength
in \(\mathcal Z_e\) of the subgroup satisfying the local
crystalline relations above \(\ell\). Define \(O_e\) as
in Theorem~\ref{thm:scalarcost717}, and, for an old place,
suppose a polynomial \(o_w\in\mathcal O[Z]\) satisfies
\[
        (u_x(\delta_w)-1)o_w(u_x(\delta_w))=0
                      \quad\text{for all }x,\qquad o_w(1)\ne0.
\]
Write \(k_0=C^d\prod_w o_w(1)\) and retain the ideal
\(\mathfrak n=(X_j^2:j\in J)\) and monomial
\(m=\prod_j X_j\) of Section~\ref{sec:scalarfitting718}.

\begin{theorem}[Simultaneous integral matrices on the characters]
\leavevmode\label{thm:charrows719}
There are finitely many \(k\)-square matrices \(B_i\) with rows
in \(\mathcal R\) and weights \(h_i\in B\) such that
\[
       \sum_i h_i\prod_{e=0}^{t-1}\det\chi_e(B_i)
           \equiv k_0^t\,\alpha m\pmod{\mathfrak n},
       \qquad
       v(\alpha)\le
          \sum_e(l_e+\len O_e)+(|J|+td)v(C).
\]
Here \(\alpha\) is a nonzero element of \(\mathcal O\).
In particular the valuation loss displayed in this statement
is independent of the number of supplementary generators
\(|H_0|\) when the other data are fixed.
\end{theorem}
\begin{proof}
Fix a character \(\chi=\chi_e\). Read all coefficients of rows
under \(\chi:\Lambda\to B\) and discard their \(W\)-coordinates.
Then apply a \(B\)-linear map to the remaining coordinates and
take the quotient by the image of \(\mathcal R\). One may also
impose further \(B\)-linear relations. Write \(Y(g)\) for the
image of \(v(g)\) in such a quotient. The \(\mathcal O\)-span
of the \(Y(g)\) is finite and killed by \(a\), since it is in
the span over \(\mathcal O\) of the images of the \(k\) standard
rows, each killed by \(a\). An \(\mathcal O\)-linear functional
on this finite span with values in \(\mathcal O/a\) gives a
function \(z(g)\) by evaluation on \(Y(g)\). It is continuous
by the continuity of \(v(g)\) and finite duality, or by extending
the functional to the span of the standard rows by that duality.

Lemma~\ref{lem:readrows719} makes \(z\) a member of \(\mathcal Z_e\).
It vanishes at the old generators since their standard rows
were discarded. At a place outside the listed ramification
\(D(g)=0\) on inertia. At an old place the twist is trivial
on inertia, and an \(\ell\)-primary homomorphism out of inertia
is detected on its tame generator. Thus \(z\) is unramified
outside \(\ell\) and the auxiliary pairs. Moreover \(\chi(H)=0\)
since \(\eta\) lies in inertia away from \(\ell\). Reading the
last relation of the lemma gives
\[
           (\chi(g)-1)\bigl(c_1z(\eta)+z(\gamma)\bigr)=0
                        \quad(g\in G).
\]
In particular every such \(z\) unramified outside \(\ell\)
satisfies \((u^{-e}-1)z(\gamma)=0\). The two auxiliary relations
give \(Y(b)=X_jY(f)\), since \(\chi(f)\equiv1\bmod a\) and
\(\epsilon_j(1)\equiv1\bmod a\).

First project just to the \(H_0\)-coordinates. The above
functions vanish at every element of \(D_0\), again by
\(v(g_i)-e_i\in\mathcal R\). Among these tests are both tame
generators at all auxiliary pairs. The functions are consequently
unramified outside \(\ell\), and the remaining tests detect them
by Lemma~\ref{lem:testpoints714}. Finite duality makes every
\(Y(g)\) zero. In particular the images of all the standard
\(H_0\)-rows vanish. This gives a matrix of rows \(V_e\) in
\(\mathcal R\) whose \(H_0\)-square part is the identity after
evaluation at \(\chi_e\).

One matrix \(V\) works for all characters. To combine two
matrices of the given sizes, let \(V_{H_0}\) denote the square
part of the first matrix before evaluation and replace it by
the rows of \(V+(1-V_{H_0})V'\). The defect of the square part
from the identity is the product
\((1-V_{H_0})(1-V'_{H_0})\), in that order. It vanishes under
any character for which either defect vanishes. Starting with
one of the \(V_e\) and repeating this operation gives the common
rows. This uses no division, even when two characters agree
modulo the maximal ideal.

For the given \(e\), discard \(W\) and send a row with remaining
parts \(b_H,b_D\), already read on \(\chi_e\), to
\[
                   b_D-b_H\,\chi_e(V_{D_0})\in B^d .
\]
Write \(N_e\) for the image of \(\mathcal R\). Add the relations
\(\mathfrak q_eB^d\), where \(\mathfrak q_e\) is the colon ideal
in the proof of Theorem~\ref{thm:allscalarfit718}.
In the resulting quotient the \(Y(g_i)\) for \(i\in D_0\)
are the \(d\) standard generators. Theorem~\ref{thm:scalarcost717}
applies to its functions with length cost at most \(l_e+\len O_e\).
Indeed their quotient by the local crystalline equations embeds
in the corresponding quotient of \(\mathcal Z_e\), and their
members unramified outside \(\ell\) satisfy the normalization
proved above. Theorem~\ref{thm:fittingcost718} now gives
\[
  d_e\equiv F_e\prod_{j\in I_e}X_j\pmod{\mathfrak q_e},\qquad
  d_e\in\operatorname{Fitt}(B^d/N_e),\qquad
  v(F_e(0))\le l_e+\len O_e+(|I_e|+d)v(C).
\]
The lifting from the quotient by \(\mathfrak q_e\) is the base
change of maximal minors used in Theorem~\ref{thm:allscalarfit718}.

Write \(d_e=\sum_i h_{ei}\det A_{ei}\), with \(h_{ei}\in B\)
and every row of \(A_{ei}\) in \(N_e\). Each such row lifts to
a row in \(\mathcal R\) whose \(\chi_e\)-evaluation vanishes on
\(H_0\) and equals the chosen row on \(D_0\). In fact a lift
before elimination has parts \(b_H,b_D\). Subtract \(b_H V\)
using the central coefficients \(b_H\in B^{|H_0|}\). This
gives precisely the asserted evaluation.

For interpolation put \(s_e=\chi_e(\gamma)=u^{-e}\) and set
\[
        Q_e=\frac{C}{\prod_{f\ne e}(s_e-s_f)}
                  \prod_{f\ne e}(\gamma-s_f)\in\mathcal O[G].
\]
The denominators divide \(C\) by
Theorem~\ref{thm:trianglebudget714}, since replacing \(u\)
by \(u^{-1}\) changes each difference by a unit. The character
\(\chi_f\) reads \(Q_e\) as \(C\) if \(f=e\), and zero otherwise.
Choose an index \(i_e\) independently for every \(e\). Build
one \(k\)-square matrix by using \(o_w(\delta_w)e_w\) on \(W\),
the common rows \(V\) on \(H_0\), and, on \(D_0\), the sum
over \(e\) of the lifted rows of \(A_{e i_e}\) multiplied on
the left by \(Q_e\). The old rows belong to \(\mathcal R\) by
the imposed polynomial identity, and the other rows belong
to it by construction. At \(\chi_e\) the \(W\)-square part
has diagonal \(o_w(1)\) with zero entries to its right. In
the remaining square the upper left block is the identity,
the lower left block is zero and the lower right block is
\(C A_{e i_e}\). The determinant is \(k_0\det A_{e i_e}\).

Weight the matrix for this tuple by \(\prod_e h_{e i_e}\)
and sum over the tuples. Distributivity gives exactly
\(k_0^t\prod_e d_e\). As in Theorem~\ref{thm:allscalarfit718},
multiplication in increasing order of the weights gives
\(\prod_e d_e\equiv m\prod_e F_e(0)\pmod{\mathfrak n}\).
Set \(\alpha=\prod_e F_e(0)\) and add the valuation bounds.
\end{proof}
\clearpage


\subsection{A trace multiplier for the projected lifts}
\label{sec:tracemult720}

Keep the finite family of representations, auxiliary pairs and
integral row module of Section~\ref{sec:relationmatrices719}.
Let \(A=\prod_x F[X_j:j\in J]\), where \(F\) is the fraction
field of \(\mathcal O\). Let \(S\subseteq A\) be any \(B\)-subalgebra
containing the coefficient families of \(\det(Z-u_x(g))\) for
every \(g\in G\). Write \(U:S[G]\to M_n(A)\) for simultaneous
evaluation. Let \(\mathcal J\) be the \(S\)-submodule of \(A\)
generated by the following trace families, for \(g\in G\) and
each of the two pairs \((b,f)\) at \(j\).
\[
 a\,\operatorname{tr} U(g),\qquad
 \operatorname{tr}U\bigl((b-1-X_j\epsilon_j(f)(f-1))g\bigr).
\]
All these families will belong to \(S\). The construction also
works if \(\mathcal J\) is enlarged.

\begin{lemma}[Integral characteristic polynomials on blocks]
\leavevmode\label{lem:blockamitsur720}
If \(D\) is any square matrix over \(S[G]\), the characteristic
polynomial of the matrix with blocks \(U(D_{ij})\) has coefficients
in \(S\). The assertion remains valid after adjoining a central
indeterminate to \(S\) and \(A\).
\end{lemma}
\begin{proof}
Use one vertex for each block index. For every term \(b g\) in
the entry \(D_{ij}\) draw an edge from \(i\) to \(j\) with
scalar \(b\in S\) and group label \(g\). For an oriented closed
walk \(w\) let \(b_w\) and \(g_w\) denote the ordered products
of its scalars and labels. Identify walks differing by cyclic
rotation. The characteristic polynomial of \(U(g_w)\) is
unchanged by rotation because the resulting group elements
are conjugate. A closed walk is primitive if it is no
positive power of a shorter walk of smaller length.
For the flattened matrix \(M=(U(D_{ij}))\), in \(A[[z]]\) one has
\[
 \det(1-zM)=
   \prod_{[w]\ {\rm primitive}}
       \det\bigl(1-z^{|w|}b_w U(g_w)\bigr).
\]
Indeed \(\log\det(1-zM)=-\sum_{h\ge1}\operatorname{tr}(M^h)z^h/h\).
This identity follows by differentiating the determinant, using
the formal geometric series for \((1-zM)^{-1}\), and taking
constant term zero for the logarithm. Expanding the matrix
trace counts based closed walks. The \(j\)-fold power of a
primitive cyclic walk \(w\) contributes \(|w|\) based terms,
and its denominator is \(j|w|\). Its contribution is exactly
the \(j\)-th logarithmic term of the corresponding factor.
All calculations with logarithms take place over the
characteristic-zero algebra \(A\). For each power of \(z\)
only finitely many factors matter, and each factor already has
coefficients in \(S\). Equating coefficients proves the result
without dividing in \(S\). The same proof applies over \(S[T]\).
This is the cycle expansion in Amitsur's determinant formula.
~\cite{IX-VII-VII-amitsur720}
\end{proof}

Write \(y_i=g_i-1\), using the \(k\) group elements of
Section~\ref{sec:relationmatrices719}, and put
\[
        \mathfrak b=\{z\in S[G]:
            \operatorname{tr}U(zh)\in\mathcal J
                    \text{ for all }h\in S[G]\}.
\]
This is a two-sided ideal by linearity and cyclicity of trace.
It contains the kernel of \(U\).

\begin{proposition}[An adjugate for the row relations]
\leavevmode\label{prop:ribetadj720}
Let \(B'\) be a \(k\)-square matrix of rows in \(\mathcal R\).
There is \(P(T)\in S[T]\) with
\[
        P(0)=\det\bigl(U(B'_{ij})\bigr),\qquad
        (g-1)P(R_0)\in\mathfrak b\quad(g\in G).
\]
For every \(R_1\in S[G]\), put
\(f=\operatorname{tr}U(R_1P(R_0))\). Then \(f\in S\), and for
every finite integral combination \(w=\sum c_g(g-1)\),
\[
                     f\det U(w)\in\mathcal J .
\]
\end{proposition}
\begin{proof}
Modulo \(\mathfrak b\), the family kernel, \(a I\), and the
auxiliary generators defining \(\mathcal D\) are zero. For
the latter two facts it suffices by \(S\)-linearity to test
traces against a single group element, which gives exactly
the defining generators of \(\mathcal J\) and their linear
combinations. Also every \(g-1\) agrees with
\(\sum_i v(g)_iy_i\) modulo the family kernel. Since an
element of the augmentation ideal is a sum of scalar multiples
of \(g-1\), the remaining relations \(I R_0\) give, in
\(S[G]/\mathfrak b\), a matrix \(D\) over \(S[G]\) such that
\[
                     B' y=D y R_0 ,
              \qquad y=(y_1,\ldots,y_k)^{\mathsf t}.
\]
Here right multiplication by \(R_0\) is applied to every entry.

Put \(M(T)=B'-TD\) and let \(P(T)\) be the determinant of
its evaluation by \(U\) as a block matrix. By
Lemma~\ref{lem:blockamitsur720}, the full characteristic
polynomial, and hence this determinant, lies over \(S[T]\).
The adjugate of the evaluated matrix is a polynomial in
that matrix with coefficients in \(S[T]\), by Cayley--Hamilton.
Applying the same polynomial to \(M(T)\) gives a matrix \(D'(T)\)
over \(S[G][T]\). The entries of
\(D'(T)M(T)-P(T)1_k\) belong coefficientwise to the kernel of
\(U\), since the ordinary adjugate identity holds after evaluation.

In the quotient by \(\mathfrak b\), let the central variable \(T\)
act by right multiplication by \(R_0\), with coefficients acting
by left multiplication. These actions commute. With this action
\(M(T)y=0\), by the row equation. The adjugate identity gives
\(y_iP(R_0)=0\) in the quotient for every \(i\). As
\(g-1\) is a linear combination of the \(y_i\) there, this
proves the first assertion.

The trace and determinant of an element of \(S[G]\) belong to
\(S\) by the lemma for one block. Write \(d=\det U(w)\).
Again the adjugate polynomial gives an element \(v\in S[G]\)
with \(U(vw)=d\,1_n\). Since \(wP(R_0)\in\mathfrak b\), also
\(R_1vwP(R_0)\in\mathfrak b\). Taking its trace against \(1\)
gives \(df\in\mathcal J\).
\end{proof}

Let \(X_0\) be a nonempty subset of the parameter set such
that at each \(x\in X_0\) the family has a direct sum
\[
           u_x=T_x\oplus\bigoplus_{e=0}^{t-1}\chi_e,\qquad
                        \chi_e=\kappa^{-e}.
\]
An isomorphism of the characteristic-zero representations
suffices to read traces and determinants in this decomposition.
Each \(T_x\) is assumed to have an integral lattice and an
inertial element \(\eta_x\) away from \(\ell\) acting by a
scalar \(y_x\) with \(y_x-1\in\mathcal O^\times\). Select a
point \(x_*\in X_0\) and take \(T=T_{x_*}\), \(\eta=\eta_{x_*}\)
as in Section~\ref{sec:relationmatrices719}, including the
full matrix span at this point.

\begin{theorem}[One integral trace multiplier with bounded coefficient]
\leavevmode\label{thm:ribettrace720}
There is \(f\in S\) such that
\(f\det U(\sum c_g(g-1))\in\mathcal J\) for all finite
families \(c_g\in\mathcal O\). For every \(x\in X_0\),
\[
                     f_x\equiv\beta_x\prod_{j\in J}X_j
                             \pmod{(X_j^2:j\in J)}
\]
with \(\beta_x\in\mathcal O\). The coefficient at the selected
point is nonzero and satisfies
\[
 \begin{split}
 v(\beta_{x_*})\le {}&
   \lambda_T+\sum_{e=0}^{t-1}l_e+
                  v(\ell)v_\ell(h_E\mathcal J_{N,t})\\
 &+t\sum_w v(o_w(1))
      +(1+|J|+2td)v(C).
 \end{split}
\]
Here \(l_e,\lambda_T\) are the scalar crystalline colengths
and the length on the main summand defined in
Section~\ref{sec:relationmatrices719}. The final term divided
by \(v(\ell)\) is at most \((4t+3)c_\ell\) on the same
interval of fields.
\end{theorem}
\begin{proof}
Fix any order of the finite set \(X_0\) and form
\[
      K_0=\prod_{x\in X_0}
          \bigl(1-(y_x-1)^{-1}(\eta_x-1)\bigr)\in\mathcal O[G].
\]
It acts as zero under every \(T_x\) on this set, since one
factor acts by zero, and as one on every \(\chi_e\), because
\(\eta_x\) has cyclotomic value one. Use \(Q_0\) from the
interpolation in Theorem~\ref{thm:charrows719} and put
\(R_1=K_0Q_0\). In the decomposition at \(x\), the only
nonzero part of \(R_1\) is the scalar \(C\) on the trivial
character \(\chi_0\). The augmentation of \(R_0\) is zero.
Consequently \(u_x(R_1R_0)=0\) for \(x\in X_0\).

Take \(B_0\) from Proposition~\ref{prop:blockrows719} and
\(B_i,h_i\) from Theorem~\ref{thm:charrows719}.
The matrix \(M_i=(1-K_0)B_0+K_0B_i\), with multipliers
on the left in every entry, again has rows in \(\mathcal R\).
Apply Proposition~\ref{prop:ribetadj720} to \(M_i\), obtaining
\(P_i\) and \(f_i=\operatorname{tr}U(R_1P_i(R_0))\).
Then \(f=\sum_i h_i f_i\in S\) has the required property
against every \(w\). At \(x\in X_0\) only the constant
coefficient \(P_i(0)\) survives multiplication by \(R_1\),
so \((f_i)_x=C P_i(0)_x\).
Reordering block coordinates by the direct summands now gives
\[
       P_i(0)_x
        =\det\bigl(T_x(B_0)\bigr)
                     \prod_{e=0}^{t-1}\det\chi_e(B_i).
\]
The first determinant is that of the flattened matrix over
\(B\). With \(k_0=C^d\prod_w o_w(1)\) and \(\alpha\) of
Theorem~\ref{thm:charrows719}, this gives
\(\beta_x=C\,\det(T_x(B_0))(0)\,k_0^t\alpha\). Indeed every
variable kills \(\prod X_j\) modulo the squares. The formula
also shows that every proper squarefree coefficient is zero
at every \(x\in X_0\). At \(x_*\), the bound for the first
determinant is Proposition~\ref{prop:blockrows719}.
Adding the valuations of the other factors and using
Theorem~\ref{thm:oddinteger677} gives the displayed inequality.
Finally \(d\le 2+|S_\ell|+2|J|\), \(|J|=(t-1)|S_\ell|\), and
\(c_\ell=(2+t|S_\ell|)b_\ell\). Substitution gives the last
bound.
\end{proof}

\begin{corollary}[One augmentation determinant detects the complement]
\leavevmode\label{cor:augdet720}
There is a single finite integral combination \(w=\sum c_g(g-1)\)
such that \(\det u_x(w)\ne0\) for every \(x\) whose representation
over \(F\) has no trivial composition factor. On every
representation with a trivial composition factor
\(\det u_x(w)=0\).
\end{corollary}
\begin{proof}
Form the finite-dimensional \(F\)-algebra which is the image
of \(F[G]\) on the direct sum of all the \(u_x\) and the trivial
representation. Its quotient by its Jacobson radical is a product
of simple Artinian algebras. The augmented character on this
product is the projection onto its factor \(F\) for the trivial
simple module. Take the element that is zero on that factor
and the identity on all the others. Lift it to \(F[G]\).
Its augmentation is zero since the augmentation factors
through the semisimple quotient. Express it as a finite sum
of multiples of \(g-1\), and clear the coefficient denominators
by a single nonzero element of \(\mathcal O\). On every
nontrivial simple subquotient the resulting element acts
by this nonzero scalar, and on each trivial one it acts by
zero. In a basis respecting a composition series its matrix
is block triangular. The product of the diagonal determinants
proves the claim.
\end{proof}

The multiplier belongs to the algebra generated by the
characteristic-polynomial families and the auxiliary variables
when this algebra is chosen as \(S\). To use its coefficient
on automorphic forms requires control of the trace relations under
passage from auxiliary level to the level of the period. The
crystalline colengths and the main-summand Selmer length in
the bound must also be estimated for those same representations.
\clearpage


\subsection{First-order transfer at a split auxiliary place}
\label{sec:auxtransfer721}

Fix one of the places \(w\) of Corollary~\ref{cor:markeddet715}.
Write \(L=E_w\), \(k_w=\mathbb F_q\), \(m=\ell^\nu\mid q-1\)
and \(D_m=\ZZ/m\ZZ\). Let \(\mathcal O\) be a finite extension
of \(\ZZ_\ell\) in the sense of valuation rings. At principal
level the group \(\Gamma=\GL_n(k_w)\) acts on the integral
lattice of forms \(M\) by operators \(\rho(g)\). All the
operators below act in \(R=\operatorname{End}_{\mathcal O}(M)\).
Congruence modulo \(a\) means difference in \(aR\).
Choose a uniformizer \(\varpi\) whose local Artin symbol is the
class of the chosen geometric Frobenius \(\phi\), and a unit
with residue \(z\) whose Artin symbol is the class of the chosen
tame generator \(\tau\). Here only the tame quotient of inertia
matters for the level characters. Local reciprocity identifies
the tame part of abelianized inertia with \(k_w^\times\). Its
chosen generator has a logarithm
\(b_0:k_w^\times\to D_m\) with \(b_0(z)=1\).
Changing a uniformizer by a unit realizes any lift of geometric
Frobenius in the local abelianization.

Let \(P\subset\Gamma\) stabilize the first coordinate line. Put
\(a(h)=b_0(h_{11})\) on \(P\), \(b(g)=b_0(\det g)\) on
\(\Gamma\), and \(d=\operatorname{diag}(z,1,\ldots,1)\).
Choose representatives \(r_i\) for the lines \(\Gamma/P\).
For an integer-valued function \(f\) on \(D_m\) write
\[
 L_a(f)=\sum_{h\in P}f(a(h))\rho(h),\qquad
 L_b(f)=\sum_{g\in\Gamma}f(b(g))\rho(g),\qquad
 T^\pm=\sum_i\rho(r_i^{\pm1}).
\]
Set \(N=L_b(1)\), \(N_P=L_a(1)\), and choose any integral
representative \(\lambda(x)\) of each residue \(x\in D_m\).

\begin{lemma}[Two weighted norms]\leavevmode\label{lem:transfer721}
As operators modulo \(m\),
\[
        T^+L_a(\lambda)T^-\equiv L_b(\lambda),
       \qquad(\rho(d)-1)L_b(\lambda)\equiv-N .
\]
For each \(g\in\Gamma\) there is an integral combination
\(Q_g\) of the \(\rho(r_i^{-1})\) such that
\[
       L_a(\lambda)T^-(\rho(g)-1)\equiv N_P Q_g\pmod m .
\]
\end{lemma}
\begin{proof}
In the group ring the coefficient at \(g\) in the first product
is \(\sum_i\lambda(a(r_{gi}^{-1}g r_i))\). Modulo \(m\) this
sum is additive in products \(g_1g_2\), by the chain rule for
the line scalars and permutation of the lines by \(g_2\).
On a transvection its value is zero since the order of a
transvection is prime to \(m\). For \(g=\operatorname{diag}
(x,1,\ldots,1)\), the first line contributes \(b_0(x)\).
On any other orbit of lines the product of the successive
scalars is the scalar by which the corresponding power of \(g\)
acts on the initial line. A vector on this line has a nonzero
coordinate after the first, which \(g\) leaves fixed. Hence this
scalar is one. Diagonals in other positions are conjugate to
these examples, and transvections and diagonals generate
\(\Gamma\) by elimination. This proves the first formula.
Left translation by \(d\) changes the weight at \(g\) from
\(\lambda(b(g))\) to \(\lambda(b(g)-1)\), proving the second.
Finally \(r_{gi}^{-1}g=h_i r_i^{-1}\) with
\(h_i=r_{gi}^{-1}g r_i\in P\).
Right translation by \(h_i\) replaces the weight at \(h\)
by \(\lambda(a(h)-a(h_i))\).
The difference is \(-a(h_i)\) modulo \(m\). Thus one can take
\(Q_g=-\sum_i\lambda(a(h_i))\rho(r_i^{-1})\).
\end{proof}

Write \(H\) for the common polynomial of
Theorem~\ref{thm:commonhecke716}, with \(H(1)=1\), and
\(p(Y)=(Y-q^{k-1})H(Y)\). On pro-Iwahori invariants put
\(\widetilde U=C_w^{-1}U_w\). Define on the whole principal
level lattice
\[
            B(f)=p(\widetilde U)L_a(f),\qquad
            \partial=-T^+B(\lambda)T^-.
\]
Every image of \(L_a(f)\) is fixed by the inverse image of
\(\ker a\) in \(\GL_n(\mathcal O_L)\). This inverse image
contains the pro-Iwahori. On these invariants \(U_w\) is the sum
of the \(q^{n-1}\) translations by
\(n(v)\operatorname{diag}(\varpi,1,\ldots,1)\), where \(n(v)\)
has upper first row \((1,v)\) and the remaining rows of the
identity. These representatives also give the double coset for
the inverse image of \(\ker a\). Indeed, conjugating its
matrices by the diagonal and intersecting with this group
imposes the single extra condition that the off-diagonal
upper first row vanish modulo \(\varpi\). This follows by
multiplying the blocks, using that its lower-left block is
already divisible by \(\varpi\). The same calculation holds
for the inverse image of \(P\). Thus \(U_w\) preserves the
first invariant space and commutes there with \(P/\ker a\):
one can lift generators by diagonal units, which permute the
row representatives and commute with the diagonal uniformizer.
It follows that \(B(1)\) is left \(P\)-invariant, and all these
operators preserve \(M\). The normalizing scalar \(C_w\) is
an \(\ell\)-adic unit.

Use a characteristic-zero decomposition of the forms by
automorphic representations. Suppose a continuous integral
representation \(u_x\) of \(G_E\) is attached to every system
occurring at this principal level. At \(w\) assume agreement of
its Frobenius-semisimplified Weil--Deligne representation with
the local Langlands parameter, including its monodromy operator.
Use geometric reciprocity and the normalization in which a
character \(\mu\) of \(L^\times\) in the normalized induction
contributes \(C_w^{-1}q^{(n-1)/2}\mu(\varpi)\) at \(\phi\)
and \(\mu(z)\) at \(\tau\). In this formula \(\mu(z)\) means
evaluation at a unit lifting \(z\), for a tame character.
Assume the normalized good Frobenius characteristic coefficients
act integrally at principal level at a conjugacy-dense set of
good places. Then the coefficient families for every \(g\in G_E\)
act integrally. Indeed these are limits in a finite-dimensional
space of scalar operators on the isotypic summands, and the
lattice \(R\) is closed. In particular the following traces are
integral operators, commuting with all local translations. Write
\(\operatorname{Tr}\) for the family of traces in these \(u_x\),
extended linearly to the group ring.

\begin{theorem}[The tame trace and its Frobenius term]
\leavevmode\label{thm:tametransfer721}
Put \(m'=m/\gcd(2,m)\). For every \(g\in G_E\) one has
\[
 \begin{split}
        N\operatorname{Tr}((\tau-1)g)&=0,\\
 \partial\operatorname{Tr}((\tau-1)g)&\equiv
        N\operatorname{Tr}\bigl((\phi-1)H(\phi)g\bigr)
                                                   \pmod {m'} .
 \end{split}
\]
The same two formulas hold with \((\tau,\phi)\) replaced by
\(((c\tau c^{-1})^{-1},(c\phi c^{-1})^{-1})\) if \(c\) is
complex conjugation and the family has a common polarization
\(\operatorname{Tr}(g)=\xi(g)\operatorname{Tr}((cgc^{-1})^{-1})\),
where \(\xi\) is an integral unit character, \(\xi(\tau)=1\)
and \(\xi(\phi)\equiv1\pmod m\).
\end{theorem}
\begin{proof}
For this proof write \(A=\operatorname{Tr}((\tau-1)g)\),
\(B_p=\operatorname{Tr}(p(\phi)(\tau-1)g)\) and
\(F_p=\operatorname{Tr}(p(\phi)g)\). We work over a splitting
field when decomposing functions into characters. On a vector
transforming under \(\Gamma\) by \(\psi\circ b\), the local
representation becomes spherical after removal of that
determinant twist. Its parameter is unramified after the same
twist and its monodromy is zero. Thus \(u_x(\tau)=\psi(1)I\).
This follows also for the original Galois representation from
Frobenius semisimplification: it preserves the restriction to
inertia and the monodromy operator. On this vector \(B_p\)
therefore agrees with \((\rho(d)-1)F_p\).
The operator \(L_b(f)\) takes values in the sum of exactly
these types, by decomposing \(f\) into characters of \(D_m\).
It commutes with \(\rho(d)\). Consequently
\[
        L_b(f)B_p=(\rho(d)-1)L_b(f)F_p,\qquad NA=0.
\]

On a vector transforming under \(P\) by a nontrivial
\(\psi\circ a\), the local parameter has one tame character
\(\mu\) on a line, with \(\mu|_{\mathcal O_L^\times}=\psi
\circ b_0\) after reduction, and an unramified complement with
zero monodromy. The indicated type is a one-dimensional space
in its irreducible local representation. These facts are the
level-one classification of Clozel, Harris and Taylor.
Their subgroup fixes a hyperplane with trivial action on its
quotient. The automorphism \(g\mapsto w_0\,{}^t g^{-1} w_0^{-1}\),
where \(w_0\) reverses the basis, exchanges this subgroup with
the group fixing the first vector and inverts the scalar
character. In their classification a spherical representation
has only the trivial character, and in the ramified case with
a nontrivial character the other \(n-1\) parameters are
unramified and spherical. The alternative involving a
Steinberg block has only the trivial character on the fixed line.
Their assertions concern smooth representations in characteristic
zero.~\cite[Lemmas 3.1.3 and 3.1.5]{IX-VII-VII-chttransfer721}

For the normalization of \(U_w\), the character space can be
computed in normalized induction from the upper parabolic of
type \((1,n-1)\), with factors \(\mu\) and the spherical
representation on \(\GL_{n-1}\). Its line is represented by a
function supported on the product of this parabolic with the
inverse image of \(P\), taking a nonzero spherical value at
the identity. The defining equivariance is consistent on the
intersection since its line character is \(\mu\) on units.
This is the line identified by the level-one classification.
Evaluating its Hecke sum at the identity gives
\(q^{n-1}q^{-(n-1)/2}\mu(\varpi)\) times that value, by
the modulus character of the parabolic. If \(n=1\), the same
calculation is evaluation of a single character.
Since \(\psi(1)\ne1\), the image of \(u_x(\tau)-1\) is the
one-dimensional \(\psi(1)\)-eigenspace and is stable under
\(\phi\). Frobenius semisimplification does not change its
eigenvalue there. On this line \(u_x(\phi)\) has exactly the
eigenvalue of \(\widetilde U\) on the form.
It follows after multiplication by \(u_x(g)\) and taking trace
that \(p(\widetilde U)A=B_p\) on this character space.

The integral linear function \(E(f)=B(f)A-L_a(f)B_p\)
now kills every difference \(f(\,\cdot-1)-f\): such a
difference has zero constant Fourier coefficient, and each
remaining Fourier term has image of nontrivial line character.
This exact equality descends to the original lattice. Apply
it to \(f(x)=\lambda(x)(\lambda(x)+1)\). Its difference is
\(-2\lambda(x)\) modulo \(m\). Hence \(2E(\lambda)\in mR\).
Freeness over \(\mathcal O\) gives \(E(\lambda)\in m'R\).
Using Lemma~\ref{lem:transfer721} and commutation of traces
with translations gives
\[
 \partial A\equiv-L_b(\lambda)B_p
      =-(\rho(d)-1)L_b(\lambda)F_p\equiv NF_p\pmod {m'} .
\]
Since \(q\equiv1\pmod m\), replacing \(p(\phi)\) by
\((\phi-1)H(\phi)\) changes the integral operator by a multiple
of \(m\). For the conjugate pair expand the polynomial in
powers of Frobenius. The map \(g\mapsto(cgc^{-1})^{-1}\)
reverses products. Cyclicity puts the power of Frobenius first
in each two-factor trace. Polarization gives exactly the same
factor \(\xi(g)\) for the two terms of the inertia difference.
For each Frobenius power the extra factor is a power of
\(\xi(\phi)\) or its inverse, hence is one modulo \(m\).
Applying the unprimed formulas with \((cgc^{-1})^{-1}\)
in place of \(g\) proves both claims.
\end{proof}

\begin{proposition}[Right translation of the weighted operator]
\leavevmode\label{prop:righttransfer721}
Suppose the lattice at principal level is the direct sum of its
coefficient lattices over a finite \(\Gamma\)-set of double
cosets, with the action giving isomorphisms of the fibers.
Let \(c_0>0\) be a multiple of all point stabilizer orders.
For each \(g\in\Gamma\) there is \(Z_g\in R\) such that
\[
        c_0\partial(\rho(g)-1)\equiv NZ_g\pmod m .
\]
Every integral operator \(K\) satisfying \(K\rho(g')=K\)
for \(g'\in\Gamma\) and \(K B(1)=0\) also satisfies \(KZ_g=0\).
\end{proposition}
\begin{proof}
For a subgroup \(H_0\subset\Gamma\) take one point of each
\(H_0\)-orbit. On its fiber put \(c_0/|\operatorname{Stab}_{H_0}|\)
times the identity, and put zero on the other fibers.
The resulting endomorphism \(\Psi_{H_0}\) satisfies
\(\sum_{h\in H_0}\rho(h)^{-1}\Psi_{H_0}\rho(h)=c_0 I\):
each point maps to the representative exactly
\(|\operatorname{Stab}_{H_0}|\) times.
In particular \(c_0 B(1)=N_P\Psi_P B(1)\) by left
\(P\)-invariance. Apply \(p(\widetilde U)\) to the last
congruence of Lemma~\ref{lem:transfer721}. This is legitimate
modulo \(m\) since the divided difference takes values in the
\(\ker a\)-fixed lattice, which is saturated, and the Hecke
sum is integral there. Multiply by \(-c_0 T^+\) on the left.
Since \(T^+N_P=N\), one obtains the congruence with
\(Z_g=-\Psi_P B(1)Q_g\). Finally
\(|P|K\Psi_P B(1)=c_0 K B(1)=0\). Cancellation of the
nonzero integer \(|P|\) in characteristic zero proves \(KZ_g=0\).
\end{proof}

For applications take \(K=D N_{\mathrm{all}}\), where
\(N_{\mathrm{all}}\) is the product of the norms at all the
auxiliary places and \(D\) acts by scalars on systems.
If the base-level support of \(D\) consists of the spherical
systems retained in Theorem~\ref{thm:commonhecke716}, then
\(K B(1)=0\). Indeed \(N_{\mathrm{all}}\) is the product of
the group orders times orthogonal projection to base level.
Testing against each spherical vector on this support after
applying \(p(\widetilde U)\) gives zero by that theorem.
For the definite quotient, a principal-level orbit stabilizer
is the image in the residue group of the stabilizer of a
base-level double coset: lifting its elements to the compact
level subgroup gives this description by the definition of a
double coset. Thus any common multiple valid at base level is
valid throughout this calculation. The local transfer uses the
same place as the cohomology test and has no dependence on its
norm beyond \(m\mid q-1\). Its automorphic application requires
the indicated characteristic-zero compatibility for the
systems at principal level.
\clearpage

\subsection{Galois representations at the auxiliary split primes}
\label{sec:splitattachment722}

Let \(E/F\), \(F=E^+\), be a field occurring in either family of
Corollary~\ref{cor:fieldrelations712}. Use the unitary group of
the identity Hermitian form in \(n\) variables over \(E\), an
algebraic coefficient representation at infinity and any compact
open level. In particular the level may be principal at every
prime in the auxiliary set, including primes chosen after the
coefficient precision has been fixed. Fix
\(\iota:\overline{\QQ}_\ell\simeq\CC\). We use the local
Artin map taking a uniformizer to geometric Frobenius. For
a complex local representation write \(\operatorname{rec}\)
for the unitary normalization of local Langlands, in which a
character of \(E_w^\times\) corresponds to that character via
reciprocity.

\begin{theorem}[Attachment on the period fields]\leavevmode
\label{thm:splitattachment722}
Every automorphic constituent \(x\) of this space has a continuous
semisimple representation \(r_x:G_E\to\GL_n(\overline{\QQ}_\ell)\).
At each split \(w\nmid\ell\) where the local Langlands parameter of
\(\pi_{x,w}\) has zero monodromy, one has
\[
 \iota\,\mathrm{WD}(r_x|_{G_{E_w}})^{\mathrm{F\text{-}ss}}
       \simeq
       \operatorname{rec}(\pi_{x,w})\otimes
                                 |\cdot|_w^{(1-n)/2}.
\]
In particular this holds at every spherical place, after any
determinant twist of a spherical representation, and on every
nontrivial line-character type of Section~\ref{sec:auxtransfer721}.
Moreover \(r_x^c\simeq r_x^\vee\otimes\kappa^{1-n}\).
For the finitely many systems at the chosen level a single
finite extension of \(\QQ_\ell\) supplies stable lattices.
\end{theorem}
\begin{proof}
The extension \(E/F\) is unramified at all finite places and
\([F:\QQ]>1\) by Corollary~\ref{cor:fieldrelations712}.
At a nonsplit finite place \(E_w/F_v\) is unramified quadratic.
The norm is onto on units. Indeed the residue-field norm is
onto, and on successive principal-unit quotients the norm
induces the surjective residue trace. Corrections at successive
powers of the uniformizer converge by completeness. Choose
\(a\in E_w^\times\) with \(a\bar a=-1\). Each two-coordinate
plane in the identity form then contains the isotropic vector
\((1,a)\) and is hyperbolic. Pairing the coordinates two at
a time gives Witt index \(\lfloor n/2\rfloor\). Hence the group
is quasi-split at these places, and it is \(\GL_n\) at each
split place. At infinity it is compact.

Labesse's cohomological base change gives
\(\Pi=\boxplus_i\Pi_i\), where \(\Pi_i\) is discrete automorphic
on \(\GL_{n_i}(\mathbb A_E)\) and conjugate self-dual.
The transfer matches the local representation at every split
place and its infinitesimal character at infinity. Compactness
at infinity ensures the discrete-series and sign hypotheses
in Labesse's corollary. The hypothesis on the degree of \(F\)
and the local hypotheses just checked apply to all our fields.
The formulation for a definite unitary group, including the
split-place condition, is also given by Guerberoff.
~\cite[Cor.~5.3]{IX-VII-VII-labesse722}\cite[Thm.~2.2]{IX-VII-VII-guerberoff722}

By the discrete-spectrum theorem of M{\oe}glin and Waldspurger,
each \(\Pi_i\) is the isobaric representation formed from one
unitary cuspidal representation \(\rho_i\) of rank \(a_i\)
and its string of norm twists, with \(a_i b_i=n_i\).
In unitary normalization the exponents in this string are
\(j=(1-b_i)/2,(3-b_i)/2,\ldots,(b_i-1)/2\).
At a finite place the parameter of the isobaric sum is the
direct sum of the parameters of its terms. This uses the
Langlands quotient defining the local isobaric sum. Write
each cuspidal factor locally in its essentially square-integrable
Langlands data, twist the data, and take their union in
Langlands order. The local correspondence takes that quotient
to the direct sum of the corresponding Weil--Deligne blocks.
This description includes residual discrete global factors.
~\cite{IX-VII-VII-mw722}

For each \(i,j\) put
\[
       \pi_{ij}=\rho_i|\det|^{j+(a_i-n)/2}.
\]
These cuspidal representations are regular algebraic in their
own ranks. At a complex embedding, the full infinitesimal
character is that of an algebraic highest weight for \(\GL_n\).
Its entries are distinct and belong to \((n-1)/2+\ZZ\), by
the half-sum of the positive roots. The infinitesimal character
of \(\rho_i|\det|^j\) contributes a submultiset. Shifting by
\((a_i-n)/2\) places its entries in \((a_i-1)/2+\ZZ\) while
preserving distinctness. Order the entries decreasingly and
subtract the half-sum of positive roots of \(\GL_{a_i}\).
Successive differences are nonnegative integers and the entries
are integers. This gives the required algebraic highest weight.
Apply the argument at every embedding, including its complex
conjugate.

Take the direct sum of the Galois representations attached to
the \(\pi_{ij}\) by Harris, Lan, Taylor and Thorne, using class
field theory in rank one. Their normalization introduces the
twist \(|\cdot|^{(1-a_i)/2}\), so the net exponent on the
parameter of \(\rho_i\) is \(j+(1-n)/2\).
At \(w\nmid\ell\), Varma's compatibility theorem gives the
semisimple Weil representation of each summand and bounds
its monodromy by the automorphic monodromy. In particular if
the latter operator is zero then the former is zero, since
every Jordan block in the upper bound has length one.
At a split place where \(\operatorname{rec}(\pi_{x,w})\) has
zero monodromy, each direct summand of its parameter has zero
monodromy. Thus Varma's result gives the displayed
Frobenius-semisimplified equality, including the zero operator.
Spherical parameters and their character twists have zero
monodromy. For the nontrivial line character this follows from
the classification used in Theorem~\ref{thm:tametransfer721}.
~\cite{IX-VII-VII-hltt722}\cite[Thms.~1--2]{IX-VII-VII-varma722}

The places of \(E\) above split places of \(F\) have Dirichlet
density one as places of \(E\). Indeed the inert prime norms
are squares of the norms downstairs, whose reciprocal squares
have convergent sum. Removing the finitely many ramified, level
and coefficient places preserves density one. At the remaining
places the polynomial of geometric Frobenius is exactly the
arithmetic Hecke polynomial in Section~\ref{sec:heckechains716}.
The conjugate place has dual unitary Satake parameters.
Multiplying by the arithmetic normalization on both sides gives
conjugate roots \(q^{n-1}/\alpha\) for the roots \(\alpha\) at \(w\).
Chebotarev density, continuity of traces and Brauer--Nesbitt
give \(r_x^c\simeq r_x^\vee\kappa^{1-n}\). The same argument
makes the representation depend only on the good Hecke system.

For completeness the image of each \(r_x\) is compact. It
is contained in \(\GL_n(K)\) for a finite coefficient extension
\(K\). Indeed \(\overline\QQ_\ell\) is a union of countably many
finite extensions (there are finitely many of each given degree
up to isomorphism and finitely many embeddings of each).
Intersection with \(\GL_n(K)\) gives closed subgroups of the
compact image. By the Baire theorem one has nonempty interior
and therefore finite index. Adjoin the entries of a finite set
of coset representatives to this \(K\). The orbit of the standard
lattice is bounded and its integral span is a stable lattice.
Taking the compositum for finitely many constituents finishes
the proof.
\end{proof}

\begin{corollary}[The normalized traces at both auxiliary places]
\leavevmode\label{cor:normalizedtransfer722}
Take the unitary theta splitting character \(\mu\) for the
source of dimension \(d\) and put \(k=n-d\). At the auxiliary
places away from its conductor use
\(C_w=q^{d/2}\mu_w(\varpi)\).
For the systems of the preceding theorem both trace identities
of Theorem~\ref{thm:tametransfer721}, including their conjugate
versions, hold with the same auxiliary primes and \(m'\).
For the orthogonality statement of
Proposition~\ref{prop:righttransfer721} one may take as support
all base-level weak systems congruent in their normalized
characteristic polynomials to the selected lift.
\end{corollary}
\begin{proof}
The splitting character restricts on \(\mathbb A_F^\times\)
to the \(d\)-th power of the quadratic character of \(E/F\).
It is therefore trivial on norms from \(E\), so \(\mu^c=\mu^{-1}\).
Its angular exponent \(h\) at a complex place has parity \(d\),
by restriction to the real element \(-1\).
The character \(\mu^{-1}|\cdot|_E^{d/2}\) is algebraic. At
each infinity place the exponents of \(z,\bar z\) are
\((d-h)/2,(d+h)/2\). Use its \(\ell\)-adic avatar \(\alpha\)
via global reciprocity. At good \(w\nmid\ell\) its value at
geometric Frobenius is \(\iota^{-1}(C_w^{-1})\), and
\(\alpha^c\alpha=\kappa^d\), since the norm character has
Frobenius value \(q^{-1}\). It is a unit-valued character,
as its Galois image is compact. For \(u_x=\alpha r_x\) the
polarization multiplier is thus \(\kappa^{1-k}\). This takes
value one at tame inertia and is one modulo \(m\) at either
Frobenius since \(m\mid q-1\).

The proof of Theorem~\ref{thm:tametransfer721} uses local
compatibility on vectors transforming by a determinant character
and on vectors transforming by a nontrivial line character.
Both are covered by Theorem~\ref{thm:splitattachment722}. On
a nontrivial line character the unique ramified line has the
eigenvalue of \(\widetilde U\), including the factor \(C_w^{-1}\).
At good split places the characteristic coefficients act
integrally because they are double-coset operators multiplied
by unit powers of \(\alpha\), signs and integral powers of
the prime norm. These Hecke operators are sums of translations
away from \(\ell\). The conjugacy density and closed-lattice
argument of Section~\ref{sec:auxtransfer721} therefore applies
on the full principal-level lattice.

On a base-level weak system \(u_x(g)\) has eigenvalue \(1\)
for all \(g\). Equality of polynomials at the spherical place
\(w\) gives \(P_x(C_w)=0\). Residual agreement with the
selected lift gives the residual polynomial required in
Theorem~\ref{thm:commonhecke716}, by its decomposition into
the main summand and the \(k\) cyclotomic characters.
The chosen auxiliary primes already make the residual
determinant on the main summand a unit at \(1\).
Consequently a single integral \(H\) at each \(w\) gives the
orthogonality for all these systems.
\end{proof}

This proves compatibility at the auxiliary places used in the
transfer, in every coefficient characteristic. The control of
integral extension classes at places above the coefficient prime
requires an estimate with the varying Hodge--Tate weights.
\clearpage

\subsection{Descent of the squarefree trace coefficient to base level}
\label{sec:topdescent723}

Use the two completely split level primes \(p_0,p_1\) of
Section~\ref{sec:levelmass678}. Keep the base field, rank \(n\)
and base level fixed for now. At auxiliary precision
\(m=\ell^\nu\), write \(M\) for the lattice at principal
level at the auxiliary places \(w_j,\ j\in J\), and \(M_0\)
for the original lattice of forms. In the model of forms as
functions valued in an integral coefficient lattice,
\(M_0=M^\Gamma\), where
\(\Gamma=\prod_j\GL_n(k_{w_j})\): invariance under the
residue group restores the original compact open level at
each \(w_j\), with the same values of the function. Let
\(R=\operatorname{End}_{\mathcal O}(M)\), and let \(\mathbb T\)
be the integral algebra of scalar families generated by
the normalized Galois characteristic coefficients of
Corollary~\ref{cor:normalizedtransfer722}. Its action commutes
with the local translations. All congruences in \(R\) in this
section are modulo \(m'=m/\gcd(2,m)\).

\begin{lemma}[A stabilizer multiple fixed at base level]
\leavevmode\label{lem:basestabilizer723}
The integer
\[
 c_0=\gcd\bigl(|\GL_n(\mathbb F_{p_0})|,
                |\GL_n(\mathbb F_{p_1})|\bigr)
\]
is a multiple of every point stabilizer order for \(\Gamma\)
on the set of principal-level double cosets. It satisfies
\(\log c_0\le n^2\min_i\log p_i\), independently of the
auxiliary places and their number. There is \(\Psi\in R\)
with \(\sum_{g\in\Gamma}\rho(g)^{-1}\Psi\rho(g)=c_0 I\).
Writing \(N=\sum_{g\in\Gamma}\rho(g)\), one has
\(N\Psi z=c_0 z\) for \(z\in M_0\).
\end{lemma}
\begin{proof}
The rational stabilizer of a base-level representative is
finite. Indeed it is discrete in the adeles and lies in a
compact subgroup at every place, including all infinite
places, with adelically compact product. Its embedding
at either marked place above \(p_i\) is faithful and has
bounded image in \(\GL_n(\QQ_{p_i})\).
The integral span of the translates of \(\ZZ_{p_i}^n\) by
this finite image is a stable lattice. On reduction of this
lattice the kernel is torsion-free, since \(p_i\) is odd.
To see this, for a nonidentity matrix \(1+p_i^a B\), \(a\ge1\)
and \(B\) nonzero modulo \(p_i\), raising to a power prime to
\(p_i\) preserves \(a\). Raising to \(p_i\) increases \(a\)
by exactly one, by binomial expansion. No positive power is
the identity. The finite rational stabilizer thus embeds in
each \(\GL_n(\mathbb F_{p_i})\). Its image in the auxiliary
quotient is the corresponding \(\Gamma\)-stabilizer of a
principal-level point. Indeed a residue translation fixes a
double coset precisely when a lift differs from a rational
stabilizer element, conjugated into the base compact group,
by an element of principal level. This gives the divisibility.
Counting matrices gives \(|\GL_n(\mathbb F_p)|\le p^{n^2}\).
Applying the orbit construction in the proof of
Proposition~\ref{prop:righttransfer721} to the whole product
gives \(\Psi\). Acting on an invariant vector in the displayed
sum and reindexing by inverses gives the last identity.
\end{proof}

At \(j\) denote the norm by \(N_j\), the operator of
Theorem~\ref{thm:tametransfer721} by \(\partial_j\), and put
\[
    N_S=\prod_{j\in S}\partial_j\prod_{j\notin S}N_j,\qquad
    \mathcal N(X)=\prod_{j\in J}(N_j+X_j\partial_j).
\]
Operators at different \(j\) commute because their translations
have disjoint local support. Thus \(N_\varnothing=N\) and
\(N_S\) is the coefficient of \(X_S=\prod_{j\in S}X_j\).
For \(f\in\mathbb T[X_j:j\in J]\) set \(f_S=[X_S]f\) and put
\(G(f)=[X_J]\mathcal N f=\sum_S N_{J\setminus S}f_S\).
Let \(\mathcal J\) be the trace ideal of
Section~\ref{sec:tracemult720} with precision \(a=m\) and the
common \(H_j\)'s of Corollary~\ref{cor:normalizedtransfer722}.

\begin{proposition}[Two identities for the deformed norms]
\leavevmode\label{prop:normfactor723}
For every \(f\in\mathcal J\) one has \(G(f)\equiv0\).
Suppose \(D\in\mathbb T\) has base-level support in the
spherical systems retained at every \(j\) in
Theorem~\ref{thm:commonhecke716}. Put \(\delta=c_0^2D\).
For \(i\ge |S|\) there is an integral endomorphism \(Y_{S,i}\)
with \(\delta^iN_S\equiv Y_{S,i}N\). One can take
\(Y_{\varnothing,i}=\delta^i\).
\end{proposition}
\begin{proof}
Compute the top coefficient modulo the squares \(X_j^2\).
A generator \(t-X_j r\) at \(j\), with \(t,r\) the inertia
trace and the Frobenius trace in Theorem~\ref{thm:tametransfer721},
satisfies
\[
   (N_j+X_j\partial_j)(t-X_jr)
       \equiv N_j t +X_j(\partial_jt-N_jr)=0
\]
in that quotient modulo \(m'\), for both conjugate choices.
Multiplication by any polynomial in \(\mathbb T[X]\) and
the other local factors preserves the equality. Generators
divisible by \(m\) also give zero. This proves the first assertion.

For the second, induct on \(|S|\), allowing every \(i\ge|S|\).
The assertion for the empty set is exact. Otherwise
\(i=e+1\) and \(e\ge|S|-1\). We show that
\(U=c_0D\delta^e N_S\) is right invariant modulo \(m'\).
For \(g\in\Gamma_j\) and \(j\notin S\), its \(N_j\) factor
gives \(U(\rho(g)-1)=0\). For \(j\in S\), Proposition
~\ref{prop:righttransfer721} gives \(Z\in R\) with
\(c_0\partial_j(\rho(g)-1)\equiv N_j Z\) and \(DNZ=0\).
The latter is the proposition with \(K=DN\), using spherical
orthogonality on its base support. Consequently
\[
 U(\rho(g)-1)\equiv
     D\delta^e N_{S\setminus\{j\}}Z
     =\delta^e N_{S\setminus\{j\}}D Z
     \equiv Y_{S\setminus\{j\},e}\,N D Z=0 .
\]
Here \(D\) has been moved past local factors before use of the
induction hypothesis. Invariance extends to \(\Gamma\), which
the individual factors generate. Multiplying the averaging
identity of Lemma~\ref{lem:basestabilizer723} on the left by
\(U\) now gives \(c_0U\equiv U\Psi N\).
Since \(c_0 U=\delta^i N_S\), this completes the induction.
\end{proof}

\begin{theorem}[Integral descent of the top coefficient]
\leavevmode\label{thm:topdescent723}
Let \(A,\delta_0\in\mathbb T\) and \(0\ne n_0\in\mathcal O\).
Assume \(A\delta_0\) has the base support in
Proposition~\ref{prop:normfactor723}. Let
\((\delta_0-n_0)f\in\mathcal J\), \(i=|J|\), and
\((n_0c_0^2)^{i+1}\mid m'\) in \(\mathcal O\).
Let \(l,y\in M_0\otimes\operatorname{Frac}(\mathcal O)\)
and \(h\in\mathcal O\) satisfy
\[
        h l\in M_0,\qquad \delta_0y=n_0y,\qquad
        A\delta_0l=n_0Ay,\qquad
        f_S y=0\quad(S\subsetneq J).
\]
Then \(h c_0 A^{i+1} f_J y\in M_0\).
\end{theorem}
\begin{proof}
Put \(n_1=n_0c_0^2,\ \delta=c_0^2 A\delta_0\).
In the commutative Hecke algebra the difference
\(\delta^{i+1}-(n_1A)^{i+1}\) is divisible by
\(\delta_0-n_0\), by the difference of powers. Every Hecke
scalar commutes with the coefficients of \(\mathcal N\).
Applying Proposition~\ref{prop:normfactor723} first to this
difference times \(f\), then to \(\delta^iN_{J\setminus S}\)
in each term of \(\delta^{i+1}G(f)\), gives
\[
   (n_1A)^{i+1}G(f)\equiv\delta^{i+1}G(f)
        \equiv\sum_{S\subseteq J}Y_{J\setminus S,i}\,N\delta f_S .
\]
Apply the congruence to the integral vector \(u=\Psi(h l)\).
By the norm identity \(Nu=h c_0 l\). Thus
\[
       N\delta f_S u=h c_0 f_S\delta l
                      =n_1 h c_0 A f_S y.
\]
Only \(S=J\) survives, and \(Y_{\varnothing,i}=\delta^i\)
acts on \(A f_J y\) as \((n_1 A)^i\). It follows that
\[
 n_1^{i+1}\bigl(A^{i+1}G(f)u-hc_0 A^{i+1}f_J y\bigr)
       \in m'M .
\]
Divide in the fraction space. Since \(m'/n_1^{i+1}\in\mathcal O\)
and \(A^{i+1}G(f)u\in M\), the second term in parentheses
lies in \(M\). It is at base level because its scalar operators
preserve \(M^\Gamma\) over the fraction field. This proves the
claim with one factor of \(c_0\).
\end{proof}

\begin{corollary}[Base separation before the choice of primes]
\leavevmode\label{cor:baseseparation723}
Fix a nonempty set \(X_0\) of split lifts among the base systems
as in Theorem~\ref{thm:ribettrace720}, and select \(x_*\in X_0\).
Use the row hypotheses of that theorem at auxiliary level.
For all \(l,y\) as above with \(hl\in M_0\), assume \(y\) is
supported on \(X_0\) and \(l,y\) have equal components on base
systems whose normalized representation has a trivial
composition factor. There is a diagonal integral Hecke operator
\(s\) at base level with \(hc_0s y\in M_0\) for all these
pairs. At \(x_*\) its eigenvalue has the bound for
\(v(\beta_{x_*})\) in Theorem~\ref{thm:ribettrace720}.
\end{corollary}
\begin{proof}
Apply Corollary~\ref{cor:augdet720} to the finite base family
alone. It gives a group algebra element \(w\) such that
\(\beta(x)=\det u_x(w)\) vanishes exactly on the base systems
with a trivial composition factor. The same expression defines
an integral Hecke element on every auxiliary family by
Lemma~\ref{lem:blockamitsur720}. Let \(V\) be the set of its
distinct nonzero values on the base family and put
\[
    \delta_0=\prod_{\lambda\in V}(\lambda-\beta),\qquad
    n_0=\prod_{\lambda\in V}\lambda .
\]
Both empty products mean one. At base level \(\delta_0\) has
value \(n_0\) precisely on the systems with a trivial factor
and vanishes elsewhere. Every such system is weak, since each
individual matrix has eigenvalue one on that subquotient.
In particular \(\delta_0y=n_0y\) and \(\delta_0l=n_0y\).
The nonzero scalar \(n_0\) is fixed before any auxiliary primes
are chosen. The number of primes \(i=(k-1)|S_\ell|\) of the
construction of Theorem~\ref{thm:auxtriangle713} is independent
of \(\nu\). Therefore choose \(\nu\) such that
\[
       \nu-v_\ell(\gcd(2,\ell^\nu))
           \ \ge\ (i+1)v_\ell(n_0c_0^2),
\]
where the valuation extends \(v_\ell(\ell)=1\). This divisibility
persists on enlarging the coefficient field to include the
auxiliary systems.

Let \(A\) be the indicator of the systems residually congruent
to \(x_*\) in all normalized characteristic coefficients.
It is integral. To prove this, for each system outside this
class select one distinguishing coefficient \(a\in\mathbb T\).
Over a common residue field of size \(Q\) form the product of
\(1-(a-a(x_*))^{Q-1}\) over the chosen coefficients. Its
values reduce to one on the class and zero off it. Its
\(\ell^u\)-th powers tend to \(A\): a unit congruent to one
has \(\ell^u\)-th powers tending to one, by the binomial
formula, while a nonunit tends to zero. The algebra \(\mathbb T\)
is an \(\mathcal O\)-submodule of the finite free module \(R\).
It is finite and closed, so it contains the limit. By
Corollary~\ref{cor:normalizedtransfer722}, \(A\delta_0\)
has the required spherical support.

Take \(f\) from Theorem~\ref{thm:ribettrace720} for this family.
It satisfies \(\beta f\in\mathcal J\). Since \(\beta\) divides
\(\delta_0-n_0\) in \(\mathbb T\), it also satisfies the
ideal condition of Theorem~\ref{thm:topdescent723}. All
proper squarefree coefficients of \(f\) vanish on \(X_0\).
That theorem now applies with \(s\) the restriction of \(Af_J\)
to base level, because \(A^{i+1}=A\). At \(x_*\) its value
is \(\beta_{x_*}\) from the trace multiplier. The scalar \(n_0\)
only determined how large a precision was needed. It disappears
from the integral vector and from the eigenvalue bound.
\end{proof}

The remaining quantities in this eigenvalue estimate include
the crystalline colengths of the row relations and the Selmer
length of the actual main summand. Their sum over the period
fields, and the occurrence of the required projected vectors
in the integral period pairing, are needed for the height
comparison in Section~\ref{sec:missing}.
\clearpage

\subsection{Exact tame rows for the character determinants}
\label{sec:tamerows725}

The factor \(t\sum_w v(o_w(1))\) in Theorem~\ref{thm:ribettrace720}
was obtained by an annihilating polynomial for every representation
at each old place. The arithmetic Frobenius at \(w\) yields a row for
all the character determinants at once. Its value in degree \(e\) is
\(q_w^{e+1}-1\), where \(q_w\) is the residue cardinality. This also
applies when an auxiliary-level representation has nonzero monodromy
at an old place.

\begin{lemma}[Tame inertia in a scalar cocycle]\leavevmode
\label{lem:scalartame725}
Let \(K\) be a nonarchimedean local field of residue cardinality
\(q\), with residue characteristic \(p\ne\ell\). Let \(C\) be a
finite commutative \(\ZZ_\ell\)-algebra and let \(\chi:G_K\to C^\times\)
be a continuous unramified character. Give a finite \(C\)-module
\(M\) the action by \(\chi\). If \(\delta\) lifts a generator of the
pro-\(\ell\) tame inertia quotient and \(\phi\) lifts arithmetic
Frobenius, the possible values \(z(\delta)\) of continuous cocycles
\(z:G_K\to M\) form exactly \(\ker(q-\chi(\phi):M\to M)\).
\end{lemma}
\begin{proof}
Inertia acts trivially, so restriction is a continuous homomorphism.
Wild inertia is pro-\(p\) and the prime-to-\(\ell\) part of tame inertia
also has zero image. In the remaining quotient Frobenius acts by the
\(q\)-th power. The cocycle identity gives
\(z(\phi\delta\phi^{-1})=\chi(\phi)z(\delta)\), whereas its restriction
to inertia gives \(qz(\delta)\).

Conversely write the quotient by wild inertia and the prime-to-\(\ell\)
part of tame inertia as \(\ZZ_\ell(1)\rtimes\widehat{\ZZ}\), using the
chosen Frobenius lift. The section \(\widehat{\ZZ}\to G\) in this tame
quotient is continuous because taking powers of the chosen element
extends from \(\ZZ\) to its profinite completion. Choose the generator
of \(\ZZ_\ell(1)\) induced by \(\delta\). Send its power \(b\) to \(b m\),
where \((q-\chi(\phi))m=0\), and give the Frobenius section value zero.
This defines a continuous cocycle since the homomorphism on inertia
intertwines conjugation with \(\chi\) on the integer powers of Frobenius
and hence on their closure.
\end{proof}

\begin{theorem}[A single group-algebra row at each old place]\leavevmode
\label{thm:jointtamerow725}
Keep the row module \(\mathcal R\), its \(k\) coordinates and the old-place
indices \(W\) from Section~\ref{sec:relationmatrices719}. At every old
place \(w\nmid\ell\) there is a row \(r_w\in\mathcal R\) with entries in
the constant group algebra \(\mathcal O[G_E]\) for which, simultaneously
for \(0\le e<t\),
\[
                     \chi_e(r_w)=(q_w^{e+1}-1)e_w .
\]
Here \(e_w\) is the standard row on the coordinate for \(\delta_w\).
The assertion requires only the finite integral family \(u_x\) used to
define the coordinates. Its members may have arbitrary inertia at \(w\).
\end{theorem}
\begin{proof}
Let \(A\) be the image of \(\mathcal O[G_E]\) in \(\prod_{e<t}\mathcal O\)
by the simultaneous evaluation \(\Xi=(\chi_e)_e\). This is an
\(\mathcal O\)-finite commutative ring. Define \(\mathcal R^0\) to be
the left \(\mathcal O[G_E]\)-module generated by the constant rows
\[
 v(g_i)-e_i,\qquad a e_i,\qquad v(gh)-v(g)-g v(h)
                       \quad(g,h\in G_E),
\]
where \(a=\ell^\nu\). They belong to \(\mathcal R\) by
Lemma~\ref{lem:readrows719}. Their joint image, and the image of their
span, are taken coordinatewise. Since \(\mathcal O[G_E]\) maps onto
\(A\), the set of joint images of rows in \(\mathcal R^0\) already is
an \(A\)-submodule \(N\) of \(A^k\). In the quotient \(M=A^k/N\) write
\(Y(g)\) for the class of \(v(g)\). The displayed rows give
\[
          Y(gh)=Y(g)+\Xi(g)Y(h),\qquad Y(g_i)=e_i\bmod N .
\]
Also \(aM=0\), so both \(M\) and \(A/aA\) are finite. The map \(Y\)
is continuous by the continuity of the original coordinate functions.
The character \(g\mapsto\Xi(g)\bmod aA\) is continuous as well.
Indeed \(A\) is a finite \(\mathcal O\)-submodule of
\(\prod_e\mathcal O\), with its subspace topology by elementary divisors.
Its unit-valued action on \(M\) is unramified at \(w\).

Apply Lemma~\ref{lem:scalartame725} with \(C=A/aA\) to this actual
cocycle. The result says that \((q_w-\Xi(\phi_w))e_w\in N\), hence
this row is the joint evaluation of an element of \(\mathcal R^0\).
Multiply a lifting row on the left by \(\phi_w^{-1}\). Since
\(\kappa(\phi_w)=q_w\), its evaluation on \(\chi_e\) is exactly
\((q_w^{e+1}-1)e_w\) as an element of \(\mathcal O^k\). The equality
holds integrally because \(N\) was defined as an image in \(A^k\).
\end{proof}

\begin{corollary}[The contribution of the two original level primes]
\leavevmode\label{cor:levelcost725}
In Theorem~\ref{thm:ribettrace720} one may choose any subset \(W_0\)
of old places away from \(\ell\), omit the hypotheses on their
polynomials and replace their contribution \(t\sum_{w\in W_0}v(o_w(1))\)
by
\[
                  \sum_{w\in W_0}\sum_{a=1}^t v(q_w^a-1).
\]
At the nodes of Corollary~\ref{cor:actualmarks714}, take \(W_0\) to
consist of the places over both original level primes \(p_0,p_1\)
different from \(\ell\), each listed once. The displayed sum is at most
\(2v(\ell)v_\ell(G_0)\). Its sum over coefficient primes, after division
by \(v(\ell)\) and multiplication by \(\log\ell\), is at most
\(2\log G_0\). These bounds are independent of the auxiliary primes
and the precision.
\end{corollary}
\begin{proof}
In the proof of Theorem~\ref{thm:charrows719}, put \(r_w\) in the row
for \(w\in W_0\) of every square matrix used in the final interpolation.
All earlier eliminations discard the old-place coordinates and use
only the fact that zero tame value forces a scalar cocycle to be
unramified there. They therefore apply without change. The \(W\)-square
block remains diagonal at every \(\chi_e\), and all its entries to the
right vanish. At that character the determinant now acquires the factor
\(\prod_{w\in W_0}(q_w^{e+1}-1)\), together with the previous factors at
the other old places. Their product over \(e\) gives the asserted
replacement in the trace multiplier. The rows on the main summand in
Proposition~\ref{prop:blockrows719} were constructed before introducing
any of the polynomials \(o_w\).

Both \(p_i\) split completely in \(E\), giving \(2f\) places of norm
\(p_i\) for \(f=[E^+:\QQ]\). Write \(d_*=rt_0\ge t\), as in
Corollary~\ref{cor:actualmarks714}. For each \(p_i\ne\ell\) the product
\(\prod_{a=1}^t(p_i^a-1)\) divides
\(\lvert\mathrm{GL}_{d_*}(\mathbb F_{p_i})\rvert\), by the ordered-basis
formula. Hence the product over the chosen \(w\) divides
\(\prod_i\lvert\mathrm{GL}_{d_*}(\mathbb F_{p_i})\rvert^{2f}=G_0^2\).
Prime factorization gives both assertions.
\end{proof}

\begin{proposition}[The local group on a marked main summand]\leavevmode
\label{prop:maintamagawa725}
At an old place \(w\nmid\ell\), suppose the main lattice \(T\) has tame
inertia and the action of a full tame generator \(\sigma\) has no
eigenvalue \(1\) on \(V=T\otimes L\). For \(M=V/T\) one has
\[
 H^1(E_w,M)=H^1_{\rm ur}(E_w,M)
          \simeq M^{I_w}/(\phi_w-1)M^{I_w},
 \qquad
 \operatorname{length}_{\mathcal O}H^1(E_w,M)
                   \le v(\det(1-T(\sigma))).
\]
Here \(v\) gives the uniformizer value one. The local finite condition
obtained as the image of \(H^1_f(E_w,V)\) in \(H^1(E_w,M)\) is zero.
\end{proposition}
\begin{proof}
Wild inertia acts trivially on \(M\) and its continuous homomorphisms
to this \(\ell\)-primary group vanish. After restriction to inertia,
a cocycle therefore descends to the procyclic tame quotient. Since
\(T(\sigma)-1\) is invertible on \(V\), it is surjective on \(V/T\).
One can subtract a coboundary to make the cocycle zero on \(\sigma\).
Continuity and density of the integer powers of \(\sigma\) then give
zero throughout tame inertia. This proves \(H^1(I_w,M)=0\).
Inflation-restriction and evaluation on the topological generator
of the unramified quotient now give the displayed isomorphism.
In detail, \(M^{I_w}\) is the kernel of \(T(\sigma)-1\) on \(M\).
Multiplication by \(T(\sigma)-1\) identifies it with the finite module
\(T/(T(\sigma)-1)T\), whose length is the stated valuation.
For a finite module over the procyclic unramified quotient every
prescribed Frobenius value gives a continuous cocycle. If the action
has period \(b\) and the module exponent is \(c\), the cocycle factors
through the quotient of order \(bc\). Coboundaries give exactly
\((\phi_w-1)M^{I_w}\). Finally \(V^{I_w}=0\). The Bloch--Kato finite
condition away from \(\ell\) is unramified cohomology and hence vanishes
on \(V\).
\end{proof}

For example take \(T=\theta R\) of the prescribed squarefree or rank-one
type in Section~\ref{sec:markedmatrix715}, with \(\theta\) unramified at
the two level primes. At every place under consideration that section
excludes \(1\) from the scaled inertia spectrum. An eigenvalue \(\rho\)
from a degree-\(b\) label has order dividing \(p_i^b-1\), \(b\le r\).
Evaluating the geometric sum of that length at \(\rho\) and at \(1\)
shows \(v(1-\rho)\le v(p_i^b-1)\). The rank-one label has order
dividing \(p_i-1\), giving the same assertion. Thus for those places
\(v(\det(1-T(\sigma)))\le\operatorname{rank}(T)\,
v(\lvert\mathrm{GL}_{d_*}(\mathbb F_{p_i})\rvert)\). In
Proposition~\ref{prop:mainselmer724} one can therefore impose the
finite local condition at every such old place at the total cost of at
most \(2\operatorname{rank}(T)v(\ell)v_\ell(G_0)\) in length. Indeed
restriction of the relaxed Selmer group to the direct sum of the local
groups has kernel the group with the finite conditions there.
Proposition~\ref{prop:maintamagawa725} bounds the length of its image.
This gives an estimate at the original level primes and leaves the
global finite Selmer group and the integral theta-period comparison
to be controlled on the same intervals.
\clearpage

\section{Symplectic deformations at strict classes}\label{nav:chapter-symplectic-classes}
\subsection{Strict classes in a four-dimensional symplectic representation}
\label{sec:symplecticstrict429}

For the twists at which two ordinary inertia characters coincide,
we use a different matrix entry. Let \(\mathcal B\) be the
finite set in Corollary~\ref{cor:unboundedprimes428} and put
\[
  n_*=[E_*:\QQ],\qquad
  L_*=2t_*+2+\max(\mathcal B\cup\{1\}),\qquad Y=N_*+L_*+2 .
\]
Every remaining small coefficient prime is at most \(L_*\).
Fix an odd prime \(\ell\le L_*\), \(2\le r<t(x)\), and a
geometric character \(\alpha\) of \(E_x\) of sign \((-1)^r\).
We allow \(\ell\) to divide the order of \(\alpha\). Set
\(\Psi=\alpha\kappa_\ell^r\) and let \(\omega_\ell\) be the
Teichm\"uller lift of the residual cyclotomic character. In
the four-dimensional representation the multiplier will
have cyclotomic exponent \(r+1\).
These are the characters in
\srcloc{Development/AutomorphicArithmetic/AutomorphicVarieties/SymplecticGSp4/DecompSigmaOrd}{57}.

\begin{lemma}[One dyadic character for all the local tests]\leavevmode
\label{lem:commondyadictwist429}
There is a primitive odd character \(\xi\) of conductor \(2^m\),
with \(2^m\le16n_*L_*(Y+1)\), such that for every odd prime
\(p\le Y\) its value at \(p\) has order greater than \(n_*L_*\).
This also holds after reduction in every odd residue
characteristic. The conductor is bounded by a fixed power of
\(V\) on the original period intervals.
\end{lemma}
\begin{proof}
Choose the smallest \(m\) with \(2^m>8n_*L_*(Y+1)\).
The units modulo \(2^m\) are the direct product generated by
\(-1\) and \(5\), of orders two and \(2^{m-2}\). Send \(-1\)
to \(-1\) and \(5\) to a primitive \(2^{m-2}\)-th root.
The resulting odd character is primitive since it is nontrivial
on the kernel of reduction modulo \(2^{m-1}\).
Choose a sign so that \(b=\pm p\equiv1\bmod4\) and let
\(a=v_2(b-1)\). Then \(2^a\le p+1\). Squaring an integer
congruent to one modulo four increases this valuation by exactly
one, since its sum with one has valuation one. Therefore \(b\)
has order \(2^{m-a}\) modulo \(2^m\). Our character is faithful
on that factor. Multiplying its value by \(-1\) preserves this
order, which is greater than two. It is at least
\(2^m/(p+1)>n_*L_*\). Reduction of roots of unity of
two-power order is injective in odd residue characteristic.
The upper bound follows from minimality of \(m\) and the bounds
for \(N_*,n_*,t_*\) on the interval.
\end{proof}

Enlarge \(O_\alpha\) unramified to an integer ring \(O\)
containing the dyadic character values. Write \(J\) for its
fraction field, \(\pi\) for a uniformizer, \(e=v_\pi(\ell)\),
and \(k\) for its residue field. We identify finite characters
of rational ideles with Galois characters in the arithmetic
Frobenius convention. Put
\[
  A=\kappa_\ell\omega_\ell^{-1}\xi,\qquad
  \mu=A\Psi,\qquad d(g)=\operatorname{diag}
           \bigl(\mu(g),\Psi(g),A(g),1\bigr).
\]
On the ordered basis \(e_0,e_1,e_2,e_3\) take the alternating
form with \(\langle e_0,e_3\rangle=\langle e_1,e_2\rangle=1\)
and other independent pairings zero. Both products of paired
diagonal entries equal \(\mu\), so \(d\) preserves the form
up to this multiplier. Let
\[
   T_\alpha=\{\ell\}\cup\{p\ne\ell:\ell\mid|\alpha(I_p)|\}.
\]
All members are odd. The inertia quotient of an abelian
extension of \(\QQ\) at two has two-power order, by local
reciprocity on \(\ZZ_2^\times\).

\begin{proposition}[The six odd roots at the specified local groups]\leavevmode
\label{prop:localsymplecticroots429}
For each \(p\in T_\alpha\) the two pairs of residual
characters of \(d\) with indices \(\{0,2\}\) and \(\{1,3\}\)
can be separated by a Frobenius lift. For each character
\[
           \gamma\in\{A,A^{-1},\Psi/A,A/\Psi,
                                  A\Psi,(A\Psi)^{-1}\}
\]
one has \(H^0(\QQ_p,k(\bar\gamma))=
H^2(\QQ_p,k(\bar\gamma))=0\). Its first cohomology has
dimension one at \(p=\ell\) and zero at the other \(p\in T_\alpha\).
\end{proposition}
\begin{proof}
At \(\ell\) take an arithmetic Frobenius lift \(f_\ell\)
with \(\kappa_\ell(f_\ell)=1\). Such a lift exists since the
local cyclotomic extensions are totally ramified, so inertia
surjects onto the cyclotomic character image. At \(p\ne\ell\)
use any arithmetic Frobenius lift \(f_p\).
The value of \(\bar A\) on \(f_p\) is \(\bar\xi(p)\).
The subgroup generated by \(\bar\Psi(f_p)\) and
\(\bar\kappa_\ell(f_p)\) has order at most \(n_*(\ell-1)\).
Indeed all its elements are roots of unity of order dividing
\(\operatorname{lcm}(n_*,\ell-1)\) in the field \(k\).
The lemma makes \(\bar A(f_p)\) lie outside that subgroup.
Here \(p\le Y\), since \(p=\ell\) or \(p\mid N_*\).
Writing \(a=\bar A(f_p)\), \(b=\bar\Psi(f_p)\), the two
pairs of eigenvalues are \(\{ab,a\}\) and \(\{b,1\}\).
An intersection would give \(a\in\{1,b,b^{-1}\}\).

Each \(\bar\gamma\) takes a value at \(f_p\) involving
\(a^{\pm1}\) times a power of \(b\). Both it and its Tate
dual consequently have nontrivial Frobenius value. The
invariant spaces are zero, and local duality proves the
assertion in degree two. The local Euler characteristic
formula gives the indicated dimensions in degree one. The
six spaces are precisely the odd root lines for the
adjoint action on the symplectic Lie algebra. Indeed the
upper entries have ratios \(A,\Psi,\Psi/A,A\Psi\), with
entries reflected across the anti-diagonal paired by the
form. The opposite entries have the inverse ratios.
At complex conjugation \(A=-1\) and \(\Psi=1\).
\end{proof}

Write \(D=(J/O)(\theta)\), \(\theta=\Psi^{-1}\), and let
\(D_h=D[\pi^h]\). Use \(S=S^{N_*}(\alpha,r)\) of
Theorem~\ref{thm:levelstrict345} with these coefficients.
Its \(O\)-length is unchanged by the unramified extension
from \(O_\alpha\), by tensoring a composition series. Henceforth
suppose \(\bar\Psi\ne1\) on \(G_\QQ\). Choose \(g_0\) with
both \(\bar\Psi(g_0)\ne1\) and \(\bar A(g_0)\ne1\). Each
kernel is proper and two proper subgroups cannot cover a
group. If neither contains the other, the product of one
element exclusive to each is in neither.

\begin{lemma}[Local splittings at a bounded additional precision]\leavevmode
\label{lem:polarizedwitness429}
Normalize cocycles by \(z(g_0)=0\). Let \(Z_h\) be the
module of continuous normalized \(\theta\)-cocycles with
values in \(D_h\) which are coboundaries on \(G_{\QQ_p}\)
for \(p\in T_\alpha\) and on inertia at every other finite
prime, all with witnesses in \(D_h\). In the normalized
divisible coefficient module put
\[
 C_\alpha=\max_{p\in T_\alpha}\min_{g\in G_{\QQ_p}}
                     v_\pi(\theta(g)-1),\qquad v_\pi(0)=+\infty .
\]
Then \(C_\alpha\le e(1+v_\ell(n_*r))\), and for every \(h\ge1\)
\[
       \pi^{C_\alpha} S[\pi^h]\ \subseteq\ Z_h\
              \subseteq\ S[\pi^h]\ \subseteq\ Z_{h+C_\alpha}.
\]
\end{lemma}
\begin{proof}
Subtracting the coboundary with witness
\(z(g_0)/(\theta(g_0)-1)\) normalizes a cocycle. Its
denominator is a unit, so this operation works in every
coefficient module in question and kills exactly the
coboundaries. In particular a divisible class killed by
\(\pi^h\) has a unique normalized representative, whose
values are all killed by \(\pi^h\).

At \(\ell\) choose an inertia element with cyclotomic value
\(1+\ell\) and raise it to the \(n_*\)-th power. This kills
its finite character and gives valuation
\(v_\pi((1+\ell)^{-rn_*}-1)=e(1+v_\ell(rn_*))\), by the
binomial formula at an odd prime. At \(p\in T_\alpha\)
different from \(\ell\), the inertia image of \(\alpha\)
contains a primitive \(\ell\)-th root of unity. The value
of \(\theta-1\) on an element mapping to such a root has
valuation \(e/(\ell-1)\). This proves the uniform bound.

If \(z\) represents a class of \(S[\pi^h]\), its restriction
at \(p\in T_\alpha\) is \((\theta-1)y_p\) for some \(y_p\in D\),
since these primes are \(\ell\) or divide \(N_*\). Evaluate at
an element attaining the finite minimum there. Multiplication
by \(\pi^h(\theta(g)-1)\) kills \(y_p\), hence
\(y_p\in D_{h+C_\alpha}\). At the other finite primes the
inertia action is through a group of order prime to \(\ell\).
If the character on inertia is trivial then the restricted
coboundary is zero. Otherwise a unit difference \(\theta(g)-1\)
makes its witness already lie in \(D_h\). Multiplying every
chosen witness by \(\pi^{C_\alpha}\) puts it in \(D_h\), which
proves the first and last inclusions. For the middle inclusion
the only further issue is the strict condition on the full
decomposition group at \(p\mid N_*\) outside \(T_\alpha\).
There the unramified subgroup for \(D\) vanishes. Its local
formula is \(D^{I_p}/(f_p-1)D^{I_p}\). With nontrivial
prime-to-\(\ell\) inertia its numerator is zero. With trivial
inertia the numerator is divisible and
\(\alpha(f_p)^{-1}p^{-r}-1\ne0\) acts surjectively.
Thus the image of the finite cocycle satisfies that strict
condition as well.
\end{proof}

\begin{theorem}[The symplectic entry carrying every strict class]\leavevmode
\label{thm:symplecticclasses429}
In the ring \(R_h=O\oplus D_h\), take \(D_h\) as a
square-zero ideal and use its underlying \(O\)-module as
coefficients for entries. Write \(E_{ij}\) for an elementary
matrix and \(X=E_{20}+E_{31}\). The rule
\[
                   z\ \longmapsto\
           \rho_z(g)=(1+z(g)X)d(g)
\]
identifies \(Z_h\) with the lifts of this form and multiplier
\(\mu\) satisfying the following conditions. The value at
\(g_0\) is \(d(g_0)\). At \(p\in T_\alpha\) the local lift
is conjugate to \(d\) by some \(1+y_pX\), \(y_p\in D_h\).
At all other finite primes the lift kills the kernel of
\(\bar d\) on inertia. Every such local conjugate at
\(p\in T_\alpha\) lies on an integral family through \(d\)
with the same ordinary or finite inertia condition.
The union of these lifts, under the inclusions of their
coefficient ideals, realizes exactly \(S\).
\end{theorem}
\begin{proof}
The image of \(X\) is isotropic, \(X^2=0\), and
\(\langle Xv,w\rangle+\langle v,Xw\rangle=0\). For the last
identity the only nontrivial pair to check is \(v=e_0,w=e_1\),
giving \(-1+1=0\), together with its reverse. Thus
\(1+yX\) is symplectic and its inverse is \(1-yX\).
Conjugation by \(d(g)\) multiplies \(X\) by
\(\theta(g)\), since the two entry ratios both equal
\(\Psi(g)^{-1}\). Matrix multiplication now gives the
group law for \(\rho_z\) precisely when
\[
                   z(gh)=z(g)+\theta(g)z(h).
\]
The \((2,0)\)-entry divided by \(\mu(g)\) recovers \(z(g)\),
so this correspondence is injective. Conjugating \(d(g)\)
by \(1-yX\) gives the entry \((\theta(g)-1)y\).
This proves the asserted equivalence at \(g_0\) and the
designated decomposition groups.

At every other finite prime the images of all the diagonal
inertia characters have order prime to \(\ell\). Indeed both
cyclotomic characters are unramified there, \(\xi\) has
two-power order, and \(\alpha(I_p)\) has order prime to
\(\ell\). Reduction on this finite diagonal image is
injective. A cocycle which is a coboundary on inertia
therefore vanishes on the kernel of \(\bar d\). Conversely
vanishing on that kernel makes the inertia cocycle descend
to the finite diagonal image. Averaging over this group
gives a coboundary over \(D_h\). This proves the converse in
the characterization as well.

At each \(p\in T_\alpha\) the family
\((1-Y'X)d(g)(1+Y'X)\) over the domain \(O[[Y']]\)
specializes to \(d\) at zero and to the given local lift
at its square-zero witness. At \(\ell\) conjugating the
ordered split flag gives its ordinary flag with exact
characters of cyclotomic exponents \(r+1,r,1,0\).
At \(p\ne\ell\) conjugation preserves the condition of
killing \(I_p\cap\ker\alpha\), since \(\xi\) is unramified
at these odd primes. The matrices preserve the two coordinate
planes indexed by \(\{0,2\}\) and \(\{1,3\}\). Their values
at \(g_0\) and \(g_0^{-1}\) are diagonal, which gives the
off-diagonal normalization for the sum
\(\rho_z(g_0)+\mu(g_0)\rho_z(g_0^{-1})\) in the source.
Conjugation by \(\operatorname{diag}(ab,b,a,1)\) multiplies
\(X\) by \(b^{-1}\), identifying the indicated torus weight.
Finally every class of \(S\) is killed by some power of
\(\pi\), by finiteness in Theorem~\ref{thm:levelstrict345}.
The two depth inclusions in the lemma identify the union
with \(S\). The trace on every one of these lifts still
equals \(\operatorname{tr}d\), since the two varying entries
are off the diagonal.
\end{proof}

\begin{corollary}[Precision at the two collision residues]\leavevmode
\label{cor:collisionprecision429}
At the primes outside \(\mathcal B\cup\{q,d_0\}\), suppose
\(\ell-1\mid2r\) and \(\bar\Psi\ne1\). At
\(r\equiv(\ell-1)/2\bmod\ell-1\) one has
\(Z_h=S[\pi^h]\) at every depth. The same holds at
\(r\equiv0\bmod\ell-1\) if \(\alpha(\ell)\ne1\).
In the other case the required additional precision is
at most \(1+v_\ell(r)\) powers of \(\ell\). These statements
refer to every scalar component on each original \(E_x\).
\end{corollary}
\begin{proof}
Here \(\ell\nmid n_*N_*\), the coefficient extension is
unramified, and \(T_\alpha=\{\ell\}\). At the half residue
the residual inertia character of \(\Psi\) is the nontrivial
quadratic character. At the zero residue the residual
character on a Frobenius lift with \(\kappa_\ell=1\) is
\(\bar\alpha(\ell)\). A nontrivial value of this prime-to-\(\ell\)
character remains nontrivial after reduction. In both cases
the minimum in the lemma is zero. Otherwise \(e=1\) and
\(v_\ell(n_*)=0\), giving the claimed bound. The integral
decomposition of Theorem~\ref{thm:levelstrict345} identifies
these scalar modules with the corresponding field components
after coefficient extension. The construction realizes the
orders at every depth. To bound them using four-dimensional
traces one still needs a comparison with irreducible
deformations of the split point.
\end{proof}


\subsection{Global localization for the six odd symplectic roots}
\label{sec:oddrootsalllevels430}

Keep \(n_*,N_*,L_*,Y\) and the range of \((x,\ell,r,\alpha)\)
in Section~\ref{sec:symplecticstrict429}. We now choose the dyadic
conductor large enough to apply the Bernoulli unit theorem in
Volume~V. This choice works simultaneously for every odd
\(\ell\le L_*\) and every indicated \(\alpha,r,x\). The group on
which the target scalar classes are measured is still the one for
\(E_x\). The additional dyadic character occurs in the other
diagonal entry \(A\). In this subsection \(\bar\Psi\) may be
trivial. Write
\[
 \mathcal R=\{A,A^{-1},\Psi/A,A/\Psi,A\Psi,(A\Psi)^{-1}\}.
\]
We count these lines with their symplectic root indices when two
characters agree modulo a coefficient prime.

\begin{lemma}[One conductor for the six Bernoulli units]\leavevmode
\label{lem:uniformoddroots430}
There is a single odd primitive dyadic character \(\xi\), of
conductor bounded by a fixed power of \(V\), for which
Proposition~\ref{prop:localsymplecticroots429} holds and for
which the following assertion holds at every indicated \(\ell\).
For each \(\Gamma\in\mathcal R\), let \(\delta\) be the
primitive Dirichlet character attached to the Teichm\"uller
lift of \(\bar\Gamma\). Both \(B_{1,\delta}\) and
\(B_{1,\delta^{-1}}\) are units at the chosen prime over \(\ell\).
\end{lemma}
\begin{proof}
Choose \(M=2^m\) to be the smallest power of two strictly larger
than
\[
 \max\{8n_*L_*(Y+1),\,2^{1+v_2(N_*)},\,L_*^5N_* n_*^2\}.
\]
Take the character on \(-1,5\) used in the proof of
Lemma~\ref{lem:commondyadictwist429}, now at conductor \(M\).
Its order at each odd prime \(p\le Y\) exceeds \(n_*L_*\)
by the same valuation calculation. This proves all the local
assertions of Proposition~\ref{prop:localsymplecticroots429}
for the present choice. The quantities in the displayed maximum
are polynomially bounded on the interval and \(M\) is at most
twice that maximum.

Write \(n'=n_*/\ell^{v_\ell(n_*)}\). Take \(D\) to be twice the
multiplicative order of \(\ell\) modulo \(n'\), with \(D=2\)
if \(n'=1\). Thus \(D\le2n_*\), \(Q=\ell^D\equiv1\bmod4\),
and both \(\ell-1\) and \(n'\) divide \(Q-1\). Factor \(\delta\)
by the Chinese remainder theorem as \(\phi\xi'\), with \(\phi\)
primitive of odd conductor \(f\) and \(\xi'\) of dyadic conductor.
The residual character of \(A\) is \(\bar\xi\), and that of
\(\Psi\) is \(\bar\alpha\,\bar\kappa_\ell^r\). The finite
character \(\alpha\) has order dividing \(n_*\) and conductor
dividing \(N_*\). Its reduction has order dividing \(n'\).
Consequently \(f\mid\ell N_*\), and the order of \(\phi\) divides
\(\operatorname{lcm}(n',\ell-1)\), hence divides \(Q-1\).
On \(1+M/2\) the dyadic part of \(\alpha\) is trivial, whereas
each of \(\xi,\xi^{-1}\) is \(-1\). Thus \(\xi'\) is primitive
of conductor \(M\). Moreover \(\delta(-1)=-1\), since
\(\Psi\) is even and \(A\) is odd.

Writing \(e_2=v_2(Q-1)\), repeated squaring gives
\[
 e_2=v_2(\ell-1)+v_2(\ell+1)+v_2(D)-1,\qquad
 2^{e_2}\le(\ell^2-1)D/2\le L_*^2n_*.
\]
Indeed the odd factor of \(D\) leaves the valuation unchanged,
the first squaring contributes \(v_2(\ell+1)\), and each
subsequent squaring contributes one. Hence
\(f4^{e_2}\le L_*^5N_*n_*^2<M\).
Theorem~\ref{thm:fourier145}, including its proof of integrality
when \(\ell\mid f\), applies to these primitive factors and to
their inverses.
\end{proof}

Let \(O\) contain the values as in the preceding subsection, let
\(\pi\) be a uniformizer, and put \(O_h=O/\pi^h\) and
\(M_h=O_h(\Gamma)\) for \(h\ge1\). Let \(T=T_\alpha\).
Define \(S_h^0(\Gamma)\subset H^1(\QQ,M_h)\) by imposing zero
restriction to \(G_{\QQ_p}\) for \(p\in T\), and zero restriction
to \(I_p\) for every finite \(p\notin T\). Let \(S_h^T(\Gamma)\)
use the same inertia conditions and allow arbitrary local
classes at \(T\). All cohomology is continuous.

\begin{proposition}[Vanishing of the strict odd root groups at every depth]\leavevmode
\label{prop:oddstrictalllevels430}
For \(\xi\) as above, \(S_h^0(\Gamma)=0\) for all six root
lines and all \(h\ge1\).
\end{proposition}
\begin{proof}
For \(h=1\) a strict cocycle vanishes on
\(I_p\cap\ker\bar\Gamma\) at every finite prime. Indeed on
inertia it is a coboundary, which is zero on that kernel.
Corollary~\ref{cor:coc145} applies by the preceding lemma.
In that corollary the cocycle restricts over
\(\QQ(\zeta_{fM})\) to an unramified class group character.
The integral Stickelberger annihilator acts on its image by a
unit, using exactly the Bernoulli unit just established.
The resulting class on the abelian cyclotomic Galois group
vanishes by commuting complex conjugation with an arbitrary
element. This proves the assertion in the residue field.

Suppose \(h>1\) and subtract a lifted global coboundary from
a strict cocycle so that its values lie in \(\pi M_h\).
Identify this submodule with \(M_{h-1}\) by multiplication
by \(\pi\). At \(p\in T\) write the original, adjusted cocycle
as \(g\mapsto(\Gamma(g)-1)y_p\) on the decomposition group.
There is a \(g\) there with \(\bar\Gamma(g)\ne1\), by
Proposition~\ref{prop:localsymplecticroots429}. Its value
forces \(y_p\in\pi M_h\). At \(p\notin T\) the inertia image
of \(\Gamma\) is finite of order prime to \(\ell\). A nontrivial
inertia image therefore supplies the same unit difference.
For a trivial image the cocycle vanishes on inertia already.
The cocycle with values in \(M_{h-1}\) consequently satisfies
all the strict conditions. Induction makes it a coboundary,
which proves the result in \(M_h\).
\end{proof}

\begin{theorem}[Global localization on an odd root line]\leavevmode
\label{thm:oddrootlocalization430}
For every \(\Gamma\in\mathcal R\), localization induces an
isomorphism of \(O_h\)-modules
\[
              S_h^T(\Gamma)\ \xrightarrow{\ \sim\ }\
                  H^1(\QQ_\ell,M_h)\simeq O_h .
\]
The localization isomorphisms commute with reduction of the
coefficients through all depths. The Selmer group of the finite
Tate dual with the local conditions orthogonal to those for
\(S_h^0(\Gamma)\) is zero.
\end{theorem}
\begin{proof}
Write \(M_h^*=\operatorname{Hom}(M_h,\mu_{\ell^\infty})\) for
the finite Tate dual, with its usual contragredient Galois
action. Its residual character on a simple constituent is
\(\bar\kappa_\ell\bar\Gamma^{-1}\). Both \(M_h\) and \(M_h^*\)
have zero invariants at each \(p\in T\), since a Frobenius
lift acts on each residual line by a scalar different from
one. Their global invariant groups vanish as well. At the real
place complex conjugation acts on \(M_h\) by \(-1\), and
\(\ell\) is odd.

Use the finite Poitou--Tate size formula as in
Proposition~\ref{prop:strictclasses365}. For local conditions
\(L_v\), the ratio of the orders of the two mutually dual
Selmer groups is the ratio of global invariant orders multiplied
by \(\prod_v |L_v|/|H^0(\QQ_v,M_h)|\), including the real
place. One may use any finite set of places containing
\(2,\ell,\infty\) and all conductor primes, with unramified
conditions elsewhere. In the strict problem the factor at
each \(p\in T\) is one. At every other finite place the
unramified subgroup is the cokernel of Frobenius minus one
on \(M_h^{I_p}\), and the local invariants are its kernel.
They have equal orders. The real factor is also one. The
preceding proposition therefore gives the trivial group for
the orthogonal conditions on \(M_h^*\).

The exact sequence comparing two choices of local conditions
in Poitou--Tate duality now gives
\[
 0\longrightarrow S_h^0(\Gamma)\longrightarrow S_h^T(\Gamma)
   \longrightarrow \bigoplus_{p\in T}H^1(\QQ_p,M_h)
   \longrightarrow
      \operatorname{Sel}_{(L^0)^\perp}(M_h^*)^\vee ,
\]
where the last symbol denotes the Pontryagin dual for the
strict choice \(L^0\). Both outside terms vanish. By local
duality \(H^2(\QQ_p,M_h)=0\) at every \(p\in T\).
The local Euler characteristic formula gives
\(|H^1(\QQ_p,M_h)|=1\) for \(p\ne\ell\) in \(T\) and
\(|H^1(\QQ_\ell,M_h)|=|O_h|\).

To determine the module structure at \(\ell\), note that
reduction from level \(h\) to \(h-1\) is surjective in
degree one because \(H^2(\QQ_\ell,M_1)=0\). The sequence
\(0\to M_{h-1}\xrightarrow{\pi}M_h\to M_1\to0\)
then identifies the kernel of reduction to \(H^1(\QQ_\ell,M_1)\)
with \(\pi H^1(\QQ_\ell,M_h)\). Here \(h>1\), and the last
reduction is surjective because \(H^2(\QQ_\ell,M_{h-1})=0\).
At level one the local group is one-dimensional.
Nakayama's lemma gives a single generator at level \(h\),
and the equality of orders makes it a free generator over
\(O_h\). Localization commutes with the maps of coefficient
modules, proving compatibility.
\end{proof}

\begin{corollary}[The odd part of the ordinary first-order problem]\leavevmode
\label{cor:ordinaryoddroots430}
At the split point \(d\) over each \(O_h\), consider fixed
multiplier first-order symplectic deformations modulo
conjugation by matrices reducing to the identity modulo the
square-zero parameter. Require the inertia restrictions away
from \(T\) and the restrictions at \(T\setminus\{\ell\}\)
to be conjugate to those of \(d\). Require a lift at \(\ell\)
of the displayed ordered symplectic flag, with its diagonal
characters fixed. On the six odd root lines the resulting
module is free of rank three over \(O_h\). Its three
coordinates are the upper triangular odd root classes at
\(\ell\).
\end{corollary}
\begin{proof}
Over \(O_h[\epsilon]/(\epsilon^2)\) write the lift in the
form \((1+\epsilon Z(g))d(g)\). The fixed multiplier condition
places \(Z(g)\) in the symplectic Lie algebra. Multiplication
gives the cocycle law for the adjoint action and conjugation
changes \(Z\) by a coboundary. The diagonal representation
preserves each matrix root line separately, also when some
residual characters agree. A lifted symplectic flag can be
put back in the original position by a symplectic basis change
reducing to the identity. Indeed lift an adapted basis for its
isotropic plane, choose dual partners for the perfect pairing,
and add elements of the plane to the partners to make their
mutual pairing zero. These operations to correct the standard
pairings can be made in the square-zero ideal. In that basis
the lower entries vanish. Conversely any upper root cocycle
preserves the original flag with the prescribed diagonal, and
coboundaries account for a change of flag position.
At \(\ell\) the allowed lower classes are thus zero and the
upper classes are unrestricted. The other primes in \(T\)
have zero first cohomology on every odd root. The localization
theorem gives one free coordinate for each of the three
upper odd roots and zero for each opposite root. This calculation
leaves the even root with character \(\Psi^{-1}\), which
carries the strict group of Theorem~\ref{thm:symplecticclasses429}.
\end{proof}

The residual vanishing condition used here also appears in
\srcloc{Development/AutomorphicArithmetic/AutomorphicVarieties/SymplecticGSp4/XvanExt}{27}.


\subsection{Extensions preserving a Lagrangian plane at a strict class}
\label{sec:strictsiegelplanes431}

Keep the choice of \(\xi\) in Lemma~\ref{lem:uniformoddroots430}.
Fix \(h\ge1\), put \(R=O/\pi^h\), and take any
\(z\in Z_h\) from Lemma~\ref{lem:polarizedwitness429}.
We use \(\bar\Psi\ne1\) as in that lemma. Identify \(D_h\)
with \(R(\Psi^{-1})\) by \(t\mapsto t/\pi^h\bmod O\) from
\(R\) to \(D_h\). In the following matrices \(z\) denotes
the cocycle in these \(R\)-coordinates. Write
\[
  P=\langle e_0,e_2\rangle,\qquad Q=\langle e_1,e_3\rangle .
\]
They are complementary Lagrangian planes for the alternating
form of Section~\ref{sec:symplecticstrict429}. The Siegel
parabolic of either plane is its stabilizer in the symplectic
similitude group.

\begin{lemma}[A representation on the two paired planes]\leavevmode
\label{lem:strictblock431}
The formula \(\rho(g)=(1+z(g)X)d(g)\), with
\(X=E_{20}+E_{31}\), defines a representation over \(R\)
with multiplier \(\mu\). It preserves both \(P\) and \(Q\).
On each \(G_{\QQ_p}\) for \(p\in T_\alpha\), and on each
\(I_p\) otherwise, an element preserving both planes
conjugates it to \(d\). With the indicated ordered bases
the two blocks are
\[
  \rho_P(g)=\begin{pmatrix}\mu(g)&0\\z(g)\mu(g)&A(g)\end{pmatrix},
  \qquad
  \rho_Q(g)=\begin{pmatrix}\Psi(g)&0\\z(g)\Psi(g)&1\end{pmatrix}.
\]
\end{lemma}
\begin{proof}
For arbitrary \(t,u\in R\) one has \((1+tX)(1+uX)=1+(t+u)X\)
because \(X^2=0\). The image of \(X\) is isotropic and \(X\)
is skew for the alternating form, so \(1+tX\) is symplectic
for every \(t\). Conjugation by \(d(g)\) acts on \(X\) by
\(\Psi(g)^{-1}\). These identities and the cocycle law give
the representation over \(R\) without a restriction on the
products of the values \(z(g)\). The block formulas follow on
the four basis vectors. Where \(z(g)=(\Psi(g)^{-1}-1)y\), one
has \(\rho(g)=(1-yX)d(g)(1+yX)\). The local witnesses in the
definition of \(Z_h\) give this equality over \(R\) at every
specified group.
\end{proof}

Let \(U_P\) be the additive group of skew endomorphisms that
kill \(P\) and map \(Q\) into \(P\). Interchanging \(P,Q\)
defines \(U_Q\). Here skew means
\(\langle tv,w\rangle+\langle v,tw\rangle=0\).
Each \(U\) has the Galois action by \(\rho\)-conjugation.
Write \(S^0(U)\) and \(S^T(U)\) for the strict and relaxed
groups defined before Proposition~\ref{prop:oddstrictalllevels430},
using these finite coefficients and \(T=T_\alpha\).

\begin{proposition}[Localization over an arbitrary strict class]\leavevmode
\label{prop:siegellocalization431}
For either \(U=U_P,U_Q\),
\[
          S^0(U)=0,\qquad
          S^T(U)\xrightarrow{\ \sim\ }H^1(\QQ_\ell,U)
                      \simeq R^3 .
\]
The middle map is localization. At each \(p\in T\setminus\{\ell\}\)
one has \(H^1(\QQ_p,U)=0\).
\end{proposition}
\begin{proof}
The pairing identifies \(Q\) equivariantly with
\(P^\vee\otimes\mu\). Under this identification the skew
condition makes \(U_P\) the symmetric tensors
\(\operatorname{Sym}^2(P)\otimes\mu^{-1}\).
In fact a map \(Q\to P\) belongs to \(U_P\) precisely when
the associated bilinear form on \(Q\) is symmetric. Here two
is invertible in \(R\), so symmetric tensors and the symmetric
square identify. The same reasoning gives
\(U_Q=\operatorname{Sym}^2(Q)\otimes\mu^{-1}\).
The line \(Re_2\subset P\) is stable and the quotient is free
on the image of \(e_0\). In the symmetric square, take first
the square of this line and then all products with this line.
This is a stable filtration, with three free rank-one successive
quotients. Their characters on \(U_P\), in order, are
\[
              A/\Psi,\qquad A,\qquad A\Psi .
\]
On \(Q\) use \(Re_3\) in the same construction. The
constituents on \(U_Q\) have the three inverse characters
in reverse order.

We show that strict vanishing passes through this filtration.
In \(0\to U'\to U''\to B\to0\) at one step, \(B\) has one
of the six characters. Project a strict global cocycle on \(U''\)
to \(B\). Proposition~\ref{prop:oddstrictalllevels430} makes
the projection a coboundary. Lifting its witness to \(U''\)
and subtracting leaves a cocycle valued in \(U'\).
At \(p\in T\) the projection to \(B\) of a local witness for
this cocycle is invariant, hence zero by the Frobenius unit
difference in Proposition~\ref{prop:localsymplecticroots429}.
At other primes the action on inertia factors through a finite
group of order prime to \(\ell\). This follows from the local
conjugacy of the block representation to the finite diagonal
inertia image. A projected local witness is inertia invariant.
Lift it to \(U''\), average its lift over the finite inertia
image, and subtract the invariant lift from the witness.
In either case one obtains a witness in \(U'\). Induction
up the filtration proves \(S^0(U)=0\).

At every \(p\in T\) the local block conjugation gives a direct
sum of the three indicated characters on \(U\). The same is
true of their local Tate duals. Both invariant groups vanish
there by Proposition~\ref{prop:localsymplecticroots429}.
Complex conjugation acts on each constituent of \(U\) by
\(-1\), so its invariant group vanishes, by the filtration.
In the finite Poitou--Tate size formula for the strict problem
all factors are therefore one, as in
Theorem~\ref{thm:oddrootlocalization430}. In particular the
dual group for the orthogonal conditions is zero.
Comparison with the relaxed local conditions gives an
isomorphism onto \(\bigoplus_{p\in T}H^1(\QQ_p,U)\).
By the local direct sum and the local calculations in that
theorem, the only nonzero term is at \(\ell\), where it is
free of rank three over \(R\).
\end{proof}

At \(\ell\) choose a witness \(y_\ell\) and transport the
displayed complete flag for \(d\) by \(1-y_\ell X\), obtaining
a flag \(\mathcal F_\ell\) for \(\rho\). For \(U=U_P\) or
\(U_Q\) let \(1+U\) denote its subgroup of the symplectic group.
Consider representations with multiplier \(\mu\) preserving the
specified plane, with the induced action on that plane and its
quotient fixed as for \(\rho\). Two are equivalent when
conjugate by \(1+U\). Impose the following local conditions.
On inertia outside \(T\) and on the full local groups in
\(T\setminus\{\ell\}\), require conjugacy to \(\rho\) by
\(1+U\). At \(\ell\) require that some \(1+U\) conjugate
preserve \(\mathcal F_\ell\) with its prescribed graded characters.

\begin{theorem}[Ordinary extensions in the two Siegel parabolics]\leavevmode
\label{thm:siegelclasses431}
For the plane \(P\) the classes of these representations form a
free \(R\)-module of rank two, with origin \(\rho\).
For \(Q\) they form a free \(R\)-module of rank one.
Localization identifies them with the allowed upper triangular
classes at \(\ell\) in each of the two unipotent radicals.
These statements hold for every \(z\in Z_h\) on the original
field and interval data.
\end{theorem}
\begin{proof}
In either parabolic, divide a representation value on the right
by \(\rho(g)\). The quotient acts trivially on the plane and
on the quotient by it, and therefore has the form \(1+u(g)\),
where \(u(g)\) kills the plane and maps into it. Its image
is isotropic. The symplectic condition consequently says exactly
that \(u(g)\) is skew. Products of two such endomorphisms
vanish, so the multiplication law of \(1+U\) is addition.
Thus \((1+u(g))\rho(g)\) is a representation exactly when
\(u\) is a cocycle on \(U\). Conjugation by \(1+U\) changes
it by a coboundary. This gives the claimed module of classes
inside \(S^T(U)\) over the actual ring \(R\).

Put the local representation at \(\ell\) back into the
diagonal basis by the chosen block conjugation. It preserves
\(P,Q\), so also preserves each of the two unipotent radicals.
The endomorphisms in \(U_P\) have independent entries at
\((0,3),(0,1),(2,1)\). The skew identity determines the
remaining entry \((2,3)\) from \((0,1)\). The first two
parameters give upper triangular endomorphisms and the third
gives a lower triangular one. In \(U_Q\) the transpose
positions give one upper and two lower parameters. In each
case the upper and lower coordinates are local Galois
submodules in the diagonal basis. A local class preserves
the required flag after the allowed conjugation exactly when
its lower coordinates are coboundaries. Its upper coordinates
then give arbitrary extensions with the same graded characters.

Proposition~\ref{prop:siegellocalization431} identifies
the global choices with these local classes. Each allowed
coordinate contributes one copy of \(R\), by
Theorem~\ref{thm:oddrootlocalization430}. This gives the two
ranks. These constructions use the original \(\Psi\) on the
same period fields because the dyadic twist occurs only in \(A\).
\end{proof}

\begin{corollary}[Characteristic polynomials throughout both parabolics]\leavevmode
\label{cor:siegeltrace431}
For every \(g\in G_\QQ\), every representation in either
module of Theorem~\ref{thm:siegelclasses431}, and every
admissible \(z\), its characteristic polynomial at \(g\) is
\[
    (t-\mu(g))(t-\Psi(g))(t-A(g))(t-1)\quad\hbox{in }R[t].
\]
A class in \(S[\pi^a]\) can be represented in this calculation
at \(h=a+C_\alpha\) for every \(a\ge1\).
\end{corollary}
\begin{proof}
Order a basis by the invariant plane and the quotient. The
resulting block triangular matrix has the two diagonal blocks
of the lemma, each triangular in its indicated basis.
Taking determinants proves the formula over \(R\). The
precision assertion is the inclusion
\(S[\pi^a]\subseteq Z_{a+C_\alpha}\) of
Lemma~\ref{lem:polarizedwitness429}. In particular all
characteristic polynomials on these two families agree at the
same depth, including at Frobenius elements. Estimating the
size of the strict group through traces still requires control
of representations outside the two parabolics.
\end{proof}


\subsection{A common Frobenius test for the rank-four representations}
\label{sec:symplecticfrobenius433}

Use \(A,\Psi,\mu=A\Psi\) and the conductor \(M\) of
Lemma~\ref{lem:uniformoddroots430}. The split diagonal is
\((\mu,\Psi,A,1)\). Its four cyclotomic exponents, in this
order, are \(r+1,r,1,0\). Frobenius is arithmetic, so the
cyclotomic character at an unramified prime \(p\ne\ell\)
has value \(p\).

\begin{proposition}[The determinants of the two parity pairs]\leavevmode
\label{prop:evenplanedeterminants433}
Let \(\rho\) be a continuous rank-four representation of
\(G_\QQ\) over a finite extension of \(\QQ_\ell\), with a stable
lattice, residual semisimplification \(\bar d^{\,\rm ss}\),
and multiplier \(\mu\) for a symplectic pairing. Suppose it is
unramified outside finitely many places. At \(\ell\) assume an
ordinary flag whose ordered diagonal characters reduce to those
of \(d\) on the decomposition group and agree with those of
\(d\) on inertia. If the global semisimplification has a
two-dimensional subrepresentation on which complex conjugation
is scalar, the determinants on it and its complementary
two-dimensional summand, in some order, have the form
\[
              \varepsilon_1\kappa_\ell^{\,r+2},\qquad
              \varepsilon_2\kappa_\ell^{\,r},
\]
with finite characters \(\varepsilon_i\). In particular at
every unramified \(p\ne\ell\), under any complex embedding of
its Frobenius eigenvalues, one has
\[
        \frac{\max_i|\lambda_i(p)|}{\min_i|\lambda_i(p)|}\ \ge\ p .
\]
The lower bound at Frobenius to the power \(s\ge1\) is \(p^s\).
\end{proposition}
\begin{proof}
Enlarge the coefficient field to define the given summand and
its complement. Passing to a stable subspace or quotient
commutes with taking the multiset of residual constituents.
Indeed intersect a stable lattice with the subspace. Its
quotient in the lattice is torsion-free, so reduction preserves
the resulting exact sequence. One can apply this argument
along a composition series, and semisimplified reduction is
independent of the chosen stable lattice by the
Brauer--Nesbitt theorem.

At complex conjugation the residual characters indexed by
\(0,2\) have value \(-1\), and those indexed by \(1,3\)
have value \(1\). As \(\ell\) is odd, the scalar summand must
take both residual characters from just one pair and the other
summand takes the other pair, with multiplicities.
On the decomposition group, the semisimplification of the
local ordinary flag has four one-dimensional constituents.
Their reductions have disjoint sets of values at the single
Frobenius lift of Proposition~\ref{prop:localsymplecticroots429}
for these two pairs. Hence the constituents on the first
of the parity pairs at this local group have indices \(0,2\),
and those on the second have indices \(1,3\). This also follows
by factoring the local Frobenius characteristic polynomial and
reducing its roots. One only needs separation between the pairs.
Determinants multiply in the local semisimplification, so their
inertia exponents are \((r+1)+1=r+2\) and \(r+0=r\), up to
finite characters.

After removing the corresponding cyclotomic power, each global
determinant has finite inertia at \(\ell\). At each other
finite \(p\), its abelian inertia image is finite. Indeed
local reciprocity identifies the abelian inertia group with
\(\ZZ_p^\times\), whose open pro-\(p\) subgroup has finite
image in the units of a finite extension of \(\QQ_\ell\)
when \(p\ne\ell\). Outside a finite set the character is
unramified. By Kronecker--Weber the product of the local unit
groups surjects onto \(G_\QQ^{\rm ab}\), giving a finite
global image. This proves the two determinant formulas.

At a good \(p\), the products of the two pairs of eigenvalues
therefore have absolute values \(p^{r+2}\) and \(p^r\).
Their ratio \(p^2\) is at most the square of the ratio of the
largest to the smallest eigenvalue magnitude. The power
assertion follows by raising the eigenvalues to that power.
The block version of this obstruction occurs in
\srcloc{Development/AutomorphicArithmetic/Foundations/GaloisRepresentations/AdmissibleResidues/PosOfMulPosOfPos}{231}.
\end{proof}

Choose once on the interval a rational prime \(p\) between
\(Y+32\) and \(2(Y+32)\), using Bertrand's postulate.
It lies outside the conductor of every \(\alpha,\xi\) under
consideration and differs from every remaining coefficient
prime \(\ell\le L_*\). Both \(p\) and \(M\) are bounded by
fixed powers of \(V\).

\begin{lemma}[Cuspidal traces at a prime of arbitrary residue degree]\leavevmode
\label{lem:cuspidaltrace433}
Suppose a rank-four unramified Frobenius matrix at \(p\) has
eigenvalues \(\lambda_i\) in a complex realization. Suppose
that, for a positive integer \(s\), the multiset
\(\{\lambda_i^s/p^{s(r+1)/2}\}\) is the unramified Satake
multiset of a cuspidal \(\mathrm{GL}_4\) representation with
unitary central character at a place of norm \(p^s\). Then
\[
 p^{(r+1)/2-9/22}\le|\lambda_i|\le p^{(r+1)/2+9/22}.
\]
In particular the preceding parity pairs cannot occur for
such a representation. Its trace \(u=\sum_i\lambda_i\) satisfies
\[
 \big|u-(1+p\zeta)(1+p^r\eta)\big|>\tfrac12p^{r+1}
\]
for all roots of unity \(\zeta,\eta\) and every \(r\ge2\).
\end{lemma}
\begin{proof}
Theorem~1 of Blomer and Brumley bounds the absolute values of
the real exponents of an unramified unitary cuspidal Satake
parameter on \(\mathrm{GL}_4\) by \(9/22\) over any number
field. At norm \(p^s\) this puts each normalized magnitude
between \((p^s)^{-9/22}\) and \((p^s)^{9/22}\).
Taking \(s\)-th roots gives the claimed estimate. The ratio
of any two magnitudes is at most \(p^{9/11}<p\), so the
determinant obstruction applies. In particular a match after
restriction to an auxiliary number field suffices at any place
above \(p\) with this Satake identity.~\cite{IX-VII-bb433}

Since \(9/22<1/2\), one has \(|u|\le4p^{(r+2)/2}\).
The split term has magnitude at least \((p-1)(p^r-1)\).
As \(r\ge2\) and \(p>32\), the difference, after division
by \(p^{r+1}\), is at least
\((1-1/p)(1-1/p^2)-4/p>1/2\).
\end{proof}

\begin{theorem}[An integer bounding the simultaneous cuspidal trace congruences]\leavevmode
\label{thm:commonsptrace433}
For each integer \(2\le r\le t_*\), take a finite multiset
\(\mathcal C_r\) of algebraic integers, stable under rational
Galois conjugation including multiplicities. Suppose every
element has magnitude at most \(4p^{(r+1)/2+9/22}\).
These multisets may in particular consist of integral traces
whose conjugates have the cuspidal realizations of the lemma.
There is a positive rational integer \(D_{\rm Sp}\) with
\[
\begin{split}
 \log D_{\rm Sp}\ \le\ 
 n_* M L_*^2\sum_{2\le r\le t_*}|\mathcal C_r|\,
         \big((r+1)\log p+\log 8\big).
\end{split}
\]
For every common field, every geometric character \(\alpha\)
under consideration, and every odd \(\ell\le L_*\), fix a
coefficient embedding and normalize \(v_\ell(\ell)=1\).
Then for each allowed \(r\),
\[
 \sum_{u\in\mathcal C_r}v_\ell
          \bigl(u-\operatorname{tr}d(\operatorname{Frob}_p)\bigr)
                     \ \le\ v_\ell(D_{\rm Sp}).
\]
This includes the split trace of every representation in the
two families of Theorem~\ref{thm:siegelclasses431} after
reduction at any depth.
\end{theorem}
\begin{proof}
Put \(m_\ell=\operatorname{lcm}(M,\ell-1)\) and form the
monic polynomial
\[
 B_r(Z)=
 \prod_{\substack{3\le\ell\le L_*\\\ell\ {\rm prime}}}
 \ \prod_{\zeta^{m_\ell}=1}\ \prod_{\eta^{n_*}=1}
                \bigl(Z-(1+p\zeta)(1+p^r\eta)\bigr).
\]
Its coefficients are algebraic integers fixed by the rational
Galois group, hence are integers. Its degree is at most
\(n_* M L_*^2\), since \(m_\ell\le M(\ell-1)\).
On the diagonal \(d\) at \(p\) the trace is
\((1+A(p))(1+\Psi(p))\), where
\(A(p)=p\,\xi(p)\omega_\ell(p)^{-1}\) and
\(\Psi(p)=p^r\alpha(p)\). The two finite factors have orders
dividing \(m_\ell\) and \(n_*\). Thus the trace for every
choice in the theorem occurs as a root of the same \(B_r\).

Take the absolute value of the product
\(\prod_r\prod_{u\in\mathcal C_r}B_r(u)\).
It is an integer by Galois stability. Every factor is nonzero
by the numerical estimate in the proof of the lemma, which
uses only the asserted magnitude bound for \(u\).
Each difference in this product has magnitude at most
\(8p^{r+1}\), by the triangle inequality and \(r\ge2\).
This proves the upper bound for its logarithm. In any
\(\ell\)-adic embedding all factors have nonnegative
valuation. Retaining the one specified root at each
\(u\in\mathcal C_r\) therefore proves the local inequality,
including for a submultiset. The reduction statement follows
from Corollary~\ref{cor:siegeltrace431}.
\end{proof}

\begin{corollary}[The size of spectra at fixed rational rank and level]\leavevmode
\label{cor:fixedlevelcount433}
Let \(\mathcal N\) be a finite list of positive levels. On the
rational \(\mathrm{GL}_4\) locally symmetric spaces of principal
level \(3N\), \(N\in\mathcal N\), consider the simultaneous
Hecke eigensystems in ordinary characteristic-zero cohomology
with irreducible rational algebraic coefficients of dominant
highest weight \(\lambda\), \(|\lambda_i|\le B\).
Even counted with their generalized eigenspace multiplicities
and summed over degrees, their total over this list is at most
\[
               C(B+1)^{10}\sum_{N\in\mathcal N}(3N)^{16}
\]
for an absolute \(C\). In particular the integer in the theorem
has \(\log D_{\rm Sp}\ll V^{C'}\) for integral cuspidal trace
lists obtained from these eigensystems with one normalized
trace per eigensystem per \(r\), if their conjugates meet the
lemma and \(B,|\mathcal N|,\max(\mathcal N\cup\{1\})\) are
polynomially bounded in \(V\).
\end{corollary}
\begin{proof}
Divide out the positive real center at infinity. The level-three
space has finitely many components and a finite cell model by
the Borel--Serre compactification. The principal level-three
subgroup is torsion-free on each component. Indeed a rational
finite-order stabilizer preserves the corresponding integral
lattice and acts trivially on its reduction modulo three, so
the congruence-kernel argument of Theorem~\ref{thm:end14}
applies. The level-\(3N\) cover has at most \((3N)^{16}\)
sheets. The cosets inject into the reduction of integral
four-by-four matrices modulo \(3N\). A lifted cell model
therefore has at most the same factor times the fixed number
of cells at level three.~\cite{IX-VII-V-bs}

The rank of the local system is
\(\prod_{i<j}(\lambda_i-\lambda_j+j-i)/(j-i)\) by the Weyl
dimension formula and is at most \((2B+4)^6\).
Cellular cochains, summed over degrees, have dimension at
most the number of cells times this rank. The simultaneous
generalized eigenspaces fit in cohomology, whose total
dimension is bounded by that of cochains. There are at most
\((2\lfloor B\rfloor+1)^4\) possible highest weights. These
estimates give the formula including eigenspace multiplicities.
Summing over at most \(t_*\) twists in the theorem preserves
a polynomial bound in \(V\).
\end{proof}

The trace estimate applies as soon as an integral spectrum of
the stated size contains the deformations in question. To
deduce the remaining strict length from it one still needs
both that spectral identification and an integral comparison
of the strict group with the sum of the trace valuations.
The source states its latter comparison with hypotheses on
finiteness and rigidity in
\srcloc{Development/AutomorphicArithmetic/AutomorphicVarieties/SymplecticGSp4/GaloisRepIneqOfAllIrreducibleNonvanishing}{152}.


\subsection{The integral equations of the negative symplectic root}
\label{sec:saturatedstrictroot434}

Write \(\theta=\Psi^{-1}\), \(D=(J/O)(\theta)\), and
\(S=S^{N_*}(\alpha,r)\) as in Section~\ref{sec:symplecticstrict429}.
Keep \(\bar\Psi\ne1\), \(T=T_\alpha\), the normalization at
\(g_0\), and \(X=E_{20}+E_{31}\).
The explicit witnesses for a local conjugation in
Lemma~\ref{lem:polarizedwitness429} may require extra
coefficient precision. We determine the integral equations
of that conjugation family and compute the resulting global
ring.

\begin{lemma}[The image of the local conjugation family]\leavevmode
\label{lem:saturatedorbit434}
For \(p\in T\), take \(\sigma_p\in G_{\QQ_p}\) minimizing
\(v_\pi(\theta(g)-1)\) and set
\[
     a_p=\theta(\sigma_p)-1,\qquad
     b_p(g)=(\theta(g)-1)/a_p .
\]
Then \(a_p\ne0\), \(b_p(g)\in O\) and \(b_p(\sigma_p)=1\).
In the matrix family \((1+z(g)X)d(g)\) the scheme-theoretic
image of conjugation by \(1-YX\) is the formal affine line
\[
                  z(g)=b_p(g)t,\qquad t=z(\sigma_p).
\]
In particular every formal local matrix equation vanishing
on these conjugates vanishes identically on this line.
\end{lemma}
\begin{proof}
The finiteness of the minimum follows from
Lemma~\ref{lem:polarizedwitness429}. The ratio \(b_p\) is
continuous and satisfies
\(b_p(gh)=b_p(g)+\theta(g)b_p(h)\), by expanding
\(\theta(gh)-1\). The displayed formula consequently gives
a representation over \(O[[t]]\). Its matrix entries generate
the parameter \(t\), by taking the \((2,0)\)-entry at
\(\sigma_p\) and dividing by \(\mu(\sigma_p)\), a unit.
The actual conjugate \((1-YX)d(g)(1+YX)\) has parameter
\(t=a_pY\). Substitution \(O[[t]]\to O[[Y]]\) by
\(t\mapsto a_pY\) is injective. The coefficient of \(Y^j\)
is the original coefficient of \(t^j\) times the nonzero
element \(a_p^j\) of the domain \(O\). The kernel of the map
from formal matrix coordinates to \(O[[Y]]\) is thus precisely
the kernel for the displayed line. This is the asserted
scheme-theoretic image. In particular its conclusion concerns
equations also over coefficient rings where \(a_p\) is a
zero divisor.
\end{proof}

For a finite length \(O\)-module \(M_0\) write
\(\mathcal Z(M_0)\) for the continuous functions
\(z:G_\QQ\to M_0\) with the cocycle law for \(\theta\),
with \(z(g_0)=0\), and satisfying the local equations
\[
\begin{aligned}
 z(g)&=b_p(g)z(\sigma_p) &&(g\in G_{\QQ_p},\ p\in T),\\
 z(g)&=0 &&(g\in I_p\cap\ker\bar d,\ p\notin T).
\end{aligned}
\]
The second line is imposed at finite places. Give \(M_0\)
the discrete topology. Put \(U=\operatorname{Hom}_O(S,J/O)\)
as an \(O\)-module. It is finite by
Theorem~\ref{thm:levelstrict345}.

\begin{proposition}[The integral coefficient module of the normalized cocycles]\leavevmode
\label{prop:strictrootmodule434}
There is a natural isomorphism for all finite length \(O\)-modules
\[
               \operatorname{Hom}_O(U,M_0)
                        \ \simeq\ \mathcal Z(M_0).
\]
In particular \(\mathcal Z((J/O)[\pi^h])=S[\pi^h]\) through
normalized cocycles, without a shift in \(h\).
\end{proposition}
\begin{proof}
The unique normalized representative \(z_s\) of \(s\in S\)
exists because \(\theta(g_0)-1\) is a unit. Its value at
each \(g\) is \(O\)-linear in \(s\). Forgetting the character
on the coefficient module, define \(x(g)\in U\) by
\(x(g)(s)=z_s(g)\). The function \(x\) is continuous since
\(S\) is finite and the individual representatives are
continuous. Where \(p\in T\), write \(z_s(g)=(\theta(g)-1)y\)
in \(D\). Evaluating at \(\sigma_p\) and multiplying by
\(b_p(g)\) shows that \(x(g)=b_p(g)x(\sigma_p)\) in \(U\).
At every other finite prime the local strict or unramified
condition also implies vanishing on \(I_p\cap\ker\bar d\).
Indeed the diagonal inertia characters there have prime-to-\(\ell\)
image, and reduction is injective on this image.
Thus \(x\in\mathcal Z(U)\).

The elements \(x(g)\) generate \(U\). Otherwise a nonzero
functional on their quotient with values in \(J/O\) would
give a nonzero \(s\in S\) with \(z_s(g)=0\) for all \(g\).
Here the finite \(O\)-linear bidual of \(S\) identifies
canonically with \(S\). This is checked on \(O/\pi^h\)
from \((J/O)[\pi^h]=\pi^{-h}O/O\), and then on direct sums.
Consequently composition with \(x\) is always injective
from the indicated Hom module.

First take \(M_0=(J/O)[\pi^h]\). A cocycle in
\(\mathcal Z(M_0)\) is a coboundary with divisible
coefficients at each \(p\in T\). Indeed choose
\(y\in J/O\) solving \(a_py=z(\sigma_p)\). This gives
\(z(g)=(\theta(g)-1)y\) on \(G_{\QQ_p}\). At the other
finite places vanishing on the indicated kernel makes the
inertia restriction descend to its finite prime-to-\(\ell\)
diagonal image. Averaging there gives a coboundary. At any
additional \(p\mid N_*\) at which the full strict condition
is required, the unramified cohomology with divisible
coefficients vanishes. If the prime-to-\(\ell\) inertia
character of \(\theta\) is nontrivial its invariant group
in \(D\) is zero. Otherwise Frobenius minus one acts
surjectively on \(D\), since the scalar
\(\alpha(p)^{-1}p^{-r}-1\) is nonzero in the unramified
case. The unramified cohomology is the cokernel of this action
on inertia invariants. Hence the global image belongs to
\(S[\pi^h]\). Conversely its normalized representative for
any \(s\in S[\pi^h]\) has values in \(M_0\), by normalization,
and the equations follow from those for \(x\). This proves the
identification at level \(h\), where the Hom isomorphism sends
evaluation at \(s\) to \(z_s\).

Every finite length \(O\)-module is a finite direct sum of
cyclic modules isomorphic to such \(M_0\). Both functors take
these direct sums to direct sums, and composition with \(x\)
commutes with all \(O\)-linear maps. This proves the stated
natural isomorphism independently of any cyclic decomposition.
\end{proof}

\begin{theorem}[The full formal ring in the negative root]\leavevmode
\label{thm:strictrootring434}
In a complete noetherian local \(O\)-algebra \(B\) with
residue field \(k\), consider representations
\((1+z(g)X)d(g)\) with \(z\) valued in \(\mathfrak m_B\).
Use the normalization and both local equations defining
\(\mathcal Z\), interpreted in the maximal-ideal topology.
This functor is represented by the completion \(R_\theta\)
of \(\operatorname{Sym}_O U\) at its ideal generated by
\(\pi\) and \(U\). Its zero section has augmentation ideal
\(I_\theta\) satisfying
\[
       I_\theta/I_\theta^2\cong U,\qquad
       R_\theta[1/\pi]=J .
\]
In particular the elementary divisors of this relative
cotangent module are exactly those of the full strict group.
\end{theorem}
\begin{proof}
The identity \(X^2=0\) and isotropy of its image allow
arbitrary coefficients \(z(g)\) in the maximal ideal of
\(B\). As in Lemma~\ref{lem:strictblock431} the representation
condition is exactly the linear cocycle law. At depth \(n\)
the coefficient module \(\mathfrak m_B/\mathfrak m_B^n\)
has finite \(O\)-length, by its filtration over the residue
field. The proposition identifies all its solutions with
\(O\)-linear maps out of \(U\). Taking the inverse limit
identifies the continuous solutions with maps
\(U\to\mathfrak m_B\). Continuity follows since each finite
quotient solution is continuous, and all equations may be
tested in the finite quotients. The universal property of
the symmetric algebra and completion identifies these maps
with continuous local \(O\)-algebra maps \(R_\theta\to B\).

To compute the two asserted invariants write
\(U=\bigoplus_{i=1}^m O/\pi^{c_i}\) with \(c_i>0\).
The ring just constructed is
\[
     O[[t_1,\ldots,t_m]]/
                (\pi^{c_1}t_1,\ldots,\pi^{c_m}t_m).
\]
The augmentation sends the variables to zero. Keeping terms
of degree one gives \(I_\theta/I_\theta^2=U\) canonically
by the inclusion of the generators. When \(m>0\), the whole
ideal of positive-degree terms is killed by
\(\pi^{\max_i c_i}\), since it is generated by the \(t_i\).
Inverting \(\pi\) gives \(J\), also when \(m=0\).
The cyclic decomposition and finite duality give the stated
elementary divisors.
\end{proof}

\begin{corollary}[The strict dual in a global matrix cotangent module]\leavevmode
\label{cor:saturatedcotangent434}
Take a framed global matrix deformation ring at \(\bar d\)
with the given multiplier. Impose the inertia kernel condition
away from \(T\), and at \(p\in T\) impose closed local matrix
equations which vanish on the formal conjugates
\((1-YX)d(g)(1+YX)\). One may also impose the diagonal value
at \(g_0\), or the weaker off-diagonal normalization in
Theorem~\ref{thm:symplecticclasses429}, and the fixed value
at complex conjugation. With augmentation at \(d\), write
\(I\) for the augmentation ideal of the resulting ring.
There is a surjection
\[
                 I/I^2\ \longrightarrow\ U .
\]
The lower matrix entries at global group elements already
surject onto this module. In particular if a quotient of the
cotangent module through which this map factors has finite
length, its length is at least \(\operatorname{length}_O S\).
\end{corollary}
\begin{proof}
Evaluate the framed matrices at the universal representation
over \(R_\theta\). The local matrices at \(p\in T\) are obtained
by substituting \(t=z(\sigma_p)\) into the local line of
Lemma~\ref{lem:saturatedorbit434}. Thus they satisfy every
one of the stipulated closed local equations. At \(p\notin T\)
the diagonal representation kills the indicated inertia
kernel and the defining equation kills its varying root entry.
The normalization at \(g_0\) is built in. At complex conjugation
\(c\), \(\theta(c)=1\) and \(c^2=1\), so \(2z(c)=0\) and
\(z(c)=0\). This verifies the additional normalizations.
The universal property of the framed matrix ring therefore
gives an augmented homomorphism to \(R_\theta\).

The \((2,0)\)-entry at \(g\), divided by \(\mu(g)\), maps
to the generator \(x(g)\) in the symmetric algebra. These
elements generate \(U\) by the proposition, proving the
surjection on relative cotangent modules and the length
bound, with the whole strict dual as coefficient module.
The corresponding square-zero derivation in the source
appears in
\srcloc{Development/AutomorphicArithmetic/AutomorphicVarieties/SymplecticGSp4/ExistsQuotientDerivationOfShaCocycle}{252}.
\end{proof}

The generic fiber of this explicit quotient consists of the
single split point for every possible list of the \(c_i\).
To upper-bound their sum one must control the relations in
the larger deformation ring comparing this cotangent direction
with its characteristic-zero points. The trace separation in
Theorem~\ref{thm:commonsptrace433} addresses the values at such
points once their spectra are identified.

\section{Galois representations and Frobenius traces}\label{nav:chapter-galois-strict-frobenius}

\subsection{The ordinary symplectic flag at a collision weight}
\label{sec:sympflag468}

Fix a period interval and the single character \(\xi\) of
Lemma~\ref{lem:uniformoddroots430}. Take an indicated prime
\(\ell\ge5\) prime to \(n_*N_*\), an indicated \(r\ge2\) with
\(\ell-1\mid2r\), and the even character \(\Psi=\alpha\kappa_\ell^r\).
Suppose \(\bar\Psi|_{G_{\QQ_\ell}}\ne1\). These hypotheses include
the faithful dyadic half-residue weights of
Section~\ref{sec:finiteexception467} whenever \(\ell\nmid n_*N_*\).
Use the unramified coefficient ring \(O\) containing the character
values, its residue field \(k\), the alternating form and
\[
  (d_0,d_1,d_2,d_3)=(A\Psi,\Psi,A,1),\qquad \mu=A\Psi
\]
from Proposition~\ref{prop:localsymplecticroots429}.
Write \(D_\ell=G_{\QQ_\ell}\) and \(I_\ell\) for its inertia group.
Coefficient rings below are complete Noetherian local \(O\)-algebras
with the given finite residue field. A frame identifies the residual
representation with the displayed diagonal one in the prescribed basis.

\begin{lemma}[A unique integral flag tested on Frobenius eigenlines]\leavevmode
\label{lem:sympflag468}
There is an arithmetic Frobenius lift \(h\in D_\ell\) having four
distinct residual eigenvalues on \(d\). For every framed symplectic
similitude lift \(\rho\) with multiplier \(\mu\), let \(P_i\) be the
spectral projector of \(\rho(h)\) lifting the \(i\)-th coordinate
projector. The following equations characterize existence of a
\(D_\ell\)-stable full flag lifting the ordered coordinate flag,
with character \(d_i|_{I_\ell}\) on its \(i\)-th graded line:
\[
 \begin{aligned}
 P_j\rho(g)P_i&=0 &&(i<j,\ g\in D_\ell),\\
 P_i(\rho(s)-d_i(s)1)P_i&=0 &&(0\le i\le3,\ s\in I_\ell).
 \end{aligned}
\]
The flag is unique and self-dual for the alternating form.
\end{lemma}
\begin{proof}
If \(\bar\Psi\) is unramified, its nontriviality gives
\(\bar\Psi(h)\ne1\) for every Frobenius lift. Otherwise multiply
a lift by an inertia element if necessary to get this inequality.
Put \(b=\bar\Psi(h)\). Its order divides
\(\operatorname{lcm}(n_*,\ell-1)\). The value \(a=\bar A(h)=\bar\xi(\ell)\)
has order greater than \(n_*(\ell-1)\) by the choice of \(\xi\).
The values \(ab,b,a,1\) are therefore distinct.

Hensel factorization of the characteristic polynomial gives unique
roots \(z_i\) lifting the four labeled eigenvalues. Every difference
\(z_i-z_j\) for \(i\ne j\) is a unit. Thus
\[
          P_i=\prod_{j\ne i}\frac{\rho(h)-z_j1}{z_i-z_j}
\]
are orthogonal idempotents with sum one, by Cayley--Hamilton. Their
images are free lines by reduction modulo the maximal ideal. The
multiplier identity conjugates the transpose inverse of \(\rho(h)\),
multiplied by \(\mu(h)\), to \(\rho(h)\). Uniqueness of the labeled
roots implies \(z_{3-i}=\mu(h)/z_i\). The form consequently pairs
each line perfectly with the line numbered \(3-i\), and all other
pairings vanish by the unit differences.

Any \(h\)-stable flag as in the statement is preserved by every
polynomial projector. Each step is the direct sum of its projector
images. These are direct summands of that step, and their ranks are
fixed by reduction to the coordinate flag. In each selected line
the corresponding summand fills the line by Nakayama's lemma.
Thus the successive steps must be the partial sums of the indicated
lines. The pairing makes these steps self-dual. The first set of
equations is exactly stability of the steps and the second prescribes
the induced inertia actions.
\end{proof}

Call the frame normalized at \(h\) if \(\rho(h)\) is diagonal.
Every framed lift admits such a frame by choosing a generator in
each eigenline reducing to its specified basis vector and rescaling
the two opposite generators to give pairings one. Choices are
related by a diagonal symplectic change of basis.

\begin{theorem}[Six integral parameters for the normalized ordinary lift]\leavevmode
\label{thm:sympordinary468}
The lifts in the preceding lemma satisfying its flag equations,
with fixed multiplier and frame normalized at \(h\), have a
representing ring
\[
              R_\ell=O[[u_1,u_2,a,x,b,c]].
\]
Under conjugation by \(\operatorname{diag}(st,t,s,1)\), the two
\(u_i\) have weight zero and \(a,x,b,c\) have respective weights
\((1,0),(0,1),(-1,1),(1,1)\). Here weight \((i,j)\) means that
evaluation after conjugation multiplies the coordinate by \(s^it^j\).
All six coordinates vanish at the split lift.
\end{theorem}
\begin{proof}
The flag is the coordinate flag in this frame. The diagonal
characters are \(d_i\) times unramified characters reducing to one.
The multiplier gives two independent ratios, since the products
on the two paired lines are prescribed. Each independent ratio is
determined by its value on \(h\), and that value can be lifted across
every surjection of Artin coefficient rings. Indeed the unramified
quotient is procyclic and the unit groups of the Artin rings here
are finite.

Consider a small surjection of such rings, with kernel \(L\) killed
by the maximal ideal. Given an upper triangular lift modulo \(L\),
first lift its two diagonal characters. Choose continuous upper
triangular symplectic similitude matrices lifting the group values
with these characters on the diagonal. Such a choice exists because
the upper unipotent symplectic group is parametrized by its four
root subgroups in order of increasing height. Equivalently one can
choose its entries \(01,12,02,03\) freely and solve the two pairing
equations for \(23,13\). These choices require only addition and
multiplication and apply in each finite quotient. A section of the
resulting finite sets supplies the continuous initial choice.
Its multiplicative defect is a two-cocycle in the strictly upper
symplectic Lie algebra tensored with \(L\). Conjugation on this
kernel factors through the residual diagonal representation.
Its four lines have characters
\(\bar A,\bar\Psi,\bar\Psi/\bar A,\bar A\bar\Psi\).
For the three odd lines the second cohomology vanishes by
Proposition~\ref{prop:localsymplecticroots429}. For \(\bar\Psi\)
the Tate dual has nontrivial inertia: the exponents of \(\bar\Psi\)
on inertia are \(0\) or \((\ell-1)/2\) modulo \(\ell-1\), neither of
which is \(1\). Local duality gives the same vanishing
\cite{IX-VII-V-milne98}. Correcting the chosen cochain by a one-cochain
therefore gives a representation lifting the original one. If
needed, conjugate by an infinitesimal upper unipotent to restore
diagonality at \(h\). On every root line its effect on the value
there is multiplication by a unit \(\bar\gamma(h)-1\), up to sign.
This conjugation retains the required diagonal characters.

In the tangent space each of the four roots contributes the
one-cocycles \(z\) with \(z(h)=0\). Coboundaries evaluate
isomorphically onto the root line at \(h\), so this is a copy of
\(H^1(D_\ell,k(\bar\gamma))\). Local invariants and second cohomology
are zero for all four lines. The local Euler characteristic formula
gives dimension one on each. There are also the two independent
unramified diagonal directions. The functor is a closed subfunctor
of framed lifts: the equations of the lemma and diagonality at \(h\)
are equations in the universal matrices. The framed ring is
Noetherian by the finite-dimensionality of local first cohomology
over the finite residual splitting field, as in
Theorem~\ref{thm:ring39}. The lifting argument gives formal
smoothness, so six parameters spanning the relative cotangent
space identify this ring with the displayed power series ring.
For example take the two varying diagonal ratios at \(h\) minus one.
For each of the four normalized tangent cocycles choose an element
on which its value is nonzero and take the corresponding upper
matrix entry. These entries have independent tangent functionals.
Conjugating the entries gives the four weights in the statement.
\end{proof}

\begin{corollary}[Identification of the component through the split lift]\leavevmode
\label{cor:sympcomponent468}
For a frame normalized at \(h\), the flag condition of
Lemma~\ref{lem:sympflag468} is equivalent to specialization from
a complete local \(O\)-domain of ordinary lifts containing the
split lift in the sense of
\srcloc{Development/AutomorphicArithmetic/AutomorphicVarieties/SymplecticGSp4/DecompSigmaOrd}{27}.
Here the ordinary equations in that domain are the polynomial
conditions at an inertia element \(s_0\) described below.
\end{corollary}
\begin{proof}
Choose \(s_0\) with cyclotomic value \(1+\ell\), possible by
total ramification of the local cyclotomic extension. The four
\(d_i(s_0)\) are the powers of \(1+\ell\) with exponents
\(r+1,r,1,0\). They are distinct in \(J=\operatorname{Frac}(O)\).
The polynomial conditions prescribe this characteristic polynomial
at \(s_0\) and the equations of Lemma~\ref{lem:sympflag468} with
each \(P_i\) replaced by
\(\prod_{j\ne i}(\rho(s_0)-d_j(s_0)1)\). These are the equations in
\srcloc{Development/AutomorphicArithmetic/AutomorphicVarieties/SymplecticGSp4/MapRepApply}{99}.

Suppose the domain \(B\) satisfies those conditions and admits an
\(O\)-specialization of its matrices to the split lift. Then \(O\)
injects into \(B\). In \(B[1/\ell]\) divide each polynomial by
\(\prod_{j\ne i}(d_i(s_0)-d_j(s_0))\) to form the projectors \(Q_i\).
Over \(\operatorname{Frac}(B)\) their partial sums give a stable
flag by the polynomial conditions. Since it is stable under \(h\),
each step is the sum of a subset of the four simple \(h\)-eigenlines.
Let \(P\) be the sum of their integral projectors and \(Q\) the
corresponding partial sum of the \(Q_i\). Equality of the two images
gives \(PQ=Q\) and \(QP=P\), identities already in \(B[1/\ell]\).
At the split specialization each labeled \(h\)-projector becomes the
labeled coordinate projector; this follows from its residual label
and Hensel uniqueness. Each \(Q_i\) also becomes the \(i\)-th
coordinate projector, by direct evaluation at \(d(s_0)\). The two
identities force the subset in the given step to be the initial
segment in the prescribed order. Thus the flag over the fraction
field is the ordered \(h\)-flag, including its inertia characters.
Its equations with the \(P_i\) are integral, hence hold in the
domain and at every specialization.

Conversely the universal normalized ring in
Theorem~\ref{thm:sympordinary468} is a domain with its split
augmentation. Over its fraction field the prescribed triangular
characters give the four distinct eigenvalues at \(s_0\).
Their ordered partial sums of eigenlines form the coordinate flag:
on each initial step the characteristic polynomial is the product
of the first diagonal factors. This proves the polynomial conditions
there, including the equations on each inertia quotient. Hence
they hold in the universal domain. Its universal representation
and augmentation supply the two specializations required in the
statement for every normalized lift.
\end{proof}


\subsection{The global symplectic equations on a surviving dyadic orbit}
\label{sec:sympglobal469}

Take the faithful characters of \(C_j\), the fields \(E_x\), primes
\(\ell\equiv5\bmod8\), and half-residue weights \(r\) in
Section~\ref{sec:spacingdyadic464}. Impose also
\(\ell\nmid n_*N_*\), \(2\le r<t(x)\), and use the interval character
\(\xi\) and coefficient extension \(O\) of Section~\ref{sec:sympflag468}.
Assume the negative strict group \(S(\alpha,r)\) is nonzero and put
\(k(r)=\operatorname{length}_O S(\alpha,r)\), after extending scalars
to \(O\). We identify unramified scalar extension throughout: it
preserves length, since a composition factor becomes the extended
residue field. Here \(k\) alone denotes the residue field of \(O\).

\begin{lemma}[The two opposite even roots]\leavevmode
\label{lem:oppositeroot469}
For the strict residual condition, zero at \(\ell\) and unramified
at every other finite prime, the groups of
\(k(\bar\Psi^{-1})\) and \(k(\bar\Psi)\) have respective dimensions
one and zero. All six odd roots of
Proposition~\ref{prop:localsymplecticroots429} have strict dimension zero.
\end{lemma}
\begin{proof}
The residual character \(\bar\Psi\) and its inverse have the same
fixed field \(L\). Their values generate a cyclic group of order
\(2^j\). Frobenius on its character values raises to \(\ell\).
The two characters lie in different Frobenius orbits: a congruence
\(\ell^i\equiv-1\bmod 2^j\) would contradict \(\ell\equiv1\bmod4\).
By Theorem~\ref{thm:commonorbits459}, restriction to \(L\) and
class reciprocity identify the two residual strict spaces with the
equivariant homomorphisms from the split class quotient \(H\) in
Section~\ref{sec:spacingdyadic464}. Explicitly one quotients its
\(\ell\)-primary class group by the prime classes above \(\ell\).
Each nonzero faithful orbit component has order at least \(\ell^f\)
where \(f=2^{j-2}\), and \(|H|<\ell^{2f}\). Thus a nonzero component
has length one and excludes the opposite component. The hypothesis
on \(S(\alpha,r)\) detects exactly the negative component by
Lemma~\ref{lem:collisionfamily464}. Extending the finite coefficient
field tensors the equivariant homomorphism spaces with \(k\): one can
compute them componentwise in a basis over the smaller field.
The assertion for the odd roots is
Proposition~\ref{prop:oddstrictalllevels430} at depth one.
\end{proof}

Let \(\Sigma\) consist of \(\ell\) and the finite conductor primes
of \(d\), including \(2\). Let \(G_\Sigma\) be the Galois group of
\(\QQ\) unramified away from \(\Sigma\) and infinity. For
\(q\in\Sigma\setminus\{\ell\}\) write
\(N_q=I_q\cap\ker\bar d\), \(H_q=G_{\QQ_q}/N_q\).
The inertia image here is finite of order prime to \(\ell\) and
\(d\) itself is trivial on \(N_q\). Consider framed lifts to
symplectic similitudes of multiplier \(\mu\), trivial on every
\(N_q\), whose restrictions at \(\ell\) are the normalized ordinary
lifts of Theorem~\ref{thm:sympordinary468}. Denote their complete
local ring by \(\mathcal B\). These conditions are closed equations
in the global framed ring, whose Noetherian construction uses the
finiteness of first cohomology with finite coefficients over the
residual splitting field, as in Theorem~\ref{thm:ring39}.

\begin{theorem}[The global equations and their two torus weights]\leavevmode
\label{thm:symprelations469}
There is a presentation
\[
 \mathcal B =
 O[[u_1,u_2,a,x,b,c,y]]/(Z_1,Z_2,E,F_1,F_2).
\]
The first six variables are the local ones in
Theorem~\ref{thm:sympordinary468}. One may take \(y\) to be a
global matrix entry of weight \((0,-1)\). The \(Z_i\) have weight
zero, \(E\) has weight \((0,1)\), and \(F_1,F_2\) have weight
\((0,-1)\). A vacant relation of a given weight may be entered
as zero.
\end{theorem}
\begin{proof}
We give the deformation count relative to \(R_\ell\).
For each torus weight take the corresponding space \(M\) of the
residual symplectic Lie algebra. It consists of two trivial lines
at weight zero, one line for each even root and one for each of
the six odd roots. Form the shifted cone \(K\) of the cochain map
\[
\begin{split}
 C^\bullet(G_\Sigma,M)\oplus\!
       \bigoplus_{q\ne\ell,\ q\in\Sigma} C^\bullet(H_q,M)
       &\longrightarrow C^\bullet(D_\ell,M)\oplus\!
       \bigoplus_{q\ne\ell,\ q\in\Sigma} C^\bullet(G_{\QQ_q},M),\\
 (z,(z_q))&\longmapsto
       (z|_{D_\ell},(z|_{G_{\QQ_q}}-\operatorname{inf}z_q)_q).
\end{split}
\]
The shift is by minus one. As in Lemma~\ref{lem:rel70}, we may
subtract a global coboundary to set the degree-zero witness at
\(\ell\) to zero and subtract local quotient coboundaries to set
the other witnesses to zero. This identifies \(H^1(K)\) with the
global cocycles vanishing on \(D_\ell\) whose restrictions at the
other \(q\) factor through \(H_q\), without further identification
of representatives. On a nonzero weight line the invariants at
\(D_\ell\) are zero by Lemma~\ref{lem:sympflag468}. Thus a class
zero there has exactly one representative vanishing there.
At \(q\ne\ell\) the finite residual inertia quotient has order
prime to \(\ell\); averaging identifies the condition of vanishing
on \(N_q\) with triviality of the restriction class on \(I_q\).
On the two trivial lines the cocycles are everywhere unramified
additive characters of \(G_\QQ\) vanishing at \(\ell\). They
vanish by class reciprocity for \(\QQ\). Hence Lemma~\ref{lem:oppositeroot469}
leaves one relative tangent dimension, on the negative even root.

In degree zero the map from the invariants in the domain of the
cochain map is injective. The last terms in Poitou--Tate duality
identify \(H^3(K)\) with the dual of \(H^0(G_\Sigma,M^*(1))\).
Indeed \(H_q\) has cohomology only in degrees zero and one by
averaging over its finite inertia image and passing to the
unramified quotient. All the indicated Tate-dual invariants vanish
already at \(\ell\). For the odd roots this follows from
Proposition~\ref{prop:localsymplecticroots429}. For the neutral
and even lines the tame inertia exponents of the lines are \(0\)
or \((\ell-1)/2\), distinct from \(1\) since \(\ell\ge5\).
The local Euler characteristics at \(q\ne\ell\) and for \(H_q\)
are zero. Those of \(G_\Sigma\) and \(D_\ell\) are respectively
\(-\dim M+\dim M^{c_\infty=1}\) and \(-\dim M\), where
\(c_\infty\) is complex conjugation. These duality and Euler
formulas apply to finite \(\ell\)-power coefficients and the
chosen \(\Sigma\) \cite{IX-VII-V-milne98}. It follows that
\[
             \dim H^2(K)=\dim H^1(K)+\dim M^{c_\infty=1}.
\]
This gives two dimensions at weight zero, one on the positive
even root, two on the negative even root and zero on the odd roots.

For clarity, the role of \(H^2(K)\) is an upper bound on relations.
In a small extension lift the global matrices as a continuous
cochain of fixed multiplier. Such set-theoretic lifts exist by
smoothness of the symplectic group, or by lifting a symplectic
basis and correcting its pairings. Lift the local quotient
representations at \(q\ne\ell\), using \(H^2(H_q,M)=0\) on all
weights. Compare these lifts and the given lift at \(\ell\) with
the global cochain. Its multiplicative defect and the local
differences define the degree-two cone class with values in the
small kernel, exactly as in Theorem~\ref{thm:pres70}. Vanishing
allows a global cochain correction giving a homomorphism. A
global infinitesimal conjugation makes its restriction at \(\ell\)
equal to the prescribed lift; this preserves factorization at
every \(q\ne\ell\). Equality at \(\ell\) also restores the
chosen frame condition at \(h\).

Choose an element on which the relative tangent cocycle on the
negative root is nonzero. The \(20\) entry there detects it,
since the root matrix is \(E_{20}+E_{31}\).
This entry as \(y\) gives a surjection
\(P=R_\ell[[y]]\twoheadrightarrow\mathcal B\) by complete Nakayama.
For its kernel \(I\), apply the obstruction to
\(P/\mathfrak m_P I\to P/I\). A nonzero linear functional on
\(I/\mathfrak m_P I\) cannot annihilate that obstruction. Otherwise
the pushed-out extension would admit a representation over \(R_\ell\)
lifting the universal one. The induced map from \(P\) would differ
from the natural one only on \(y\), by a square-zero kernel
element. Every element of \(I\) has zero residual linear
coefficient in \(y\) by minimality, so the two maps would agree
on \(I\), contradicting the chosen functional. The continuous
cochain construction for this finite-kernel extension follows by
choosing a section on finitely many open cosets as in the proof of
Theorem~\ref{thm:pres70}.

These dimension bounds apply separately to each weight. Indeed
the torus acts by diagonal conjugation on all entries and on the
cochains. A homogeneous relation in a matrix entry of weight
\(\lambda\) has scalar transformation law \(s^{\lambda_1}t^{\lambda_2}\).
In the universal obstruction a defect in that entry has this same
law. Equivalently the dual of the relation space maps into the
corresponding cohomology space with the dual convention for the
torus on linear functionals. Weight separation here uses the
split algebraic torus over the coefficient ring: in each finite
jet retain coefficients of the Laurent monomials in formal \(s,t\).
This gives exact projections, also when characters agree on some
finite group of residue-field units. Taking inverse limits lifts
a weight basis of \(I/\mathfrak m_P I\) to homogeneous relations.
Nakayama's lemma gives the asserted presentation.
\end{proof}

\begin{corollary}[Elimination of the diagonal variables and first torus invariants]\leavevmode
\label{cor:firsttorus469}
Let the first torus mean conjugation with \(t=1\). Its invariant
ring in \(\mathcal B\) has a presentation
\[
 \mathcal B^{s}=O[[x,v,w,y]]/(E,F_1,F_2),\qquad v=ab,\quad w=bc .
\]
The weights of \(x,v,w,y\) for the remaining torus are
\(1,1,2,-1\). In this formula \(E,F_1,F_2\) denote the relations
after the following elimination.
\end{corollary}
\begin{proof}
Impose \(a=x=b=c=y=0\). Every universal off-diagonal matrix
entry has nonzero torus weight and therefore vanishes in this
quotient, on which the torus is trivial. Conversely, the five
variables were chosen from off-diagonal entries, so imposing
their vanishing represents exactly the diagonal lifts. Their four
character ratios to \(d_i\) are unramified at all finite places.
At \(\ell\) this is the flag condition. Elsewhere inertia acts
through a finite prime-to-\(\ell\) quotient and the ratios take
values in the principal congruence units, a pro-\(\ell\) group.
Global reciprocity for \(\QQ\) forces every ratio to be one. One
can see this on each Artin quotient: a finite abelian extension
of \(\QQ\) unramified at all finite primes is trivial. Thus the
diagonal quotient is \(O\) via the split augmentation.

Only \(Z_1,Z_2\) can survive after these five variables are set
to zero. Their restrictions generate \((u_1,u_2)\), so their
linear parts in the two diagonal variables form an invertible
matrix modulo the maximal ideal. The formal inverse function
theorem eliminates \(u_1,u_2\) over \(O\) by \(Z_1=Z_2=0\);
alternatively solve successively modulo powers of the maximal
ideal using that matrix. This substitution respects both weights
by uniqueness. The remaining relations all have first weight zero.
In the power series ring on \(a,x,b,c,y\) a monomial has first
weight zero precisely when its exponent on \(b\) is the sum of its
exponents on \(a,c\). It has a unique expression as a monomial in
\(x,ab,bc,y\). This substitution identifies the corresponding
complete power series rings on neutral monomials, since the
generator degrees are between one and two. Projection of
\(E H+F_1H_1+F_2H_2\) to first weight zero simply projects the
three coefficients to first weight zero. Exactness of projection
on each finite jet, then completion, proves the stated formula.
\end{proof}


\subsection{The integral Frobenius congruence giving the strict length}
\label{sec:symptrace470}

Continue with the surviving faithful dyadic character and the same
weight, prime and field as in Section~\ref{sec:sympglobal469}. Write
\(\rho:G_\Sigma\to\operatorname{GSp}_4(\mathcal B)\) for its universal
normalized lift, and \(t(g)=\operatorname{tr}\rho(g)\),
\(t_0(g)=\operatorname{tr}d(g)\). Let \(T\) be the closure of the
\(O\)-subalgebra generated by the \(t(g)\), in the maximal-ideal
topology of \(\mathcal B\). We determine an integral quotient of
\(T\) directly from the negative root cohomology.

\begin{lemma}[Three products of entries generate the trace algebra]\leavevmode
\label{lem:invarianttrace470}
In Corollary~\ref{cor:firsttorus469} set
\[
     U=xy,\qquad V=vy,\qquad W=wy^2,\qquad
     P=O[[U,V,W]],\qquad D=(U,V,W).
\]
Write \(F_i=yf_i(U,V,W)\) and \(e(U,V,W)=yE\). Then \(e\in D\)
and the full torus invariant ring in \(\mathcal B\) equals \(T\).
Its presentation, including all integral relations, is
\[
                      T=P/\big(e,\ D(f_1,f_2)\big).
\]
The augmentation \(T\to O\) is evaluation at \(U=V=W=0\).
\end{lemma}
\begin{proof}
For a monomial \(x^iv^jw^ky^l\), second weight zero means
\(l=i+j+2k\). Thus every neutral monomial is uniquely a product
\(U^iV^jW^k\). The two adic topologies on these power series are
equivalent because the three monomial degrees in \(x,v,w,y\)
are \(2,2,3\). Every series of weight \(-1\) is \(y\) times a
neutral series. Multiplication by \(y\) takes the series of weight
\(+1\) onto \(D\): it gives precisely the neutral series with no
constant monomial, as can also be seen from \(U=yx,V=yv,W=y(wy)\).
Projection of the ideal \((E,F_1,F_2)\) to degree zero therefore
gives exactly \((e,D(f_1,f_2))\). Exactness follows by projecting
the coefficients of an expression in the homogeneous generators
term by term, as in Corollary~\ref{cor:firsttorus469}.

Every trace is invariant. To prove the converse, observe first
that the coefficients of the characteristic polynomial of \(\rho(h)\)
belong to \(T\) by Newton's identities on \(t(h^i)\), \(1\le i\le4\),
since \(\ell\ge5\). Its four labeled roots belong to \(T\) as well:
start with lifts in \(O\) of the simple residual roots and apply
Newton iteration in the complete ambient ring. Inverses of elements
of \(T\) with nonzero residue belong to \(T\), by subtracting their
residues lifted to units in \(O\) and using a geometric series.
Every iterate is therefore in \(T\), and convergence puts the limit
there. The polynomials in \(\rho(h)\) giving \(P_i\) in
Lemma~\ref{lem:sympflag468} have coefficients in \(T\). These
projectors are \(E_{ii}\) since the frame is normalized. It follows
that every cyclic product of matrix entries belongs to \(T\).
Indeed it is the trace of
\(P_{i_0}\rho(g_1)P_{i_1}\rho(g_2)\cdots P_{i_{m-1}}\rho(g_m)\),
with entries successively indexed \(i_0i_1,\ldots,i_{m-1}i_0\).
Expand the \(P_i\) as polynomials; each trace in the expansion is
the trace of a group element.

Choose \(a,x,b,c\) as entries with indices \(01,02,12,03\),
respectively, using a group element with nonzero tangent evaluation
on the associated root for each choice. Take \(y=\rho(g_y)_{20}\)
as in Theorem~\ref{thm:symprelations469}. The products \(U=xy\)
and \(V=aby\) already follow the cycles of indices \(020\) and
\(0120\). In \(W=bcy^2\) use the alternating form on one copy
of \(y\). The multiplier identity gives
\[
          \rho(g_y)_{20}=-\mu(g_y)\rho(g_y^{-1})_{31}.
\]
The remaining entries \(03,31,12,20\) form a cycle. Thus
\(U,V,W\in T\). Since \(T\) is closed, it contains every power
series in these entries. Their value at the split lift is zero.
\end{proof}

\begin{theorem}[The exact integral quotient of the exceptional strict length]\leavevmode
\label{thm:exactstricttrace470}
Define \(C=P/(e,f_1,f_2)\), a quotient of the actual trace algebra.
There is an isomorphism of \(O\)-algebras
\[
               C/DC = O/\ell^{\,k(r)}O,\qquad
                    k(r)=\operatorname{length}_O S(\alpha,r).
\]
In particular the formula holds at the possible weight of length
above twenty-five in Corollary~\ref{cor:leveltwentyfive467}.
It requires no finiteness hypothesis on \(C\) as an \(O\)-module.
\end{theorem}
\begin{proof}
Put \(O_m=O/\ell^mO\) for \(m\ge1\). In the presentation from
Theorem~\ref{thm:symprelations469}, after the diagonal elimination,
consider maps to \(O_m[\epsilon]/(\epsilon^2)\) given by
\[
              a=x=b=c=0,\qquad y=\epsilon z_0,\quad z_0\in O_m .
\]
The neutral diagonal variables eliminated by \(Z_1,Z_2\) take their
split values: they were expressed as neutral series and all nonconstant
neutral terms vanish under this substitution. Since the only possible
linear term has weight \((0,-1)\), \(E\) vanishes and the two remaining
relations require exactly \(f_i(0)z_0=0\) for \(i=1,2\). Conversely
these equations suffice to give the specified map. Its matrices have
the form
\[
       \rho_\epsilon(g)=(1+\epsilon z(g)X)d(g),\qquad
                              X=E_{20}+E_{31},
\]
where \(d\) is reduced modulo \(\ell^m\). Indeed the only possible
varying entries have the indicated weight; the multiplier equation
on their infinitesimal coefficients gives the displayed one-dimensional
symplectic root space. The group law is precisely
\[
           z(gg')=z(g)+\Psi^{-1}(g) z(g').
\]
The upper triangular condition at \(\ell\) forces \(z|_{D_\ell}=0\).
Away from \(\ell\) the framed equations force \(z|_{N_q}=0\) at
the conductor primes and impose unramifiedness at the other primes.
Conversely any such cocycle gives these matrices and this map:
all the locally chosen upper coordinates are zero, the diagonal is
fixed, and the chosen \(20\) entry gives its value of \(y\).

We verify that these are exactly the classes in \(S(\alpha,r)[\ell^m]\).
View \(O_m(\Psi^{-1})\) as the \(\ell^m\)-torsion of
\((\operatorname{Frac}(O)/O)(\Psi^{-1})\) by division by \(\ell^m\).
On both modules \(\Psi^{-1}(h)-1\) is a unit. Every global class
has exactly one cocycle representative with \(z(h)=0\), obtained
by subtracting a coboundary. If its class is killed by \(\ell^m\),
the normalized representative is killed by \(\ell^m\), because
the normalized representative of the zero class is zero. For a
strict class the cocycle on \(D_\ell\) is a coboundary, whose
value zero at \(h\) forces its witness to be zero. At \(q\ne\ell\)
a coboundary on inertia restricts to zero on \(N_q\), since
\(\Psi\) is trivial there. Conversely vanishing on \(N_q\) implies
that the inertia restriction factors through a group of order prime
to \(\ell\). Averaging there makes it a coboundary, including for
finite coefficients. This proves the claimed identification with all
the specified cocycles. The correspondence is \(O_m\)-linear, as
is seen from the square-zero formula for the matrices.

By Corollary~\ref{cor:dyadicorbit459}, \(S(\alpha,r)\) is a finite
cyclic \(O\)-module of length \(k(r)\). Consequently the simultaneous
annihilator of \(f_1(0),f_2(0)\) in \(O_m\) has length
\(\min(m,k(r))\) for every \(m\). For an ideal of the discrete valuation
ring \(O\) these lengths determine its valuation, including the zero
ideal by letting \(m\) grow. They give
\((f_1(0),f_2(0))=\ell^{k(r)}O\). Setting \(U,V,W\) to zero in \(C\)
and using \(e(0)=0\) proves the assertion.
\end{proof}

\begin{corollary}[Generation by good-prime trace differences]\leavevmode
\label{cor:frobideal470}
For \(q\notin\Sigma\) write \(\delta_q=t(\operatorname{Frob}_q)-
t_0(\operatorname{Frob}_q)\). The ideal of \(T\) generated by all
these \(\delta_q\) is \(DT\). Some set of at most three such primes
generates this ideal. For any such generating set \(\mathcal Q\),
\[
                  C/(\delta_q:q\in\mathcal Q)C
                      \simeq O/\ell^{k(r)}O .
\]
\end{corollary}
\begin{proof}
The topology induced from \(\mathcal B\) on \(T\) is its maximal-ideal
topology. To see this directly, \(P\), with finite residue field, is
compact in its maximal-ideal topology as the inverse limit of finite
quotients. The continuous surjection from \(P\) onto the closed
subalgebra in the lemma is thus a quotient map onto a Hausdorff
space. In particular all ideals of \(T\) are closed, either by
completeness and Noetherianity or by expressing each as the image
of a finite product of compact copies of \(T\).

Both traces are continuous class functions on \(G_\Sigma\). The
Chebotarev density theorem in every finite quotient of this group
implies that the conjugates of good-prime Frobenius elements are
dense. Thus the closed ideal generated by the \(\delta_q\) contains
\(t(g)-t_0(g)\) for every \(g\). In the quotient by that ideal all
traces are scalars, and the scalars form a closed subring, being
the continuous image of the compact ring \(O\). They generate the
whole quotient by the definition of \(T\). The split augmentation
identifies that quotient with \(O\), proving the ideal equality.
Since \(DT\) is generated by \(U,V,W\), its quotient modulo the
maximal ideal times \(DT\) has dimension at most three. Select a
basis from the images of the \(\delta_q\) and apply Nakayama's
lemma. The last assertion now follows from the theorem. The
choice by Nakayama gives an existence assertion about the primes;
their sizes do not enter the proof.
\end{proof}


\subsection{The finite trace algebra and the determinant at Frobenius}
\label{sec:sympdet471}

Use \(P,D,e,f_1,f_2,T,C\) and the actual strict length \(k(r)>0\)
of Section~\ref{sec:symptrace470}. Put \(J=\operatorname{Frac}(O)\).
In \(R=P/(e)\) write \(D_R\) for the ideal of the three variables and
\(I=(f_1,f_2)R\). Thus \(T=R/D_R I,\ C=R/I,\ R/D_R=O\).
We first determine what finiteness is needed to compare the
strict group with characteristic-zero traces.

\begin{lemma}[Finiteness and the component separate from the split trace]\leavevmode
\label{lem:finitetrace471}
The following conditions are equivalent: \(T\) is finite as an
\(O\)-module, \(C\) is finite as an \(O\)-module, and \(C/\ell C\)
has dimension zero. Regardless of these conditions,
\[
        T[1/\ell]\simeq J\times C[1/\ell].
\]
If the conditions hold, then \(C\) is a nonzero finite free
complete intersection over \(O\).
\end{lemma}
\begin{proof}
The kernel of \(T\to C\) is \(I/D_R I\), generated over
\(R/D_R=O\) by the two images of \(f_1,f_2\). This proves the
first equivalence. If \(C/\ell C\) has dimension zero, it is an
Artinian local quotient of \(k[[U,V,W]]\), hence finite over \(k\).
Then some power of the maximal ideal of \(C\) is contained in
\(\ell C\), and the two adic topologies are equivalent. Lift a
\(k\)-basis and express any element successively modulo powers
of \(\ell\) in these lifts. Completeness of \(O\) gives an
expression over \(O\). The converse follows by reduction.
By Theorem~\ref{thm:exactstricttrace470},
\(R/(D_R+I)=O/\ell^{k(r)}O\); thus \(D_R\) and \(I\) are comaximal
after inverting \(\ell\). The Chinese remainder theorem gives
the product formula.

Under finiteness the ideal \((e,f_1,f_2,\ell)\subset P\) is
primary to the maximal ideal. All four elements lie in that ideal
since \(k(r)>0\). The regular local ring \(P=O[[U,V,W]]\) is
Cohen--Macaulay of dimension four. The four elements, in the
displayed order, form a system of parameters and hence a regular
sequence \cite{IX-VII-cm471}. Thus multiplication by \(\ell\) on \(C\)
is injective. A finite torsion-free module over a discrete valuation
ring is free; the quotient by \(D\) is nonzero. This proves the
last assertion.
\end{proof}

\begin{theorem}[Torsion and the congruence with the split trace]\leavevmode
\label{thm:congruenceorder471}
Assume the finiteness conditions of Lemma~\ref{lem:finitetrace471}.
There is a surjective ring map with square-zero kernel
\[
       T\longrightarrow O\times_{\,O/\ell^{k(r)}O} C .
\]
The map from \(C\) in the fiber product is its quotient by \(DC\).
The kernel is the full \(O\)-torsion of \(T\) and is isomorphic
as an \(O\)-module to \(O/\ell^{k(r)}O\). In particular the
scalars occurring in the first factor with second coordinate zero
in the torsion-free trace algebra form exactly \(\ell^{k(r)}O\).
\end{theorem}
\begin{proof}
The elementary exact sequence for two ideals gives
\(R/(D_R\cap I)=(R/D_R)\times_{R/(D_R+I)}(R/I)\).
For surjectivity, lift a compatible pair to two elements of \(R\).
Their difference lies in \(D_R+I\), so subtracting a suitable
element of each ideal gives a common lift. Hence the claimed
kernel is \((D_R\cap I)/(D_R I)\), and its square is zero.

The regular sequence from the lemma makes \(f_1,f_2\) a regular
sequence in \(R\). The relations between them are exactly the
multiples of \((-f_2,f_1)\): reduce a relation modulo \(f_1\),
use injectivity of \(f_2\) there and then cancel \(f_1\) in \(R\).
Right exactness of tensor product over \(R\) now gives
\[
 I/D_R I \simeq
    O^2/O\big(-f_2(0),f_1(0)\big).
\]
Its map to \(R/D_R=O\) has image \((f_1(0),f_2(0))=\ell^{k(r)}O\)
and kernel exactly \((D_R\cap I)/(D_R I)\). Write
\((f_1(0),f_2(0))=\ell^{k(r)}(a,b)\) as a pair of elements, so
\(aO+bO=O\). The solutions in \(O^2\) of \(au+bv=0\) are precisely
the multiples of \((-b,a)\), as one of \(a,b\) is a unit. Quotienting
by its \(\ell^{k(r)}\)-multiple yields the claimed module. The
fiber product lies in the torsion-free module \(O\oplus C\);
therefore this kernel accounts for all \(O\)-torsion. Its removal
gives the asserted exact scalar ideal by the definition of the
fiber product.
\end{proof}

\begin{corollary}[A determinant bound for the full exceptional length]\leavevmode
\label{cor:strictdet471}
Under the same finiteness hypothesis let \(z\in DC\). For instance
one can take a good-prime trace difference \(\delta_q\). Suppose
that \(z\) is nonzero at every geometric point of \(C[1/\ell]\).
For each \(J\)-embedding point \(\sigma:C[1/\ell]\to\overline J\)
let \(m_\sigma\) be the length of its local factor after extension
to \(\overline J\). Then, with \(v_\ell(\ell)=1\),
\[
 \begin{aligned}
 k(r)&\le \operatorname{length}_O(C/zC)
      =v_\ell\big(\det\nolimits_O(z\mid C)\big)\\
     &=\sum_{\sigma}m_\sigma\,v_\ell\big(\sigma(z)\big).
 \end{aligned}
\]
Some integral linear combination of the at most three trace
differences of Corollary~\ref{cor:frobideal470} always satisfies
the nonvanishing hypothesis.
\end{corollary}
\begin{proof}
As \(z\in DC\), the quotient \(C/zC\) surjects onto
\(O/\ell^{k(r)}O\) by Theorem~\ref{thm:exactstricttrace470}.
Over \(\overline J\), multiplication by \(z\) on each local
Artin factor acts on the graded pieces of its maximal-ideal
filtration by the scalar \(\sigma(z)\). Its determinant on the
generic fiber is therefore \(\prod_\sigma\sigma(z)^{m_\sigma}\).
Nonvanishing makes multiplication on the finite free module
\(C\) injective; diagonalizing the inclusion of lattices over
\(O\) gives exactly the length displayed.

Let \(z_1,z_2,z_3\) be the chosen differences, padding with zeros
if necessary. They generate \(DC\), which generates the unit ideal
after inverting \(\ell\), since \(C/DC\) has finite length.
Thus at every geometric point at least one of their values is
nonzero. At a point the polynomial
\(\sigma(z_1)+X\sigma(z_2)+X^2\sigma(z_3)\) has at most two
roots. If there are \(n\) points, some integer \(m\) between
zero and \(2n\) avoids all the roots. Take
\(z=z_1+mz_2+m^2z_3\).
\end{proof}

There is a specific use of this bound on the original interval.
At the common prime \(p\) of Theorem~\ref{thm:commonsptrace433},
suppose the traces at the geometric points, counted with the
\(m_\sigma\), occur under the coefficient embedding in the integral
multiset \(\mathcal C_r\) of that theorem. Its archimedean trace
estimate ensures \(t(\operatorname{Frob}_p)-t_0(\operatorname{Frob}_p)\)
is nonzero at every such point. Taking \(z=\delta_p\) in the
corollary and applying that theorem gives
\(k(r)\le v_\ell(D_{\rm Sp})\). In particular a common such list of
the polynomially bounded size in Corollary~\ref{cor:fixedlevelcount433}
would bound these lengths on the period interval. The determinant
counts the multiplicity of the actual trace algebra at each point.
Establishing finiteness of its special fiber and identifying those
multiplicities and traces on the indicated fields are further
arithmetic requirements.



\subsection{Galois representations realizing the trace coordinates}
\label{sec:tracereps473}

Keep the family and rings of Section~\ref{sec:symptrace470}.
An integral specialization will mean a continuous local \(O\)-algebra
map to a complete discrete valuation ring \(A'\) with finite residue
field containing \(k\). The coefficient map may also factor through
\(k\), so the statements apply on the special fiber. Write
\((U',V',W')\) for the values of the three trace coordinates in \(A'\).

\begin{lemma}[Realization over a quartic extension]\leavevmode
\label{lem:realizetrace473}
Every integral specialization of \(C\) has the trace of a genuine
global normalized symplectic lift after adjoining a fourth root of
a uniformizer of \(A'\). Its local component is the one in
Corollary~\ref{cor:sympcomponent468}.
\end{lemma}
\begin{proof}
Let \(s\) be a fourth root of the uniformizer and work in the
integers \(A''\) of the resulting finite extension of the fraction
field. Choose the following values in the presentation after
elimination of \(u_1,u_2\) in Corollary~\ref{cor:firsttorus469}:
\[
             y=s,\quad b=s,\quad x=U'/s,\quad
                     a=V'/s^2,\quad c=W'/s^3 .
\]
All belong to the maximal ideal: \(U',V',W'\) have valuation at
least that of \(s^4\), with zero allowed. Thus all formal series
converge. The products \(xy,aby,bcy^2\) have precisely the
required values. The equations \(F_i=y f_i\) hold because the
specialization kills the two \(f_i\). Since \(yE=e\) and \(s\ne0\)
in a domain, \(E=0\) as well. The eliminated diagonal variables
are supplied by their unique series. This gives a continuous
specialization of \(\mathcal B\), hence actual matrices with their
required local and inertia conditions. Their traces agree with
the given map on \(T\) since \(U,V,W\) topologically generate it.
If \(V'=W'=0\) we may instead take \(a=b=c=0\), with the same
\(x,y\) and identical verification.
\end{proof}

\begin{theorem}[The rank of nonzero specializations]\leavevmode
\label{thm:rankdichotomy473}
Suppose \((U',V',W')\ne(0,0,0)\).
\begin{enumerate}
\item If \(V'\ne0\) or \(W'\ne0\), the constructed representation
is absolutely irreducible of dimension four.
\item If \(V'=W'=0\), there is a realization as the direct sum of
two absolutely irreducible planes. Complex conjugation acts on
these planes as \(-1\) and \(+1\), respectively. With the indexing
\(\{0,2\},\{1,3\}\) their determinants are precisely the images
of \(A^2\Psi,\Psi\), respectively. In particular both planes are
even.
\end{enumerate}
The second case cannot occur at any nonzero integral specialization
if the unit condition in Theorem~\ref{thm:unitray472} holds.
\end{theorem}
\begin{proof}
Work over an algebraic closure of the fraction field of \(A''\).
The fixed element \(h\) has four distinct eigenvalues since their
residues are distinct. A subspace stable under the group must be a
sum of some of its four eigenlines, by the polynomials \(P_i(h)\).
For each nonzero entry with indices \(ij\) the presence of line
\(j\) in this sum forces the presence of line \(i\). The symplectic
identity supplies another nonzero entry at \((3-j,3-i)\) in the
matrix of the inverse group element.

If \(V'\ne0\), the entries \(01,12,20\) are nonzero at their
respective chosen elements. They connect \(0,1,2\) by a directed
cycle. Reflecting \(20\) and \(01\) gives entries \(31,23\),
which join \(3\) to that cycle in both directions. If \(W'\ne0\),
the entries \(03,12,20\) are nonzero. Reflection of \(20\) gives
\(31\), and the cycle of indices \(0,3,1,2,0\) uses all four
lines. Thus in both cases there is no nonempty proper stable
subset of the eigenlines.

If only \(U'\ne0\), use \(a=b=c=0\) in the second construction
of the lemma. Only first-torus-weight-zero variables are left,
so every entry of nonzero first weight vanishes. The matrices
preserve the two coordinate planes in the statement. The entries
\(02,20\) are nonzero because \(xy=U'\); reflection yields the
opposite pair \(13,31\) on the other plane. The same eigenline
argument proves absolute irreducibility of both. The alternating
form identifies either plane with the dual of the other twisted
by \(\mu\).

At a complex conjugation the matrices on the two planes reduce
to \(-1,+1\) respectively and square to one. Since \(2\) is a
unit, the projectors onto the opposite signs reduce to zero and
therefore vanish. At \(\ell\) the determinant inertia characters
of the planes are \(d_0d_2\) and \(d_1d_3\) by the ordinary flag.
At \(q\ne\ell\) the determinants factor on inertia through its
finite prime-to-\(\ell\) image under \(\bar d\). Thus the ratio
of each determinant to the corresponding product of the \(d_i\)
is unramified there: it reduces to one and the principal
congruence units of \(A''\) form a pro-\(\ell\) group. These
ratios are everywhere unramified abelian characters congruent to
one, hence trivial by reciprocity for \(\QQ\). This proves the
determinant formulas in both characteristics. Under the unit
condition \(U=H(V,W)\) with \(H(0,0)=0\) by
Corollary~\ref{cor:eliminateupper472}. This excludes a nonzero
triple with \(V'=W'=0\).
\end{proof}

\begin{corollary}[A characteristic-\(\ell\) representation detecting nonfiniteness]\leavevmode
\label{cor:finitecurve473}
The ring \(T\) is infinite as an \(O\)-module if and only if there
is an integral specialization of \(C\) to the integers of a finite
extension of \(k((s_1))\) for which the trace of some good-prime
Frobenius is transcendental over \(k\). Every such failure of
finiteness has a witness of one of the two types in
Theorem~\ref{thm:rankdichotomy473}. With a nonzero unit projection
as in Theorem~\ref{thm:unitray472}, it has a rank-four absolutely
irreducible witness. When \(C\) is finite over \(O\), every geometric
point entering the determinant in Corollary~\ref{cor:strictdet471}
also has a realization of one of the two types, in characteristic
zero.
\end{corollary}
\begin{proof}
If \(T\) is infinite, \(C/\ell C\) has positive dimension by
Lemma~\ref{lem:finitetrace471}. Take a prime ideal with quotient
of dimension one, for example the penultimate member of a chain
of maximal length ending at the maximal ideal. Its quotient \(B'\)
is a one-dimensional complete local domain over \(k\). At least
one of \(U,V,W\) has nonzero image there, since killing all three
leaves only \(k\). Choose a nonzero parameter \(s_1\) in its
maximal ideal. The quotient by \(s_1\) has finite \(k\)-dimension,
so \(B'\) is finite over \(k[[s_1]]\): lift a basis and expand
successive remainders in powers of \(s_1\), using completeness
and the equivalence of the parameter and maximal-ideal topologies.
The coefficient map on \(k[[s_1]]\) is injective since \(B'\) is
a domain and \(s_1\) is nonzero. Its fraction field is a finite
extension of \(k((s_1))\). The valuation ring of that local field
contains \(B'\) by integrality, and the inclusion is local since
some power of the maximal ideal of \(B'\) lies in \(s_1B'\).
It therefore gives the required integral specialization.

By Corollary~\ref{cor:frobideal470} the images of good-prime
differences generate the ideal of the three coordinates. At
least one such difference is nonzero. Its image has positive
valuation, because it lies in the maximal ideal. Every nonzero
element algebraic over the finite field \(k\) has finite
multiplicative order and valuation zero. The split trace is in
\(k\) in characteristic \(\ell\), so the corresponding Frobenius
trace is transcendental. Lemma~\ref{lem:realizetrace473} and the
theorem now give the claimed matrices. Conversely, if \(T\) is
finite, \(C/\ell C\) is Artinian local with residue field \(k\).
A map to a characteristic-\(\ell\) domain kills its nilpotent
maximal ideal. All trace values are then in \(k\).

Finally, under finiteness, each characteristic-zero geometric
point has values in a finite extension of \(J\), integral on
\(C\). Its map to the integer ring is local and continuous,
since the maximal ideal of \(C/\ell C\) is nilpotent. It cannot
kill \(U,V,W\) simultaneously as \(C/DC=O/\ell^{k(r)}O\).
The lemma and theorem therefore apply to it as well.
\end{proof}


\subsection{Irreducibility after restriction to the period fields}
\label{sec:strongtrace474}

Use the genuine representations constructed in
Theorem~\ref{thm:rankdichotomy473}. We work either in characteristic
zero at a nonzero triple or in characteristic \(\ell\) at one of
the specializations with transcendental trace in
Corollary~\ref{cor:finitecurve473}. The fields of definition of
the matrices are local fields with finite residue fields. For a
rank-four irreducible realization write \(V\) for its representation
space. For the block realization, take \(V\) to be either one of
its irreducible planes. In discussing subspaces, extend the
coefficient field to an algebraic closure.

\begin{lemma}[Exclusion of permuted direct summands]\leavevmode
\label{lem:noblocks474}
There is no decomposition of \(V\) as the direct sum of at least
two nonzero subspaces permuted by the global Galois group.
\end{lemma}
\begin{proof}
The semisimplified residual characters of \(V\) form the full set
\(\{1,\tau,\bar A,\bar A\tau\}\) in rank four, or one of its
prescribed pairs in rank two. No ratio of distinct members has
order two. The ratio \(\tau^{\pm1}\) has order \(2^j\ge8\);
each of the other possible ratios is \(\bar A^{\pm1}\tau^i\)
for \(i=0,\pm1\). At the lift \(h\) of Lemma~\ref{lem:sympflag468}
the order of \(\bar A(h)\) is a power of two strictly larger
than the order of \(\tau(h)\) and larger than two. Multiplying
it by such a power of \(\tau(h)\) preserves its order, since
\(\tau(h)\) also has two-power order. Thus no nontrivial
quadratic character fixes the residual multiset under twisting:
it would have to take a member to a different member with that
quadratic ratio.

If a direct decomposition were permuted, irreducibility would
make the permutation action transitive. The number of summands
is then \(2\) or \(4\). The permutation action of the compact
image is continuous: all subspaces involved are defined over
some finite extension of the coefficient local field, and their
finite orbit in the Grassmannian is discrete. The kernel of
matrix reduction is pro-\(\ell\), by its successive principal
congruence quotients. Since \(\ell\ge5\) it maps trivially to
the relevant symmetric group. The permutation image therefore
factors through the abelian residual matrix image. A faithful,
transitive abelian permutation group acts regularly, since a
stabilizer fixes every translate of its point. Its order is
the number of summands, so it has a nontrivial quadratic
character. Multiply vectors in each summand by the value of this
character on the group element labeling the summand. The resulting
operator exhibits \(V\) as isomorphic to its quadratic twist.
Reduction would give the same isomorphism on residual
semisimplifications. Indeed isomorphic representations have the
same characteristic polynomials, and reducing these in stable
lattices identifies their residual semisimplifications by
Brauer--Nesbitt. The quadratic character stays nontrivial since
\(\ell\) is odd. This contradicts the first paragraph.
\end{proof}

\begin{theorem}[Strong irreducibility at the nonconstant traces]\leavevmode
\label{thm:strongtrace474}
In both characteristics \(V\) remains absolutely irreducible on
every open subgroup of \(G_\QQ\). In particular this holds on
\(G_{E_x}\) for the original period field and on every further
finite extension of that field.
\end{theorem}
\begin{proof}
It suffices to consider an open normal subgroup \(H\): every open
subgroup contains its open normal core. The restriction to \(H\)
is semisimple even in characteristic \(\ell\). Indeed choose a
simple \(H\)-submodule. Its conjugates are simple \(H\)-submodules
and span a nonzero global subrepresentation, hence all of \(V\).
A sum of simple submodules is semisimple. Its isotypic subspaces
are permuted by the global action. By Lemma~\ref{lem:noblocks474}
there is only one type.

In characteristic zero \(H\) contains a positive power of the
inertia element \(s_0\) in Corollary~\ref{cor:sympcomponent468}.
Its four eigenvalues in rank four are the corresponding powers
of \(1+\ell\) with exponents \(r+1,r,1,0\); they stay distinct
under any positive common power. On either plane the two
corresponding eigenvalues stay distinct. An isotypic sum with
multiplicity greater than one would give every eigenvalue of
each element of \(H\) with at least that multiplicity. Thus the
restriction is irreducible.

In characteristic \(\ell\) first exclude scalar action of \(H\).
The determinant in rank four is \(\mu^2\), and the determinants
of the planes are given in Theorem~\ref{thm:rankdichotomy473}.
All have values in a finite field in this characteristic. If a
finite-index subgroup acts by scalars, its scalar image is finite,
because their powers of degree \(\dim V\) lie in this finite set.
The whole image of \(V\) would be finite. For a plane the other
plane is its dual times \(\mu\), which also has finite image as
a character. Thus in either case the original four-dimensional
image would be finite. All eigenvalues of its matrices would
then be algebraic over the finite field, contrary to the
transcendental trace of the chosen specialization. This settles
rank two and the case of a one-dimensional simple type in rank four.

The only remaining case is rank four with \(V|_H=W\otimes M\),
where \(W\) is simple of dimension two and \(M\) is a multiplicity
space of dimension two with trivial \(H\)-action. The algebra
of endomorphisms commuting with \(H\) is
\(1\otimes\operatorname{End}(M)\), by Schur's lemma. Conjugation
by \(G_\QQ\) acts on it through a finite group, since \(H\) acts
trivially. Each automorphism of a full matrix algebra over an
algebraically closed field is inner. Thus we have a finite
projective action on \(M\). If \(B_g\) represents this projective
action, multiplication of the matrix of \(g\) by \(1\otimes B_g^{-1}\)
on the left centralizes \(1\otimes\operatorname{End}(M)\);
hence the matrix of \(g\) has form \(A_g\otimes B_g\).
In particular any fixed line, or any permuted pair of distinct
lines, in \(M\) gives a stable subspace, or permuted summands,
in \(V\).

The image of the principal congruence kernel in this finite
projective group is a normal \(\ell\)-subgroup. If nontrivial
it has a central element of order \(\ell\). In characteristic
\(\ell\) a nonidentity projective transformation of order \(\ell\)
on a line has a single fixed point: its matrix up to scale is a
nontrivial unipotent Jordan block. Centrality forces the whole
\(\ell\)-subgroup to fix this point, its unique common fixed
point. Normality then forces the full projective group to fix
the point, contradicting irreducibility of \(V\). The projective
action must therefore factor through the abelian residual matrix
image of two-power order. If trivial, it fixes every line in
\(M\). Otherwise take a nonidentity element. It has two distinct
eigenlines in characteristic \(\ell\), and every element commuting
with it projectively permutes this pair. This contradicts
Lemma~\ref{lem:noblocks474} or irreducibility if each is fixed.
The proof is complete.
\end{proof}

\begin{corollary}[One plane and its prescribed scalar twist]\leavevmode
\label{cor:twoplane474}
In the two-plane case let \(Q\) denote the plane with indices
\(\{1,3\}\). The four-dimensional representation is isomorphic to
\(Q\oplus(A\otimes Q)\). Its trace at \(g\) is
\((1+A(g))\operatorname{tr}Q(g)\). In characteristic \(\ell\)
it is \(2\operatorname{tr}Q(g)\) on the kernel of \(\bar A\).
\end{corollary}
\begin{proof}
The alternating form pairs the two planes perfectly with multiplier
\(\mu\), identifying the other plane with \(Q^\vee\otimes\mu\).
The exterior product pairing in dimension two identifies
\(Q^\vee\) with \(Q\otimes(\det Q)^{-1}\).
Theorem~\ref{thm:rankdichotomy473} gives \(\det Q=\Psi\),
and \(\mu=A\Psi\). This proves the representation identity
and both trace formulas.
\end{proof}


\subsection{Simultaneous Frobenius parameters on the period interval}
\label{sec:splitparameters475}

Fix one of the intervals \([V,V^{3/2}]\) in Section~\ref{sec:missing}.
Take the algebras \(T_i,C_i\) for the surviving dyadic data of
Section~\ref{sec:sympglobal469} occurring at its points. The
subscript \(i\) indexes the distinct choices of coefficient prime,
dyadic layer, character and weight; write \(O_i,k_i,\ell_i,r_i\)
for their coefficient rings, residue fields, primes and weights.
This list is finite: the layer and weight indices and the primes
are bounded on the interval and each layer has finitely many
characters. We use the single \(\xi\) for that interval. There
are also only finitely many distinct \(E_x\) there, since the
auxiliary quadratic fields are fixed and the dyadic layers have
bounded indices. Choose any finite Galois extension \(F/\QQ\)
containing all these period fields and all the residual splitting
fields of \(d_i\). A good prime in this subsection avoids the
union of all the indicated ramified primes. Put
\(\delta_{i,p}=t_i(\operatorname{Frob}_p)-t_{0,i}(\operatorname{Frob}_p)\).

\begin{lemma}[Separation of traces on every open coset]\leavevmode
\label{lem:opencoset475}
For a characteristic-zero geometric point of a finite \(C_i\), or
a characteristic-\(\ell_i\) curve specialization as in
Corollary~\ref{cor:finitecurve473} with nonzero trace-coordinate
triple, the locus
\[
        \{g\in G_F:\operatorname{tr}\rho_i(g)=
                                     \operatorname{tr}d_i(g)\}
\]
is closed and has empty interior in \(G_F\). This holds for the
realizations of either rank type.
\end{lemma}
\begin{proof}
Continuity gives closedness. At a nonzero triple in characteristic
\(\ell_i\) some good-prime trace differs from its split value,
by Corollary~\ref{cor:frobideal470}. The difference belongs to
the maximal ideal of a local field of positive characteristic
and is nonzero, so that trace is transcendental over \(k_i\).
Theorem~\ref{thm:strongtrace474} therefore applies. Suppose the
equality held on a coset \(gH\) of an open subgroup \(H\subset G_F\).
In rank four the restriction of \(\rho_i\) to \(H\) is absolutely
irreducible and distinct from all the one-dimensional constituents
of \(d_i\). In the two-plane case in characteristic zero, the
absolutely irreducible restrictions \(Q|_H,(A\otimes Q)|_H\)
are distinct. Otherwise their determinants would give \(A^2|_H=1\).
But \(H\) contains a positive power of the inertia element \(s_0\)
and \(A(s_0)=1+\ell_i\). In the two-plane case in positive
characteristic \(A=1\) on \(G_F\), so the trace is
\(2\operatorname{tr}Q\) there by Corollary~\ref{cor:twoplane474}.

Over the algebraic closure of the coefficient field, for any
finite list of pairwise inequivalent simple representations
\(\pi_a\) of \(H\), the linear span of the tuples
\((\pi_a(h))_a,\ h\in H\), is the product of their full matrix
algebras. Indeed each simple image of the group algebra is a
full matrix algebra by Burnside's theorem. Its kernel is a
maximal two-sided ideal, distinct for distinct simple modules,
since a full matrix algebra has just one simple module up to
isomorphism. To select one factor, for each other kernel choose
an element in it congruent to one modulo the selected kernel.
Their product selects that factor, and arbitrary matrices in it
are then obtained by multiplication. This proves the span assertion.

An equality of traces for every \(gh\), \(h\in H\), would be
a linear identity on that product of matrix algebras. Select
the component of \(\rho_i\) in rank four or of \(Q\) in rank two
and put zero on the other components, including each distinct
diagonal character. On the selected matrix \(X\) the identity
would require \(\operatorname{tr}(\pi(g)X)=0\) for every \(X\),
with an extra factor \(2\) in the repeated-plane case. The
matrix \(\pi(g)\) is invertible, the trace pairing is nondegenerate,
and \(2\) is nonzero. This contradiction excludes every such coset.
\end{proof}

\begin{theorem}[At most three common split primes give finite parameter maps]\leavevmode
\label{thm:splitparameters475}
There are distinct good rational primes \(p_1,\ldots,p_d\), all
splitting completely in \(F\), with
\[
            d\le\max_i\dim(C_i/\ell_i C_i)\le3,
\]
such that both \(C_i\) and \(T_i\) are finite modules over
\(O_i[[X_1,\ldots,X_d]]\) via \(X_a\mapsto\delta_{i,p_a}\),
for every \(i\). Use \(d=0\) for an empty list of data or if all
special fibers have dimension zero. The bound on \(d\) is two
if every datum satisfies the unit condition of
Theorem~\ref{thm:unitray472}. For this common choice \(T_i\) is
finite over \(O_i\) if and only if all the \(\delta_{i,p_a}\)
are nilpotent in \(C_i/\ell_i C_i\).
\end{theorem}
\begin{proof}
Begin with the rings \(B_i=C_i/\ell_i C_i\), quotients of the
three-variable power series rings in Lemma~\ref{lem:invarianttrace470}.
Under the unit condition, Corollary~\ref{cor:eliminateupper472}
gives a two-variable presentation. We successively quotient each
\(B_i\) by trace differences at the same primes.

At one step consider all the positive-dimensional irreducible
components of the spectra of the current quotients. For each,
the associated complete local domain admits a one-dimensional
domain quotient. Passing to the integers of the fraction field
of that quotient as in Corollary~\ref{cor:finitecurve473} gives
an integral specialization with some trace coordinate nonzero.
Indeed the three coordinates and the residue field generate the
ambient special fiber, so they cannot all vanish in a domain of
dimension one. Lemma~\ref{lem:realizetrace473} provides the matrices.
On \(G_F\) the complement of each trace-equality locus for these
representations is open and dense by Lemma~\ref{lem:opencoset475}.
There are finitely many components since the list of rings is finite
and each is Noetherian. The intersection of finitely many dense
open sets is dense and open. A nonempty open subset of \(G_F\)
thus satisfies all the desired inequalities.

The extension \(F/\QQ\) is finite Galois. In the chosen open
set take a coset \(gN\) of an open normal subgroup of \(G_\QQ\)
contained in \(G_F\). Chebotarev for the finite quotient
\(G_\QQ/N\) gives infinitely many rational primes with the
conjugacy class of \(g\). At any such good prime the splitting
in \(F\) is complete. Its trace values satisfy all the inequalities,
which are conjugation invariant. Choose one new prime. Its
difference is nonzero on each selected curve and hence outside
the minimal prime of every positive-dimensional component.
Quotienting by it reduces each positive dimension by at least
one: any positive-length chain of primes containing the element
can be extended downward by a strictly smaller minimal prime
of the original ring. The same argument bounds a surviving
zero-dimensional quotient by zero. Repetition proves the asserted
number of choices.

Write \(B=O_i[[X_1,\ldots,X_d]]\). The ideal
\(\mathfrak m_B C_i\) now has zero-dimensional quotient with
finite residue field, hence is primary to the maximal ideal.
Lift a residue-field basis of \(C_i/\mathfrak m_B C_i\).
It generates \(C_i\) over \(B\): express a remainder at each
stage in monomials in \(\ell_i,\delta_{i,p_1},\ldots,\delta_{i,p_d}\)
of the next total degree times the basis lifts. The coefficients
converge in \(B\) and the remainders tend to zero in \(C_i\),
since its two adic topologies agree. The kernel of \(T_i\to C_i\)
has two \(O_i\)-generators by Lemma~\ref{lem:finitetrace471},
so the same module-finiteness follows for \(T_i\).

If each chosen difference is nilpotent on the special fiber, its
parameter map from \(k_i[[X_1,\ldots,X_d]]\) factors through a
quotient killing a positive power of each variable. That quotient
is finite over \(k_i\). Thus \(C_i/\ell_i C_i\) is finite over
\(k_i\), the condition in Lemma~\ref{lem:finitetrace471}.
Conversely, if \(T_i\) is finite over \(O_i\), the maximal ideal
of the Artinian local special fiber of \(C_i\) is nilpotent,
and contains all the differences.
\end{proof}

\begin{corollary}[Simultaneous determinant separation over the period fields]\leavevmode
\label{cor:splitdet475}
There are infinitely many good rational primes \(p\) splitting
completely in \(F\) for which every \(i\) with \(C_i\) finite
over \(O_i\) simultaneously satisfies
\[
    k(r_i)\ \le\
    v_{\ell_i}\!\left(\det\nolimits_{O_i}
                 (\delta_{i,p}\mid C_i)\right)<\infty .
\]
When any parameter prime is needed in the theorem, its first
choice can satisfy these inequalities as well.
\end{corollary}
\begin{proof}
Each of the finitely many finite algebras \(C_i\) has finitely
many characteristic-zero geometric points. Impose the open
dense conditions from Lemma~\ref{lem:opencoset475} for every
one of these points on \(G_F\), along with any conditions from
the first parameter step. The same open-coset choice and
Chebotarev argument work simultaneously. For each selected
prime the generic determinants are nonzero, so
Corollary~\ref{cor:strictdet471} gives the inequality with the
full geometric multiplicities. This proves existence on the
same fields throughout the fixed interval. An estimate uniform
in the location of the interval would additionally have to
control the sizes of the determinants or their trace spectra;
the qualitative prime choice gives no such estimate.
\end{proof}


\subsection{Algebraic monodromy of the exceptional trace representations}
\label{sec:monodromy476}

Let \(V\) be a representation in Theorem~\ref{thm:strongtrace474},
over an algebraic closure of its coefficient local field. Its
algebraic monodromy group means the reduced Zariski closure \(G\)
of the Galois image. In rank four it is contained in the group of
similitudes of the specified alternating form. In rank two we
use the determinant as the similitude character.

\begin{lemma}[The residual eigenvalues exclude a symmetric cube]\leavevmode
\label{lem:fourweights476}
In rank four the group \(G\) does not normalize a conjugate of
the image of \(\mathrm{SL}_2\) on its third symmetric power.
\end{lemma}
\begin{proof}
In characteristic zero or characteristic at least five, the
ordinary symmetric cube has image isomorphic to \(\mathrm{SL}_2\).
Indeed the induced map on the Lie algebra is nonzero and injective,
and the center acts faithfully. A finite normal kernel with zero
Lie algebra lies in the center, so the kernel is trivial and the
map is separable onto its image. An algebraic automorphism of
\(\mathrm{SL}_2\) over an algebraically closed field is inner:
in the root datum description there is no diagram automorphism
of type \(A_1\). The centralizer on this irreducible four-dimensional
module consists of scalars. Thus every element of the normalizer
is a scalar times a symmetric cube. Its four eigenvalues, with
multiplicities in some order, form a geometric progression,
by applying the symmetric cube to the two eigenvalues of a
two-by-two upper triangular matrix.

Apply this to \(\rho(h)\) with \(h\) from Lemma~\ref{lem:sympflag468}.
All eigenvalues are integral units and reduction would preserve
a geometric progression, since the common ratio is a quotient
of two of these units. Their actual reductions are
\(\{1,b,a,ab\}\) with \(b=\tau(h)\) of order dividing \(2^j\le n_*\)
and \(a=\bar\xi(\ell)\) of order greater than \(n_*L_*\).
Modulo the subgroup generated by \(b\) these become the two
values \(1,a\langle b\rangle\), each repeated twice. The order
of the second value exceeds two, since it is at least
\(\operatorname{ord}(a)/|\langle b\rangle|>L_*\).
A four-term geometric progression with exactly two distinct
values must have common ratio of order two. The quotient of
its distinct values then has order two, a contradiction.
\end{proof}

\begin{theorem}[The connected algebraic monodromy groups]\leavevmode
\label{thm:monodromy476}
For a rank-four witness, the derived group of \(G^\circ\) is
\(\mathrm{Sp}_4\). In characteristic zero \(G=\mathrm{GSp}_4\);
in characteristic \(\ell\) its identity component is
\(\mathrm{Sp}_4\). On either irreducible plane the derived group
is \(\mathrm{SL}_2\), with full \(\mathrm{GL}_2\) in
characteristic zero and identity component \(\mathrm{SL}_2\)
in characteristic \(\ell\).
\end{theorem}
\begin{proof}
The inverse image of \(G^\circ\) in the profinite Galois group
is open, since there are finitely many closed components and the
representation is continuous in the local-field topology.
Theorem~\ref{thm:strongtrace474} makes the action of \(G^\circ\)
irreducible. Its unipotent radical fixes a nonzero vector, by
triangularization, and normality makes the fixed subspace
invariant. The radical consequently acts trivially and is
trivial in this matrix group. Thus \(G^\circ\) is reductive.
Its central torus acts by scalars by Schur's lemma, so its
connected semisimple derived group \(H\) acts irreducibly as
well. We use the usual classification of reduced root systems,
the central covering by the simply connected group, and the
highest-weight classification of its simple rational modules
\cite{IX-VII-VI-milne122}.

In rank four, \(H\subset\mathrm{Sp}_4\) has rank at most two
and dimension at most ten. The rank-zero case cannot act
irreducibly. A simple group of type \(G_2\) has dimension
fourteen and is excluded. Type \(C_2=B_2\) already has
dimension ten and a connected closed subgroup of that dimension
in \(\mathrm{Sp}_4\) is the whole group. In type \(A_2\) lift
the representation to \(\mathrm{SL}_3\). Self-duality for the
invariant form requires its highest weight \((a,b)\) in
fundamental weights to satisfy \(a=b\), as duality takes
\((a,b)\) to \((b,a)\). A nonzero \((a,a)\) is regular, with
six distinct Weyl conjugates, each a weight of the module.
This exceeds dimension four in either characteristic. In type
\(A_1\times A_1\) the irreducible representation on the
simply connected cover is a tensor product. Both factors act
nontrivially, so both modules have dimension two. Each has
an invariant alternating determinant form because a connected
semisimple group has trivial characters. The tensor of the two
forms is symmetric and nondegenerate. Schur's lemma makes
the space of invariant bilinear forms one-dimensional, which
excludes an alternating form in odd characteristic or
characteristic zero.

It remains to exclude type \(A_1\). In characteristic zero the
four-dimensional simple module of \(\mathrm{SL}_2\) is its
third symmetric power. In characteristic \(\ell\) write a
highest weight \(n\) in base \(\ell\), with digits \(n_i\).
The simple module has dimension \(\prod_i(n_i+1)\) and is
the tensor of the corresponding Frobenius twists of the
\(n_i\)-th symmetric powers \cite{IX-VII-steinberg476}. In this rank
one can also see the formula directly. Multiply the digit
polynomials in \(X^{\ell^i},Y^{\ell^i}\) to obtain the span
of \(X^{n-j}Y^j\) with \(0\le j_i\le n_i\) in each digit.
The monomials have distinct torus weights. From any one of them
the substitution \(Y\mapsto Y+X\) and projection to the highest
weight yield \(X^n\). From \(X^n\), substitution \(X\mapsto X+Y\)
and weight projection recover every allowed monomial, since the
product of the binomial coefficients of its digits is nonzero.
This module is simple of highest weight \(n\), giving the formula.
In dimension four there is either a single digit three or two
digits one. In the latter case the two-dimensional forms give the
same symmetric-form obstruction as above. In the former case the
Frobenius-twisted symmetric cube has the same algebraic image as
the symmetric cube, by surjectivity of Frobenius on algebraically
closed field points. In either characteristic the full group
would normalize this image, since it normalizes \(H\).
Lemma~\ref{lem:fourweights476} excludes it. This proves the
rank-four assertion about \(H\).

In rank two \(H\subset\mathrm{SL}_2\) is nontrivial and of rank
at most one. Its type is \(A_1\), so dimension three forces
equality with \(\mathrm{SL}_2\). In characteristic zero each
of the prescribed similitude characters has infinite image:
their exponents on the chosen principal inertia element are
\(r+1\) in rank four and \(r+2,r\) on the planes. An infinite
subgroup is Zariski dense in the multiplicative group. Since
\(G\) already contains the kernel of the similitude, it is the
full similitude group. In characteristic \(\ell\) the similitude
character has finite image by its explicit formula. Its connected
kernel then gives exactly the asserted identity component.
\end{proof}

\begin{corollary}[Eigenprojectors on the cyclotomic kernel of a period field]\leavevmode
\label{cor:cyclomonodromy476}
For any finite extension \(F/\QQ\), including \(E_x\), and any
nonzero \(X\in\operatorname{End}(V)\), there is an element of
\(G_{F(\mu_{\ell^\infty})}\) whose matrix has distinct eigenvalues
and for which some eigenprojector \(P\) satisfies
\(\operatorname{tr}(PX)\ne0\).
\end{corollary}
\begin{proof}
The image of \(G_F\) has finite index in the original image, so
its Zariski closure has the same identity component. The closure
of the subgroup generated by its commutators contains the
derived group of this component: the commutator map is algebraic,
and the image of a dense set of pairs is dense in the image of
the map. All the commutators preserve the alternating form with
multiplier one and lie in the cyclotomic kernel. Their subgroup
is therefore dense in \(\mathrm{Sp}_4\) or \(\mathrm{SL}_2\)
by the theorem. The defining representation of this connected
group is irreducible, so its matrices span the full matrix
algebra. The function \(g\mapsto\operatorname{tr}(gX)\) is
nonzero somewhere by nondegeneracy of the trace pairing.
The open locus of distinct eigenvalues is nonempty, by taking
a generic element of the diagonal torus. The two nonempty open
loci intersect in the irreducible group and hence meet the
dense commutator subgroup. For such a matrix \(g\), the spectral
identity \(g=\sum_i z_iP_i\) shows that at least one
\(\operatorname{tr}(P_iX)\) is nonzero.
\end{proof}


\subsection{Split Frobenius congruences and theta tests on one interval}
\label{sec:jointprimes484}

Fix one interval \([V,V^{3/2}]\) of Section~\ref{sec:missing}.
Use its full finite list of exceptional dyadic trace algebras
\(C_i,T_i\) from Theorem~\ref{thm:splitparameters475}. Choose
any finite collection of the signed elliptic product tests over
its period fields as in Corollary~\ref{cor:projectsymbol483}.
Enlarge the common finite Galois field \(F\) of
Theorem~\ref{thm:splitparameters475} to contain the field \(J\)
of Theorem~\ref{thm:commonsplit483} for these tests. Write
\(\Gamma=G_F\) and use Haar probability on \(\Gamma\). For a
split prime \(q\), the trace difference \(\delta_{i,q}\) means
the universal trace of Frobenius less its split diagonal trace.
We use \(\ell_i\) for the coefficient prime of \(C_i\).
The auxiliary testing prime \(q\) can be different from all the
\(\ell_i\).

\begin{lemma}[A zero set of traces has measure zero]\leavevmode
\label{lem:tracezero484}
For each trace representation in Lemma~\ref{lem:opencoset475},
the set of \(g\in\Gamma\) where its trace agrees with the split
diagonal trace has Haar measure zero. Among the rational primes
in \([X,2X]\) splitting completely in \(F\), the proportion
giving that equality tends to zero as \(X\) tends to infinity.
Both assertions hold for the union over any fixed finite list
of the indicated representations.
\end{lemma}
\begin{proof}
Map \(\Gamma\) to matrices by the direct sum of the representation
and its diagonal comparison. Write \(G\) for the reduced Zariski
closure of this image over an algebraic closure of its coefficient
local field. The inverse image of \(G^\circ\) is an open subgroup
\(\Gamma_0\). Indeed there are finitely many components and their
inverse images are closed. The image of \(\Gamma_0\) is dense
in \(G^\circ\), by density of the original image and the fact that
each component is Zariski open as well as closed.

For any proper closed subvariety \(Y\) of the connected group
\(G^\circ\), its inverse image in \(\Gamma_0\) has measure zero.
We prove this by induction on dimension, treating irreducible
components separately. Distinct translates of an irreducible
\(Y\) have intersections of smaller dimension. Such intersections
have zero measure by induction, with empty intersection in
dimension zero. If \(Y\) had positive measure, only finitely
many distinct translates by the image of \(\Gamma_0\) could have
positive measure, since every finite collection would have union
of measure its cardinality times that positive measure. Thus the
stabilizer has finite index in that image. Its Zariski closure
is all of \(G^\circ\), since an irreducible group cannot be a
finite union of proper closed subsets. For any \(y\in Y\), the
inverse image of \(Y\) under \(g\mapsto gy\) is closed and
contains this dense stabilizer. Hence \(G^\circ y\subset Y\),
contradicting properness. The preimages used here are measurable
since the polynomials defining a given closed subvariety are
continuous over the finite coefficient extensions containing
their coefficients.

Trace equality cuts out a closed subvariety on each component.
If it contained a component, equality would hold on a coset
of \(\Gamma_0\), against Lemma~\ref{lem:opencoset475}. Translating
each component to the identity component proves the measure
assertion. For the statement about primes, the closed equality
set is stable under conjugacy by \(G_{\mathbb Q}\) because both
traces are restrictions of class functions on that group.
It is contained in conjugacy-stable unions of cosets, inside
\(\Gamma\), of open normal subgroups of \(G_{\mathbb Q}\) with
arbitrarily small relative measure. To see this, first take an
open neighborhood of suitably small measure by regularity of
Haar measure. By compactness, cosets of one open normal subgroup
meeting the closed set can all be taken inside that neighborhood.
Chebotarev in the resulting finite quotient gives the relative
measure of this union as the limiting proportion in \([X,2X]\)
among the split primes. Taking arbitrarily small neighborhoods
gives the result. A finite union preserves measure zero and the
limiting bound.
\end{proof}

\begin{theorem}[Common primes for the trace parameters and the theta tests]\leavevmode
\label{thm:jointprimes484}
Fix \(0<\delta<1/4\). The at most three primes of
Theorem~\ref{thm:splitparameters475} may be chosen successively
with the following additional properties. After each fixed
prefix, a proportion tending to one of the primes in
\([X,2X]\) splitting completely in \(F\) are eligible for
the next parameter, as \(X\to\infty\). At each chosen prime
\(q\) every marked elliptic factor satisfies
\(\operatorname{ord}(H_{q,l})>X^{1/4-\delta}\).
The prime separates the split trace at all characteristic-zero
geometric points of all \(C_i\) finite over \(O_i\).
If there is no parameter to choose, the latter two properties
still hold with relative proportion tending to one.
\end{theorem}
\begin{proof}
At a given step, take the one-dimensional domain specialization
from each current positive-dimensional component as in
Theorem~\ref{thm:splitparameters475}. Add to this finite list
all geometric points of each finite generic algebra. Apply the
lemma. Outside a set of relative proportion tending to zero,
the difference at \(q\) is nonzero on every selected
specialization. In particular it lies outside each indicated
minimal prime. Its quotient reduces the greatest positive
dimension by at least one. The finite-module conclusion then
follows from the same complete local parameter argument as in
Theorem~\ref{thm:splitparameters475}. Theorem~\ref{thm:commonsplit483}
excludes only \(O(X^{1-2\delta})\) further rational primes for
all the elliptic factors together. This is a vanishing proportion
of the split primes for fixed \(F\). All finite bad places,
previously chosen primes and primes at which a specified nonzero
coefficient array becomes zero can also be omitted. Thus
Corollary~\ref{cor:projectsymbol483} applies at each chosen
prime with its stated rank and degree bounds.
\end{proof}

For a geometric point \(\sigma\) of
\(C_i[1/\ell_i]\) with \(C_i\) finite and nonzero, put
\[
 a_{i,\sigma}=\min_{g\in\Gamma}
 v_{\ell_i}\bigl(t_{i,\sigma}(g)-t_{0,i}(g)\bigr),
 \qquad v_{\ell_i}(\ell_i)=1 .
\]
These depths are finite by Lemma~\ref{lem:opencoset475} and
nonnegative by integrality. Each minimum is attained since the
values lie in a discrete valuation group in a finite extension.

\begin{lemma}[Primes attaining the individual trace depths]\leavevmode
\label{lem:mintraces484}
For each such geometric point, at least \(1/20\) of \(\Gamma\)
in Haar probability attains its minimum. For a list of \(s\ge1\)
points, at most \(\lfloor20\log s\rfloor+1\) good split primes
suffice to attain every minimum, with the prime allowed to depend
on the point. Every selected prime can also satisfy the theta
order bounds and the simultaneous separation in the theorem.
\end{lemma}
\begin{proof}
Fix a point and write \(f\) for its trace difference on \(\Gamma\).
Divide by an element of valuation \(a_{i,\sigma}\) and reduce,
obtaining a nonzero function \(\bar f\) on \(\Gamma\) over a
finite residue field. The left translates of \(f\) before
reduction span a space of dimension at most \(4^2+4=20\)
because the trace of a translate of a four-dimensional representation
is linear in its sixteen matrix entries and the split
representation contributes at most four one-dimensional
characters. The same dimension bound holds for the translates
of \(\bar f\). Otherwise \(21\) of them would give an invertible
evaluation matrix on some \(21\) group elements, as point
evaluations span the dual of a finite-dimensional space of
functions. Lifting its normalized entries would contradict the
characteristic-zero bound.

Choose a basis of at most \(20\) left translates of \(\bar f\).
At each \(g\in\Gamma\) some translate takes the nonzero value of
\(\bar f\) at a point of its support. Thus at least one basis
translate is nonzero at \(g\). The basis supports cover \(\Gamma\)
and have equal measure, proving the \(1/20\) bound.
The condition of attaining each minimum is clopen and invariant
under \(G_{\mathbb Q}\)-conjugacy. Given the uncovered points,
the expected number of their minima attained by a single element
of \(\Gamma\) is at least one twentieth of their number.
Some nonempty clopen set attains at least this many.
Chebotarev gives positive relative density of corresponding
split primes, also in dyadic intervals. Removing the exceptional
sets of relative proportion tending to zero in the theorem is
permitted. Each step leaves at most \(19/20\) of the preceding
number uncovered. After the stated number of steps the remaining
number is less than one, since \((19/20)^j\le e^{-j/20}\).
\end{proof}

\begin{corollary}[The least determinant valuation on the common splitting field]\leavevmode
\label{cor:strictmindet484}
For each of the finite algebras \(C_i\) above, let \(m_{i,\sigma}\)
be its generic geometric multiplicities and let \(k_i\) be the
strict length in Theorem~\ref{thm:congruenceorder471}.
After a finite unramified extension of \(O_i\), there is an
integral linear combination \(z_i\) of the trace differences
at the primes of Lemma~\ref{lem:mintraces484} such that
\[
    k_i\ \le\ v_{\ell_i}\det(z_i\mid C_i)
           =\sum_\sigma m_{i,\sigma}a_{i,\sigma}.
\]
This is the minimum determinant valuation among elements of the
ideal generated by all good split-prime trace differences, allowing
finite unramified scalar extension. The primes in the lemma may
be chosen once for the union of the geometric points across
the interval.
\end{corollary}
\begin{proof}
Write the selected differences in \(C_i\) as \(z_1,\ldots,z_b\).
For each point divide the tuple \((\sigma(z_j))_j\) by
a scalar of valuation \(a_{i,\sigma}\) and reduce. The resulting
row is nonzero by the lemma. In the algebraic closure of the
residue field there is a vector outside the kernels of all the
rows, since their defining linear forms have nonzero product and the
field is infinite. Its coordinates lie in a finite residue
extension. Lift them to integers \(c_j\) in the corresponding
unramified extension. Fix embeddings into a common algebraic
closure at \(\ell_i\). Then \(z_i=\sum_j c_jz_j\) has valuation
exactly \(a_{i,\sigma}\) at each extended geometric point, since
its normalized residue there is nonzero.

Finite freeness of \(C_i\), the local generic multiplicities, and
the quotient \(C_i/DC_i=O_i/\ell_i^{k_i}O_i\) of
Theorem~\ref{thm:exactstricttrace470} all persist under unramified
scalar extension, using the new coefficient ring for the quotient.
Each \(z_j\) belongs to \(DC_i\). Hence the determinant and
length argument of Corollary~\ref{cor:strictdet471} gives the
inequality and the displayed equality with multiplicities.
At any point, every split-prime difference has valuation at
least \(a_{i,\sigma}\) by definition. Multiplication by an
element of the finite integral algebra and addition preserve
that lower bound. The same determinant formula proves minimality.
\end{proof}

Both conclusions use primes splitting in the same field \(F\).
The local calculation still requires finiteness of the outstanding
trace algebras and an upper estimate for the sum of their congruence
depths with their geometric multiplicities. The theta ranks
at these primes concern the mixed derivatives of
Corollary~\ref{cor:projectsymbol483}. For the fixed final period
and its archimedean determinant, the comparison of
Section~\ref{sec:missing} remains open.


\subsection{Finite quotients at common split primes}
\label{sec:artiniansplit490}

Fix an entire period interval \(I=[V,V^{3/2}]\). Use the same
exceptional dyadic data, test fields, and finite collection of
elliptic factors as in Section~\ref{sec:jointprimes484}. In particular
\(F/\mathbb Q\) is finite Galois and contains every period field
and the auxiliary fields for these tests on \(I\). For datum \(i\)
write \(O_i\) for its unramified coefficient ring, \(\ell_i\) for
its uniformizer, \(\kappa_i\) for its finite residue field, and
\(k_i>0\) for its strict length. Its actual exceptional trace
algebra and augmentation ideal satisfy
\[
 C_i=O_i[[U,V',W]]/(e_i,f_{i1},f_{i2}),\qquad
 D_i=(U,V',W)C_i,\qquad C_i/D_i=O_i/\ell_i^{k_i}O_i
\]
by Theorem~\ref{thm:exactstricttrace470}. We write \(V'\) for
the middle trace coordinate to distinguish it from the endpoint
of \(I\). At a good rational prime \(q\) write
\(\delta_{i,q}=t_i(\operatorname{Frob}_q)-
t_{0,i}(\operatorname{Frob}_q)\in D_i\).

\begin{lemma}[A trace on a curve in either characteristic]\leavevmode
\label{lem:mixedtrace490}
Let \(B\) be any positive-dimensional complete local domain
quotient of \(C_i\). There is a continuous local map from \(B\)
to the integers of a local field of characteristic \(0\) or
\(\ell_i\) such that the specialized trace differs from the
split trace on a set of full Haar measure in \(G_F\).
The proportion of split primes in \([X,2X]\) detecting this
difference tends to one as \(X\to\infty\).
\end{lemma}
\begin{proof}
Take a one-dimensional domain quotient \(B'\) of \(B\). For
instance, take the penultimate prime in a chain of maximal length
ending in the maximal ideal. The quotient is complete since
the original ring is complete Noetherian local. The three trace
coordinates cannot all vanish in \(B'\), since \(C_i/D_i\) is
Artinian. If \(\ell_i=0\) in \(B'\), choose a nonzero parameter
of \(B'\). The quotient by the parameter is finite-dimensional
over \(\kappa_i\). Lifting a basis successively in its powers
shows that \(B'\) is finite over the power series ring in the
parameter over \(\kappa_i\), as in
Corollary~\ref{cor:finitecurve473}. Its fraction field is a
local field with finite residue field and its valuation ring
contains \(B'\). Some power of the maximal ideal lies in the
parameter ideal, so the inclusion is local and continuous.
Some good-prime trace difference has nonzero image, since those
differences generate \(D_i\) by Corollary~\ref{cor:frobideal470}.
That image has positive valuation, hence is transcendental over
the finite constant field. The diagonal trace takes values in a
finite field, so the representation satisfies the transcendental
trace hypothesis of Theorem~\ref{thm:strongtrace474}.

If \(\ell_i\ne0\) in \(B'\), it is a parameter since \(B'\) is a
one-dimensional local domain. The same successive lifting with
this parameter makes \(B'\) finite over \(O_i\). Its fraction
field is a finite extension of \(\operatorname{Frac}(O_i)\) and
the inclusion into its integers is local and continuous for the
same reason. The nonzero coordinate triple gives a genuine
representation by Lemma~\ref{lem:realizetrace473} and
Theorem~\ref{thm:rankdichotomy473}. The characteristic-zero
statement of Theorem~\ref{thm:strongtrace474} applies to this
specialization. Only the one-dimensional quotient was used to
obtain its finite coefficient field.

In either case the trace cannot agree with the split trace on any
open coset in \(G_F\). Indeed the rank-four representation restricts
irreducibly to each open normal subgroup. In the two-plane case the
two irreducible planes there are distinct in characteristic zero
by the determinant argument of Lemma~\ref{lem:opencoset475}, and
the repeated plane in characteristic \(\ell_i\) has multiplicity
two prime to \(\ell_i\). On a coset, the trace pairing on a
nontrivial matrix block therefore cannot cancel with the diagonal
characters, by the group-algebra argument in that lemma.
Lemma~\ref{lem:tracezero484} now applies: on each component of
the algebraic monodromy group the equality cuts out a proper
closed subvariety. Its inverse image has measure zero, and
Chebotarev in finite quotients gives the stated proportion.
\end{proof}

\begin{theorem}[At most three split tests bound every exceptional strict length]\leavevmode
\label{thm:artiniansplit490}
For the whole list on \(I\) there are distinct good primes
\(q_1,\ldots,q_h\) splitting completely in \(F\), with
\[
       h\le \max_i\dim C_i\le3 ,
\]
such that each quotient
\[
 A_i=C_i/(\delta_{i,q_1},\ldots,\delta_{i,q_h})
\]
is a finite ring. If every datum satisfies the unit condition of
Corollary~\ref{cor:eliminateupper472}, one can take \(h\le2\).
For an empty list take \(h=0\). At each successive choice, for
any fixed \(0<\eta<1/4\), a proportion tending to one of split
primes in \([X,2X]\) are eligible and satisfy
\(\operatorname{ord}(H_{q,l})>X^{1/4-\eta}\) for every elliptic
factor \(l\) in the specified tests. For
\[
 h_i=\min\{a\ge0:\ell_i^a1=0\text{ in }A_i\}
\]
one has \(k_i\le h_i\) and \(A_i/D_iA_i=O_i/\ell_i^{k_i}O_i\),
where \(D_iA_i\) denotes the image ideal.
\end{theorem}
\begin{proof}
The ambient power series ring over \(O_i\) is a regular local
domain of dimension four. Some defining relation is nonzero,
since on setting the three coordinates to zero their values
generate \(\ell_i^{k_i}O_i\). Quotienting by a nonzero element
reduces the dimension by at least one, by adjoining the zero
prime below any chain of primes containing it. Hence
\(\dim C_i\le3\). Under the unit condition the two-variable
presentation over \(O_i\) in Corollary~\ref{cor:eliminateupper472}
has the same quotient \(O_i/\ell_i^{k_i}O_i\) at its origin.
The same argument gives dimension at most two.

At each step take every positive-dimensional minimal component
of every current quotient \(C_i/(\delta_{i,q_a}:a\text{ already chosen})\).
There are finitely many since the rings are Noetherian and the
list on \(I\) is finite. Its domain is a complete local quotient
of \(C_i\). Choose a specialization by the lemma for each one.
For a relative proportion tending to one of the split primes
all the resulting traces differ from the diagonal traces.
The difference at such a prime lies outside the minimal prime
of each listed component. Theorem~\ref{thm:commonsplit483}
permits imposing the elliptic order conditions at the same time.
The finitely many bad or previous primes can be omitted.

If a chain of primes of positive length survives the new quotient,
the corresponding chain upstairs begins at a prime containing the
new difference. A minimal prime below it belongs to a
positive-dimensional component and cannot contain the difference.
It extends the lifted chain downwards by one. Hence the maximum
positive dimension drops by at least one at each step until
all quotients have dimension zero. Each quotient is local
Noetherian of dimension zero, hence Artinian, with finite residue
field. Its maximal-ideal filtration now shows its underlying set
is finite. All the chosen differences belong to \(D_i\), which
proves the quotient identity and the inequality on the additive
order of one. This argument also applies to an exceptional trace
algebra having a positive-dimensional special fiber.
\end{proof}

\begin{corollary}[Finite precision and an explicit bound from initial forms]\leavevmode
\label{cor:jetslength490}
Let \(P=O_i[[u_1,\ldots,u_m]]\), where \(m=3\) or \(m=2\) in
the reduced presentation. Let \(J\) be the ideal of the displayed
relations defining \(C_i\) and lifts of the \(h\) chosen differences.
Write \(\mathfrak n=(\ell_i,u_1,\ldots,u_m)\), so \(A_i=P/J\).
For \(N\ge1\) compute
\[
       b_i(N)=\operatorname{length}_{O_i}P/(J+\mathfrak n^N).
\]
For some finite \(N\) one has \(b_i(N+1)=b_i(N)\).
At the first such \(N\), \(P/(J+\mathfrak n^N)=A_i\), and its
largest elementary-divisor exponent over \(O_i\) equals \(h_i\).
In particular \(k_i\le h_i\le N\).

There is also a sufficient test using leading forms alone. If the
initial ideal of \(J\) in \(\kappa_i[L,X_1,\ldots,X_m]\) contains
\(m+1\) homogeneous elements of positive degrees \(a_0,\ldots,a_m\)
with no common projective zero over \(\overline{\kappa_i}\), then
\[
                  k_i\le h_i\le 1+\sum_{j=0}^m(a_j-1).
\]
\end{corollary}
\begin{proof}
Modulo \(\mathfrak n^N\), a monomial \(u^\alpha\) of degree
\(|\alpha|<N\) has coefficients in \(O_i/\ell_i^{N-|\alpha|}O_i\).
These are independent additive coordinates. The truncated ideal
of relations is the \(O_i\)-span of all monomial multiples of
the given generators of \(J\) through this precision. Thus each
\(b_i(N)\) uses only finite coefficient matrices over \(O_i\).
For the maximal ideal \(\mathfrak m\) of \(A_i\), successive
lengths differ by \(\operatorname{length}(\mathfrak m^N/
\mathfrak m^{N+1})\). If this vanishes, Nakayama's lemma applied to
\(\mathfrak m^N\) gives \(\mathfrak m^N=0\). Such a stage occurs
by the theorem. Multiplication by \(\ell_i^a\) kills the whole ring
precisely when it kills one. The structure theorem for finite
modules over a discrete valuation ring identifies its least such
exponent with the largest elementary-divisor exponent. Since
\(\ell_i\in\mathfrak m\), this exponent is at most \(N\).

The associated graded ring of \(P\) is the indicated polynomial
ring, with \(L\) the image of \(\ell_i\); this follows also from
the monomial description. The assumed homogeneous ideal generated
by the \(m+1\) elements has radical the irrelevant maximal ideal
by the homogeneous Nullstellensatz and extension of the residue
field. They form a system of parameters and hence a regular
sequence in the polynomial ring localized at that ideal, since it
is Cohen--Macaulay. The homogeneous sequence is regular also in
the graded ring: a nonzero homogeneous class cannot disappear on
localizing at the irrelevant maximal ideal, as multiplication
by a polynomial with nonzero constant term preserves its lowest
nonzero degree. Successively using the exact sequence for
multiplication by each homogeneous nonzerodivisor gives the
Hilbert series of their quotient as
\(\prod_{j=0}^m(1+t+\cdots+t^{a_j-1})\). The graded ring of
\(A_i\) is a quotient of this ring since it is the graded ring
of \(P\) modulo the entire initial ideal of \(J\).
It vanishes in degree \(1+\sum_j(a_j-1)\). Nakayama as above
proves the bound.
\end{proof}

The theorem gives finite upper bounds on the same \(E_x\) and
the same \(I\) used in the theta tests. To deduce the estimate of
Section~\ref{sec:missing} one still needs bounds for their weighted
sum as \(V\) increases. The stabilization exponent in the corollary
is allowed to vary with the datum and the interval.


\subsection{Finite trace quotients preserving Frobenius spectra}
\label{sec:spectralquot499}

Fix the interval \(I=[V,V^{3/2}]\) of
Section~\ref{sec:jointprimes484}. Retain its finite list of
exceptional dyadic trace algebras \(C_i\), their unramified
coefficient rings \(O_i\) with uniformizers \(\ell_i\) and
residue fields \(\kappa_i\), and the ideals \(D_i\) with
\(C_i/D_i=O_i/\ell_i^{k_i}O_i\). Write \(t_{i,q}\) and
\(t_{0,i,q}\) for the universal and diagonal Frobenius traces.
The same finite Galois field \(F\) contains the period fields
\(E_x\), the residual splitting fields, and the fields of the
selected elliptic tests on \(I\).

\begin{lemma}[Avoidance of all finite-field trace values]\leavevmode
\label{lem:constants499}
From any positive-dimensional complete local domain quotient
of \(C_i/\ell_i C_i\) there is a local specialization to the
valuation ring of a local field \(L\) of characteristic \(\ell_i\)
with the following property. The proportion of good primes
in \([X,2X]\) split in \(F\) for which the specialized trace
is algebraic over \(\kappa_i\) tends to zero. In particular it
avoids the roots of any given nonzero polynomial over
\(\kappa_i\) with relative proportion tending to one, even
when the polynomial and its degree depend on the prime.
\end{lemma}
\begin{proof}
Take a one-dimensional domain quotient and its local-field
specialization as in Corollary~\ref{cor:finitecurve473}.
The image of some trace difference is nonzero and has positive
valuation, since the residue at the closed point is diagonal
and \(C_i/D_i\) is Artinian. The procedure in
Lemma~\ref{lem:mixedtrace490} gives the associated
representation with a transcendental trace. Algebraic elements
of \(L\) over the finite field \(\kappa_i\) form a finite field:
every nonzero such element has finite multiplicative order,
hence has valuation zero. The residue map is injective on these
constants since the same is true of their nonzero differences.

On \(G_F\) the trace is either that of an absolutely irreducible
four-dimensional representation or twice the trace of an
absolutely irreducible plane. Theorem~\ref{thm:strongtrace474}
gives irreducibility on every open subgroup, and
Corollary~\ref{cor:twoplane474} gives the factor two since the
residual twist is trivial on \(G_F\). This factor is nonzero.
The trace cannot equal a constant on any open coset \(gH\)
in \(G_F\). Shrink \(H\) to an open normal subgroup. Over an
algebraic closure the span of the matrices \(\rho(h)\) on the
irreducible factor \(W\) is \(\operatorname{End}(W)\) by
Burnside's theorem. The span of \(\rho(h)-1\) is a nonzero
two-sided ideal there: left or right multiplication by
\(\rho(h')\) is a difference of two group matrices, and
the representation of \(H\) is nontrivial. This ideal is
therefore all of \(\operatorname{End}(W)\). Constancy on the
coset would force the trace pairing with the invertible matrix
\(\rho(g)\) to vanish on the whole algebra, which is impossible.

In the reduced algebraic monodromy group, each specified
constant trace fiber is thus a proper closed subset on every
component. The measure-zero argument in
Lemma~\ref{lem:tracezero484} applies to this polynomial equation:
the inverse image of a proper closed subset of the identity
component under a compact Zariski-dense group has measure zero,
and translation gives the result on its other components.
Take the finite union over the constant field in \(L\).
Chebotarev in the finite quotients used in that lemma gives
the stated limiting relative proportion. Every root in \(L\)
of every nonzero polynomial over \(\kappa_i\) belongs to this
same finite field.
\end{proof}

\begin{theorem}[Common split primes retaining prescribed integral spectra]\leavevmode
\label{thm:spectralquot499}
For every good split prime \(q\) and each \(i\), prescribe any
monic polynomial \(P_{i,q}\in O_i[X]\) satisfying
\(P_{i,q}(t_{0,i,q})\in\ell_i^{k_i}O_i\).
At most three common primes \(q_1,\ldots,q_s\) suffice to make
all the algebras
\[
 A_i=C_i/\big(P_{i,q_j}(t_{i,q_j}):1\le j\le s\big)
\]
finite as \(O_i\)-modules, with
\(A_i/D_iA_i=O_i/\ell_i^{k_i}O_i\). At each stage a relative
proportion tending to one of the split primes is eligible;
one may impose the elliptic order bounds of
Theorem~\ref{thm:jointprimes484} on them. Every specialization
whose traces meet the prescribed equations factors through
the quotient. In particular one can retain simultaneously any
fixed finite collection of integral characteristic-zero
specializations.

For the original algebra \(C_i\) itself, finiteness over \(O_i\)
is equivalent to the following positive-density condition:
on a set of split good primes of positive upper relative
proportion in dyadic intervals, the image of \(t_{i,q}\) in
\((C_i/\ell_i C_i)_{\rm red}\) is algebraic over \(\kappa_i\).
No degree bound for these individual algebraic relations
is required.
\end{theorem}
\begin{proof}
At the start and after any fixed prefix of selected primes,
look at the positive-dimensional minimal components of the
current special fibers for all \(i\). This is a finite list
by Noetherianity. Choose a specialization from each by
Lemma~\ref{lem:constants499}. The reductions of the \(P_{i,q}\)
are nonzero because they are monic. Thus, outside a set of
relative proportion tending to zero, every new equation is
nonzero on each of its listed components, regardless of its
degree. The elliptic order constraint has the same relative
proportion one by Theorem~\ref{thm:commonsplit483}. A
positive-length chain of primes in a new special fiber lifts
to a chain upstairs to which one can prepend a minimal prime:
the new equation was outside every positive-dimensional minimal
component. Hence the maximal positive dimension decreases.
Initially each special fiber is a quotient of
\(\kappa_i[[U,V',W]]\). At most three steps make all of them
Artinian.

Each \(A_i\) is complete local and its residue modulo \(\ell_i\)
is finite-dimensional over \(\kappa_i\). A power of its maximal
ideal lies in \(\ell_i A_i\), so its maximal-ideal and
\(\ell_i\)-adic topologies are equivalent. Lift a basis of
\(A_i/\ell_i A_i\) and expand any element successively modulo
the powers of \(\ell_i\). Completeness of \(O_i\) makes the
lifts generate \(A_i\) over \(O_i\). By the hypothesis at the
diagonal trace, all imposed equations belong to \(D_i\). This
proves the assertion about the strict quotient. For finitely
many chosen specializations into finite extensions of
\(\operatorname{Frac}(O_i)\), one may take
\(P_{i,q}\) to be \(X-t_{0,i,q}\) times the product of the
minimal polynomials of their trace values. These are monic
integral polynomials since a discrete valuation ring is
integrally closed. The quotient then retains each chosen
specialization by the universal property of a quotient ring.

If \(C_i\) is finite, its reduced special fiber is its finite
residue field, and the trace condition holds at every good
prime. Otherwise its special fiber has positive dimension
by Lemma~\ref{lem:finitetrace471}. Choose one of the
specializations of the lemma on such a component. It kills
the nilradical. Every prime whose reduced universal trace
is algebraic gives a trace in the finite constant field of
this fixed local field, a set of relative proportion tending
to zero. This contradicts positive upper relative proportion.
\end{proof}

\begin{corollary}[The strict length with finite-module torsion]\leavevmode
\label{cor:strictvertical499}
For these \(A_i\) put \(U_i=A_i[\ell_i^\infty]\) and
\(A_i^{\rm f}=A_i/U_i\), and let \(b_i\) be the least \(b\ge0\)
with \(\ell_i^bU_i=0\). A further common good split prime \(q\)
can be chosen so that multiplication by
\(z_i=t_{i,q}-t_{0,i,q}\) on \(A_i^{\rm f}\) has nonzero
determinant for every \(i\). Then
\[
 k_i\le b_i+
       v_{\ell_i}\!\left(\det\big(z_i\mid A_i^{\rm f}\big)\right),
             \qquad v_{\ell_i}(\ell_i)=1 .
\]
The determinant is one for a zero free quotient. Its valuation
is the sum of the branch valuations with the dimensions of the
local generic algebras as multiplicities.
\end{corollary}
\begin{proof}
The generic fibers are finite algebras and have only finitely
many geometric points. Each point gives an integral local
specialization of \(C_i\) into a local field of characteristic
zero: the image of \(A_i\) is integral and its maximal ideal
has a power in \(\ell_i A_i\). It cannot annihilate \(D_i\)
since \(C_i/D_i\) is killed by \(\ell_i^{k_i}\). The
characteristic-zero trace separation in
Lemma~\ref{lem:mixedtrace490}, or its open-coset proof applied
to this point, holds with relative proportion one.
Consequently one good split prime separates all these points
from their diagonal traces. In the generic algebra \(z_i\)
is a unit. Multiplication on the finite free \(O_i\)-module
\(A_i^{\rm f}\) is injective, and the length of its cokernel
is the valuation of the determinant by elementary divisors.
Over an algebraic closure its action on each local Artinian
factor is its residue value plus a nilpotent, which proves the
multiplicity statement.

The image of \(U_i\) in the cyclic module \(A_i/D_iA_i\) has
length at most \(b_i\) because \(\ell_i^{b_i}\) kills it.
The remaining quotient is precisely \(A_i^{\rm f}/D_iA_i^{\rm f}\).
Since \(z_i\in D_i\), its length is at most the cokernel length
just computed. Adding the two contributions proves the bound.
\end{proof}
The polynomials permit direct comparisons with fixed integral
Hecke spectra whenever the relevant trace identities are
available. To deduce the uniform strict cohomology estimate
one still has to control the torsion exponents and the
determinant valuations with their weights on \(I\), or establish
the estimate by another argument.





\subsection{Tame inertia and the integral matrix algebra at the level primes}
\label{sec:markedmatrix715}

On the shorter family of Section~\ref{sec:levelmass678}, put
\(Q=(q_0d)^{s+6}\), where \(18480\mid q_0\) and \(2^s=m^{22}\).
The two level primes \(p_0,p_1\) split completely and
\(Q^2\mid p_i-1\). The principal rank is \(r=3^a\). At the
sufficiently large nodes \(a\ge1\) and \(m(4r+2)<2^s\).
Indeed \(m\ge\sqrt x\) and \(r\le3x^5\), as in
Corollary~\ref{cor:oddwindow677}. For any fixed coefficient
prime \(\ell\), one of \(p_0,p_1\) differs from \(\ell\). Over
that prime choose a place with a primitive label, as required
by the projection in Section~\ref{sec:levelmass678}. We will
also need the rank-one case of a label \(X-1\). Write \(q\)
for the norm of the chosen place.

A tame character over the residue field of characteristic \(q\)
can be read in characteristic zero by fixing compatible names
for roots of unity of order prime to \(q\). These identifications
preserve orders and \(q\)-th powers. Reduction of such a root
in residue characteristic \(\ell\ne q\) discards precisely
the \(\ell\)-part of its order. In fact roots of order prime to
\(\ell\) reduce injectively by Hensel's lemma, and a residue
field of characteristic \(\ell\) has no nonidentity
\(\ell\)-power roots of unity.

\begin{lemma}[The residual orbit of a primitive label]\leavevmode
\label{lem:tameorbit715}
Suppose \(q\equiv1\bmod 9\), \(r=3^a>1\), and \(Q^2\mid q-1\).
Let \(\xi\) generate the cyclic group of order \(q^r-1\).
If \(\beta^Q=1\), then \(\lambda=\beta\xi\) again has order
\(q^r-1\). For every \(\ell\ne q\) the \(r\) elements
\(\lambda,\lambda^q,\ldots,\lambda^{q^{r-1}}\) have distinct
reductions in characteristic \(\ell\). All their powers of
exponent \((q^r-1)/(q-1)\) are the same primitive
\((q-1)\)-st root of unity.
\end{lemma}
\begin{proof}
Set \(L=q^r-1\). Inside the group generated by \(\xi\),
\(\beta\) is a power of \(\xi^{L/Q}\). Each prime dividing
\(L\) divides \(L/Q\), because \(Q^2\mid L\). Thus
\(1+kL/Q\) is prime to \(L\) for every integer \(k\), proving
the primitivity of \(\lambda\).

For residual distinctness it suffices that \(q\) have order
\(r\) modulo a divisor of \(L\) prime to \(\ell\). If
\(\ell\ne3\), take the full \(3\)-part of \(L\). For
\(v_3(q-1)=b\) one has \(v_3(q^{3^j}-1)=b+j\), since
\(A^2+A+1\) has \(3\)-adic valuation one for \(A\equiv1\bmod3\).
Thus \(q\) has order \(3^a\) modulo \(3^{b+a}\). If
\(\ell=3\), put \(B=q^{r/3}\). The integer \(B^2+B+1>3\)
has \(3\)-adic valuation one, hence a prime divisor \(p\ne3\).
It divides \(q^r-1\), and it cannot divide \(B-1\) as their
common divisors with \(B^2+B+1\) divide three. The order
of \(q\bmod p\) divides \(r\) and fails to divide \(r/3\),
so equals \(r\). This settles every coefficient prime.
Finally raising a generator of order \(L\) to \(L/(q-1)\)
produces a generator of order \(q-1\), fixed by raising to
the power \(q\).
\end{proof}

Use the following local meaning of the prescribed type. On a
tame inertia generator \(\sigma\) at \(w\mid q\), the
inertial characters of a rank-\(r\) primitive label have
values \(\lambda^{q^j}\), \(0\le j<r\), where the roots of
the primitive label have been multiplied by the fixed
character value \(\beta\) of order dividing \(Q\).
The local Weil--Deligne inertia representation is tame
with these characters. In the rank-one case its character
value is \(\omega^{-j}\), where \(\omega\) has order \(Q\)
and \(1\le j\le m\). These are the two marked types used
here. A further twist unramified at the marked place
preserves them.

\begin{theorem}[A scalar and integral matrix units at a marked place]
\leavevmode\label{thm:markedmatrix715}
Let \(R\) be an \(\ell\)-adic representation with one of the
two marked local types above, and write \(d'=r\) or \(d'=1\)
for its dimension. After extending the coefficient field to
contain the type values, let \(\mathcal L\) be any invariant
lattice over its ring of integers \(\mathcal O\). The following
hold on \(\mathcal L\).

There is a power \(\eta\) of tame inertia on which \(R\) acts
by a scalar whose residual order is greater than \(4r+2\).
One has \(\kappa(\eta)=1\). For any arithmetic Frobenius
\(\Phi\) in the decomposition group at this place,
\(R(\Phi^{d'})\) is scalar, with
\(\kappa(\Phi^{d'})=q^{d'}\). The \(\mathcal O\)-algebra
generated by \(R(\sigma)\) and \(R(\Phi)\) is
\(\operatorname{End}_{\mathcal O}(\mathcal L)\).
\end{theorem}
\begin{proof}
The monodromy operator of the Weil--Deligne representation
commutes with inertia. In the primitive case the inertia
characters have distinct one-dimensional eigenspaces in
characteristic zero. Nilpotence makes monodromy zero on each
line. The same conclusion follows from dimension one in the
other case. Thus actual inertia on \(R\) is the specified
finite tame action.

In the primitive case let \(\eta=\sigma^{(q^r-1)/(q-1)}\).
Lemma~\ref{lem:tameorbit715} makes the action scalar of order
\(q-1\). In rank one take \(\eta=\sigma\) itself, with scalar
order \(Q/\gcd(Q,j)\). Choose \(p\in\{2,3\}\) different from
\(\ell\). As \(p^s\mid Q\), in the first case the order after
reduction is at least \(p^s\). In the second it is at least
\(p^s/j\). Both lower bounds exceed \(4r+2\) by the numerical
inequality at the start of the subsection. All tame inertia
here fixes the \(\ell\)-power roots of unity since \(q\ne\ell\).

Write \(A=R(\sigma)\), and denote its distinct eigenvalues
by \(\lambda_i\). Their pairwise differences are units in
\(\mathcal O\) by the lemma. The polynomial operators
\[
       P_i=\prod_{j\ne i}\frac{A-\lambda_j I}{\lambda_i-\lambda_j}
\]
therefore belong to \(\mathcal O[A]\). They give a direct sum
of rank-one lattices \(\mathcal L_i=P_i\mathcal L\).
Conjugation by arithmetic Frobenius raises a tame generator
to its \(q\)-th power in the tame inertia quotient. Hence
\(R(\Phi)\) permutes these lattices in one cycle, in the
order of the orbit or its reverse according to how the
idempotents are indexed. Each transition is an isomorphism
because \(\mathcal L\) is invariant under \(\Phi^{-1}\) as
well. In generators of the lines the transition factors are
units. After \(d'\) steps their product, independent of the
starting line, is the scalar of \(R(\Phi^{d'})\). For each
pair \(i,j\) some power \(k\ge0\) maps \(\mathcal L_i\) to
\(\mathcal L_j\). Then \(P_jR(\Phi)^kP_i\) generates the
rank-one module of maps between these two summands, extended
by zero on the others. Their span is the full matrix algebra.
For rank one the same statements follow with \(P_1=I\).
Finally arithmetic Frobenius at \(q\ne\ell\) has cyclotomic
value \(q\), proving the stated cyclotomic identity.
\end{proof}

\begin{corollary}[Unit determinants for the twisted main summand]
\leavevmode\label{cor:markeddet715}
Suppose \(R\) is a global representation of the prescribed
type and \(\theta\) is a character unramified at both marked
level primes. For each \(\ell\), the same auxiliary primes as
in Theorem~\ref{thm:trianglebudget714} can be chosen with
\[
            \det(1-(\theta R)(\phi_j))\in\mathcal O^\times
                 \qquad\text{for every }j ,
\]
after a finite extension of coefficients containing the values
in use. Here \(\theta\) has unit values, as holds for the
finite-order auxiliary twists. The bound on the total
valuation cost is the one in Corollary~\ref{cor:actualmarks714}.
The determinant is also a unit for geometric Frobenius.
\end{corollary}
\begin{proof}
By construction of the marked projection there is a primitive
label over each of the two level primes for the principal
rank, and the rank-one datum supplies its indicated character
over each. Choose the level prime different from \(\ell\).
The inertia element of Theorem~\ref{thm:markedmatrix715}
has the same scalar on \(\theta R\), since \(\theta\) is
unramified there. Its residual order exceeds the dimension.
Corollary~\ref{cor:scalarfrob713} therefore imposes the
determinant condition as part of the original prime selection.
The test count and the tame valuations in that selection do
not depend on the imposed unit condition. For an invertible
integral matrix \(T\),
\(\det(1-T^{-1})=(-1)^{d'}\det(T)^{-1}\det(1-T)\), giving the
last assertion.
\end{proof}

There is also a determinant bound at the other marked places
away from \(\ell\), where the labels can be any allowed
squarefree polynomials in the principal rank. For a root
from a factor of degree \(d''\), the scaled inertia eigenvalue
\(\rho\) has order dividing \(q^{d''}-1\). Indeed the scaling
factor belongs to \(\mathbb F_q^\times\). This \(\rho\) is
different from \(1\), since otherwise the unscaled label
would have a root in the subgroup of order \(Q\).
The same exclusion holds in rank one by the lower bound
on the scalar order. For a nontrivial root \(\rho\) of order
dividing \(L=q^{d''}-1\), the polynomial \(1+X+\cdots+X^{L-1}\)
vanishes at \(\rho\) and takes value \(L\) at \(1\). Thus
\(1-\rho\) divides \(L\) in the valuation ring. Every such
\(L\) divides the level group order \(G_0\) of
Corollary~\ref{cor:actualmarks714}. Multiplying the bounds
over the eigenvalues gives
\(v_\ell(\det(1-R(\sigma)))\le d'v_\ell(G_0)\) at any
such place. If \(R^c\simeq R^\vee\otimes\kappa^k\), the
matrix \(R(c\sigma^{-1}c^{-1})\) has the same eigenvalues
as \(R(\sigma)\), since \(\kappa(\sigma)=1\). This gives the same bound there.

These arguments evaluate the prescribed local types on an
invariant lattice. To use them for projection onto a global
space of lifts one must also identify which Hecke eigenspaces
carry the required Galois representations and control the
congruences between those eigenspaces.
\clearpage

\section{Ordinary symplectic equations and trace multiplicities}\label{nav:chapter-ordinary-symplectic-equations}

\subsection{The ordinary equations on the strict symplectic root}
\label{sec:strictordinary545}

The four-dimensional construction uses polynomial equations to
describe an ordinary local representation. We compute these equations
on the matrix entries that carry a strict cohomology class. The
calculation also applies when the residual character is trivial on
the entire coefficient-prime decomposition group.

Let \(p\ge5\), let \(O/\ZZ_p\) be a finite unramified extension
of discrete valuation rings, and write \(J=\operatorname{Frac}(O)\).
For \(N\ge1\) set \(D_N=(J/O)[p^N]\). With the scalar action of
a character understood, this module is isomorphic to \(O/p^NO\).
Consider the local characters of
Proposition~\ref{prop:localsymplecticroots429},
\[
    d=(A\Psi,\Psi,A,1),\qquad \theta=\Psi^{-1},\qquad
    X=E_{20}+E_{31}.
\]
Here indices run from zero to three. Write \(G_p=G_{\QQ_p}\)
and \(I_p\) for inertia. A continuous cocycle \(z\) for
\(D_N(\theta)\) gives matrices \(\rho_z(g)=(1+z(g)X)d(g)\)
over \(O\oplus D_N\), in which products of two elements of
the second summand are zero. This is the entry used in
Theorem~\ref{thm:symplecticclasses429}.

Fix \(s_0\in I_p\) for which the four units
\(\lambda_i=d_i(s_0)\) are distinct in \(J\). Set
\[
       B_i(T)=\prod_{k\ne i}(T-\lambda_k),\qquad
       \Delta_i=B_i(\lambda_i),\qquad t_0=\theta(s_0)-1,\qquad
       C=\frac{\Delta_0\Delta_2}{t_0}.
\]
The denominator \(t_0\ne0\) divides \(\Delta_0\) in \(O\), since
\(\lambda_2-\lambda_0=\lambda_0 t_0\).
The ordinary condition at \(s_0\) prescribes the characteristic
polynomial \(\prod_i(T-\lambda_i)\) for \(M=\rho_z(s_0)\) and
the equations
\begin{align*}
       B_j(M)\rho_z(g)B_i(M)&=0
              &&(g\in G_p,\ i<j),\\
       B_i(M)(\rho_z(s)-d_i(s)1)B_i(M)&=0
              &&(s\in I_p,\ 0\le i\le3).
\end{align*}
These are the original local equations in
\srcloc{Development/AutomorphicArithmetic/AutomorphicVarieties/SymplecticGSp4/MapRepApply}{99}.

\begin{proposition}[The obstruction to the ordinary equations at finite depth]
\leavevmode\label{prop:rootpolynomial545}
The ordinary condition on \(\rho_z\) is equivalent to the single
family of equations
\begin{equation}
     C\,\bigl(t_0 z(g)-(\theta(g)-1)z(s_0)\bigr)=0
                             \qquad(g\in G_p)
     \label{eq:rootcommutator545}
\end{equation}
in \(D_N\). In particular the class of the restriction of \(z\)
in \(H^1(G_p,D_N(\theta))\) is killed by \(\Delta_0\Delta_2\) whenever
the ordinary condition holds. Every local coboundary satisfies
the condition at the original depth.
\end{proposition}
\begin{proof}
The two coordinate planes with ordered bases \((e_0,e_2)\)
and \((e_1,e_3)\) are preserved. On either plane, with indices
\((i,j)\) in this order, the matrix at \(s_0\) is
\[
      \begin{pmatrix}
          \lambda_i&0\\ \lambda_i z(s_0)&\lambda_j
      \end{pmatrix},\qquad \lambda_j/\lambda_i=\theta(s_0).
\]
For any polynomial \(B\) its value on this matrix has diagonal
\((B(\lambda_i),B(\lambda_j))\) and lower-left entry
\[
       \lambda_i z(s_0)
       \frac{B(\lambda_i)-B(\lambda_j)}{\lambda_i-\lambda_j}.
\]
The divided difference is an integral polynomial expression
in the two units, as follows first for each monomial by summing
a finite geometric progression and then by linearity. If
\(k\notin\{i,j\}\), \(B_k\) vanishes at both diagonal entries
and \(B_k(M)\) is zero on this plane. On the plane in question
\(B_i(M)\) has only its first column and \(B_j(M)\) only its
second row. Multiplying the two matrices with \(\rho_z(g)\)
between them gives just the entry
\[
    \bigl(B_j(M)\rho_z(g)B_i(M)\bigr)_{ji}
       =d_i(g)\,\frac{\Delta_i\Delta_j}{t_0}
            \bigl(t_0z(g)-(\theta(g)-1)z(s_0)\bigr).
\]
This formula uses integral multiples of the elements of \(D_N\).
The diagonal inertia equations vanish on both planes. Indeed,
on the indicated plane the preceding first column has second
column zero, so sandwiching \(\rho_z(s)-d_i(s)1\) between
two copies of \(B_i(M)\) gives zero. With the operator having
only its second row, the same calculation uses its second
diagonal entry \(d_j(s)\). The characteristic polynomial
already has the required value by triangularity.

Put \(\mu=\lambda_0\lambda_3=\lambda_1\lambda_2\).
Replacing a pair of units \(u,v\) by \(\mu/u,\mu/v\)
preserves the valuation of \(u-v\). Thus
\(v_p(\Delta_0)=v_p(\Delta_3)\) and \(v_p(\Delta_2)=v_p(\Delta_1)\).
The two factors \(\Delta_i\Delta_j/t_0\) from the two planes
differ by a unit. This proves the equivalence. Equation
\eqref{eq:rootcommutator545} writes \(\Delta_0\Delta_2 z(g)\)
as the coboundary with witness \(Cz(s_0)\) at every \(g\).
Conversely, substituting \(z(g)=(\theta(g)-1)y\) makes its
left hand side zero for any \(y\in D_N\).
\end{proof}

\begin{theorem}[The cost of replacing a strict local condition]
\leavevmode\label{thm:strictordinary545}
Use a field \(E_x\) and a character \(\alpha\) from the intervals
of Theorem~\ref{thm:conductor232}. Suppose \(p\ge5\),
\(p\nmid n_*N_*\), \(2\le r<R=\lfloor V^{1203/2}\rfloor\),
\(p-1\mid2r\), and \(\alpha(c)=(-1)^r\). Use
\(\Psi=\alpha\kappa_p^r\) and the auxiliary \(A\) of the
symplectic construction. At finite depth \(N\), fix the same
global coefficient module \(D_N(\theta)\) and the same local
conditions everywhere away from \(p\). Let \(S_N\) be the
resulting group with zero localization at \(p\), and \(P_N\)
the group with the displayed ordinary equations there for the
source inertia element \(\kappa_p(s_0)=1+p\). Then
\[
  S_N\subset P_N,\qquad
  \operatorname{length}_O(P_N/S_N)
       \le 2\bigl(6+2v_p(r)+v_p(r-1)+v_p(r+1)\bigr).
\]
The estimate is uniform in \(N\) and in the unramified
coefficient extension used for the auxiliary character.
\end{theorem}
\begin{proof}
At \(p\), the finite conductor characters in \(d\) are
unramified, apart from the Teichmüller factor of the cyclotomic
character. This factor is one on \(s_0\). It follows that
\(\lambda_i\) are the powers of \(1+p\) with exponents
\((r+1,r,1,0)\), in that order. For \(a\ne b\) the valuation
of \((1+p)^a-(1+p)^b\) is \(1+v_p(a-b)\), by the binomial
theorem at the odd prime \(p\). Consequently
\[
     H:=v_p(\Delta_0\Delta_2)=6+2v_p(r)+v_p(r-1)+v_p(r+1).
\]
The criterion \eqref{eq:rootcommutator545} is linear in
the cocycle and is invariant under addition of a coboundary.
It thus defines a subgroup on cohomology. Localization
identifies \(P_N/S_N\) with a subgroup of
\(H^1(G_p,D_N(\theta))[p^H]\) by the proposition.

We bound this last group uniformly in \(N\). Put
\(k=O/pO\) and \(T_\theta=O(\theta)\). The residual dual
character \(\bar\kappa_p\bar\theta^{-1}\) is nontrivial on
inertia. Its exponent is \(1+r\), which at the two residues
equals \(1\) or \(1+(p-1)/2\) modulo \(p-1\), both nonzero
for \(p\ge5\). The divisible dual coefficient therefore has
no invariants, since some inertia element acts by a scalar
whose difference from one is a unit. Local Tate duality gives
\(H^2(G_p,T_\theta)=0\). Here the trace dual of \(O\) over
\(\ZZ_p\) is \(O\) since \(O\) is unramified. The
multiplication sequences now give
\[
       H^1(G_p,D_N(\theta))\simeq
          H^1(G_p,T_\theta)/p^NH^1(G_p,T_\theta).
\]
In particular its quotient modulo \(p\) is
\(H^1(G_p,k(\bar\theta))\). For this residual module local
duality gives zero in degree two as well, and the local Euler
characteristic formula gives dimension
\(1+\dim_k H^0(G_p,k(\bar\theta))\le2\) in degree one.
These local duality and Euler characteristic formulas are
also used in Proposition~\ref{prop:localsymplecticroots429}.
The finite group on the left is consequently generated by at
most two elements over \(O\), by Nakayama's lemma. Its
\(p^H\)-torsion has length at most \(2H\), as follows by
writing its elementary divisors over the discrete valuation
ring. This proves the estimate at every depth, including
the case \(\bar\Psi|_{G_p}=1\).
\end{proof}

\begin{corollary}[A simultaneous bound on the period fields]
\leavevmode\label{cor:collisionloss545}
At every \(x\in[V,V^{3/2}]\), sum the excess lengths
\(\operatorname{length}_O(P_N/S_N)\log p\) over the primes
and characters of the theorem and the colliding weights
\(2\le r<t(x)\). Count each geometric character of
\(\operatorname{Gal}(E_x/\QQ)\) once in its prescribed sign.
The sum is \(O(n_*R\log(3R))\), even when the finite depths
are chosen separately for each summand.
\end{corollary}
\begin{proof}
There are at most \(n_*=[E_*:\QQ]\) characters at every pair.
For one prime the number of weights with \(p-1\mid2r\) is at
most \(2R/(p-1)\), and such primes satisfy \(p\le2R+1\).
The constant term of the preceding bound, summed with \(\log p\),
therefore costs \(O(R\log(3R))\) before the character count.
Indeed \(\sum_{p\le y}\log p/(p-1)\ll\log(3y)\). On each
range \(m<p\le2m\) the product of the primes divides
\(\binom{2m}{m}\le4^m\). Taking \(m\) through powers of two
proves this estimate by a dyadic partition. For each \(r\ge2\),
even the sum over all primes of
\((2v_p(r)+v_p(r-1)+v_p(r+1))\log p\) equals
\(\log(r^2(r-1)(r+1))\ll\log(3R)\). Sum over the weights
and multiply by \(n_*\). The bound is polynomial in \(V\)
on the original family, hence \(o(\exp(x^\varepsilon/2))\)
uniformly there.
\end{proof}

This comparison measures the extra finite classes admitted by
the polynomial ordinary condition on the specified root.
Bounding the total number of classes satisfying that condition
still requires a global argument. The projected integral
theta-period estimate is also required on these fields in
Section~\ref{sec:missing}.


\subsection{Ordinary flags at finite precision}
\label{sec:finiteflag546}

Let \(O\) be a complete discrete valuation ring with uniformizer
\(\pi\), fraction field \(J\), and normalized valuation \(v\).
Write \(L=O^d\) and \(L_N=L/\pi^N L\).
Let \(I\) be a subgroup of a group \(G\), choose \(s_0\in I\),
and give characters \(\chi_i:I\to O^\times\) for \(0\le i<d\).
The units \(\lambda_i=\chi_i(s_0)\) are assumed pairwise
distinct in \(J\). Put
\[
    f(T)=\prod_i(T-\lambda_i),\quad
    B_i(T)=\prod_{j\ne i}(T-\lambda_j),\quad
    h_i=v(B_i(\lambda_i)),\quad h=\max_i h_i .
\]
An ordered ordinary flag for a representation on \(L_N\)
is a full flag of direct summands, stable under \(G\),
such that \(I\) acts by \(\chi_i\) on the \(i\)-th successive
line, in the given order. We first allow its reduction
to be any flag with these characters.

\begin{theorem}[Recovering a flag from the polynomial equations]
\leavevmode\label{thm:finiteflag546}
Let \(\rho_N:G\to\operatorname{GL}(L_N)\) and
\(M=\rho_N(s_0)\). An ordered ordinary flag gives the
identities over \(O/\pi^N\)
\begin{align*}
   \det(T-M)&=f(T),\\
   B_j(M)\rho_N(g)B_i(M)&=0 &&(i<j,\ g\in G),\\
   B_i(M)(\rho_N(s)-\chi_i(s)1)B_i(M)&=0 &&(s\in I).
\end{align*}
Conversely, if these equations hold and \(N>3h\), the
representation has an ordered ordinary flag modulo
\(\pi^{N-3h}\). If \(d\) is even, \(L\) has its standard
perfect alternating form, \(\rho_N\) has a prescribed
similitude character \(\mu:G\to O^\times\) modulo \(\pi^N\),
and \(\lambda_i\lambda_{d-1-i}=\mu(s_0)\), then the recovered
flag can be taken self-dual. For an existing pair of ordered
ordinary flags at a depth \(m>h\), their reductions already
agree modulo \(\pi^{m-h}\).
\end{theorem}
\begin{proof}
Number the steps starting with zero and write \(F_i\) for
the step ending with the line labeled \(i\), with \(F_{-1}=0\).
Given a flag, choose an adapted basis. The matrix \(M\) is
upper triangular with diagonal \(\lambda_i\) at the stated
precision. The factors of \(B_i(T)\) indexed by \(j>i\)
give the characteristic polynomial of \(M\) on \(L_N/F_i\).
The factors indexed by \(j<i\) give its characteristic
polynomial on \(F_{i-1}\). Cayley--Hamilton therefore gives
\(B_i(M)L_N\subset F_i\) and \(B_i(M)F_{i-1}=0\).
Stability of each step and the character on each quotient
now imply the two matrix equations.

For the converse take a matrix \(A\) over \(O\) lifting \(M\).
Its characteristic polynomial \(f_A\) satisfies
\(f_A\equiv f\bmod\pi^N\). At each \(\lambda_i\) its value
has valuation at least \(N\) and its derivative has valuation
exactly \(h_i\). Since \(N>2h_i\), Hensel's lemma gives a
root \(\lambda_i'\) of \(f_A\) with
\(v(\lambda_i'-\lambda_i)\ge N-h_i\). For \(i\ne j\) one
has \(v(\lambda_i-\lambda_j)\le h\), while \(N-h_i>h\).
The lifted roots are distinct and exhaust the roots of \(f_A\).
Every difference \(\lambda_i'-\lambda_j'\) has the valuation
of the corresponding original difference. Set
\[
       B_i'(T)=\prod_{j\ne i}(T-\lambda_j'),\qquad
       P_i=B_i'(A)/B_i'(\lambda_i').
\]
The prime here distinguishes the new polynomial. Its value in
the denominator has valuation \(h_i\). By interpolation the
\(P_i\) are the pairwise annihilating idempotents projecting
onto the labeled eigenlines in \(L\otimes J\).

The coefficients of \(B_i'-B_i\) lie in \(\pi^{N-h}O\).
For each \(g\) take any integral lift \(U_g\) of \(\rho_N(g)\).
The polynomial matrix equations first give divisibility by
\(\pi^N\) with \(B_j(A),B_i(A)\) as the two outside factors.
Replacing both by the primed polynomials retains divisibility
by \(\pi^{N-h}\), since all matrices involved at this stage
are integral. Dividing by the two scalar denominators gives,
with \(k=N-3h\),
\[
  P_j U_g P_i\in\pi^k\operatorname{End}_O(L)\quad(j>i),
   \qquad
  P_i(U_s-\chi_i(s)1)P_i\in\pi^k\operatorname{End}_O(L)
                                  \quad(s\in I).
\]
This uses a lift of each individual matrix. Their products
need only represent the group law at depth \(N\), where it
holds by hypothesis.

Intersect the successive sums of the labeled \(J\)-eigenlines
with \(L\) to form a saturated flag \(F_i\). Over a discrete
valuation ring each such step and each step in the preceding
quotient is a direct summand. For \(v_0\in F_i\) the part of
\(U_g v_0\) in the sum of the eigenlines above \(i\) is
\(\sum_{a\le i<j}P_j U_gP_a v_0\in\pi^kL\). After subtracting
it, the result is integral and belongs to \(F_i\). This
proves stability at depth \(k\). For \(s\in I\), project
\((U_s-\chi_i(s)1)v_0\) onto the sum of the eigenlines
numbered \(j\ge i\). All its terms are covered by the first
matrix divisibility except \(a=j=i\), which is covered by
the second. Thus the result belongs to \(F_{i-1}+\pi^kL\).
This gives the claimed inertia action on the reduction.

In the symplectic case choose the lift \(A\) to have multiplier
exactly \(\mu(s_0)\). Start with lifts of the first and last
columns in a symplectic basis, then scale the last by a
unit congruent to one modulo \(\pi^N\) to make their
pairing exact. Project the remaining columns onto their
orthogonal complement. All subtracted cross pairings were
already zero modulo \(\pi^N\). Repeating with each successive
pair gives the desired lift. Its eigenvalues are preserved
as a set by \(z\mapsto\mu(s_0)/z\). The disjoint Hensel
neighborhoods show that this map takes \(\lambda_i'\)
to \(\lambda_{d-1-i}'\). The perfect form pairs each
eigenline with its opposite and is zero between the other
lines. The resulting flag over \(J\) is self-dual. Taking
intersections with the unimodular \(L\) gives the same
orthogonal-complement identities in \(L\). They descend to
the reduction since the steps are direct summands.

Finally consider two flags at depth \(m\). Lift a basis
adapted to each flag to a basis of \(L\). In this basis
lift the matrix of \(M\) at depth \(m\) to an upper triangular
matrix with diagonal exactly \(\lambda_i\). In the original
basis the two matrices \(A_1,A_2\) agree modulo \(\pi^m\)
and both have characteristic polynomial \(f\). Their
projectors at \(i\) are \(B_i(A_\nu)/B_i(\lambda_i)\)
and differ by an endomorphism in
\(\pi^{m-h_i}\operatorname{End}_O(L)\). The sum of the
first \(i+1\) eigenlines for each matrix is exactly the
span over \(J\) of the first \(i+1\) adapted basis vectors,
by triangularity and the distinct diagonal entries. For
an integral vector in that span for \(A_1\), project it
onto the corresponding sum for \(A_2\). Its change lies
in \(\pi^{m-h}L\), so its image is also integral. Interchanging
the two matrices gives equality of the two steps after
reduction to depth \(m-h\). The proof of this last assertion
only uses the one matrix and its ordered actions on the
successive lines.
\end{proof}

\begin{theorem}[Sharp precision for the locally trivial residual character]
\leavevmode\label{thm:doubleflag546}
Take the four characters and the source element \(s_0\) of
Theorem~\ref{thm:strictordinary545}. The polynomial test
gives a self-dual ordinary flag at depth \(N-3h>0\), where
\[
       h=3+v_p(r)+\max\{v_p(r-1),v_p(r+1)\}.
\]
Suppose now that \(p-1\mid r\) and \(\alpha(p)=1\).
Write \(c=1+v_p(r)\). For any two ordinary flags lifting
the prescribed coordinate flag of a framed local representation
with residual diagonal \(\bar d\), equality holds modulo
\(p^{m-c}\) if the flags exist at depth \(m>c\). This loss
of \(c\) powers of \(p\) is attained on the split
representation over the same local field.
\end{theorem}
\begin{proof}
For the first assertion the four exponents at \(s_0\) are
\((r+1,r,1,0)\). Summing \(1+v_p(a-b)\) over the other
three exponents for each \(a\) gives the displayed maximum.
The similitude pairings of the four characters have the
prescribed product. Apply Theorem~\ref{thm:finiteflag546} to
the equations with the original inertia element.

For the sharper assertion, put \(O_m=O/p^mO\). Choose
an arithmetic Frobenius lift \(g\) with \(\kappa_p(g)=1\)
as in Proposition~\ref{prop:localsymplecticroots429}.
The residual action of \(g\) has eigenvalue \(\bar A(g)\ne1\)
at indices \(0,2\) and eigenvalue one at \(1,3\). The
auxiliary character in Lemma~\ref{lem:uniformoddroots430}
supplies this inequality at the indicated primes. Hensel
factorization of the characteristic polynomial of \(\rho_m(g)\)
over \(O_m\) separates the two relatively prime residual
factors of degree two. Bézout's identity for their monic
lifts and Cayley--Hamilton give complementary polynomial
idempotents. Write \(U,W\) for their images in the stated
order. Both are free of rank two. These summands are determined
by the representation and \(g\).

Every step of an ordinary flag is stable under these
idempotents because it is stable under \(\rho_m(g)\).
Its parts in \(U,W\) are direct summands, with ranks read
off from the prescribed residual flag. The same argument
applied to the quotient by the step shows that its two parts
are direct summands of \(U,W\) themselves. It follows that
the successive proper steps are
\[
             U_1,\qquad U_1\oplus W_1,\qquad U\oplus W_1
\]
for lines \(U_1\subset U,\ W_1\subset W\).
Compress \(\rho_m(s_0)\) to \(U,W\) using the respective
polynomial idempotents. In a basis for the displayed flag
ordered by indices \(0,1,2,3\), the matrix of \(\rho_m(s_0)\)
is upper triangular with entries \(\lambda_i\) on the diagonal.
Consequently its two compressions are upper triangular on
\((U_1,U/U_1)\) and \((W_1,W/W_1)\), with labeled entries
\((\lambda_0,\lambda_2)\) and \((\lambda_1,\lambda_3)\).
This calculation only compresses one matrix and requires no
invariance of the two full summands under inertia.

In each pair the difference of the entries has valuation
\(v_p((1+p)^r-1)=c\). Apply the two-dimensional, single-matrix
uniqueness argument at the end of Theorem~\ref{thm:finiteflag546}
to each compression. The two possibilities for each line
agree at depth \(m-c\). The three displayed steps therefore
agree there as well.

For sharpness use the split diagonal representation \(d\).
The finite character \(\alpha\) is unramified at \(p\) and
equals one on Frobenius, so it is trivial locally. Since
\(p-1\mid r\), the image of \(\theta=\kappa_p^{-r}\)
is contained in \(1+p^cO\) on all of \(G_p\), by the
binomial valuation formula. With \(X=E_{20}+E_{31}\) the
matrix \(1+p^{m-c}X\) is symplectic and commutes modulo
\(p^m\) with every \(d(g)\), since \(d(g)Xd(g)^{-1}=\theta(g)X\).
It sends the coordinate flag to a second self-dual ordinary
flag with the same ordered characters and the same reduction
modulo \(p\). The first line is now generated by
\(e_0+p^{m-c}e_2\). The lines are distinct modulo
\(p^{m-c+1}\), which proves the assertion.
\end{proof}

For a residual flag fixed in advance the reduction of a
recovered flag must be checked against that choice. Both
theorems concern the coefficient-prime local equations.
The global number of ordinary classes and its relation to
the integral period in Section~\ref{sec:missing} remain to
be bounded.


\subsection{The global equations for the two even roots}
\label{sec:generaleven547}

Fix one of the period fields \(E_x\) on the interval of
Section~\ref{sec:missing}. Let \(\alpha\) be a character of this field,
\(2\le r<t(x)\), \(\alpha(c_\infty)=(-1)^r\), and let \(\ell\ge5\)
be prime with \(\ell\nmid n_*N_*\) and \(\ell-1\mid 2r\).
Write \(\Psi=\alpha\kappa_\ell^r\), and suppose
\(\bar\Psi|_{D_\ell}\ne1\), where \(D_\ell=G_{\QQ_\ell}\).
At the half residue this holds on inertia. At the zero residue it
means \(\alpha(\ell)\ne1\). Use the interval character \(\xi\) of
Proposition~\ref{prop:oddstrictalllevels430}, the resulting
\(A=\kappa_\ell\omega_\ell^{-1}\xi\), and the unramified coefficient
ring \(O\) containing the values of the finite characters. Put
\(J=\operatorname{Frac}(O)\), \(k=O/\ell O\). The split
representation and its multiplier are
\[
 d=\operatorname{diag}(A\Psi,\Psi,A,1),\qquad \mu=A\Psi.
\]
For a character \(\gamma\) and a coefficient module \(V(\gamma)\),
write \(S(V(\gamma))\) for the kernel in \(H^1(\QQ,V(\gamma))\)
of restriction to \(D_\ell\) and to inertia at every other finite
prime. Define
\[
 a_\pm=\dim_k S\bigl(k(\bar\Psi^{\pm1})\bigr),\qquad
 S^\pm=S\bigl((J/O)(\Psi^{\pm1})\bigr).
\]
Thus \(S^-=S(\alpha,r)\) after extending coefficients. The integer
\(a_-\) may be greater than one, and both residual groups may occur.

Choose \(h\in D_\ell\) as in Lemma~\ref{lem:sympflag468}.
Let \(\Sigma\) consist of \(\ell,2\) and the finite conductor primes
of \(d\). Let \(\mathcal B\) represent the framed global lifts of
Section~\ref{sec:sympglobal469}, now for the data just specified.
A lift has multiplier \(\mu\), is unramified off \(\Sigma\), kills
\(N_q=I_q\cap\ker\bar d\) for \(q\in\Sigma\setminus\{\ell\}\),
and at \(\ell\) has the prescribed ordinary flag and is diagonal at
\(h\). Take \(u_1,u_2,a,x_0,b,c\) for the local parameters of
Theorem~\ref{thm:sympordinary468}, writing \(x_0\) for its \(02\)
entry to distinguish it from the interval parameter \(x\).

\begin{theorem}[The global presentation at both colliding residues]
\leavevmode\label{thm:allglobalroots547}
Over \(R_\ell=O[[u_1,\allowbreak u_2,a,x_0,b,c]]\) the ring \(\mathcal B\)
has \(a_+\) further parameters \(P_i\) of torus weight \((0,1)\)
and \(a_-\) parameters \(Y_j\) of weight \((0,-1)\). Its ideal
of relations admits generators of the following weights, with zero
generators allowed:
\[
\begin{array}{c|c|c}
 \text{weight}&\text{generators}&\text{number}\\ \hline
 (0,0)&Z_\nu&2\\
 (0,1)&E_\nu&a_++1\\
 (0,-1)&F_\nu&a_-+1 .
\end{array}
\]
One can take the \(P_i,Y_j\) to be \(02,20\) entries, respectively,
of global matrices. Eliminating \(u_1,u_2\) gives
\[
 \mathcal B^{s}=
 O[[x_0,v,w,P_1,\ldots,P_{a_+},Y_1,\ldots,Y_{a_-}]]
       /(E_\nu,F_\nu),\qquad v=ab,\quad w=bc,
\]
where the superscript denotes invariants under the first factor of
\(\operatorname{diag}(st,t,s,1)\). The remaining weights are \(1\) on
\(x_0,v,P_i\), \(2\) on \(w\), and \(-1\) on \(Y_j\).
\end{theorem}
\begin{proof}
The four residual eigenvalues at \(h\) are distinct by
Lemma~\ref{lem:sympflag468}. Consequently the ordinary local ring
with this normalization is precisely \(R_\ell\). We compute the
relative deformation spaces separately for each torus weight \(M\)
in the symplectic Lie algebra. They are the cohomology of the
degree-shifted mapping cone in Theorem~\ref{thm:symprelations469},
where the maps compare a global cochain with the fixed local
cochain at \(\ell\) and with cochains through \(G_{\QQ_q}/N_q\)
at the other \(q\in\Sigma\). Here the image of \(I_q\) in that
quotient has order prime to \(\ell\), as each diagonal character
on inertia has this property.

In degree one one may set the local witness at \(\ell\) to zero
by subtracting a global coboundary. The witnesses at the other
primes can then be absorbed in local quotient coboundaries. On
each root line the remaining representative is unique because
the eigenvalue ratio at \(h\) differs from one by a unit. Its
vanishing on \(N_q\) is equivalent to the inertia class being zero,
by averaging on \(I_q/N_q\). Thus the relative tangent dimensions
on the even roots are \(a_+\) and \(a_-\). On the odd roots they
are zero by Proposition~\ref{prop:oddstrictalllevels430}, since
\(T_\alpha=\{\ell\}\). On the diagonal the cocycles are everywhere
unramified additive characters of \(G_\QQ\), hence zero by class
reciprocity for \(\QQ\).

In the same cone \(H^0=0\), by injectivity on degree-zero
invariants, and \(H^3=0\), by Poitou--Tate duality. For the
latter assertion the dual global invariant space vanishes already
at \(\ell\). For the neutral and even weights the tame exponents
are \(0\) or \((\ell-1)/2\), both distinct from \(1\) modulo
\(\ell-1\). For the odd roots use
Proposition~\ref{prop:localsymplecticroots429}. The quotient local
groups at \(q\ne\ell\) have \(\ell\)-cohomological dimension at
most one by inertia averaging. Their Euler characteristics and
those of \(G_{\QQ_q}\) are zero. Subtracting the Euler characteristic
at \(\ell\) from the global one now gives
\[
       \dim H^2=\dim H^1+\dim M^{c_\infty=1}.
\]
This adds two to the neutral dimension and one to each even root
dimension. It adds zero to the odd dimensions.

To realize the relative parameters, consider the linear
functionals that evaluate a normalized cocycle at all elements
of \(G_\Sigma\). They separate the cocycles. On each finite
dimensional root space a subset of these evaluations is a basis
of its dual. Multiplication by the corresponding unit diagonal
entry of \(d(g)\) identifies these evaluations with the indicated
matrix entries on relative tangents. Complete Nakayama gives a
surjection from \(R_\ell[[P_i,Y_j]]\). The obstruction calculation
in the proof of Theorem~\ref{thm:symprelations469} bounds the
number of minimal relations of each weight by the corresponding
\(\dim H^2\) just computed. Indeed the defects of lifted cochains
define its cone obstruction. A functional on the minimal relation
space killing this obstruction would lift the universal matrices
across its square-zero extension while fixing \(R_\ell\). Varying
any of the \(P_i,Y_j\) by elements of the small kernel preserves
every relation there, since their relative linear terms vanish
modulo the maximal ideal. This contradicts a nonzero functional.
Taking Laurent coefficients of the two formal torus variables on
finite jets gives homogeneous generators, with the same bounds
after passage to the limit.

Imposing zero on all off-diagonal parameters leaves the diagonal
deformations. Each diagonal character differs from the corresponding
one of \(d\) by an everywhere unramified character reducing to
one, hence equals it by class reciprocity. Only \(Z_1,Z_2\) survive
on \(O[[u_1,u_2]]\). They generate \((u_1,u_2)\), and the formal
inverse function theorem eliminates these variables equivariantly.
A monomial in the remaining variables has first weight zero
exactly when its exponent of \(b\) is the sum of the exponents
of \(a\) and \(c\). It has a unique expression using \(v,w\)
in their place. All remaining equations have first weight zero.
Projection onto that weight on their coefficients identifies the
invariant ideal, giving the asserted presentation.
\end{proof}

\begin{theorem}[The integral equations at every finite depth]
\leavevmode\label{thm:integralroots547}
In the presentation of Theorem~\ref{thm:allglobalroots547} keep
only the variables \(Y_j\) and set the other variables equal to
zero. The resulting equations are the rows of a matrix
\(B_-\in\operatorname{Mat}_{a_-+1,a_-}(O)\) acting on the column
of \(Y_j\). Keeping only \(P_i\), including the condition \(x_0=0\),
gives a matrix \(B_+\in\operatorname{Mat}_{a_++1,a_+}(O)\).
Put \(K_\pm=\operatorname{coker}(B_\pm^{\,t})\). At every
\(m\ge1\) there are \(O\)-linear isomorphisms
\[
 S^\pm[\ell^m]\simeq\operatorname{Hom}_O(K_\pm,O/\ell^mO).
\]
In particular, the ideal of the \(a_-\)-rowed minors of \(B_-\)
is \(\ell^{\,\operatorname{length}_O S^-}O\). For \(a_-=0\)
the minor of size zero equals one.
\end{theorem}
\begin{proof}
With only the \(Y_j\) present, a monomial of weight \(-1\) is a
single \(Y_j\) and no monomial has weight \(+1\). This leaves
precisely \(a_-+1\) linear equations over \(O\). The other
restriction gives \(a_++1\) linear equations by the same weight
count. The eliminated diagonal parameters restrict to zero in
both cases since they are neutral and their split values are zero.

Set \(O_m=O/\ell^mO\) and substitute \(\epsilon z_j\),
\(z_j\in O_m\), for the retained variables in
\(O_m[\epsilon]/(\epsilon^2)\). Such substitutions in the full
ring are in bijection with the kernel of \(B_\pm:O_m^{a_\pm}
\to O_m^{a_\pm+1}\). In the universal matrices their only
possible first-order entries have the selected weight.
The multiplier equation gives the root vector \(X_-=E_{20}+E_{31}\)
or \(X_+=E_{02}+E_{13}\) for the form
\(\langle e_0,e_3\rangle=\langle e_1,e_2\rangle=1\).
The matrices are \((1+\epsilon z(g)X_\pm)d(g)\), and the
group law says \(z(gg')=z(g)+\Psi^{\pm1}(g)z(g')\).
All six local parameters are split, so \(z\) vanishes on
\(D_\ell\). At \(q\in\Sigma\setminus\{\ell\}\) it vanishes on
\(N_q\). It is unramified away from \(\Sigma\). Conversely each
such cocycle defines a lift with just these parameter values.
The retained coordinates evaluate the cocycle multiplied by
fixed units, giving an \(O_m\)-linear correspondence.

Division by \(\ell^m\) identifies \(O_m\) with \((J/O)[\ell^m]\).
Every divisible-coefficient class has a unique representative
zero at \(h\), since \(\Psi^{\pm1}(h)-1\) is a unit. The
representative of a class killed by \(\ell^m\) therefore takes
values in the indicated finite module. For a strict class it is
zero on \(D_\ell\), because a coboundary there vanishing at \(h\)
has zero witness. At other conductor primes \(N_q\) acts trivially
on the entire coefficient module. A coboundary on inertia vanishes
on \(N_q\), and conversely averaging on the prime-to-\(\ell\)
quotient makes any such inertia restriction a coboundary.
These facts identify exactly the cocycles counted above with
\(S^\pm[\ell^m]\). The definition of the transpose cokernel now
gives the claimed Hom groups.

Theorem~\ref{thm:levelstrict345} gives finiteness of \(S^-\) on
these fields. Extension to \(O\) is unramified and preserves its
module length. A free summand of \(K_-\) would make the lengths of
its Hom groups unbounded with \(m\). Thus \(K_-\) is torsion.
For a cyclic summand \(O/\ell^bO\) the Hom length is
\(\min(b,m)\). Taking \(m\) larger than all elementary divisors
shows \(\operatorname{length}_O K_-=\operatorname{length}_O S^-\).
In Smith normal form the ideal of maximal minors of \(B_-\) is
generated by the product of these elementary divisors. This proves
the formula, including the case of the empty presentation.
\end{proof}

\begin{corollary}[A bound for the sizes of the presentation matrices]
\leavevmode\label{cor:rootranks547}
Let \(L\) be the fixed field of \(\ker\bar\Psi\), and let \(d_L=[L:\QQ]\).
On the whole indicated part of the interval,
\[
 a_+,a_-\ \le\ \frac{\log|\Cl(L)|}{\log\ell}
 \ \le\ \frac{d_L\log(\ell N_*\zeta(2))}{\log\ell},
 \qquad d_L\le 2n_* .
\]
\end{corollary}
\begin{proof}
The cyclic group \(\operatorname{Gal}(L/\QQ)\) has order prime to
\(\ell\). For either character, restriction injects its residual
strict group into homomorphisms from the class group of \(L\) to
the additive group of \(k\). Indeed restriction turns the cocycle
into an everywhere unramified homomorphism, zero at primes over
\(\ell\), and the kernel of restriction vanishes by averaging.
The \(k\)-dimension of the full space of homomorphisms is at most
\(\log|\Cl(L)|/\log\ell\). The character cutting out \(L\) is a
product of a character of order at most \(n_*\) and a character
of order at most two. Its conductor divides \(\ell N_*\). The
conductor--discriminant formula gives \(|D_L|\le(\ell N_*)^{d_L}\).
Minkowski's bound supplies an integral ideal of norm at most
\(\sqrt{|D_L|}\) in each ideal class, since \(L\) is totally
real. The number of ideals of norm at most \(H\ge1\) is at most
\(H^2\zeta_L(2)\), and the Euler factors give
\(\zeta_L(2)\le\zeta(2)^{d_L}\). Substitution proves the stated
inequality. It controls the dimensions of the matrices. The
valuation of their minors still enters the length in
Theorem~\ref{thm:integralroots547}.
\end{proof}


\subsection{The trace algebra and the negative root module}
\label{sec:roottrace548}

Use the fields, intervals, coefficient primes and characters of
Section~\ref{sec:generaleven547}. Put \(p=a_++2\), \(m=a_-\), and
write \(X_1,\ldots,X_p\) for \(x_0,P_1,\ldots,P_{a_+},v\), in
that order. Here \(p\) counts coordinates; the coefficient prime
is \(\ell\). With weights \(1,2,-1\) on \(X_i,w,Y_j\), respectively,
the first torus invariant algebra has the presentation
\[
 B=\mathcal B^s=B_0/(E_\nu,F_\nu),\qquad
 B_0=O[[X_1,\ldots,X_p,w,Y_1,\ldots,Y_m]].
\]
The weights of \(E_\nu,F_\nu\) are \(+1,-1\). A weight space of a
completed ring means the series transforming by the specified
character on every finite jet. Let \(T\) be the closed algebra in
\(\mathcal B\) generated over \(O\) by the traces of the universal
representation \(\rho\). Write \(D\) for its split augmentation
ideal and \(k_-=\operatorname{length}_O S(\alpha,r)\).

\begin{theorem}[Trace generators and integral relations]
\leavevmode\label{thm:allroottrace548}
The remaining torus invariants in \(B\) form \(T\). They are
topologically generated over \(O\) by
\[
 U_{ij}=X_iY_j\quad(1\le i\le p,\ 1\le j\le m),\qquad
 W_{jk}=wY_jY_k\quad(1\le j\le k\le m).
\]
Write \(W_{jk}=W_{kj}\) when \(j>k\). The relations among these
monomials in \(B_0\) are generated by
\[
\begin{aligned}
 U_{ij}U_{ht}&=U_{it}U_{hj},&
 U_{ij}W_{kt}&=U_{ik}W_{jt},&
 W_{ij}W_{kt}&=W_{ik}W_{jt},
\end{aligned}
\]
with indices in the ranges of their respective symbols. To obtain
\(T\), impose also the neutral equations
\[
       Y_j E_\nu=0,\qquad X_i F_\nu=0,\qquad wY_j F_\nu=0 .
\]
Each left-hand side is expressed in the monomial generators. The
images of all \(U_{ij},W_{jk}\) generate \(D\).
\end{theorem}
\begin{proof}
For a neutral monomial assign one of its negative factors to
each \(X\) and two to each \(w\). Equality of weights uses up
every \(Y\). Two monomials in the resulting symbols have the
same image precisely when the multiplicities of every \(X_i\),
of \(w\), and of each \(Y_j\) agree. Match the \(X\)-factors
by type and match the \(w\)-factors in any order. The allocations
of the \(Y\)'s now differ by a permutation of the occupied slots.
A transposition exchanges two single slots, a single and a double
slot, or slots of two double factors. The displayed identities
give these exchanges, with symmetry handling an exchange in one
double factor. Grouping polynomial terms with the same image
shows that these binomials generate the whole polynomial kernel
over \(O\). The identities persist on completion. Indeed for a
fixed image the number of symbols is the total number of its
\(X\)'s and \(w\)'s; their degrees as original monomials are two
or three. Thus the two filtrations, including powers of \(\ell\),
are cofinal. The finitely generated binomial ideal is closed in
the complete Noetherian ring.

Every monomial of weight \(-1\) has a \(Y_j\) which leaves a
neutral factor when removed. Each monomial of weight \(+1\)
has either an \(X_i\) or a factor \(wY_j\) with the same property.
Indeed, in the absence of an \(X_i\), the number of \(Y\)'s is
twice the number of \(w\)'s minus one. Grouping series by a
chosen such factor proves the two generation statements also
for power series. Project an expression in the homogeneous ideal
\((E_\nu,F_\nu)\) onto weight zero. Its coefficients on the
\(E_\nu\) have weight \(-1\) and those on the \(F_\nu\) have
weight \(+1\). This gives exactly the listed additional equations.

It remains to recover these invariant functions from traces.
Newton's identities applied to \(\operatorname{tr}\rho(h^j)\)
for \(1\le j\le4\) give the coefficients of the characteristic
polynomial of \(\rho(h)\), since \(\ell\ge5\). Newton iteration
from its four distinct residual roots gives their labeled lifts
in \(T\): each denominator has unit residue, its inverse belongs
to the closed algebra by a geometric series, and the iterations
converge. The four coordinate projectors \(e_{ii}\) are therefore
polynomials in \(\rho(h)\) over \(T\). Inserting such projectors
between matrices and taking the trace shows that any cyclic
product of matrix entries belongs to \(T\). Expanding the
projectors expresses it in traces of group elements.

The entries for \(x_0,P_i\) have indices \(02\), and those for
\(Y_j\) have indices \(20\). Thus their pairwise products are
cycles. Since \(v=ab\) uses \(01,12\), its product with \(Y_j\)
is a cycle as well. For \(W_{jk}=bcY_jY_k\) use the multiplier
identity
\[
  \rho(g_j)_{20}=-\mu(g_j)\rho(g_j^{-1})_{31},
     \qquad Y_j=\rho(g_j)_{20}.
\]
The indices \(03,31,12,20\) now form a cycle. All generators
belong to \(T\), and every trace is invariant under conjugation.
The invariant series converge in the original topology. This
proves equality with the closed trace algebra. Evaluating at
the split representation kills exactly its ideal of monomials
of positive degree in the displayed generators.
\end{proof}

\begin{theorem}[The Fitting ideal at the split trace]
\leavevmode\label{thm:tracefitting548}
Let \(M\) be the weight \(-1\) space of \(B=\mathcal B^s\),
viewed as a module over \(T\). It is generated by the \(Y_j\).
Its zeroth Fitting ideal \(I=\operatorname{Fitt}_0^T(M)\)
satisfies
\[
    DI=0,\qquad M/DM\simeq K_-,\qquad
    (I+D)/D=\ell^{k_-}O\ \subset\ T/D=O .
\]
In particular \(C=T/I\) is defined by the module itself and
\(C/DC\simeq O/\ell^{k_-}O\).
\end{theorem}
\begin{proof}
Before imposing any equations in \(B_0\), extract the coefficients
of the monomials \(Y_1,\ldots,Y_m\) from a series of weight \(-1\).
This maps the weight space onto \(O^m\). Its kernel consists of
the products of this space with the neutral augmentation ideal:
subtracting the linear terms, each remaining monomial is \(Y_j\)
times a nonconstant neutral monomial. Such expressions for series
follow by distributing the monomials among the finitely many
generators. In the quotient \(B\), the additional relations in
this \(O^m\) come exactly from the coefficients of \(Y_j\) in
the \(F_\nu\). A multiplier of an \(F_\nu\) has weight zero and
only its scalar constant can contribute. A multiplier of an
\(E_\nu\) has weight \(-2\); every such monomial has at least two
\(Y\)'s and gives no linear term. The resulting rows are those
of \(B_-\) in Theorem~\ref{thm:integralroots547}. Exact projection
onto weight \(-1\) commutes with taking the homogeneous quotient,
so \(M/DM=\operatorname{coker}(B_-^t)\).

The module is finite over the Noetherian ring \(T\), hence
finitely presented. Its zeroth Fitting ideal is generated by
the maximal minors of any relation matrix on a finite generating
set. This operation commutes with base change by tensoring the
presentation. Such minors annihilate the module, by applying
the adjugate to each square subsystem. The convention for zero
generators gives the unit ideal for the zero module. Applying
base change to \(T\to O\) and using
Theorem~\ref{thm:integralroots547} gives the stated image of \(I\).
Since \(I\) kills every \(Y_j\) in \(B\), it kills their products
\(U_{ij}\) and \(W_{jk}\), which generate \(D\). Taking the
quotient by \(I+D\) gives the last assertion.
\end{proof}

\begin{corollary}[Separation over the fraction field and good-prime traces]
\leavevmode\label{cor:tracesplit548}
There is a surjection
\[
 T\twoheadrightarrow O\times_{\,O/\ell^{k_-}O} C .
\]
Its kernel has square zero and is killed by \(\ell^{k_-}\). In
particular \(T[1/\ell]=J\times C[1/\ell]\), and \(T\) is finite
over \(O\) if and only if \(C\) is finite over \(O\). At most
\[
        pm+\frac{m(m+1)}2
\]
good-prime trace differences \(\delta_q=\operatorname{tr}
\rho(\operatorname{Frob}_q)-\operatorname{tr}d(\operatorname{Frob}_q)\),
with \(q\notin\Sigma\), generate \(D\). For such a choice the
quotient of \(C\) by these differences is \(O/\ell^{k_-}O\).
\end{corollary}
\begin{proof}
The difference of two representatives of a compatible pair lies
in \(D+I\); subtracting appropriate elements of these two ideals
gives a common lift. The kernel is \(D\cap I\), whose square
vanishes by \(DI=0\). Lift the scalar \(\ell^{k_-}\) to an element
of \(I\) using Theorem~\ref{thm:tracefitting548}. Its difference
from that scalar is in \(D\). Both terms annihilate \(D\cap I\),
so \(\ell^{k_-}\) does too. Inverting \(\ell\) gives the product.
As \(T\) is Noetherian and \(DI=0\), the ideal \(I\) is finitely
generated over \(T/D=O\). The sequence with quotient \(C\)
gives the finiteness equivalence.

The trace algebra is a continuous image of a formal power series
ring in the stated number of variables over \(O\). The latter
is compact, so its quotient topology agrees with the subspace
topology in \(\mathcal B\), which is Hausdorff. Ideals in \(T\)
are closed. By Chebotarev in the finite quotients of \(G_\Sigma\),
the ideal of the good-prime differences contains the differences
of traces at every group element. In its quotient all traces are
scalars. The scalar subring is compact and closed, and the traces
topologically generate, so this quotient is \(O\) by the split
augmentation. Thus these differences generate \(D\). Their
images span \(D/\mathfrak m_TD\), of dimension at most the
number of monomials generating \(D\). Take a basis from these
images and apply Nakayama's lemma. The final quotient follows
from Theorem~\ref{thm:tracefitting548}.
\end{proof}


\subsection{Integral realization of Frobenius traces}
\label{sec:realizeall549}

Keep the arithmetic data and the indexing of
Theorem~\ref{thm:allroottrace548}; in particular \(X_p=v=ab\).
Let \(H\) be a complete discrete valuation ring with finite residue
field and a local coefficient map from \(O\). We allow equal or
mixed characteristic. Fix a continuous local \(O\)-algebra map
\(\sigma:T\to H\), and write \(\nu\) for the valuation on \(H\)
normalized by \(\nu(H^\times)=0\) and \(\nu(\varpi_H)=1\).
Define three minima of valuations of its trace coordinates:
\[
\begin{split}
 \nu_2&=\min_{i<p,\,j}\nu\bigl(\sigma(U_{ij})\bigr),\\
 \nu_3&=\min_j\nu\bigl(\sigma(U_{pj})\bigr),\qquad
 \nu_4=\min_{j,k}\nu\bigl(\sigma(W_{jk})\bigr).
\end{split}
\]
The value of zero and the minimum of an empty set are \(+\infty\).

\begin{theorem}[The exact ramification index for an integral realization]
\leavevmode\label{thm:ramification549}
Let \(H'/H\) be the valuation rings in a finite extension of
fraction fields, with ramification index \(e\). There is a continuous
local \(\mathcal B\to H'\) inducing \(\sigma\) on traces if and only if
\[
                      e\nu_i\ge i\qquad(i=2,3,4).
\]
In particular the least possible index is
\(\max(1,\lceil2/\nu_2\rceil,\lceil3/\nu_3\rceil,
\lceil4/\nu_4\rceil)\), taking \(1/\infty=0\). It is achieved by
a totally tamely ramified extension of index at most four.
\end{theorem}
\begin{proof}
In an integral local realization every chosen off-diagonal
parameter lies in the maximal ideal. The coordinates in the
three minima are products of respectively two, three and four
such entries: they are \(X_iY_j\) for \(i<p\), \(abY_j\), and
\(bcY_jY_k\). Their valuations in \(H'\) are therefore at least
two, three and four. This proves necessity.

If all coordinates vanish, take the split representation. Otherwise
form the matrix over \(H\) whose rows are the vectors
\((\sigma(U_{ij}))_j\) and \((\sigma(W_{kj}))_j\). All \(2\)-rowed
minors vanish by Theorem~\ref{thm:allroottrace548}, including for
one row of each kind by the symmetry of \(W\). Choose a nonzero
row and an index \(j_0\) where its entries have minimal valuation.
Divide it by its \(j_0\)-th entry to obtain
\(\lambda=(\lambda_j)\in H^m\) with \(\lambda_{j_0}=1\). All rows
are multiples of this one over the fraction field, hence each
row is its own \(j_0\)-th entry times \(\lambda\). In particular
symmetry gives
\[
 \sigma(U_{ij})=u_i\lambda_j,\qquad
 \sigma(W_{jk})=w_0\lambda_j\lambda_k,\qquad
 u_i=\sigma(U_{ij_0}),\quad w_0=\sigma(W_{j_0j_0}).
\]
Thus \(\nu_2=\min_{i<p}\nu(u_i)\), \(\nu_3=\nu(u_p)\), and
\(\nu_4=\nu(w_0)\).

Take any uniformizer \(s\) of \(H'\). In the full global
presentation after elimination of the neutral variables set
\[
 Y_j=s\lambda_j,\quad X_i=u_i/s\ (i<p),\quad
 b=s,\quad a=u_p/s^2,\quad c=w_0/s^3 .
\]
The inequalities put each assigned value in the maximal ideal.
They give \(v=u_p/s\), \(w=w_0/s^2\) and all the prescribed trace
coordinates. Every relation \(E_\nu\) evaluates to zero since
\(Y_{j_0}E_\nu\) is an invariant relation and \(Y_{j_0}=s\ne0\).
Some \(X_i\) or \(wY_{j_0}\) has nonzero value because the trace
coordinates are not all zero. Its product with each \(F_\nu\)
is an invariant relation, so \(F_\nu\) evaluates to zero too.
The neutral substitutions converge, giving a continuous local
map on all of \(\mathcal B\). If \(u_p=w_0=0\), the same
verification permits \(a=b=c=0\).

Every finite \(\nu_i\) is a positive integer. The least integer
in the assertion is consequently at most four and is prime to
the residual characteristic \(\ell\ge5\). Adjoining a root of
\(Z^e-\varpi_H\) gives a totally ramified extension of that index
by Eisenstein's criterion. This also proves attainability in equal
characteristic.
\end{proof}

\begin{theorem}[The two-plane locus and its Frobenius determinants]
\leavevmode\label{thm:purityblocks549}
Suppose the trace coordinates of \(\sigma\) are not all zero.
If some \(\sigma(U_{pj})\) or \(\sigma(W_{jk})\) is nonzero, its
realization above is absolutely irreducible. If all these coordinates
are zero, it has a realization as the direct sum \(V_-\oplus V_+\)
of two absolutely irreducible planes with
\[
 c_\infty|_{V_\pm}=\pm1,\qquad
 \det V_-=A^2\Psi,\qquad \det V_+=\Psi .
\]
In characteristic zero take any finite extension \(F/\QQ\), a
place above a rational prime \(q\notin\Sigma\) of residue degree
\(f\), and an embedding of the eigenvalue field into \(\CC\).
For the four eigenvalues of arithmetic Frobenius on these planes
let \(M_q,m_q\) be the largest and smallest of their complex
absolute values. Then \(M_q/m_q\ge q^f\).
\end{theorem}
\begin{proof}
The eigenvalues of \(h\) are distinct in the fraction field
because they are distinct in the residue field. Any invariant
subspace over an algebraic closure is a sum of coordinate
eigenlines. A nonzero entry \(ij\) forces any such sum containing
line \(j\) to contain line \(i\). The multiplier identity supplies
also a nonzero entry \(3-j,3-i\) at the inverse group element.
If some \(\sigma(U_{pj})\ne0\), entries \(01,12,20\) all occur.
They connect the first three lines in both directions, and their
reflections connect the fourth in both directions. If some
\(\sigma(W_{jk})\ne0\), entries \(03,12,20\) occur, and reflection
provides \(31\). Their cycle visits all four lines. Neither set
of edges permits a proper nonempty invariant subset.

In the remaining case set \(a=b=c=0\) as in the proof of
Theorem~\ref{thm:ramification549}. Every entry of nonzero first
torus weight vanishes. The two planes with indices \(\{0,2\}\)
and \(\{1,3\}\) are therefore invariant. A nonzero remaining
\(\sigma(U_{ij})\) gives entries \(02,20\) and their reflections
\(13,31\), proving absolute irreducibility on each plane.
Complex conjugation reduces to the stated scalars on these
planes. Its idempotent for the opposite sign reduces to zero
and hence is zero, since 2 is a unit. At \(\ell\) the ordinary
flag gives the determinant inertia characters \(d_0d_2,d_1d_3\).
At \(q\ne\ell\) the inertia quotients have order prime to
\(\ell\); a determinant ratio reducing to one is trivial on
these groups since the principal congruence units form a
pro-\(\ell\) group. The two ratios to the displayed characters
are therefore everywhere unramified global characters, hence
trivial by class reciprocity for \(\QQ\).

The determinants at the indicated Frobenius have absolute values
\(q^{f(r+2)}\) and \(q^{fr}\). Indeed each finite character
value has absolute value one, and
\(\kappa_\ell(\operatorname{Frob})=q^f\). On \(V_-\) the larger
absolute value is at least \(q^{f(r+2)/2}\). On \(V_+\) the
smaller is at most \(q^{fr/2}\). Dividing proves the bound.
In particular, any Frobenius estimate giving relative spread
strictly less than \(q^f\) at even one such place excludes
this two-plane realization.
\end{proof}


\subsection{Strict length and Frobenius multiplicities}
\label{sec:mixedmultiplicity550}

Fix the period field, character, weight and coefficient prime of
Section~\ref{sec:generaleven547}. Use \(T,D,I,C=T/I\) of
Theorem~\ref{thm:tracefitting548}. Let \(\overline J\) be an
algebraic closure of \(J\), with \(v_\ell(\ell)=1\). When \(T\)
is finite over \(O\), let \(\sigma\) range over the geometric
points of \(C[1/\ell]\), and set
\[
 n_\sigma=\dim_{\overline J}(C\otimes_O\overline J)_\sigma .
\]
Here the subscript denotes the local Artin factor. This convention
counts the several embeddings of each residue field separately.

\begin{theorem}[A strict-length bound with the generic multiplicities]
\leavevmode\label{thm:lengthmultiplicity550}
If \(T\) is finite over \(O\) and \(d\in D\) is nonzero at all
the \(\sigma\), then
\[
 \operatorname{length}_O S(\alpha,r)
       \ \le\ 2\sum_\sigma n_\sigma v_\ell(\sigma(d)).
\]
Such a \(d\) exists in the integral span of the good-prime trace
differences. The assertion includes finite trace algebras with
\(\ell\)-torsion and with nonreduced generic fiber.
\end{theorem}
\begin{proof}
Write \(m=a_-\), \(p=a_++2\). When \(m=0\), the strict group
vanishes by Theorem~\ref{thm:integralroots547}. Assume \(m\ge1\).
Take the direct sum \(A\) of the integer weight spaces of the
ambient power series ring \(B_0=O[[X_i,w,Y_j]]\) from
Section~\ref{sec:roottrace548}. Its degree-zero ring \(A_0\) is
the complete monomial algebra in \(U_{ij},W_{jk}\). Each weight
space is finite over \(A_0\). Indeed, remove factors \(X_iY_j\)
and \(wY_jY_k\) from any monomial. What remains is purely
positive, purely negative, a scalar, or \(w^bY_j\) with \(b>0\).
There are only finitely many such remainders in a fixed weight.
Grouping the corresponding terms of a series proves the claim.
In particular \(A\) is generated over its Noetherian neutral
ring by the \(X_i,w,Y_j\). Taking the quotient by \(E_\nu,F_\nu\)
in \(A\) gives exactly the weight spaces of the actual completed
solution ring. To check this for one weight, project the coefficients
of an expression in the finitely many homogeneous relations to
the required weights. Its neutral quotient is \(T\).

In the notation of Sections~\ref{sec:upper35} and
\ref{sec:lower36}, the numbers of negative variables of weights
\(-1,-2\) are \(m,0\). The numbers of positive variables of
weights \(1,2\) are \(p,1\); the positive equation counts are
\(p-1,0\), and there are \(m+1\) negative equations of weight
\(-1\). There are no neutral variables or equations in the formal
chart. The chart map is coefficient extension into the same power
series ring over \(J\). Its cotangent is freely generated by the
listed variables: all monomials of ordinary degree at least two
in one weight can be grouped by two of their factors to put the
corresponding series in the square of the augmentation ideal.
Thus the negative cotangent over \(O\), modulo the linear parts
of the negative equations, is \(K_-\) of
Theorem~\ref{thm:integralroots547}. It has finite length \(k_-\).
The rational negative spanning condition follows. The split
point is the unique geometric fixed point, since killing all
nonzero weights kills all the \(U_{ij},W_{jk}\) and leaves the
scalars.

We give the change to the multiplicity calculation in
Theorem~\ref{thm:full36}. A mixed point has both a nonzero
\(Y_j\) and a nonzero \(X_i\) or \(w\), so some neutral monomial
\(g=U_{ij}\) or \(W_{jj}\) is nonzero. Conversely, on a principal
open defined by any \(U_{ij}\) or \(W_{jk}\) we can write
\(g=fY_j\) with \(f\) of weight \(1\). Inverting \(g\) in the
graded solution algebra makes \(Y_j\) invertible. Every homogeneous
element there is a neutral element times a power of \(Y_j\), and
the grading makes a sum of distinct powers direct. Consequently
this localized algebra is
\[
                       T[1/g][Y_j,Y_j^{-1}].
\]
Over the generic fiber the opens cover precisely the nonsplit
factor \(C[1/\ell]\), by Corollary~\ref{cor:tracesplit548}.
There is one scaling orbit over each \(\sigma\); the localized
weight pieces at that geometric invariant character all have
dimension \(n_\sigma\). These statements can first be checked
over a finite extension splitting the finite generic algebra.

Lift \(d\) to the ambient neutral ring with split value zero.
In the construction of Lemma~\ref{lem:support35}, the modules
in large multidegrees are finite over \(O\) and have geometric
generic support only at mixed characters. Its proof uses
finiteness of the neutral quotient, uniqueness of the fixed
point and rational spanning on the negative side, all verified
above. For completeness, at the split character a relevant
prime in the construction would give a purely negative
solution with some \(Y_j\ne0\). Setting all positive variables
to zero leaves \(B_-Y=0\), and full column rank over \(J\)
excludes it. At the other characters a product of opposite
weights is invertible, exactly the localization just described.

At such a character the upper estimate of
Theorem~\ref{thm:upper35} uses the neutral piece of the equations
with the \(m+1\) negative and \(p-1\) positive weight-one
equations squared. Before squaring, every localized weight piece
has dimension \(n_\sigma\). Squaring one equation \(h\) of
weight \(b\) at most doubles a common bound for those dimensions.
In fact for the ring \(L\) with the other current equations imposed
the sequence
\[
 (L/(h))_{j-b}\xrightarrow{\ h\ }(L/(h^2))_j
                    \longrightarrow (L/(h))_j\longrightarrow0
\]
is exact and remains so on localization. The resulting upper
bound for the local dimension in each sufficiently large
multidegree is \(2^{m+p} n_\sigma\). On a finite \(O\)-module
the difference of cokernel and kernel lengths of multiplication
by \(d\) equals its determinant valuation on the torsion-free
quotient, since the two lengths cancel on its torsion. On the
geometric local factors of this quotient the valuation is the
sum of the dimensions times \(v_\ell(\sigma(d))\). These
valuations are nonnegative by integrality. This gives the upper
bound \(2^{m+p}\sum_\sigma n_\sigma v_\ell(\sigma(d))\) for the
same difference used in Theorem~\ref{thm:full36}.

For the lower bound the integral chart and the equations just
checked are exactly those in its proof. More explicitly, take
Smith normal form of \(B_-\) with diagonal entries
\(\ell^{c_1},\ldots,\ell^{c_m}\), so \(\sum c_i=k_-\).
In the three graded modules of that proof the factors stripped
from the corresponding squared equations contribute \(2\sum c_i\).
The remaining finite difference of the monomial count is
\(2^{m-1+(p-1)}\), by the calculation at the end of
Section~\ref{sec:lower36} with no negative variable of weight
\(-2\) and one positive variable of weight \(2\). Thus the
lower bound is \(2^{m+p-1}k_-\). The lower calculation requires
only the chart, the linear terms, and finite length after
imposing \(d\) on the large pieces. The preceding upper calculation
gave this last property even at a multiple point. Comparing and
dividing gives the claimed inequality.

Finally the finitely many good-prime differences in
Corollary~\ref{cor:tracesplit548} generate \(D\), and \(DC[1/\ell]\)
is the unit ideal. At every \(\sigma\) at least one difference
is nonzero. Use these differences as the coefficients in order
of a polynomial in a new variable. Each of the finitely many
specialized polynomials has only finitely many roots, so some
integer gives a nonzero value at all the points. The empty set
of points allows \(d=0\); in the lower calculation one can use
the ambient lift adjusted by \(zF_i\) as in the proof of
Theorem~\ref{thm:full36}.
\end{proof}

\begin{corollary}[A global representation forced by a strict class]
\leavevmode\label{cor:pointalternative550}
If \(S(\alpha,r)\ne0\), at least one of the following global
ordinary representations exists, with the conductor and multiplier
conditions of Section~\ref{sec:generaleven547}:
\begin{enumerate}
\item a representation over the integers of a finite extension of
\(J\) with nonsplit trace;
\item a representation over the integers of a finite extension of
\(k((t))\) with a good-prime Frobenius trace transcendental over \(k\).
\end{enumerate}
Each has one of the two representation types in
Theorem~\ref{thm:purityblocks549}.
\end{corollary}
\begin{proof}
If \(T\) is finite, absence of any geometric point of \(C[1/\ell]\)
would give zero on the right in
Theorem~\ref{thm:lengthmultiplicity550}. Thus some such point
exists. It takes values in the integers of a finite extension
of \(J\). Its specialization map is local and continuous,
since the maximal ideal of the finite local algebra \(T/\ell T\)
is nilpotent. It cannot kill \(D\), since \(C/DC\) is torsion.
Apply Theorem~\ref{thm:ramification549}.

If \(T\) is infinite over \(O\), its special fiber has positive
dimension. Indeed, a dimension-zero special fiber is Artinian,
so a power of the maximal ideal lies in \(\ell T\). The two adic
topologies then agree, and lifting a finite basis modulo \(\ell\)
by successive approximation makes \(T\) finite over \(O\).
Take a prime quotient of the special fiber of dimension one.
It is a complete local domain over \(k\). For any nonzero
parameter \(t\) in its maximal ideal the same approximation
applied modulo \(t\) makes it finite over \(k[[t]]\); the map
from this discrete valuation ring is injective. Inclusion into
the valuation ring of its fraction field gives a local map from
\(T\) to the integers of a finite extension of \(k((t))\). The
image of \(D\) is nonzero because the quotient by \(D\) of the
special fiber is \(k\). Some good-prime difference therefore
has nonzero image of positive valuation. This element is
transcendental over the finite field, since each nonzero element
algebraic over \(k\) is a root of unity. The split trace is
constant in \(k\), so the corresponding Frobenius trace is also
transcendental. Apply Theorem~\ref{thm:ramification549} in equal
characteristic and then Theorem~\ref{thm:purityblocks549}.
\end{proof}

\begin{corollary}[The contribution at a common Frobenius prime]
\leavevmode\label{cor:intervalmultiplicity550}
Use the common prime and multisets of Theorem~\ref{thm:commonsptrace433}.
At any triple \((\ell,r,\alpha)\) under consideration suppose
\(T\) is finite and the multiset of nonsplit Frobenius traces,
with the \(n_\sigma\) multiplicities and the chosen coefficient
embedding, occurs in \(\mathcal C_r\). Then
\(\operatorname{length}_O S(\alpha,r)\le2v_\ell(D_{\rm Sp})\).
For all such triples on one \(E_x\) their total scalar length,
summed over geometric characters and weighted by \(\log\ell\),
is at most \(2n_*t_*\log D_{\rm Sp}\).
\end{corollary}
\begin{proof}
The size bound on every member of \(\mathcal C_r\) in
Theorem~\ref{thm:commonsptrace433} separates it from the split
trace. Apply Theorem~\ref{thm:lengthmultiplicity550} to the
trace difference at that same prime. The integer product in
Theorem~\ref{thm:commonsptrace433} bounds the sum of its
valuations with the indicated multiplicities, giving the first
inequality. At each \(\ell\) there are at most \(n_*t_*\)
character and weight pairs. Summing and using
\(\sum_\ell v_\ell(D_{\rm Sp})\log\ell=\log D_{\rm Sp}\) proves
the second. Establishing finiteness and the asserted occurrence
of traces throughout the interval still requires arithmetic
information about the global representations.
\end{proof}


\subsection{The two marks in the symplectic trace algebra}
\label{sec:twomarks557}

On the fields \(E_x\) and intervals of Section~\ref{sec:missing},
take the four characters \(d=(A\Psi,\Psi,A,1)\), multiplier
\(\mu=A\Psi\), and the auxiliary character \(\xi\) of
Lemma~\ref{lem:uniformoddroots430}. Let \(r\ge2\),
\(\ell\ge5\), \(\ell\nmid n_*N_*\), and \(\ell-1\mid 2r\).
In the first theorem it suffices that \(\bar\Psi\ne1\) globally.
The coefficient ring \(O\) is the unramified integer ring
containing the finite character values. Use any complete
Noetherian local \(O\)-algebra \(R\) with the same finite
residue field and a continuous matrix lift
\(\rho:G_{\QQ}\to\operatorname{GSp}_4(R)\) of \(\bar d\) with
multiplier \(\mu\). Write \(u^*\) for the adjoint of a matrix
under the fixed alternating pairing; it sends \(\rho(g)\) to
\(\mu(g)\rho(g^{-1})\).

Choose an arithmetic Frobenius lift \(s\in D_\ell\) with
\(\kappa_\ell(s)=1\), and take \(g_0\in G_{\QQ}\) with
\((\bar A(g_0)-1)(\bar\Psi(g_0)-1)\ne0\). Such a \(g_0\)
exists because two proper subgroups cannot cover a group.
By the choice of \(\xi\), \(\bar A(s)\) is outside the subgroup
generated by \(\bar\Psi(s)\). The residual eigenvalues of \(s\)
at positions \(0,2\) therefore form a set disjoint from those
at \(1,3\); repetitions within each set are allowed. Define
\[
       \phi=\lim_{j\to\infty}s^{\,\ell^{j!}},
       \qquad H=\rho(g_0)+\mu(g_0)\rho(g_0^{-1}).
\]
The limit exists in the profinite cyclic closure of \(s\).
On its finite quotients the exponents are eventually zero on
the \(\ell\)-part and one modulo the prime-to-\(\ell\) part.
The marking in
\srcloc{Development/AutomorphicArithmetic/AutomorphicVarieties/SymplecticGSp4/DecompSigmaOrd}{140}
requires that \(\rho(\phi)\) preserve the two coordinate planes
\(02,13\) and that \(H_{02}=H_{20}=0\).

\begin{theorem}[Neutral entries and traces with the two marks]
\leavevmode\label{thm:twomarks557}
There is a symplectic change of frame congruent to the identity
that gives the indicated marking. Any two resulting marked
bases differ by diagonal factors with product one on each
opposite pair. For a marked lift let \(T_0\subset R\) be the
\(O\)-algebra generated by the traces of \(\rho(g)\). Let
\(M_0\subset R\) be the \(O\)-algebra generated by the products
of matrix entries of total weight zero under conjugation by
\(\operatorname{diag}(st,t,s,1)\), where \(s,t\) here are
independent torus variables. There are \(z\in M_0\) and
\(\Delta=z^2\in T_0\), both units of \(R\), such that
\[
       M_0[1/\Delta]=T_0[1/\Delta,z],\qquad
                      \overline{M_0}=\overline{T_0}.
\]
Closures use the maximal-ideal topology of \(R\). In particular,
\(M_0\) is finite as an \(O\)-module exactly when
\(\overline{T_0}\) is. In this case \(M_0=\overline{T_0}\).
\end{theorem}
\begin{proof}
Let \(m\) be the order of the residual diagonal matrix of \(s\).
It is prime to \(\ell\). The congruence kernel in every
Artinian quotient of \(R\) is an \(\ell\)-group, by its
filtration by powers of the maximal ideal. Hence
\(\rho(\phi)^m=1\), and \(\rho(\phi)\) reduces to \(\bar d(s)\).
The \(m\)-th roots occurring in the residue field have
Teichm\"uller lifts in \(O\); here \(m\) divides the order of the
multiplicative group of that field. Interpolating at these
roots, with unit differences, gives a projector \(P\) over \(O\)
as a polynomial in \(\rho(\phi)\): it takes value one on the
two eigenvalues indexed by \(0,2\), and zero on the others.
Unoccupied eigenspaces are zero by Nakayama's lemma. The two
summands \(U=P R^4\), \(U'=(1-P)R^4\) are free of rank two.
They are isotropic and dual to each other. Indeed two
eigenspaces can pair only when their eigenvalues multiply to
\(\mu(\phi)\); the paired eigenvalues lie in opposite sets, as
one checks after reduction. Thus \(P^*=1-P\).

Now \(H=H^*\). The endomorphism \(H_U=PHP\) on \(U\) has
distinct residual eigenvalues \(\bar A(g_0)\bar\Psi(g_0)+1\)
and \(\bar A(g_0)+\bar\Psi(g_0)\). Hensel factorization of
its quadratic characteristic polynomial gives two unique
rank-one eigensummands labeled by \(0,2\). The adjoint
projectors give the dual eigensummands on \(U'\) labeled by
\(3,1\). Choose generators lifting the residual basis in \(U\)
and the paired dual generators in \(U'\). This is the desired
symplectic basis. In any marked basis the two planes must be
\(U,U'\), since the projector is a polynomial and the two
residual spectra are disjoint. The two off-diagonal equations
then select exactly the eigensummands of \(H_U\). This proves
the uniqueness assertion about bases.

Write \(s_0',s_2'\) for the \(H_U\)-eigenvalues and put
\(z=s_0'-s_2'\). The symbols \(s_i'\) in this calculation are
scalars. In the marked frame their difference is a difference
of diagonal entries of \(H\), so belongs to \(M_0\).
Since the polynomial \(P\) has coefficients in \(O\),
\(\operatorname{tr}H_U,\operatorname{tr}(H_U^2)\) belong to
\(T_0\): expand the matrix products in elements of the group.
Here the full trace agrees with the trace on \(U\) because
\(PHP\) vanishes on \(U'\). Therefore
\[
             \Delta=2\operatorname{tr}(H_U^2)
                        -(\operatorname{tr}H_U)^2=z^2\in T_0 .
\]
Its residue is nonzero. The two coordinate projectors inside
\(P\) are \((H_U-s_2'P)/z\) and its complement in \(P\), with
\(s_2'=(\operatorname{tr}H_U-z)/2\). Taking adjoints gives the
other two. Since \(\rho(g)^*=\mu(g)\rho(g^{-1})\), the four
projectors belong to the span of the matrices \(\rho(g)\)
over \(T_0[1/\Delta,z]\). Tracing products with these projectors
inserted shows that the product of the entries on any directed
cycle belongs to this scalar ring.

For a neutral product regard an entry \(ij\) as an arrow
from \(j\) to \(i\). If \(\delta_i\) is the number of incoming
arrows minus the number of outgoing arrows at \(i\), neutrality
and \(\sum_i\delta_i=0\) give
\((\delta_0,\delta_1,\delta_2,\delta_3)=(d,-d,-d,d)\).
The multiplier identity permits replacing an arrow \(j\to i\)
by \(3-i\to3-j\) for the inverse group element, at the cost
of a signed unit in \(O\). For \(d>0\) there are at least
\(2d\) arrows from \(\{1,2\}\) to \(\{0,3\}\). Replacing \(d\)
of them reduces every imbalance to zero; one replacement
subtracts one at both \(0,3\) and adds one at both \(1,2\).
For \(d<0\) use \(-d\) arrows in the opposite direction.
The resulting finite directed graph is a union of directed
cycles, by following arrows until returning to a vertex and
then removing the cycle. This gives the inclusion of \(M_0\)
in \(T_0[1/\Delta,z]\). Diagonal entries have weight zero, so
\(T_0\subset M_0\), proving the localized equality.

A unit of \(R\) in \(\overline{T_0}\) has its inverse there
by a geometric series after lifting its residue to \(O\).
Starting with a lift of the residue of \(z\), Newton iteration
for the square root of \(\Delta\) in the closed ring gives \(z\):
the two candidate squares agree and the sum of the two roots
with this residue is a unit. This proves the equality of closures.
If \(M_0\) is finite over \(O\) it is compact, hence closed in
\(R\). Conversely \(M_0\) is a submodule of \(\overline{T_0}\)
and \(O\) is Noetherian. This proves the finiteness assertion.
\end{proof}

\begin{proposition}[The local component in the two ordinary markings]
\leavevmode\label{prop:actualcomponent557}
Impose in addition \(\bar\Psi|_{D_\ell}\ne1\). Choose the
separating Frobenius \(h\) of Lemma~\ref{lem:sympflag468}.
On a framed local lift with fixed multiplier, the following
conditions are equivalent over every coefficient ring \(R\)
as above.
\begin{enumerate}
\item The lift has the ordered ordinary flag of
Lemma~\ref{lem:sympflag468}, with graded inertia characters \(d_i\).
\item The lift is a specialization of a lift over a complete
local \(O\)-domain that also specializes to the split lift
\(d|_{D_\ell}\) over \(O\). The domain lift satisfies the
ordinary polynomial equations at the marked inertia element
of Theorem~\ref{thm:finiteflag546}.
\end{enumerate}
These are the local component maps in
\srcloc{Development/AutomorphicArithmetic/AutomorphicVarieties/SymplecticGSp4/DecompSigmaOrd}{25}
and its coefficient-prime condition on line~126.
\end{proposition}
\begin{proof}
For the first implication choose a symplectic frame reducing
to the given residual frame in which the separating \(h\)
is diagonal. The flag consists of the successive labeled
eigenlines. Theorem~\ref{thm:sympordinary468} realizes the
normalized matrices by specializing \(O[[u_1,u_2,a,x,b,c]]\),
with split specialization at the origin. To allow the inverse
frame change, adjoin ten formal variables for \(X\) in the
symplectic Lie algebra. The matrix \((1+X)(1-X)^{-1}\)
is symplectic and parametrizes every symplectic change \(C\)
congruent to one: \(X=(C-1)(C+1)^{-1}\) is skew adjoint and
has entries in the maximal ideal. This uses invertibility of
two. Conjugate the universal local lift by this matrix. The
resulting ring is a power series domain through the split
point. The ordinary flag gives the polynomial equations over
this domain by the forward implication in
Theorem~\ref{thm:finiteflag546}, equivalently by Cayley--Hamilton
on its successive submodules and quotients.

Conversely let \(B\) be such a domain and \(K\) its fraction
field. The split specialization \(B\to O\) is the identity
on \(O\), so \(K\) has characteristic zero. At the specified
inertia element the four values \(d_i\) are distinct in \(K\):
its principal cyclotomic value has infinite order and the
exponents \(r+1,r,1,0\) are distinct. The polynomial equations
therefore give the ordered stable flag over \(K\) by ordinary
spectral interpolation, with the stated inertia actions.
Independently, \(\rho_B(h)\) has four distinct residual roots,
so its four labeled spectral projectors are integral over \(B\)
and belong to \(\operatorname{End}_B(B^4)\) by Hensel
factorization. Every step of the flag over \(K\) is a sum of
some of these eigenlines. For the resulting permutation, its
stability equations, and its inertia equations on each graded
line, are matrix identities in these integral projectors.
They hold over \(B\) because \(B\hookrightarrow K\). Specialize
to the split lift over \(O\), then pass to \(\operatorname{Frac}O\).
Its four distinct labeled inertia eigenvalues force the
permutation to be the original one. The equations over \(B\)
thus give the required ordered direct summands after
specialization to \(R\), including when \(R\) has nilpotents.
\end{proof}

\begin{corollary}[The finiteness condition for the strict-length inequality]
\leavevmode\label{cor:finmark557}
For the triples of Proposition~\ref{prop:actualcomponent557}
let \(T\) be the trace ring of Theorem~\ref{thm:lengthmultiplicity550}.
Its finiteness over \(O\) is equivalent to the following
arithmetic assertion for the same characters. For every framed
global lift with multiplier \(\mu\), the two marks above, the
split local component at \(\ell\), and the given inertia kernels
away from \(\ell\), \(M_0\) is finite over \(O\).
\end{corollary}
\begin{proof}
At a finite \(q\ne\ell\) the residual diagonal inertia image
has prime-to-\(\ell\) order since \(\ell\nmid n_*N_*\) and
\(\xi\) has two-power order. Its kernel is the imposed inertia
kernel in both descriptions; no tame \(\ell\)-part of
\(\alpha\) is present. Given a lift in the two marks, diagonalize
\(\rho(h)\) by a symplectic frame reducing to the identity.
Proposition~\ref{prop:actualcomponent557} identifies the local
condition. The resulting lift factors through the universal
ring of Theorem~\ref{thm:allglobalroots547}. If \(T\) is finite
over \(O\), its image is compact and contains every trace.
Theorem~\ref{thm:twomarks557} proves the stated finiteness
for \(M_0\). Conversely take the universal lift in the
\(h\)-diagonal marking and apply the basis construction of
Theorem~\ref{thm:twomarks557} over its own coefficient ring.
It satisfies the two marks and the split local component.
Conjugation preserves all of its traces, so the asserted
finiteness for this one lift makes \(T\) finite.
\end{proof}

The arithmetic assertion in the corollary is the finiteness
clause spelled out in
\srcloc{Development/AutomorphicArithmetic/AutomorphicVarieties/SymplecticGSp4/DecompSigmaOrd}{172}
for this datum. The conductor choices in the source's individual
cases are specified separately in
\srcloc{Development/AutomorphicArithmetic/AutomorphicVarieties/AutomorphicCases/SymplecticGSp4}{15}.
The equivalence proves the change of marking for a fixed datum.
An arithmetic bound on its trace algebra and its Frobenius
multiplicities is still required in the interval argument.


\subsection{Infinitesimal traces at the ordinary symplectic points}
\label{sec:inftan558}

Fix \(E_x,\alpha,r,\ell,O,J,\mathcal B,T\) as in
Section~\ref{sec:generaleven547}. In particular \(\ell\ge5\),
\(\ell\nmid n_*N_*\), \(\ell-1\mid 2r\), and the four characters
are separated at the element \(h\in D_\ell\). Assume \(T\) is
finite over \(O\). Let \(\sigma\) be a geometric point of the
generic fiber with an absolutely irreducible rank-four realization
in Theorem~\ref{thm:purityblocks549}. Take a finite extension
\(L/J\) that splits the generic trace algebra and contains the
coefficients of an integral realization of \(\sigma\). Such an
extension exists by Theorem~\ref{thm:ramification549}. Write \(V\)
for its representation space over \(L\) and
\(\mathfrak s=\{X\in\operatorname{End}_L(V):X^*=-X\}\), with
Galois acting by conjugation. Write \(\mathfrak m_\sigma\) for the
maximal ideal of the local factor of \(T\otimes_O L\) at \(\sigma\).

\begin{lemma}[Lifting a tangent of the trace algebra]
\leavevmode\label{lem:rankfourtangent558}
Every \(L\)-algebra map from the local factor to
\(L[\epsilon]/(\epsilon^2)\) with reduction \(\sigma\) is the
trace of a continuous symplectic-similitude lift of \(V\) with
the original multiplier. The lift preserves the ordered local
flag, its graded inertia characters, and the inertia kernels at
the other primes. After fixing the integral realization, one can
choose the lift so that its first derivative depends linearly
on the tangent vector.
\end{lemma}
\begin{proof}
Inside \(\mathcal B\) take the \(T\)-algebra \(A_0'\) generated by
the entries of all universal matrices. Each generator is
homogeneous for the two torus variables in
Theorem~\ref{thm:allglobalroots547}. Finite sums of its weight
spaces are direct. Indeed after separating the Laurent characters
in the torus action modulo each power of the maximal ideal, a
relation between distinct weights is zero in each weight at every
depth. The piece of weight zero is \(T\), by
Theorem~\ref{thm:allroottrace548}.

Draw an arrow \(j\to i\) if a matrix entry in row \(i\), column
\(j\) has nonzero value at the chosen realization of \(V\).
If a vertex cannot reach another, its reachable vertices span
a nonzero proper invariant subspace. Absolute irreducibility
therefore makes the graph strongly connected. The product of
entries along a path has weight equal to the difference of the
endpoint weights and is nonzero at the realization. The four
vertex weights are \((1,1),(0,1),(1,0),(0,0)\). Choose path
products \(x_+,x_-,y_+,y_-\) of respective weights
\((1,0),(-1,0),(0,1),(0,-1)\). Their product \(b\in T\)
satisfies \(\sigma(b)\ne0\). All four become units on inverting
\(b\). Multiplication by suitable powers of \(x_+\) and \(y_+\)
then makes every homogeneous element have weight zero. A fraction
of weight zero still has numerator of weight zero as \(b\) is
neutral. Directness of the grading gives an isomorphism
\[
     A_0'[1/b]=T[1/b][x_+^{\pm1},y_+^{\pm1}].
\]

Extend the given map on \(T\) through this Laurent algebra by
sending \(x_+,y_+\) to their constant values in the chosen
realization. Reduction agrees with that realization on the
whole algebra. Applying the extension to each universal entry
gives the lift. Multiplication, the pairing, zeros below the
local flag, its diagonal inertia actions, and triviality on the
specified inertia kernels are identities between these entries.
They hold after applying the map. To check continuity, choose
for each entry weight a fixed product of the four path elements
of opposite weight. It has invertible value in the dual numbers.
Multiplying the entry by this element gives a continuous
\(T\)-valued function of the group element. Indeed \(T\) is
compact, so its topology as a finite \(O\)-module agrees with
its subspace topology in the Hausdorff ring \(\mathcal B\).
The map on this finite \(O\)-module is continuous. Dividing by
the fixed invertible value proves continuity of each lifted
entry. Holding the two Laurent variables constant extends every
derivation linearly in the original derivation, which proves the
last assertion.
\end{proof}

\begin{proposition}[The finite local condition at the coefficient prime]
\leavevmode\label{prop:finiteordinary558}
At an integral point of the specified global chart, the
ordinary representation and every continuous first-order lift
with the same ordered flag and graded inertia characters are
crystalline over \(K=\QQ_\ell(\mu_\ell)\). For a first-order lift,
crystallinity concerns its underlying eight-dimensional
\(L\)-representation. Its self-extension class at \(\ell\)
belongs to \(H^1_f(\QQ_\ell,\operatorname{End}_L V)\).
\end{proposition}
\begin{proof}
On \(D_\ell\) the diagonal characters of the point are
\(d_i u_i\), where \(d=(A\Psi,\Psi,A,1)\) and the \(u_i\) are
unramified characters congruent to one. Those of a first-order
lift differ further by unramified characters in \(1+\epsilon L\).
Write \(w=(r+1,r,1,0)\) for the successive cyclotomic exponents.
Over \(K\) the tame character \(\omega_\ell\) is trivial.
Both \(\alpha\) and \(\xi\) are unramified at \(\ell\). Thus on
each graded part the representation over \(K\) is
\(\kappa_\ell^{w_i}\) times an unramified representation. In
the first-order case the unramified part has two identical
character factors. Continuous unramified representations are
crystalline of weight zero, so all graded parts are crystalline.

Let \(s\in D_\ell\) be the arithmetic Frobenius with
\(\kappa_\ell(s)=1\) used in Section~\ref{sec:twomarks557}.
It lies in \(G_K\) and induces Frobenius on the residue field,
since \(K/\QQ_\ell\) is totally ramified. Whenever \(i<j\) and
\(w_i-w_j=1\), the ratio of the unramified character factors
on the two graded parts is nontrivial. For the pairs \(01\) and
\(23\) its value at \(s\) reduces to \(\bar\xi(s)\). For the
remaining possible pair \(12\), which occurs when \(r=2\), it
reduces to \(\bar\Psi(s)/\bar\xi(s)\). Neither residue equals
one because \(\bar\xi(s)\notin\langle\bar\Psi(s)\rangle\).

Suppose an initial part \(P\) of the flag is crystalline over
\(K\), and let \(Q\) be the next graded part. Set
\(W=\operatorname{Hom}_L(Q,P)\). This is crystalline, and all
its cyclotomic exponents are positive. In the convention where
\(\kappa_\ell\) has Hodge--Tate weight \(-1\), this gives
\(\operatorname{Fil}^0D_{\rm dR}(W)=0\).
Also \(H^0(K,W^\vee(1))=0\). Indeed a character factor of the
dual twist coming from a difference greater than one is
nontrivial on an open subgroup of inertia. At a difference equal
to one, nontriviality follows from the Frobenius calculation.
The finite filtrations on \(P,Q\), refined by the two scalar
factors on a dual-number block, give a filtration on \(W^\vee(1)\).
Vanishing of invariants on every factor implies vanishing on
the whole module by left exactness.

The Bloch--Kato dimension formula and the local Euler
characteristic formula with Tate duality now give
\[
 \dim_L H^1_f(K,W)
   =[K:\QQ_\ell]\dim_L W+\dim_L H^0(K,W)
   =\dim_L H^1(K,W).
\]
Thus every extension in the flag has finite local class.
For crystalline end terms this says precisely that the middle
term is crystalline. Indeed splitting the extension after
tensoring with \(B_{\rm cris}\) gives the required crystalline
periods by the comparison isomorphisms for the two end terms.
Induction proves the assertion for both the point and its
first-order lift.~\cite{IX-VII-VI-bk117,IX-VII-V-milne98}

The exact self-extension over \(K\) therefore splits over
\(B_{\rm cris}\). One may take a linear splitting on the
crystalline period spaces and extend scalars by their comparison
isomorphisms. The class over \(\QQ_\ell\) with coefficients in
\(B_{\rm cris}\otimes_{\QQ_\ell}\operatorname{End}_L V\)
restricts to zero over \(K\). Corestriction multiplies it by
\([K:\QQ_\ell]\), which is invertible over \(L\). This gives
the asserted finite local condition over \(\QQ_\ell\).
\end{proof}

\begin{theorem}[A Selmer bound for the tangent at an irreducible trace]
\leavevmode\label{thm:rankfourselmer558}
There is an \(L\)-linear injection
\[
       \operatorname{Hom}_L(\mathfrak m_\sigma/
             \mathfrak m_\sigma^2,L)
          \lhook\joinrel\longrightarrow H^1_f(\QQ,\mathfrak s).
\]
Here the finite global group uses the unramified condition away
from \(\ell\) and the kernel of scalar extension to \(B_{\rm cris}\)
at \(\ell\). If this group vanishes, the multiplicity
\(n_\sigma\) in Theorem~\ref{thm:lengthmultiplicity550} is one.
\end{theorem}
\begin{proof}
Write a lifted matrix as \((1+\epsilon z(g))\rho(g)\). The
group law makes \(z\) a continuous one-cocycle, and the fixed
multiplier gives \(z(g)^*=-z(g)\). At each finite prime
\(q\ne\ell\) the cocycle is zero on the prescribed kernel of
inertia. Its restriction to inertia factors through a finite
group and is a coboundary by averaging. At \(\ell\) the class
is finite in \(\operatorname{End}_L V\) by
Proposition~\ref{prop:finiteordinary558}. The equivariant
projection \(X\mapsto (X-X^*)/2\) is a left inverse to the
inclusion of \(\mathfrak s\). Applying it after scalar extension
to \(B_{\rm cris}\) gives the same local condition in \(\mathfrak s\).

These operations and the choice in
Lemma~\ref{lem:rankfourtangent558} are linear on derivatives.
If the cohomology class is zero, conjugation by a constant
first-order change of basis makes all the matrices equal to
the original ones. Thus every trace has zero derivative. The
\(O\)-algebra generated by the traces is an \(O\)-submodule of
the finite module \(T\), hence finite and compact. Its closure
is \(T\), so the traces generate \(T\) algebraically. The
derivative therefore vanishes on all of \(T\), proving
injectivity. If the right side is zero, Nakayama's lemma applied
to the maximal ideal of the local Artin algebra gives
\(\mathfrak m_\sigma=0\). This factor is \(L\). At a nonsplit
point an element of the augmentation ideal \(D\) is invertible.
The identity \(D\,\operatorname{Fitt}_0^T(M)=0\) in
Theorem~\ref{thm:tracefitting548} identifies the local factor
of \(C\) with this same field. Its multiplicity is one.
\end{proof}


\subsection{Irreducibility and multiplicity on the period fields}
\label{sec:enormous559}

Take the triples of Section~\ref{sec:generaleven547}. In mixed
characteristic let \(V\) be the absolutely irreducible four-dimensional
representation of Theorem~\ref{thm:purityblocks549}, or either of
the irreducible planes in its second case. We work over a finite
extension of \(J\) over which the matrices admit the integral
model of that theorem. All statements about invariant subspaces
allow extension to the algebraic closure.

\begin{proposition}[Irreducibility over every finite extension]
\leavevmode\label{prop:strongall559}
For every finite extension \(F/\QQ\), the restriction of \(V\)
to \(G_F\) is absolutely irreducible.
\end{proposition}
\begin{proof}
Write \(n=\dim V\). There are two local elements with different
roles. At \(h\) the \(n\) residual diagonal characters have
pairwise distinct values \(b_1,\ldots,b_n\). Choose an inertia
element \(s_0\) in the kernels of the finite characters with
\(\kappa_\ell(s_0)\) of infinite order. Such elements exist
because inertia maps onto the principal cyclotomic units after
passage to a subgroup of finite index. The eigenvalues of
\(s_0\) are the powers with the prescribed exponents from
\(r+1,r,1,0\). Every positive power of \(s_0\) therefore has
simple spectrum. The ordinary flag supplies a \(D_\ell\)-stable
line, also for either plane by intersection with the coordinate
flag.

Suppose the global group permutes a direct-sum decomposition
into \(m>1\) nonzero subspaces. It acts transitively on them by
irreducibility. These finitely many subspaces are defined over
a finite extension of the coefficient field. Their orbit in its
Grassmannian is finite and Hausdorff, so the permutation action
is continuous. The kernel of reduction on the compact matrix
image is pro-\(\ell\). It acts trivially on the summands because
\(\ell\ge5\) and \(m\le4\). Thus the permutation image is a
quotient of the residual diagonal group and is abelian. A faithful
transitive abelian permutation group is regular, since a point
stabilizer fixes every translate of that point.

Take a power of \(s_0\) fixing the summands. Each summand is a
sum of its eigenlines by simple spectrum. The \(D_\ell\)-stable
line is one such eigenline, so lies in a unique summand.
Consequently \(D_\ell\) fixes that summand and acts trivially on
all summands by regularity. Choose \(g\) inducing a nonidentity
permutation. Every \(gh^j\) permutes the summands without a fixed
one, so its trace is zero. On reduction, writing \(a_i\) for the
value at \(g\) of the \(i\)-th residual character, we obtain
\[
                  \sum_{i=1}^n a_i b_i^j=0
                         \qquad(0\le j<n).
\]
The Vandermonde determinant is a unit since the \(b_i\) are
distinct. This forces \(a_i=0\) for all \(i\), whereas character
values are units. This rules out any permuted decomposition.

For an open normal subgroup \(U\) of \(G_\QQ\), choose a simple
\(U\)-submodule of \(V\). Its global conjugates are simple
\(U\)-submodules and span \(V\). Thus \(V|_U\) is semisimple.
Its isotypic components would give a permuted decomposition if
there were several. Within its single isotypic component a
multiplicity greater than one would repeat every eigenvalue of
every element of \(U\). A positive power of \(s_0\) lies there
and has simple spectrum. This proves absolute irreducibility on
\(U\), and the normal core gives it on any open subgroup.
\end{proof}

\begin{theorem}[The cyclotomic image at the two colliding residues]
\leavevmode\label{thm:enormousall559}
For every finite \(F/\QQ\), the image of
\(G_{F(\mu_{\ell^\infty})}\) on \(V\) is enormous in the sense
of Newton and Thorne. More explicitly, after a finite extension
of coefficients there are \(n^2\) elements of this image with
distinct eigenvalues on each element, whose matrices span
\(\operatorname{End}(V)\).
\end{theorem}
\begin{proof}
Let \(G\) be the Zariski closure of the global image over the
algebraic closure of the coefficient field, and let \(G^\circ\)
be its identity component. The intersection of the image of
\(G_F\) with \(G^\circ\) is Zariski dense in \(G^\circ\).
Indeed it has finite index in the global image, and finitely
many translates of a proper closed subset cannot cover a
connected algebraic group. The preimage of \(G^\circ\) is
open in the profinite group, since algebraic components are
closed in the local field topology and there are finitely many.
Thus \(G^\circ\) acts absolutely irreducibly by
Proposition~\ref{prop:strongall559}. Its unipotent radical
has a nonzero fixed vector. Normality and irreducibility make
all vectors fixed, so this radical is trivial. The connected
reductive group \(G^\circ\) is the product of its central torus
and its connected derived group \(H\). Schur's lemma makes the
central torus scalar, hence \(H\) acts absolutely irreducibly.
A power of \(s_0\) as in the proposition has image in \(G^\circ\)
with simple spectrum. Removing its scalar central factor gives
an element of \(H\) with simple spectrum.

Write \(D\) for the Zariski closure of the image of the indicated
cyclotomic kernel. It contains \(H\). In fact the commutators of
elements of \(G_F\) lie in the kernel of the cyclotomic character.
Apply the commutator morphism to the dense subset just used
inside \(G^\circ\times G^\circ\). Closedness of \(D\) puts all
commutators of \(G^\circ\) in \(D\), which contains the closed
group they generate. Thus \(D^\circ\) contains the connected
group \(H\). It follows that \(D^\circ\) acts irreducibly and
that the nonvanishing locus of the characteristic-polynomial
discriminant in \(D^\circ\) is nonempty. The actual kernel
elements in this open locus are Zariski dense there and hence
in \(D^\circ\). Burnside's theorem shows that the linear span
of \(D^\circ\) is the whole matrix algebra, so the actual
simple-spectrum elements also span it. Select a basis and
extend coefficients once to split its characteristic polynomials.

For any nonzero matrix \(X\), nondegeneracy of the trace
pairing gives a selected matrix \(M\) with
\(\operatorname{tr}(MX)\ne0\). Write \(M=\sum_a aP_a\), where
\(P_a\) is its spectral projector at the eigenvalue \(a\).
Some \(\operatorname{tr}(P_a X)\) is nonzero. Applying this to
a nonzero vector of each simple submodule of the endomorphism
space gives precisely the condition in Definition~2.23 of
Newton and Thorne.~\cite{IX-VII-VI-nt119}
\end{proof}

\begin{theorem}[Multiplicity one at potentially automorphic rank-four points]
\leavevmode\label{thm:automult559}
Suppose \(T\) is finite over \(O\). At a rank-four point
\(\sigma\), suppose that on some totally real or CM number
field \(F\) its representation \(V\) is attached to a regular
algebraic cuspidal polarized automorphic representation of
\(\mathrm{GL}_4(\mathbb A_F)\), in the Galois normalization of
Newton and Thorne. Then \(n_\sigma=1\) in the strict-length
inequality of Theorem~\ref{thm:lengthmultiplicity550}.
\end{theorem}
\begin{proof}
Write \(F^+=F\) in the totally real case, and take the maximal
totally real subfield otherwise. The multiplier of the automorphic
pairing on \(G_F\) agrees with \(\mu|_{G_F}\). To see this, in
the CM case the conjugate of \(V|_{G_F}\) is isomorphic to
\(V|_{G_F}\), since \(V\) is defined on \(G_\QQ\). In either
case the two pairings make \(V|_{G_F}\) isomorphic to its twist
by the ratio \(\nu\) of the multipliers. Determinants give
\(\nu^4=1\). On the open kernel of \(\nu\), the intertwiner is
scalar by Proposition~\ref{prop:strongall559}. It intertwines
on all of \(G_F\), so \(\nu=1\).

Over a totally real \(F\) the automorphic pairing is consequently
a scalar multiple of the original alternating pairing, by Schur's
lemma. The fixed-multiplier Lie algebra \(\mathfrak s\) is a
direct summand of the Lie algebra of its similitude group. Indeed
\(X\mapsto(X-X^*)/2\) is an equivariant projection onto it.
By Theorem~\ref{thm:enormousall559} the cyclotomic image is
enormous. Theorem~5.3 of Newton and Thorne makes the finite
Selmer group of this similitude Lie algebra over \(F\) zero,
hence also that of \(\mathfrak s\).

Over a CM field choose a complex conjugation \(c\in G_{F^+}\)
and put \(C=\rho(c)\). Let \(J_0\) be the matrix of the original
alternating pairing. Since \(C^2=1\) and \(\mu(c)=-1\), the
matrix \(B=J_0 C\) is symmetric. It pairs \(V\) with its
conjugate over \(G_F\), since
\[
     \rho(g)^{\mathsf T} B\,\rho(cgc^{-1})=\mu(g)B
                           \qquad(g\in G_F).
\]
The equality of multipliers and irreducibility give uniqueness
up to scalar. In the polarized adjoint on \(\operatorname{End}(V)\)
the action of \(c\) is \(X\mapsto -B^{-1}X^{\mathsf T}B\).
On \(\mathfrak s\) this is \(X\mapsto CXC^{-1}\), precisely
the original action of \(c\). The self-adjoint matrices are
stable under this operation and under \(G_F\), giving an
invariant complement. Theorem~5.2 of Newton and Thorne and
Theorem~\ref{thm:enormousall559} give vanishing of the finite
Selmer group of the polarized adjoint over \(F^+\).
The direct summand \(\mathfrak s\) has the same vanishing.
Both theorems use the unramified local condition away from
\(\ell\) and the finite crystalline-period condition at
\(\ell\).~\cite{IX-VII-VI-nt119}

Restriction takes \(H^1_f(\QQ,\mathfrak s)\) into
\(H^1_f(F^+,\mathfrak s)\) by functoriality of these conditions.
Corestriction composed with restriction is multiplication by
\([F^+:\QQ]\), a nonzero scalar in the coefficient field. The
global group over \(\QQ\) is therefore zero. The multiplicity
assertion follows from Theorem~\ref{thm:rankfourselmer558}.
\end{proof}

Thus a rank-four point equipped with such a cuspidal companion
contributes its Frobenius valuation once. Finiteness of the
trace algebra, existence of companions throughout the interval,
and the treatment of the two-plane points are further inputs to
the strict-length bound.

\clearpage
\section{References}\label{nav:bibliography-IX}
\begingroup\def\section*#1{}\begin{thebibliography}{99}\small
\bibitem{IX-VII-V-milne98} J. S. Milne, \emph{Arithmetic Duality Theorems}, 2nd ed., BookSurge, 2006. \href{https://www.jmilne.org/math/Books/ADTnot.pdf}{Part I, Sections 2 and 4}.
\bibitem{IX-VII-VI-milne122} J. S. Milne, \href{https://www.jmilne.org/math/Books/iAG2022.pdf}{\emph{Algebraic Groups}}, Cambridge University Press, 2017, chapters on reductive groups, rank one and root data.
\bibitem{IX-VII-brumer330} A. Brumer, \emph{On the units of algebraic number fields}, Mathematika \textbf{14} (1967), 121--124. \href{https://doi.org/10.1112/S0025579300003703}{Original paper}.
\bibitem{IX-VII-geraghty363} D. Geraghty, \emph{Modularity lifting theorems for ordinary Galois representations}, Math. Ann. \textbf{373} (2019), 1341--1427, Section 2. \href{https://doi.org/10.1007/s00208-018-1742-4}{Published paper}. Proposition and lemma numbers refer to the \href{https://citeseerx.ist.psu.edu/document?doi=8faa2d675c94ae78d0ac8ebb3f9d4f69ede42d44\&repid=rep1\&type=pdf}{original manuscript}.
\bibitem{IX-VII-I-ac123} J. Arthur and L. Clozel, \href{https://www.claymath.org/library/cw/arthur/pdf/30.pdf}{\emph{Simple Algebras, Base Change, and the Advanced Theory of the Trace Formula}}, Annals of Mathematics Studies 120, Princeton University Press, 1989, Chapter III.
\bibitem{IX-VII-lgp361} T. Barnet-Lamb, T. Gee, D. Geraghty and R. Taylor, \emph{Local-global compatibility for \(l=p\), II}, Ann. Sci. \'Ec. Norm. Sup\'er. (4) \textbf{47} (2014), 165--179, Theorem A. \href{https://www.numdam.org/articles/10.24033/asens.2212/}{Published paper}.
\bibitem{IX-VII-V-shin} S. W. Shin, \emph{Galois representations arising from some compact Shimura varieties}, Ann. of Math. \textbf{173} (2011), 1645--1741. \href{https://annals.math.princeton.edu/2011/173-3/p09}{Theorem 1.2}.
\bibitem{IX-VII-VI-nt119} J. Newton and J. A. Thorne, \href{https://arxiv.org/pdf/1912.11265}{\emph{Adjoint Selmer groups of automorphic Galois representations of unitary type}}, J. Eur. Math. Soc. \textbf{25} (2023), 1919--1967, Theorem~A and Section~2.5.
\bibitem{IX-VII-VII-tadicsph716} M. Tadi\'c, \emph{Spherical unitary dual of general linear group over non-Archimedean local field}, Ann. Inst. Fourier \textbf{36} (1986), no. 2, 47--55, Theorem 3.2. \href{https://www.numdam.org/articles/10.5802/aif.1046/}{Published paper}.
\bibitem{IX-VII-VII-im716} N. Iwahori and H. Matsumoto, \emph{On some Bruhat decomposition and the structure of the Hecke rings of \(p\)-adic Chevalley groups}, Publ. Math. IH\'ES \textbf{25} (1965), 5--48. \href{https://doi.org/10.1007/BF02684396}{Published paper}.
\bibitem{IX-VII-VII-bzgeom716} I. N. Bernstein and A. V. Zelevinsky, \emph{Induced representations of reductive \(p\)-adic groups. I}, Ann. Sci. \'Ecole Norm. Sup. \textbf{10} (1977), 441--472, geometric lemma 2.12. \href{https://www.math.tau.ac.il/~bernstei/Publication_list/publication_texts/BZ-induced-1.pdf}{Authors' text}.
\bibitem{IX-VII-VII-amitsur720} S. A. Amitsur, \emph{On the characteristic polynomial of a sum of matrices}, Linear and Multilinear Algebra \textbf{8} (1980), no.~3, 177--182. \href{https://doi.org/10.1080/03081088008817315}{Published paper}.
\bibitem{IX-VII-VII-chttransfer721} L. Clozel, M. Harris and R. Taylor, \emph{Automorphy for some \(\ell\)-adic lifts of automorphic mod \(\ell\) Galois representations}, Publ. Math. IH\'ES \textbf{108} (2008), 1--181, \S3.1. \href{https://virtualmath1.stanford.edu/~rltaylor/twugfin.pdf}{Authors' text}.
\bibitem{IX-VII-bb433} V. Blomer and F. Brumley, \emph{On the Ramanujan conjecture over number fields}, Ann. of Math. (2) \textbf{174} (2011), 581--605, Theorem 1. \href{https://annals.math.princeton.edu/2011/174-1/p18}{Published paper}.
\bibitem{IX-VII-V-bs} A. Borel and J.-P. Serre, \emph{Corners and arithmetic groups}, Comment. Math. Helv. \textbf{48} (1973), 436--491. \href{https://www.e-periodica.ch/cntmng?pid=com-001:1973:48::32}{Original paper}.
\bibitem{IX-VII-cm471} The Stacks Project Authors, \emph{The Stacks Project}, Algebra, Lemma 10.104.2, \href{https://stacks.math.columbia.edu/tag/02JN}{Tag 02JN}.
\bibitem{IX-VII-steinberg476} R. Steinberg, \emph{Representations of algebraic groups}, Nagoya Math. J. \textbf{22} (1963), 33--56. \href{https://doi.org/10.1017/S0027763000011016}{Theorem 1.1}.
\bibitem{IX-VII-VI-bk117} S. Bloch and K. Kato, \href{https://virtualmath1.stanford.edu/~conrad/BSDseminar/refs/BKTamagawa.pdf}{\emph{L-functions and Tamagawa numbers of motives}}, in \emph{The Grothendieck Festschrift}, Vol.~I, Progr. Math. \textbf{86} (1990), 333--400, Section~3.
\bibitem{IX-VII-VII-labesse722} J.-P. Labesse, \emph{Changement de base CM et séries discrètes}, in \emph{On the Stabilization of the Trace Formula}, International Press, 2011, 429--470. \href{https://webusers.imj-prg.fr/~michael.harris/CBCM.pdf}{Author's text}.
\bibitem{IX-VII-VII-guerberoff722} L. Guerberoff, \emph{Modularity lifting theorems for Galois representations of unitary type}, Compos. Math. \textbf{147} (2011), 1022--1058. \href{https://doi.org/10.1112/S0010437X10005154}{Published paper}.
\bibitem{IX-VII-VII-mw722} C. M{\oe}glin and J.-L. Waldspurger, \emph{Le spectre résiduel de \(\mathrm{GL}(n)\)}, Ann. Sci. École Norm. Sup. \textbf{22} (1989), 605--674. \href{https://www.numdam.org/articles/10.24033/asens.1595/}{Published paper}.
\bibitem{IX-VII-VII-hltt722} M. Harris, K.-W. Lan, R. Taylor and J. Thorne, \emph{On the rigid cohomology of certain Shimura varieties}, Res. Math. Sci. \textbf{3} (2016), no.~37. \href{https://doi.org/10.1186/s40687-016-0078-5}{Published paper}.
\bibitem{IX-VII-VII-varma722} I. Varma, \emph{Local-global compatibility for regular algebraic cuspidal automorphic representations when \(\ell\ne p\)}, Forum Math. Sigma \textbf{12} (2024), e21. \href{https://doi.org/10.1017/fms.2024.7}{Published paper}.
\end{thebibliography}\endgroup
\end{document}
